% Generated mechanically from frozen input p04/ja/divisors.tex.
% Do not edit this generated file; edit the bound locale source instead.

% OK, start here.
%

\setcounter{chapter}{30}
\chapter{因子}



\phantomsection
\label{divisors-section-phantom}


\section{序論}
\label{divisors-section-introduction}

\noindent
本章では、因子の定義などに関わるいくつかの非常に基本的な問題を扱う。
基本的な参考文献は \cite{EGA} である。



\section{随伴点}
\label{divisors-section-associated}

\noindent
$R$ を環、$M$ を $R$-加群とする。
素イデアル $\mathfrak p \subset R$ が $M$ に {\it 随伴する} とは、
$M$ の元で、その零化イデアルが $\mathfrak p$ であるものが存在することをいう。
Algebra, Definition \ref{algebra-definition-associated} を参照せよ。
\cite[IV Definition 3.1.1]{EGA} に従い、スキーム上の準連接層に対する
随伴点を次のように定義する。

\begin{definition}
\label{divisors-definition-associated}
$X$ をスキームとする。
$\mathcal{F}$ を $X$ 上の準連接層とする。
\begin{enumerate}
\item $x \in X$ が $\mathcal{F}$ に {\it 随伴する} とは、極大イデアル
$\mathfrak m_x$ が $\mathcal{O}_{X, x}$-加群
$\mathcal{F}_x$ に随伴することをいう。
\item $\text{Ass}(\mathcal{F})$ または $\text{Ass}_X(\mathcal{F})$ と書くものを、
$\mathcal{F}$ の随伴点全体の集合とする。
\item {\it $X$ の随伴点} とは、$\mathcal{O}_X$ の随伴点のことである。
\end{enumerate}
\end{definition}

\noindent
これらの定義は、$X$ が局所 Noether であり、
$\mathcal{F}$ が有限型である場合に最も有用である。
たとえば、$X$ の既約成分の生成点が $X$ に随伴しないこともあり得る。
Example \ref{divisors-example-no-associated-prime} を参照せよ。
非 Noether の場合には弱随伴点を用いる方が便利なことがある。
Section \ref{divisors-section-weakly-associated} を参照せよ。
ここで、スキーム論的概念とアフィン開集合上の代数的概念との関係を示す。
ただし、この対応が完全に機能するのは局所 Noether スキームの場合に限られる。

\begin{lemma}
\label{divisors-lemma-associated-affine-open}
$X$ をスキームとし、$\mathcal{F}$ を $X$ 上の準連接層とする。
$\Spec(A) = U \subset X$ をアフィン開集合とし、
$M = \Gamma(U, \mathcal{F})$ とおく。
$x \in U$ とし、対応する素イデアルを $\mathfrak p \subset A$ とする。
\begin{enumerate}
\item $\mathfrak p$ が $M$ に随伴するならば、
$x$ は $\mathcal{F}$ に随伴する。
\item $\mathfrak p$ が有限生成ならば、逆も成り立つ。
\end{enumerate}
特に $X$ が局所 Noether ならば、上のようなすべての組
$(\mathfrak p, x)$ に対して同値
$$
\mathfrak p \in \text{Ass}(M) \Leftrightarrow x \in \text{Ass}(\mathcal{F})
$$
が成り立つ。
\end{lemma}

\begin{proof}
これは Algebra, Lemma \ref{algebra-lemma-associated-primes-localize}
から従うが、次のように直接論じることもできる。
$\mathfrak p$ が $M$ に随伴すると仮定する。
このとき $m \in M$ が存在し、その零化イデアルは
$\mathfrak p$ である。
局所化は完全なので、$\mathfrak pA_{\mathfrak p}$ は
$m/1 \in M_{\mathfrak p}$ の零化イデアルである。
$M_{\mathfrak p} = \mathcal{F}_x$
（Schemes, Lemma \ref{schemes-lemma-spec-sheaves}）であるから、
$x$ は $\mathcal{F}$ に随伴する。

\medskip\noindent
逆に、$x$ が $\mathcal{F}$ に随伴し、
$\mathfrak p$ が有限生成であると仮定する。
$x$ が $\mathcal{F}$ に随伴するので、元
$m' \in M_{\mathfrak p}$ が存在し、その零化イデアルは
$\mathfrak pA_{\mathfrak p}$ である。
$m' = m/f$ と書く。ただし $f \in A$, $f \not \in \mathfrak p$ とする。
零化イデアル $I$ を $m$ のものとする。これは
$A$ のイデアルであり、
$IA_{\mathfrak p} = \mathfrak pA_{\mathfrak p}$ を満たす。したがって
$I \subset \mathfrak p$ かつ $(\mathfrak p/I)_{\mathfrak p} = 0$ である。
$\mathfrak p$ は有限生成なので、ある
$g \in A$, $g \not \in \mathfrak p$ が存在して
$g(\mathfrak p/I) = 0$ となる。
従って $gm$ の零化イデアルは $\mathfrak p$ であり、結論を得る。

\medskip\noindent
$X$ が局所 Noether ならば、$A$ は Noether であり
（Properties, Lemma \ref{properties-lemma-locally-Noetherian}）、
$\mathfrak p$ は常に有限生成である。
\end{proof}

\begin{lemma}
\label{divisors-lemma-ass-support}
$X$ をスキームとする。
$\mathcal{F}$ を準連接 $\mathcal{O}_X$-加群とする。このとき
$\text{Ass}(\mathcal{F}) \subset \text{Supp}(\mathcal{F})$ である。
\end{lemma}

\begin{proof}
定義から直ちに従う。
\end{proof}

\begin{lemma}
\label{divisors-lemma-ses-ass}
$X$ をスキームとする。
$0 \to \mathcal{F}_1 \to \mathcal{F}_2 \to \mathcal{F}_3 \to 0$
を $X$ 上の準連接層の短完全列とする。このとき
$\text{Ass}(\mathcal{F}_2) \subset
\text{Ass}(\mathcal{F}_1) \cup \text{Ass}(\mathcal{F}_3)$
かつ
$\text{Ass}(\mathcal{F}_1) \subset \text{Ass}(\mathcal{F}_2)$
である。
\end{lemma}

\begin{proof}
任意の点 $x \in X$ に対し、茎の列
$0 \to \mathcal{F}_{1, x} \to \mathcal{F}_{2, x} \to \mathcal{F}_{3, x} \to 0$
は $\mathcal{O}_{X, x}$-加群の短完全列である。
したがって補題は Algebra, Lemma \ref{algebra-lemma-ass} から従う。
\end{proof}

\begin{lemma}
\label{divisors-lemma-finite-ass}
$X$ を局所 Noether スキームとする。
$\mathcal{F}$ を連接 $\mathcal{O}_X$-加群とする。
このとき $\text{Ass}(\mathcal{F}) \cap U$ は、
任意の準コンパクト開集合 $U \subset X$ に対して有限である。
\end{lemma}

\begin{proof}
Noether 環上の有限加群の随伴素イデアル全体は有限であるからである。
Algebra, Lemma \ref{algebra-lemma-finite-ass} を参照せよ。
スキームから代数へ移すには、$U$ が有限個のアフィン開集合の和であり、
各開集合が Noether 環のスペクトルであること
（Properties, Lemma \ref{properties-lemma-locally-Noetherian}）、
$\mathcal{F}$ がこの環上の有限加群に対応すること
（Cohomology of Schemes, Lemma \ref{coherent-lemma-coherent-Noetherian}）、
および最後に Lemma \ref{divisors-lemma-associated-affine-open} を用いればよい。
\end{proof}

\begin{lemma}
\label{divisors-lemma-ass-zero}
$X$ を局所 Noether スキームとし、$\mathcal{F}$ を準連接
$\mathcal{O}_X$-加群とする。このとき
$$
\mathcal{F} = 0 \Leftrightarrow \text{Ass}(\mathcal{F}) = \emptyset.
$$
\end{lemma}

\begin{proof}
$\mathcal{F} = 0$ ならば、定義により
$\text{Ass}(\mathcal{F}) = \emptyset$ である。
逆に $\text{Ass}(\mathcal{F}) = \emptyset$ ならば、
Algebra, Lemma \ref{algebra-lemma-ass-zero} により
$\mathcal{F} = 0$ である。
スキームから代数へ移すには、任意のアフィン開集合へ制限し、
Lemma \ref{divisors-lemma-associated-affine-open} を用いればよい。
\end{proof}

\begin{example}
\label{divisors-example-no-associated-prime}
$k$ を体とする。環 $R = k[x_1, x_2, x_3, \ldots]/(x_i^2)$ は局所環であり、
その極大イデアル $\mathfrak m$ は局所冪零である。
$R$ の元で、その零化イデアルが $\mathfrak m$ となるものは存在しない。
したがって $\text{Ass}(R) = \emptyset$ であり、
$X = \Spec(R)$ は随伴点をもたないスキームの例である。
\end{example}

\begin{lemma}
\label{divisors-lemma-restriction-injective-open-contains-ass}
$X$ を局所 Noether スキームとし、$\mathcal{F}$ を準連接
$\mathcal{O}_X$-加群とする。$U \subset X$ が開であり
$\text{Ass}(\mathcal{F}) \subset U$ ならば、
$\Gamma(X, \mathcal{F}) \to \Gamma(U, \mathcal{F})$ は単射である。
\end{lemma}

\begin{proof}
$s \in \Gamma(X, \mathcal{F})$ を、$U$ 上への制限が零となる切断とする。
$\mathcal{F}' \subset \mathcal{F}$ を写像
$\mathcal{O}_X \to \mathcal{F}$ の像とする。この写像は
$s$ により定まる。このとき
$\text{Supp}(\mathcal{F}') \cap U = \emptyset$ である。一方、
Lemma \ref{divisors-lemma-ses-ass} により
$\text{Ass}(\mathcal{F}') \subset \text{Ass}(\mathcal{F})$ である。
さらに
$\text{Ass}(\mathcal{F}') \subset \text{Supp}(\mathcal{F}')$
（Lemma \ref{divisors-lemma-ass-support}）でもあるから、
$\text{Ass}(\mathcal{F}') = \emptyset$ を得る。
従って Lemma \ref{divisors-lemma-ass-zero} により $\mathcal{F}' = 0$ である。
\end{proof}

\begin{lemma}
\label{divisors-lemma-minimal-support-in-ass}
$X$ を局所 Noether スキームとする。
$\mathcal{F}$ を準連接 $\mathcal{O}_X$-加群とする。
$x \in \text{Supp}(\mathcal{F})$ を $\mathcal{F}$ の台の点で、
$\text{Supp}(\mathcal{F})$ の他の点の特殊化ではないものとする。
このとき $x \in \text{Ass}(\mathcal{F})$ である。
特に、$X$ の任意の既約成分の生成点は $X$ の随伴点である。
\end{lemma}

\begin{proof}
$x \in \text{Supp}(\mathcal{F})$ なので、加群 $\mathcal{F}_x$ は零でない。
したがって Algebra, Lemma \ref{algebra-lemma-ass-zero} により
$\text{Ass}(\mathcal{F}_x) \subset \Spec(\mathcal{O}_{X, x})$
は空でない。一方、仮定により
$\text{Supp}(\mathcal{F}_x) = \{\mathfrak m_x\}$ である。
また
$\text{Ass}(\mathcal{F}_x) \subset \text{Supp}(\mathcal{F}_x)$
（Algebra, Lemma \ref{algebra-lemma-ass-support}）なので、
$\mathfrak m_x$ は $\mathcal{F}_x$ に随伴し、結論を得る。
\end{proof}

\noindent
次の補題は More on Algebra, Lemma
\ref{more-algebra-lemma-check-injective-on-ass} の類似である。

\begin{lemma}
\label{divisors-lemma-check-injective-on-ass}
$X$ を局所 Noether スキームとし、
$\varphi : \mathcal{F} \to \mathcal{G}$ を準連接
$\mathcal{O}_X$-加群の写像とする。
任意の $x \in X$ に対して次の少なくとも一方が成り立つと仮定する。
\begin{enumerate}
\item $\mathcal{F}_x \to \mathcal{G}_x$ は単射である、または
\item $x \not \in \text{Ass}(\mathcal{F})$ である。
\end{enumerate}
このとき $\varphi$ は単射である。
\end{lemma}

\begin{proof}
仮定から $\text{Ass}(\Ker(\varphi)) = \emptyset$ が従い、
したがって Lemma \ref{divisors-lemma-ass-zero} により $\Ker(\varphi) = 0$ である。
\end{proof}

\begin{lemma}
\label{divisors-lemma-check-isomorphism-via-depth-and-ass}
$X$ を局所 Noether スキームとし、
$\varphi : \mathcal{F} \to \mathcal{G}$ を準連接
$\mathcal{O}_X$-加群の写像とする。
$\mathcal{F}$ は連接であり、任意の $x \in X$ に対して
次のいずれか一方が成り立つと仮定する。
\begin{enumerate}
\item $\mathcal{F}_x \to \mathcal{G}_x$ は同型である、または
\item $\text{depth}(\mathcal{F}_x) \geq 2$ かつ
$x \not \in \text{Ass}(\mathcal{G})$ である。
\end{enumerate}
このとき $\varphi$ は同型である。
\end{lemma}

\begin{proof}
これは More on Algebra, Lemma
\ref{more-algebra-lemma-check-isomorphism-via-depth-and-ass}
をスキームの言葉へ移したものである。
\end{proof}




\section{射と随伴点}
\label{divisors-section-morphisms-associated}

\noindent
$f : X \to S$ をスキームの射とする。
$\mathcal{F}$ を $\mathcal{O}_X$-加群の層とする。
$s \in S$ が点であるとき、$\mathcal{F}_s$ で、
$\mathcal{O}_{X_s}$-加群、すなわち $\mathcal{F}$ を射
$i_s : X_s \to X$ によって引き戻して得られるものを表すのがしばしば便利である。
ここで $X_s$ は $f$ の $s$ 上のスキーム論的ファイバーである。
式で書けば
$$
\mathcal{F}_s = i_s^*\mathcal{F}
$$
である。もちろん、この記法は $\mathcal{F}$ の点 $x \in X$ における茎について
既に用いている記法と、$f = \text{id}_X$ の場合に衝突する。
しかし、次の補題の定式化のように、この記法はしばしば便利である。

\begin{lemma}
\label{divisors-lemma-bourbaki}
$f : X \to S$ をスキームの射とする。
$\mathcal{F}$ を $X$ 上の準連接層で $S$ 上平坦なものとする。
$\mathcal{G}$ を $S$ 上の準連接層とする。このとき
$$
\text{Ass}_X(\mathcal{F} \otimes_{\mathcal{O}_X} f^*\mathcal{G})
\supset
\bigcup\nolimits_{s \in \text{Ass}_S(\mathcal{G})}
\text{Ass}_{X_s}(\mathcal{F}_s)
$$
であり、$S$ が局所 Noether ならば等号が成り立つ
（記法 $\mathcal{F}_s$ については上を参照）。
\end{lemma}

\begin{proof}
$x \in X$ とし、$s = f(x) \in S$ とする。
$B = \mathcal{O}_{X, x}$, $A = \mathcal{O}_{S, s}$,
$N = \mathcal{F}_x$, $M = \mathcal{G}_s$ とおく。
$\mathcal{F} \otimes_{\mathcal{O}_X} f^*\mathcal{G}$ の $x$ における茎は
$B$-加群 $M \otimes_A N$ に等しいことに注意する。したがって
$x \in \text{Ass}_X(\mathcal{F} \otimes_{\mathcal{O}_X} f^*\mathcal{G})$
であることと
$\mathfrak m_B$ が $\text{Ass}_B(M \otimes_A N)$ に属することは同値である。
同様に、$s \in \text{Ass}_S(\mathcal{G})$ かつ
$x \in \text{Ass}_{X_s}(\mathcal{F}_s)$ であることと
$\mathfrak m_A \in \text{Ass}_A(M)$ かつ
$\mathfrak m_B/\mathfrak m_A B \in
\text{Ass}_{B \otimes \kappa(\mathfrak m_A)}(N \otimes \kappa(\mathfrak m_A))$
であることは同値である。従って補題は
Algebra, Lemma \ref{algebra-lemma-bourbaki-fibres} から従う。
\end{proof}





\section{埋め込まれた点}
\label{divisors-section-embedded}

\noindent
$R$ を環、$M$ を $R$-加群とする。
素イデアル $\mathfrak p \subset R$ が $M$ の
{\it 埋め込まれた随伴素イデアル}であるとは、それが $M$ の随伴素イデアルで、
$M$ の随伴素イデアルの中で極小でないことをいう。
Algebra, Definition \ref{algebra-definition-embedded-primes} を参照せよ。
\cite[IV Definition 3.1.1]{EGA} に従い、スキーム上の準連接層に対する
埋め込まれた随伴点を次のように定義する。

\begin{definition}
\label{divisors-definition-embedded}
$X$ をスキームとする。
$\mathcal{F}$ を $X$ 上の準連接層とする。
\begin{enumerate}
\item $\mathcal{F}$ の {\it 埋め込まれた随伴点}とは、
$\mathcal{F}$ の随伴点の中で極大でないもの、すなわち
$\mathcal{F}$ の別の随伴点の特殊化であるものをいう。
\item $x$ を $X$ の点とする。これが {\it 埋め込まれた点}であるとは、
$x$ が $\mathcal{O}_X$ の埋め込まれた随伴点であることをいう。
\item $X$ の {\it 埋め込まれた成分}とは、既約閉集合
$Z = \overline{\{x\}}$ で、その生成点 $x$ が
$X$ の埋め込まれた点であるものをいう。
\end{enumerate}
\end{definition}

\noindent
Noether の場合に $\mathcal{F}$ が連接ならば、次が成り立つ。

\begin{lemma}
\label{divisors-lemma-embedded}
$X$ を局所 Noether スキームとする。
$\mathcal{F}$ を連接 $\mathcal{O}_X$-加群とする。このとき
\begin{enumerate}
\item $\text{Supp}(\mathcal{F})$ の既約成分の生成点は
$\mathcal{F}$ の随伴点であり、
\item $\mathcal{F}$ の随伴点が埋め込まれていることと、
$\text{Supp}(\mathcal{F})$ の既約成分の生成点でないことは同値である。
\end{enumerate}
特に、$X$ の埋め込まれた点とは、$X$ の随伴点であって
$X$ の既約成分の生成点でないものである。
\end{lemma}

\begin{proof}
この場合 $Z = \text{Supp}(\mathcal{F})$ が閉であることを思い出そう。
Morphisms, Lemma \ref{morphisms-lemma-support-finite-type} を参照せよ。
また $Z$ の既約成分の生成点は $\mathcal{F}$ の随伴点である。
Lemma \ref{divisors-lemma-minimal-support-in-ass} を参照せよ。
最後に Lemma \ref{divisors-lemma-ass-support} により
$\text{Ass}(\mathcal{F}) \subset Z$ である。
これらの結果と、$Z$ がソバー位相空間であり、従って
$Z$ の任意の点が $Z$ の生成点の特殊化であるという事実を合わせると、
(1) と (2) が従う。
\end{proof}

\begin{lemma}
\label{divisors-lemma-S1-no-embedded}
$X$ を局所 Noether スキームとする。
$\mathcal{F}$ を $X$ 上の連接層とする。このとき次は同値である。
\begin{enumerate}
\item $\mathcal{F}$ は埋め込まれた随伴点をもたず、
\item $\mathcal{F}$ は性質 $(S_1)$ をもつ。
\end{enumerate}
\end{lemma}

\begin{proof}
これは Algebra, Lemma \ref{algebra-lemma-criterion-no-embedded-primes} と
上の Lemma \ref{divisors-lemma-associated-affine-open} を合わせたものである。
\end{proof}

\begin{lemma}
\label{divisors-lemma-noetherian-dim-1-CM-no-embedded-points}
$X$ を次元 $\leq 1$ の局所 Noether スキームとする。次は同値である。
\begin{enumerate}
\item $X$ は Cohen--Macaulay であり、
\item $X$ は埋め込まれた点をもたない。
\end{enumerate}
\end{lemma}

\begin{proof}
Lemma \ref{divisors-lemma-S1-no-embedded} と定義から従う。
\end{proof}

\begin{lemma}
\label{divisors-lemma-scheme-theoretically-dense-contain-embedded-points}
$X$ を局所 Noether スキームとし、$U \subset X$ を開部分スキームとする。
次は同値である。
\begin{enumerate}
\item $U$ は $X$ においてスキーム論的に稠密である
（Morphisms, Definition \ref{morphisms-definition-scheme-theoretically-dense}）。
\item $U$ は $X$ において稠密であり、$U$ は $X$ のすべての
埋め込まれた点を含む。
\item $U$ は $X$ のすべての随伴点を含む。
\end{enumerate}
\end{lemma}

\begin{proof}
問題は $X$ 上局所的なので、$X = \Spec(A)$ で
$A$ が Noether 環であると仮定してよい。このとき $U$ は準コンパクトであり
（Properties, Lemma \ref{properties-lemma-immersion-into-noetherian}）、
従って
$U = D(f_1) \cup \ldots \cup D(f_n)$
である（Algebra, Lemma \ref{algebra-lemma-qc-open}）。
この状況で $U$ が $X$ においてスキーム論的に稠密であることと
$A \to A_{f_1} \times \ldots \times A_{f_n}$ が単射であることは同値である。
Morphisms, Example \ref{morphisms-example-scheme-theoretic-closure} を参照せよ。
条件 (3) を代数へ移すと、任意の随伴素イデアル
$\mathfrak p$（$A$ のもの）に対して、ある $i$ が存在し、
$f_i \not \in \mathfrak p$ となるという意味になる。

\medskip\noindent
(1)、すなわち $A \to A_{f_1} \times \ldots \times A_{f_n}$ が単射であると仮定する。
$a \in A$ の零化イデアルが素イデアル $\mathfrak p$ であるとする。
$a$ は $A_{f_i}$ の零でない元へ写り、そのような $i$ が存在するので、
$f_i \not \in \mathfrak p$ である。従って (3) が成り立つ。

\medskip\noindent
$U$ が $\Spec(A)$ で稠密であることと、すべての極小素イデアルを
含むことは同値である。$A$ の随伴素イデアル全体は、
極小素イデアル全体と埋め込まれた随伴素イデアル全体との和に等しい
（Algebra, Proposition
\ref{algebra-proposition-minimal-primes-associated-primes} を参照）。
従って (2) と (3) は同値である。

\medskip\noindent
(3) を仮定する。これは、任意の随伴素イデアル
$\mathfrak p$（$A$ のもの）が $A_{f_i}$ の素イデアルに対応し、
そのような $i$ が存在することを意味する。
このとき
$A \to A_{f_1} \times \ldots \times A_{f_n}$ は単射である。実際、
Algebra, Lemma \ref{algebra-lemma-zero-at-ass-zero} により
$A \to \prod_{\mathfrak p \in \text{Ass}(A)} A_\mathfrak p$
は単射である。このようにして (3) から (1) が従う。
\end{proof}

\begin{lemma}
\label{divisors-lemma-remove-embedded-points}
$X$ を局所 Noether スキームとする。
$\mathcal{F}$ を $X$ 上の連接層とする。連接部分層全体の集合
$$
\{
\mathcal{K} \subset \mathcal{F}
\mid
\text{Supp}(\mathcal{K})\text{ is nowhere dense in }\text{Supp}(\mathcal{F})
\}
$$
は最大元 $\mathcal{K}$ をもつ。
$\mathcal{F}' = \mathcal{F}/\mathcal{K}$ とおくと、次が成り立つ。
\begin{enumerate}
\item $\text{Supp}(\mathcal{F}') = \text{Supp}(\mathcal{F})$。
\item $\mathcal{F}'$ は埋め込まれた随伴点をもたない。
\item 稠密開集合 $U \subset X$ が存在し、
$U \cap \text{Supp}(\mathcal{F})$ は $\text{Supp}(\mathcal{F})$ で稠密であり、
$\mathcal{F}'|_U \cong \mathcal{F}|_U$ である。
\end{enumerate}
\end{lemma}

\begin{proof}
これは Algebra, Lemmas \ref{algebra-lemma-remove-embedded-primes} および
\ref{algebra-lemma-remove-embedded-primes-localize} から従う。
$U$ は $\mathcal{F}$ の埋め込まれた随伴点全体の閉包の補集合として
取れることに注意する。
\end{proof}

\begin{lemma}
\label{divisors-lemma-no-embedded-points-endos}
$X$ を局所 Noether スキームとする。
$\mathcal{F}$ を埋め込まれた随伴点をもたない連接
$\mathcal{O}_X$-加群とする。次のようにおく。
$$
\mathcal{I}
=
\Ker(\mathcal{O}_X
\longrightarrow
\SheafHom_{\mathcal{O}_X}(\mathcal{F}, \mathcal{F})).
$$
これは埋め込まれた点をもたない閉部分スキーム
$Z \subset X$ を定める連接イデアル層である。さらに
連接層 $\mathcal{G}$ が $Z$ 上に存在し、
(a) $\mathcal{F} = (Z \to X)_*\mathcal{G}$、
(b) $\mathcal{G}$ は埋め込まれた随伴点をもたず、
(c) $\text{Supp}(\mathcal{G}) = Z$（集合として）
を満たすものが存在する。
\end{lemma}

\begin{proof}
いくつかの主張は Cohomology of Schemes, Lemma
\ref{coherent-lemma-coherent-support-closed} の証明で既に見た。
残りは Algebra, Lemma \ref{algebra-lemma-no-embedded-primes-endos}
から従う。
\end{proof}



\section{弱随伴点}
\label{divisors-section-weakly-associated}

\noindent
$R$ を環、$M$ を $R$-加群とする。
素イデアル $\mathfrak p \subset R$ が $M$ に {\it 弱随伴する} とは、
ある元 $m$ が $M$ に存在し、$\mathfrak p$ が $m$ の零化イデアルを含む
素イデアルの中で極小となるものが存在することをいう。
Algebra, Definition \ref{algebra-definition-weakly-associated} を参照せよ。
$R$ が極大イデアル $\mathfrak m$ をもつ局所環ならば、
$\mathfrak m$ が $M$ に弱随伴することと、元 $m \in M$ が存在し、
その零化イデアルの根基が $\mathfrak m$ であることは同値である。
Algebra, Lemma \ref{algebra-lemma-weakly-ass-local} を参照せよ。

\begin{definition}
\label{divisors-definition-weakly-associated}
$X$ をスキームとする。
$\mathcal{F}$ を $X$ 上の準連接層とする。
\begin{enumerate}
\item $x \in X$ が $\mathcal{F}$ に {\it 弱随伴する} とは、極大イデアル
$\mathfrak m_x$ が $\mathcal{O}_{X, x}$-加群
$\mathcal{F}_x$ に弱随伴することをいう。
\item $\text{WeakAss}(\mathcal{F})$ で
$\mathcal{F}$ の弱随伴点全体の集合を表す。
\item {\it $X$ の弱随伴点}とは、$\mathcal{O}_X$ の弱随伴点のことである。
\end{enumerate}
\end{definition}

\noindent
この場合、任意のアフィン開集合上で、この定義は上で定義した
弱随伴素イデアルと正確に対応する。正確には次の通りである。

\begin{lemma}
\label{divisors-lemma-weakly-associated-affine-open}
$X$ をスキームとし、$\mathcal{F}$ を $X$ 上の準連接層とする。
$\Spec(A) = U \subset X$ をアフィン開集合とし、
$M = \Gamma(U, \mathcal{F})$ とおく。
$x \in U$ とし、対応する素イデアルを $\mathfrak p \subset A$ とする。
次は同値である。
\begin{enumerate}
\item $\mathfrak p$ は $M$ に弱随伴し、
\item $x$ は $\mathcal{F}$ に弱随伴する。
\end{enumerate}
\end{lemma}

\begin{proof}
Algebra, Lemma \ref{algebra-lemma-weakly-ass-local} から従う。
\end{proof}

\begin{lemma}
\label{divisors-lemma-weakly-ass-support}
$X$ をスキームとする。
$\mathcal{F}$ を準連接 $\mathcal{O}_X$-加群とする。このとき
$$
\text{Ass}(\mathcal{F}) \subset \text{WeakAss}(\mathcal{F}) \subset
\text{Supp}(\mathcal{F}).
$$
\end{lemma}

\begin{proof}
定義から直ちに従う。
\end{proof}

\begin{lemma}
\label{divisors-lemma-ses-weakly-ass}
$X$ をスキームとする。
$0 \to \mathcal{F}_1 \to \mathcal{F}_2 \to \mathcal{F}_3 \to 0$
を $X$ 上の準連接層の短完全列とする。このとき
$\text{WeakAss}(\mathcal{F}_2) \subset
\text{WeakAss}(\mathcal{F}_1) \cup \text{WeakAss}(\mathcal{F}_3)$
かつ
$\text{WeakAss}(\mathcal{F}_1) \subset \text{WeakAss}(\mathcal{F}_2)$
である。
\end{lemma}

\begin{proof}
任意の点 $x \in X$ に対し、茎の列
$0 \to \mathcal{F}_{1, x} \to \mathcal{F}_{2, x} \to \mathcal{F}_{3, x} \to 0$
は $\mathcal{O}_{X, x}$-加群の短完全列である。
したがって補題は Algebra, Lemma \ref{algebra-lemma-weakly-ass}
から従う。
\end{proof}

\begin{lemma}
\label{divisors-lemma-weakly-ass-zero}
$X$ をスキームとする。
$\mathcal{F}$ を準連接 $\mathcal{O}_X$-加群とする。このとき
$$
\mathcal{F} = (0) \Leftrightarrow \text{WeakAss}(\mathcal{F}) = \emptyset
$$
である。
\end{lemma}

\begin{proof}
Lemma \ref{divisors-lemma-weakly-associated-affine-open} および
Algebra, Lemma \ref{algebra-lemma-weakly-ass-zero}
から従う。
\end{proof}

\begin{lemma}
\label{divisors-lemma-restriction-injective-open-contains-weakly-ass}
$X$ をスキームとし、$\mathcal{F}$ を準連接
$\mathcal{O}_X$-加群とする。$U \subset X$ が開であり
$\text{WeakAss}(\mathcal{F}) \subset U$ ならば、
$\Gamma(X, \mathcal{F}) \to \Gamma(U, \mathcal{F})$
は単射である。
\end{lemma}

\begin{proof}
$s \in \Gamma(X, \mathcal{F})$ を、$U$ 上への制限が零となる切断とする。
$\mathcal{F}' \subset \mathcal{F}$ を写像
$\mathcal{O}_X \to \mathcal{F}$ の像とする。この写像は
$s$ により定まる。このとき
$\text{Supp}(\mathcal{F}') \cap U = \emptyset$ である。一方、
Lemma \ref{divisors-lemma-ses-weakly-ass} により
$\text{WeakAss}(\mathcal{F}') \subset \text{WeakAss}(\mathcal{F})$ である。
さらに
$\text{WeakAss}(\mathcal{F}') \subset \text{Supp}(\mathcal{F}')$
（Lemma \ref{divisors-lemma-weakly-ass-support}）なので、
$\text{WeakAss}(\mathcal{F}') = \emptyset$ を得る。
従って Lemma \ref{divisors-lemma-weakly-ass-zero} により $\mathcal{F}' = 0$ である。
\end{proof}

\begin{lemma}
\label{divisors-lemma-minimal-support-in-weakly-ass}
$X$ をスキームとする。
$\mathcal{F}$ を準連接 $\mathcal{O}_X$-加群とする。
$x \in \text{Supp}(\mathcal{F})$ を $\mathcal{F}$ の台の点で、
$\text{Supp}(\mathcal{F})$ の他の点の特殊化ではないものとする。このとき
$x \in \text{WeakAss}(\mathcal{F})$ である。
特に、$X$ の任意の既約成分の生成点は
$\mathcal{O}_X$ に弱随伴する。
\end{lemma}

\begin{proof}
$x \in \text{Supp}(\mathcal{F})$ なので、加群 $\mathcal{F}_x$ は零でない。
したがって
$\text{WeakAss}(\mathcal{F}_x) \subset \Spec(\mathcal{O}_{X, x})$
は Algebra, Lemma \ref{algebra-lemma-weakly-ass-zero} により空でない。
一方、仮定により
$\text{Supp}(\mathcal{F}_x) = \{\mathfrak m_x\}$ である。
また
$\text{WeakAss}(\mathcal{F}_x) \subset \text{Supp}(\mathcal{F}_x)$
であるから
（Algebra, Lemma \ref{algebra-lemma-weakly-ass-support}）、
$\mathfrak m_x$ は $\mathcal{F}_x$ に弱随伴し、結論を得る。
\end{proof}

\begin{lemma}
\label{divisors-lemma-ass-weakly-ass}
$X$ をスキームとする。
$\mathcal{F}$ を準連接 $\mathcal{O}_X$-加群とする。
$\mathfrak m_x$ が $\mathcal{O}_{X, x}$ の有限生成イデアルならば
$$
x \in \text{Ass}(\mathcal{F}) \Leftrightarrow
x \in \text{WeakAss}(\mathcal{F}).
$$
である。特に $X$ が局所 Noether ならば
$\text{Ass}(\mathcal{F}) = \text{WeakAss}(\mathcal{F})$ である。
\end{lemma}

\begin{proof}
Algebra, Lemma \ref{algebra-lemma-ass-weakly-ass} を参照せよ。
\end{proof}

\begin{lemma}
\label{divisors-lemma-weakass-pushforward}
$f : X \to S$ をスキームの準コンパクトかつ準分離な射とする。
$\mathcal{F}$ を準連接 $\mathcal{O}_X$-加群とする。
$s \in S$ を $f$ の像に属さない点とする。このとき
$s$ は $f_*\mathcal{F}$ に弱随伴しない。
\end{lemma}

\begin{proof}
$f' : X' \to \Spec(\mathcal{O}_{S, s})$ を、
$f$ を射 $g : \Spec(\mathcal{O}_{S, s}) \to S$ で基底変換して得られるものとして考え、
もう一方の射影を $g' : X' \to X$ と書く。このとき
$$
(f_*\mathcal{F})_s = (g^*f_*\mathcal{F})_s = (f'_*(g')^*\mathcal{F})_s
$$
である。最初の等号は $g$ が $s$ における局所環上に同型を誘導することから、
二つ目は平坦基底変換
（Cohomology of Schemes, Lemma
\ref{coherent-lemma-flat-base-change-cohomology}）から従う。
もちろん $s \in \Spec(\mathcal{O}_{S, s})$ は $f'$ の像に属さない。
従って $S$ は局所環 $(A, \mathfrak m)$ のスペクトルであり、
$s$ は $\mathfrak m$ に対応すると仮定してよい。
Schemes, Lemma \ref{schemes-lemma-push-forward-quasi-coherent} により
層 $f_*\mathcal{F}$ は準連接であり、$A$-加群 $M$ に対応するとしよう。
$s$ は $f$ の像に属さないので、
$X = \bigcup_{a \in \mathfrak m} f^{-1}D(a)$
は開被覆である。$X$ は準コンパクトなので、
$a_1, \ldots, a_n \in \mathfrak m$ で
$X = f^{-1}D(a_1) \cup \ldots \cup f^{-1}D(a_n)$
となるものを取れる。従って
$$
M \to M_{a_1} \oplus \ldots \oplus M_{a_r}
$$
は単射である。このため、任意の非零元 $m$（茎 $M_\mathfrak p$ のもの）に対し、
ある $i$ が存在して $a_i^n m$ はすべての $n \geq 0$ について
非零となる。従って $\mathfrak m$ は $M$ に弱随伴しない。
\end{proof}

\begin{lemma}
\label{divisors-lemma-check-injective-on-weakass}
$X$ をスキームとし、
$\varphi : \mathcal{F} \to \mathcal{G}$ を準連接
$\mathcal{O}_X$-加群の写像とする。任意の $x \in X$ に対して
次の少なくとも一方が成り立つと仮定する。
\begin{enumerate}
\item $\mathcal{F}_x \to \mathcal{G}_x$ は単射である、または
\item $x \not \in \text{WeakAss}(\mathcal{F})$ である。
\end{enumerate}
このとき $\varphi$ は単射である。
\end{lemma}

\begin{proof}
仮定から $\text{WeakAss}(\Ker(\varphi)) = \emptyset$ が従い、
従って Lemma \ref{divisors-lemma-weakly-ass-zero} により $\Ker(\varphi) = 0$ である。
\end{proof}

\begin{lemma}
\label{divisors-lemma-depth-2-hartog}
$X$ を局所 Noether スキームとし、$\mathcal{F}$ を連接
$\mathcal{O}_X$-加群とする。$j : U \to X$ を、
$x \in X \setminus U$ に対して $\text{depth}(\mathcal{F}_x) \geq 2$
となる開部分スキームとする。このとき
$$
\mathcal{F} \longrightarrow j_*(\mathcal{F}|_U)
$$
は同型であり、従って
$\Gamma(X, \mathcal{F}) \to \Gamma(U, \mathcal{F})$
も同型である。
\end{lemma}

\begin{proof}
Lemma \ref{divisors-lemma-check-isomorphism-via-depth-and-ass} が
補題に表示された写像へ適用できることを示す。
$x \in X$ とする。$x \in U$ ならば
$j_*(\mathcal{F}|_U)|_U = \mathcal{F}|_U$ なので、この写像は茎上の同型である。
$x \in X \setminus U$ ならば
$x \not \in \text{Ass}(j_*(\mathcal{F}|_U))$ である
（Lemmas \ref{divisors-lemma-weakass-pushforward} および
\ref{divisors-lemma-weakly-ass-support}）。
$\text{depth}(\mathcal{F}_x) \geq 2$ を仮定しているので、証明は完了する。
\end{proof}

\begin{lemma}
\label{divisors-lemma-weakass-reduced}
$X$ を被約スキームとする。このとき $X$ の弱随伴点は、
$X$ の既約成分の生成点にちょうど一致する。
\end{lemma}

\begin{proof}
Algebra, Lemma \ref{algebra-lemma-reduced-weakly-ass-minimal} から従う。
\end{proof}



\section{射と弱随伴点}
\label{divisors-section-morphisms-weakly-associated}

\begin{lemma}
\label{divisors-lemma-weakly-ass-reverse-functorial}
$f : X \to S$ をスキームのアフィン射とする。
$\mathcal{F}$ を準連接 $\mathcal{O}_X$-加群とする。このとき
$$
\text{WeakAss}_S(f_*\mathcal{F}) \subset f(\text{WeakAss}_X(\mathcal{F}))
$$
である。
\end{lemma}

\begin{proof}
$X$ と $S$ はアフィンであると仮定してよい。このとき $X \to S$ は環準同型
$A \to B$ から得られる。$\mathcal{F} = \widetilde M$ と書く。
ただしこれはある $B$-加群 $M$ に対するものとする。
Lemma \ref{divisors-lemma-weakly-associated-affine-open} により、
$\mathcal{F}$ の弱随伴点は $M$ の弱随伴素イデアルと正確に対応する。
同様に、$f_*\mathcal{F}$ の弱随伴点は、$M$（$A$-加群としてのもの）の
弱随伴素イデアルと正確に対応する。
従って補題は Algebra, Lemma
\ref{algebra-lemma-weakly-ass-reverse-functorial} から従う。
\end{proof}

\begin{lemma}
\label{divisors-lemma-ass-functorial-equal}
$f : X \to S$ をスキームのアフィン射とする。
$\mathcal{F}$ を準連接 $\mathcal{O}_X$-加群とする。
$X$ が局所 Noether ならば
$$
f(\text{Ass}_X(\mathcal{F})) =
\text{Ass}_S(f_*\mathcal{F}) =
\text{WeakAss}_S(f_*\mathcal{F}) =
f(\text{WeakAss}_X(\mathcal{F}))
$$
である。
\end{lemma}

\begin{proof}
$X$ と $S$ はアフィンであると仮定してよい。このとき $X \to S$ は環準同型
$A \to B$ から得られる。$X$ は局所 Noether なので、環 $B$ は Noether である。
Properties, Lemma \ref{properties-lemma-locally-Noetherian} を参照せよ。
$\mathcal{F} = \widetilde M$ と書く。ただしこれはある $B$-加群 $M$ に対するものとする。
Lemma \ref{divisors-lemma-associated-affine-open} により、
$\mathcal{F}$ の随伴点は $M$ の随伴素イデアルと正確に対応し、
$M$（$A$-加群としてのもの）の任意の随伴素イデアルは
$f_*\mathcal{F}$ の随伴点となる。従って包含
$$
f(\text{Ass}_X(\mathcal{F})) \subset \text{Ass}_S(f_*\mathcal{F})
$$
は Algebra, Lemma \ref{algebra-lemma-ass-functorial-Noetherian}
から従う。また包含
$$
\text{Ass}_S(f_*\mathcal{F}) \subset \text{WeakAss}_S(f_*\mathcal{F})
$$
は Lemma \ref{divisors-lemma-weakly-ass-support} から従う。さらに包含
$$
\text{WeakAss}_S(f_*\mathcal{F}) \subset f(\text{WeakAss}_X(\mathcal{F}))
$$
は Lemma \ref{divisors-lemma-weakly-ass-reverse-functorial} から従う。
外側の集合は Lemma \ref{divisors-lemma-ass-weakly-ass} により等しいので、
すべての箇所で等号が成り立つ。
\end{proof}

\begin{lemma}
\label{divisors-lemma-weakly-associated-finite}
$f : X \to S$ をスキームの有限射とする。
$\mathcal{F}$ を準連接 $\mathcal{O}_X$-加群とする。このとき
$\text{WeakAss}(f_*\mathcal{F}) = f(\text{WeakAss}(\mathcal{F}))$ である。
\end{lemma}

\begin{proof}
$X$ と $S$ はアフィンであると仮定してよい。このとき $X \to S$ は有限環準同型
$A \to B$ から得られる。$\mathcal{F} = \widetilde M$ と書く。
ただしこれはある $B$-加群 $M$ に対するものとする。
Lemma \ref{divisors-lemma-weakly-associated-affine-open} により、
$\mathcal{F}$ の弱随伴点は $M$ の弱随伴素イデアルと正確に対応する。
同様に、$f_*\mathcal{F}$ の弱随伴点は、$M$（$A$-加群としてのもの）の
弱随伴素イデアルと正確に対応する。
従って補題は Algebra, Lemma
\ref{algebra-lemma-weakly-ass-finite-ring-map} から従う。
\end{proof}

\begin{lemma}
\label{divisors-lemma-weakly-ass-pullback}
$f : X \to S$ をスキームの射とし、$\mathcal{G}$ を準連接
$\mathcal{O}_S$-加群とする。$x \in X$ とし、$s = f(x)$ とする。
$f$ が $x$ において平坦で、点 $x$ がファイバー $X_s$ の生成点であり、
$s \in \text{WeakAss}_S(\mathcal{G})$ ならば
$x \in \text{WeakAss}(f^*\mathcal{G})$ である。
\end{lemma}

\begin{proof}
$A = \mathcal{O}_{S, s}$, $B = \mathcal{O}_{X, x}$,
$M = \mathcal{G}_s$ とする。
$m \in M$ を、その零化イデアル
$I = \{a \in A \mid am = 0\}$ の根基が $\mathfrak m_A$ となる元とする。
$m \otimes 1$ の零化イデアルは $I B$ である。これは
$A \to B$ が忠実平坦だからである。従って $\sqrt{I B} = \mathfrak m_B$ を示せば十分である。
これは、$B/\mathfrak m_AB$ の極大イデアルが局所冪零であること
（Algebra, Lemma \ref{algebra-lemma-minimal-prime-reduced-ring} を参照）
と、$\sqrt{I} = \mathfrak m_A$ という仮定から従う。
いくつかの詳細は省略する。
\end{proof}

\begin{lemma}
\label{divisors-lemma-weakly-ass-change-fields}
$K/k$ を体拡大とし、$X$ を $k$ 上のスキームとする。
$\mathcal{F}$ を準連接 $\mathcal{O}_X$-加群とする。
$y \in X_K$ の像を $x \in X$ とする。
$y$ が引き戻し $\mathcal{F}_K$ の弱随伴点ならば、
$x$ は $\mathcal{F}$ の弱随伴点である。
\end{lemma}

\begin{proof}
これは Algebra, Lemma \ref{algebra-lemma-weakly-ass-change-fields}
をスキームの言葉へ移したものである。
\end{proof}

\noindent
次は、順像の深さがしばしば 2 以上となることを示す簡単な補題である。

\begin{lemma}
\label{divisors-lemma-depth-pushforward}
$f : X \to S$ をスキームの準コンパクトかつ準分離な射とする。
$\mathcal{F}$ を準連接 $\mathcal{O}_X$-加群とする。
$s \in S$ とする。
\begin{enumerate}
\item $s \not \in f(X)$ ならば、$s$ は $f_*\mathcal{F}$ に弱随伴しない。
\item $s \not \in f(X)$ かつ $\mathcal{O}_{S, s}$ が Noether ならば、
$s$ は $f_*\mathcal{F}$ に随伴しない。
\item $s \not \in f(X)$ であり、$(f_*\mathcal{F})_s$ が有限
$\mathcal{O}_{S, s}$-加群で、かつ $\mathcal{O}_{S, s}$ が Noether ならば
$\text{depth}((f_*\mathcal{F})_s) \geq 2$ である。
\item $\mathcal{F}$ が $S$ 上平坦で、$a \in \mathfrak m_s$ が非零因子ならば、
$a$ は $(f_*\mathcal{F})_s$ 上の非零因子である。
\item $\mathcal{F}$ が $S$ 上平坦で、$a, b \in \mathfrak m_s$ が正則列ならば、
$a$ は $(f_*\mathcal{F})_s$ 上の非零因子であり、
$b$ は $(f_*\mathcal{F})_s/a(f_*\mathcal{F})_s$ 上の非零因子である。
\item $\mathcal{F}$ が $S$ 上平坦で、$(f_*\mathcal{F})_s$ が有限
$\mathcal{O}_{S, s}$-加群ならば
$\text{depth}((f_*\mathcal{F})_s) \geq
\min(2, \text{depth}(\mathcal{O}_{S, s}))$ である。
\end{enumerate}
\end{lemma}

\begin{proof}
(1) は Lemma \ref{divisors-lemma-weakass-pushforward} である。
(2) は (1) と Lemma \ref{divisors-lemma-ass-weakly-ass} から従う。

\medskip\noindent
(3) の証明。深さが $\geq 2$ であることを示すには
$\Hom_{\mathcal{O}_{S, s}}(\kappa(s), (f_*\mathcal{F})_s) = 0$ および
$\Ext^1_{\mathcal{O}_{S, s}}(\kappa(s), (f_*\mathcal{F})_s) = 0$
を示せば十分である。
Algebra, Lemma \ref{algebra-lemma-depth-ext} を参照せよ。
完全列
$0 \to \mathfrak m_s \to \mathcal{O}_{S, s} \to \kappa(s) \to 0$
を用いると、写像
$$
\Hom_{\mathcal{O}_{S, s}}(\mathcal{O}_{S, s}, (f_*\mathcal{F})_s)
\to
\Hom_{\mathcal{O}_{S, s}}(\mathfrak m_s, (f_*\mathcal{F})_s)
$$
が同型であることを示せば十分である。
平坦基底変換
（Cohomology of Schemes, Lemma
\ref{coherent-lemma-flat-base-change-cohomology}）により、
$S$ を $\Spec(\mathcal{O}_{S, s})$ で置き換え、
$X$ を $\Spec(\mathcal{O}_{S, s}) \times_S X$ で置き換えてよい。
$\mathfrak m \subset \mathcal{O}_S$ で $s$ のイデアル層を表す。このとき
$$
\Hom_{\mathcal{O}_{S, s}}(\mathfrak m_s, (f_*\mathcal{F})_s) =
\Hom_{\mathcal{O}_S}(\mathfrak m, f_*\mathcal{F}) =
\Hom_{\mathcal{O}_X}(f^*\mathfrak m, \mathcal{F})
$$
である。最初の等号は $S$ が閉点 $s$ をもつ局所スキームだからであり、
二つ目は準連接加群上の $f^*, f_*$ に対する随伴から従う。
しかし $s \not \in f(X)$ なので $f^*\mathfrak m = \mathcal{O}_X$ である。
議論を逆向きにたどれば、求める等号を得る。

\medskip\noindent
(4), (5), (6) の証明には平坦基底変換
（Cohomology of Schemes, Lemma
\ref{coherent-lemma-flat-base-change-cohomology}）を用いて、
$S$ が $\mathcal{O}_{S, s}$ のスペクトルである場合へ帰着する。
このとき非零因子 $a \in \mathcal{O}_{S, s}$ は短完全列
$$
0 \to \mathcal{O}_S \xrightarrow{a} \mathcal{O}_S \to
\mathcal{O}_S/a \mathcal{O}_S \to 0
$$
を定める。$\mathcal{F}$ は $S$ 上平坦なので、完全列
$$
0 \to \mathcal{F} \xrightarrow{a} \mathcal{F} \to
\mathcal{F}/a\mathcal{F} \to 0
$$
を得る。順像を取ると完全列
$$
0 \to f_*\mathcal{F} \xrightarrow{a} f_*\mathcal{F} \to
f_*(\mathcal{F}/a\mathcal{F})
$$
を得る。これにより (4) が証明され、また
$f_*\mathcal{F}/ af_*\mathcal{F} \subset f_*(\mathcal{F}/a\mathcal{F})$
であることが分かる。$b$ が
$\mathcal{O}_{S, s}/a\mathcal{O}_{S, s}$ 上の非零因子ならば、
全く同じ議論により
$b : \mathcal{F}/a\mathcal{F} \to \mathcal{F}/a\mathcal{F}$
は単射である。順像を取ると
$$
b : f_*(\mathcal{F}/a\mathcal{F}) \to f_*(\mathcal{F}/a\mathcal{F})
$$
は単射であり、従って
$b : f_*\mathcal{F}/ af_*\mathcal{F} \to f_*\mathcal{F}/ af_*\mathcal{F}$
も単射である。これで (5) が証明された。(6) は (4), (5) と定義から従う。
\end{proof}










\section{相対随伴点集合}
\label{divisors-section-relative-assassin}

\noindent
$A \to B$ を環準同型とし、$N$ を $B$-加群とする。素イデアル
$\mathfrak q \subset B$ が $N$ の $B/A$ 上の相対随伴素イデアル集合に属するとは、
$\mathfrak q$ が $N \otimes_A \kappa(\mathfrak p)$ の随伴素イデアルであることをいう。
ここで $\mathfrak p = A \cap \mathfrak q$ である。
Algebra, Definition \ref{algebra-definition-relative-assassin} を参照せよ。
以下では、スキームの射上の準連接層に対する相対随伴点集合を定義する。

\begin{definition}
\label{divisors-definition-relative-assassin}
$f : X \to S$ をスキームの射とする。
$\mathcal{F}$ を準連接 $\mathcal{O}_X$-加群とする。
{\it $\mathcal{F}$ の、$X$ における $S$ 上の相対随伴点集合}とは、集合
$$
\text{Ass}_{X/S}(\mathcal{F}) =
\bigcup\nolimits_{s \in S} \text{Ass}_{X_s}(\mathcal{F}_s)
$$
をいう。ここで $\mathcal{F}_s = (X_s \to X)^*\mathcal{F}$ は、
$\mathcal{F}$ を $f$ の $s$ におけるファイバーへ制限したものである。
\end{definition}

\noindent
ここでも、この概念は $f$ のファイバーが局所 Noether で、
$\mathcal{F}$ が有限型である場合に最もよく用いられるという注意がある。
一般の場合には、次節で定義する相対弱随伴点集合を用いるべきであろう。
アフィン開集合上でスキーム論的概念を代数的概念と結び付けよう。
この対応が完全に成り立つのは、ファイバーが局所 Noether である
スキームの射についてだけであることに注意せよ。

\begin{lemma}
\label{divisors-lemma-relative-assassin-affine-open}
$f : X \to S$ をスキームの射とする。
$\mathcal{F}$ を $X$ 上の準連接層とする。
$U \subset X$ と $V \subset S$ を、$f(U) \subset V$ を満たすアフィン開集合とする。
$U = \Spec(A)$, $V = \Spec(R)$ と書き、
$M = \Gamma(U, \mathcal{F})$ と置く。
$x \in U$ とし、$\mathfrak p \subset A$ を対応する素イデアルとする。このとき
$$
\mathfrak p \in \text{Ass}_{A/R}(M) \Rightarrow
x \in \text{Ass}_{X/S}(\mathcal{F})
$$
である。すべてのファイバー $X_s$（$f$ のファイバー）が局所 Noether ならば、
$\mathfrak p \in \text{Ass}_{A/R}(M) \Leftrightarrow
x \in \text{Ass}_{X/S}(\mathcal{F})$
が上のすべての対 $(\mathfrak p, x)$ について成り立つ。
\end{lemma}

\begin{proof}
集合 $\text{Ass}_{A/R}(M)$ は Algebra, Definition
\ref{algebra-definition-relative-assassin} で定義される。
対 $(\mathfrak p, x)$ を一つ取り、$s = f(x)$ と置く。
$\mathfrak r \subset R$ を $\mathfrak p$ の下にある素イデアル、
すなわち $s$ に対応する素イデアルとする。
$\mathfrak p' \subset A \otimes_R \kappa(\mathfrak r)$ を、
その逆像が $\mathfrak p$ となる素イデアル、すなわち $x$ をそのファイバー
$X_s$ の点と見たときに対応する素イデアルとする。
このとき $\mathfrak p \in \text{Ass}_{A/R}(M)$ であることと、
$\mathfrak p'$ が $M \otimes_R \kappa(\mathfrak r)$ の随伴素イデアルであることは同値である。
Algebra, Lemma \ref{algebra-lemma-compare-relative-assassins} を参照せよ。
環 $A \otimes_R \kappa(\mathfrak r)$ は $U_s$ に対応し、
加群 $M \otimes_R \kappa(\mathfrak r)$ は準連接層
$\mathcal{F}_s|_{U_s}$ に対応することに注意せよ。
従って Lemma \ref{divisors-lemma-associated-affine-open} により、
$x$ は $\mathcal{F}_s$ の随伴点である。
$\mathfrak p'$ が有限生成ならば逆向きも成り立ち、
これにより最後の文が正しいことが分かる。
\end{proof}

\begin{lemma}
\label{divisors-lemma-base-change-relative-assassin}
$f : X \to S$ をスキームの射とする。
$\mathcal{F}$ を準連接 $\mathcal{O}_X$-加群とする。
$g : S' \to S$ をスキームの射とする。基底変換図
$$
\xymatrix{
X' \ar[d] \ar[r]_{g'} & X \ar[d] \\
S' \ar[r]^g & S
}
$$
を考え、$\mathcal{F}' = (g')^*\mathcal{F}$ と置く。
$x' \in X'$ を、その像が $x \in X$, $s' \in S'$, $s \in S$ となる点とする。
$f$ は局所有限型であると仮定する。このとき
$x' \in \text{Ass}_{X'/S'}(\mathcal{F}')$ であることと、
$x \in \text{Ass}_{X/S}(\mathcal{F})$ であり、かつ $x'$ が
$\Spec(\kappa(s') \otimes_{\kappa(s)} \kappa(x))$ の不可約成分の
生成点に対応することは同値である。
\end{lemma}

\begin{proof}
ファイバーの射 $X'_{s'} \to X_s$ を考える。
$X_{s'} = X_s \times_{\Spec(\kappa(s))} \Spec(\kappa(s'))$
なので、これは平坦射である。さらに $\mathcal{F}'_{s'}$ はこの射による
$\mathcal{F}_s$ の引き戻しである。$X_s$ は Noether スキーム
$\Spec(\kappa(s))$ 上局所有限型なので、$X_s$ は局所 Noether である。
Morphisms, Lemma \ref{morphisms-lemma-finite-type-noetherian} を参照せよ。
従って Lemma \ref{divisors-lemma-bourbaki} を適用でき、
$$
\text{Ass}_{X'_{s'}}(\mathcal{F}'_{s'}) =
\bigcup\nolimits_{x \in \text{Ass}(\mathcal{F}_s)} \text{Ass}((X'_{s'})_x).
$$
を得る。従って補題を証明するには、ファイバー $(X'_{s'})_x$、
すなわち射 $X'_{s'} \to X_s$ の $x$ 上のファイバーの随伴点が
その生成点であることを示せば十分である。
スキームとして
$(X'_{s'})_x = \Spec(\kappa(s') \otimes_{\kappa(s)} \kappa(x))$
であることに注意せよ。
Algebra, Lemma \ref{algebra-lemma-tensor-fields-CM} により、環
$\kappa(s') \otimes_{\kappa(s)} \kappa(x)$ は Noether Cohen--Macaulay 環である。
従ってその随伴素イデアルは極小素イデアルである。
Algebra, Proposition \ref{algebra-proposition-minimal-primes-associated-primes}
（極小素イデアルは随伴素イデアルである）および
Algebra, Lemma \ref{algebra-lemma-criterion-no-embedded-primes}
（埋め込まれた素イデアルはない）を参照せよ。
\end{proof}

\begin{remark}
\label{divisors-remark-base-change-relative-assassin}
Lemma \ref{divisors-lemma-base-change-relative-assassin} と同じ記号と仮定の下では、常に
$(g')^{-1}(\text{Ass}_{X/S}(\mathcal{F})) \supset
\text{Ass}_{X'/S'}(\mathcal{F}')$
である。射 $S' \to S$ が局所準有限ならば、実際
$$
(g')^{-1}(\text{Ass}_{X/S}(\mathcal{F}))
=
\text{Ass}_{X'/S'}(\mathcal{F}')
$$
である。これは、この場合、体拡大 $\kappa(s')/\kappa(s)$ が常に有限だからである。
実際、任意の射 $g : S' \to S$ に対して、すべての体拡大
$\kappa(s')/\kappa(s)$ が代数的ならば、より一般にこのことが成り立つ。この場合、
$\kappa(s') \otimes_{\kappa(s)} \kappa(x)$ のすべての素イデアルが
極大（かつ極小）素イデアルだからである。
Algebra, Lemma \ref{algebra-lemma-integral-over-field} を参照せよ。
\end{remark}




\section{相対弱随伴点集合}
\label{divisors-section-relative-weak-assassin}

\begin{definition}
\label{divisors-definition-relative-weak-assassin}
$f : X \to S$ をスキームの射とする。
$\mathcal{F}$ を準連接 $\mathcal{O}_X$-加群とする。
{\it $\mathcal{F}$ の、$X$ における $S$ 上の相対弱随伴点集合}とは、集合
$$
\text{WeakAss}_{X/S}(\mathcal{F}) =
\bigcup\nolimits_{s \in S} \text{WeakAss}(\mathcal{F}_s)
$$
をいう。ここで $\mathcal{F}_s = (X_s \to X)^*\mathcal{F}$ は、
$\mathcal{F}$ を $f$ の $s$ におけるファイバーへ制限したものである。
\end{definition}

\begin{lemma}
\label{divisors-lemma-relative-weak-assassin-assassin-finite-type}
$f : X \to S$ を局所有限型のスキームの射とする。
$\mathcal{F}$ を準連接 $\mathcal{O}_X$-加群とする。このとき
$\text{WeakAss}_{X/S}(\mathcal{F}) = \text{Ass}_{X/S}(\mathcal{F})$ である。
\end{lemma}

\begin{proof}
これは $f$ のファイバーが局所 Noether スキームであり、
局所 Noether スキーム上では随伴点と弱随伴点が一致するためである。
Lemma \ref{divisors-lemma-ass-weakly-ass} を参照せよ。
\end{proof}

\begin{lemma}
\label{divisors-lemma-relative-weak-assassin-finite}
$f : X \to S$ をスキームの射とする。
$i : Z \to X$ を有限射とする。
$\mathcal{F}$ を準連接 $\mathcal{O}_Z$-加群とする。このとき
$\text{WeakAss}_{X/S}(i_*\mathcal{F}) =
i(\text{WeakAss}_{Z/S}(\mathcal{F}))$
である。
\end{lemma}

\begin{proof}
$i_s : Z_s \to X_s$ をファイバー間に誘導される射とする。このとき
$(i_*\mathcal{F})_s = i_{s, *}(\mathcal{F}_s)$ である。
これは Cohomology of Schemes, Lemma \ref{coherent-lemma-affine-base-change}
と $i$ がアフィンであることから従う。従って
Lemma \ref{divisors-lemma-weakly-associated-finite} を適用すれば結論を得る。
\end{proof}





\section{Fitting イデアル}
\label{divisors-section-fitting-ideals}

\noindent
この節は More on Algebra, Section
\ref{more-algebra-section-fitting-ideals} における議論の続きである。
$S$ をスキームとし、$\mathcal{F}$ を有限型準連接
$\mathcal{O}_S$-加群とする。この状況では Fitting イデアル
$$
0 = \text{Fit}_{-1}(\mathcal{F}) \subset \text{Fit}_0(\mathcal{F}) \subset
\text{Fit}_1(\mathcal{F}) \subset \ldots \subset \mathcal{O}_S
$$
を、次の性質によって特徴づけられる準連接イデアルの列として構成できる。
アフィン開集合 $U = \Spec(A)$（$S$ のもの）を任意に取り、
$\mathcal{F}|_U$ が $A$-加群 $M$ に対応するならば、
$\text{Fit}_i(\mathcal{F})|_U$ はイデアル $\text{Fit}_i(M) \subset A$
に対応する。この構成は well-defined であり、準連接イデアル層となる。
実際、$f \in A$ ならば、第 $i$ Fitting イデアル、すなわち
$M_f$ を $A_f$-加群と見たときのものは、
More on Algebra, Lemma \ref{more-algebra-lemma-fitting-ideal-basics} により
$\text{Fit}_i(M) A_f$ に等しい。

\medskip\noindent
別の方法として、$\mathcal{F}$ の局所表示を用いて Fitting イデアルを構成できる。
すなわち、$U \subset X$ が開で、
$$
\bigoplus\nolimits_{i \in I} \mathcal{O}_U \to
\mathcal{O}_U^{\oplus n} \to \mathcal{F}|_U \to 0
$$
が $\mathcal{F}$ の $U$ 上での表示ならば、
$\text{Fit}_r(\mathcal{F})|_U$ は、その表示の最初の矢印を定める行列の
$(n - r) \times (n - r)$ 小行列式によって生成される。
これは上の構成と両立する。環上の加群の Fitting イデアルは、
まさにこの方法で定義されるからである。いくつかの詳細は省略する。

\begin{lemma}
\label{divisors-lemma-base-change-fitting-ideal}
$f : T \to S$ をスキームの射とする。
$\mathcal{F}$ を有限型準連接 $\mathcal{O}_S$-加群とする。このとき
$f^{-1}\text{Fit}_i(\mathcal{F}) \cdot \mathcal{O}_T =
\text{Fit}_i(f^*\mathcal{F})$
である。
\end{lemma}

\begin{proof}
More on Algebra, Lemma
\ref{more-algebra-lemma-fitting-ideal-basics} の (3) から直ちに従う。
\end{proof}

\begin{lemma}
\label{divisors-lemma-fitting-ideal-of-finitely-presented}
$S$ をスキームとする。
$\mathcal{F}$ を有限表示 $\mathcal{O}_S$-加群とする。このとき
$\text{Fit}_r(\mathcal{F})$ は有限型準連接イデアルである。
\end{lemma}

\begin{proof}
More on Algebra, Lemma
\ref{more-algebra-lemma-fitting-ideal-basics} の (4) から直ちに従う。
\end{proof}

\begin{lemma}
\label{divisors-lemma-on-subscheme-cut-out-by-Fit-0}
$S$ をスキームとする。
$\mathcal{F}$ を有限型準連接 $\mathcal{O}_S$-加群とする。
$Z_0 \subset S$ を $\text{Fit}_0(\mathcal{F})$ によって切り出される閉部分スキームとする。
$Z \subset S$ を $\mathcal{F}$ のスキーム論的台とする。このとき
\begin{enumerate}
\item 閉部分スキームとして $Z \subset Z_0 \subset S$ である。
\item 閉部分集合として $Z = Z_0 = \text{Supp}(\mathcal{F})$ である。
\item 有限型準連接 $\mathcal{O}_{Z_0}$-加群 $\mathcal{G}_0$ で
$$
(Z_0 \to X)_*\mathcal{G}_0 = \mathcal{F}.
$$
を満たすものが存在する。
\end{enumerate}
\end{lemma}

\begin{proof}
$Z$ は局所的には $\mathcal{F}$ の零化イデアルによって切り出されることを思い出そう。
Morphisms, Definition \ref{morphisms-definition-scheme-theoretic-support}
を参照せよ（これは $Z$ を定義するために Morphisms, Lemma
\ref{morphisms-lemma-scheme-theoretic-support} を用いる）。
従って More on Algebra, Lemma
\ref{more-algebra-lemma-fitting-ideal-basics} の (6) により、
スキーム論的に $Z \subset Z_0$ である。一方、Morphisms, Lemma
\ref{morphisms-lemma-scheme-theoretic-support} により集合論的に
$Z = \text{Supp}(\mathcal{F})$ であり、More on Algebra, Lemma
\ref{more-algebra-lemma-fitting-ideal-basics} の (7) により集合論的に
$Z_0 = Z$ である。最後に (3) のような $\mathcal{G}_0$ を得るには、
Morphisms, Lemma \ref{morphisms-lemma-scheme-theoretic-support} における
$\mathcal{G}$（$Z$ 上のもの）を用いて
$\mathcal{G}_0 = (Z \to Z_0)_*\mathcal{G}$ と置いてもよいし、
Morphisms, Lemma \ref{morphisms-lemma-i-star-equivalence} と、
More on Algebra, Lemma \ref{more-algebra-lemma-fitting-ideal-basics} の (6) により
$\text{Fit}_0(\mathcal{F})$ が $\mathcal{F}$ を零化することを用いてもよい。
\end{proof}

\begin{lemma}
\label{divisors-lemma-fitting-ideal-generate-locally}
$S$ をスキームとし、$\mathcal{F}$ を有限型準連接
$\mathcal{O}_S$-加群とする。$s \in S$ とする。このとき、
$\mathcal{F}$ が $r$ 個の元により生成されるような $s$ の近傍が存在することと
$\text{Fit}_r(\mathcal{F})_s = \mathcal{O}_{S, s}$ であることは同値である。
\end{lemma}

\begin{proof}
More on Algebra, Lemma
\ref{more-algebra-lemma-fitting-ideal-generate-locally} から直ちに従う。
\end{proof}

\begin{lemma}
\label{divisors-lemma-fitting-ideal-finite-locally-free}
$S$ をスキームとし、$\mathcal{F}$ を有限型準連接
$\mathcal{O}_S$-加群とする。$r \geq 0$ とする。次は同値である。
\begin{enumerate}
\item $\mathcal{F}$ は階数 $r$ の有限局所自由加群である。
\item $\text{Fit}_{r - 1}(\mathcal{F}) = 0$ かつ
$\text{Fit}_r(\mathcal{F}) = \mathcal{O}_S$ である。
\item $\text{Fit}_k(\mathcal{F}) = 0$ が $k < r$ に対して成り立ち、
$\text{Fit}_k(\mathcal{F}) = \mathcal{O}_S$ が $k \geq r$ に対して成り立つ。
\end{enumerate}
\end{lemma}

\begin{proof}
More on Algebra, Lemma
\ref{more-algebra-lemma-fitting-ideal-finite-locally-free} から直ちに従う。
\end{proof}

\begin{lemma}
\label{divisors-lemma-locally-free-rank-r-pullback}
$S$ をスキームとし、$\mathcal{F}$ を有限型準連接
$\mathcal{O}_S$-加群とする。$\mathcal{F}$ の Fitting イデアルによって定義される閉部分スキーム
$$
S = Z_{-1} \supset Z_0 \supset Z_1 \supset Z_2 \ldots
$$
は次の性質を持つ。
\begin{enumerate}
\item 共通部分 $\bigcap Z_r$ は空である。
\item 規則
$(\Sch/S)^{opp} \to \textit{Sets}$
$$
T \longmapsto
\left\{
\begin{matrix}
\{*\} & \text{if }\mathcal{F}_T\text{ is locally generated by }
\leq r\text{ sections} \\
\emptyset & \text{otherwise}
\end{matrix}
\right.
$$
によって定義される関手は、開部分スキーム $S \setminus Z_r$ により表現される。
\item 規則によって定義される関手
$F_r : (\Sch/S)^{opp} \to \textit{Sets}$
$$
T \longmapsto
\left\{
\begin{matrix}
\{*\} & \text{if }\mathcal{F}_T\text{ locally free rank }r\\
\emptyset & \text{otherwise}
\end{matrix}
\right.
$$
は、$Z_{r - 1} \setminus Z_r$（$S$ の局所閉部分スキーム）により表現される。
\end{enumerate}
$\mathcal{F}$ が有限表示ならば、$Z_r \to S$,
$S \setminus Z_r \to S$, $Z_{r - 1} \setminus Z_r \to S$ は有限表示である。
\end{lemma}

\begin{proof}
(1) は、任意のアフィン開集合 $U$ 上で、ある整数 $n$ が存在して
$\text{Fit}_n(\mathcal{F})|_U = \mathcal{O}_U$ となるため正しい。
実際、$n$ として $\mathcal{F}$ の $U$ 上での生成元の個数を取れる。
More on Algebra, Section \ref{more-algebra-section-fitting-ideals} を参照せよ。

\medskip\noindent
任意の射 $g : T \to S$ に対し、Lemmas
\ref{divisors-lemma-base-change-fitting-ideal} および
\ref{divisors-lemma-fitting-ideal-generate-locally} により、
$\mathcal{F}_T$ が局所的に $\leq r$ 個の切断で生成されることと
$\text{Fit}_r(\mathcal{F}) \cdot \mathcal{O}_T = \mathcal{O}_T$
であることは同値である。これで (2) が証明された。

\medskip\noindent
任意の射 $g : T \to S$ に対し、Lemmas
\ref{divisors-lemma-base-change-fitting-ideal} および
\ref{divisors-lemma-fitting-ideal-finite-locally-free} により、
$\mathcal{F}_T$ が階数 $r$ の自由加群であることと
$\text{Fit}_r(\mathcal{F}) \cdot \mathcal{O}_T = \mathcal{O}_T$ かつ
$\text{Fit}_{r - 1}(\mathcal{F}) \cdot \mathcal{O}_T = 0$
であることは同値である。これで (3) が証明された。

\medskip\noindent
$\mathcal{F}$ が有限表示であると仮定する。このとき射
$Z_r \to S$ はそれぞれ有限表示である。実際、
$\text{Fit}_r(\mathcal{F})$ は有限型である
（Lemma \ref{divisors-lemma-fitting-ideal-of-finitely-presented} および
Morphisms, Lemma \ref{morphisms-lemma-closed-immersion-finite-presentation}）。
これは $Z_{r - 1} \setminus Z_r$ が $Z_r$ における逆コンパクトな開集合であることを含意する
（Properties, Lemma \ref{properties-lemma-quasi-coherent-finite-type-ideals}）。
従って射 $Z_{r - 1} \setminus Z_r \to Z_r$ も有限表示である。
\end{proof}

\noindent
Lemma \ref{divisors-lemma-locally-free-rank-r-pullback} にもかかわらず、
次の補題は $\mathcal{F}$ が有限型準連接加群である場合には成り立たない。
すなわち層別化はなお存在するが、それが一般に関手 $F_{flat}$ を表すことは正しくない。

\begin{lemma}
\label{divisors-lemma-finite-presentation-module}
$S$ をスキームとする。$\mathcal{F}$ を有限表示
$\mathcal{O}_S$-加群とする。Lemma
\ref{divisors-lemma-locally-free-rank-r-pullback} と同様に
$S = Z_{-1} \supset Z_0 \supset Z_1 \supset \ldots$ とする。
$S_r = Z_{r - 1} \setminus Z_r$ と置く。このとき
$S' = \coprod_{r \geq 0} S_r$ は関手
$$
F_{flat} : \Sch/S \longrightarrow \textit{Sets},\quad\quad
T \longmapsto
\left\{
\begin{matrix}
\{*\} & \text{if }\mathcal{F}_T\text{ flat over }T\\
\emptyset & \text{otherwise}
\end{matrix}
\right.
$$
を表現する。さらに、$\mathcal{F}|_{S_r}$ は階数 $r$ の局所自由加群であり、
射 $S_r \to S$ と $S' \to S$ は有限表示である。
\end{lemma}

\begin{proof}
$g : T \to S$ を、引き戻し $\mathcal{F}_T = g^*\mathcal{F}$ が平坦となる
スキームの射と仮定する。このとき $\mathcal{F}_T$ は有限表示の平坦な
$\mathcal{O}_T$-加群である。従って $\mathcal{F}_T$ は有限局所自由である。
Properties, Lemma \ref{properties-lemma-finite-locally-free} を参照せよ。
ゆえに $T = \coprod_{r \geq 0} T_r$ と書け、$\mathcal{F}_T|_{T_r}$ は局所自由かつ
階数 $r$ である。このことから
$$
F_{flat} = \coprod\nolimits_{r \geq 0} F_r
$$
を得る。これは $\Sch/S$ 上の Zariski 層の圏における等式であり、
$F_r$ は Lemma \ref{divisors-lemma-locally-free-rank-r-pullback} のものである。
従って $F_{flat}$ は
$\coprod_{r \geq 0} (Z_{r - 1} \setminus Z_r)$ により表現される。
ここで $Z_r$ は Lemma \ref{divisors-lemma-locally-free-rank-r-pullback} のものである。
他の主張もこの補題から従う。
\end{proof}

\begin{example}
\label{divisors-example-not-fp-MB}
$R = \prod_{n \in \mathbf{N}} \mathbf{F}_2$ とし、$I \subset R$ を、
ほとんどすべての成分が零である元 $a = (a_n)_{n \in \mathbf{N}}$ のイデアルとする。
$\mathfrak m$ を $I$ を含む極大イデアルとする。このとき
$M = R/\mathfrak m$ は有限平坦 $R$-加群である。実際、$R$ は絶対平坦である
（More on Algebra, Lemma
\ref{more-algebra-lemma-product-fields-absolutely-flat}）。
$S = \Spec(R)$, $\mathcal{F} = \widetilde{M}$ と置く。
Lemma \ref{divisors-lemma-locally-free-rank-r-pullback} の閉部分スキームは
$S = Z_{-1}$, $Z_0 = \Spec(R/\mathfrak m)$, $Z_i = \emptyset$（$i > 0$）である。
しかし $\text{id} : S \to S$ は
$(S \setminus Z_0) \amalg Z_0$ を経由しない。これは $\mathfrak m$ が
$S$ の孤立していない点だからである。従って Lemma
\ref{divisors-lemma-finite-presentation-module} は有限型加群については成り立たない。
\end{example}






\section{射の特異点集合}
\label{divisors-section-singular-locus-morphism}

\noindent
$f : X \to S$ をスキームの有限型射とする。点の集合 $U$、すなわち
$f$ が滑らかである点全体は
$X$ の開集合である（Morphisms, Definition
\ref{morphisms-definition-smooth} による）。
多くの状況では、標準的な閉部分スキーム $\text{Sing}(f) \subset X$ であって、
その補集合が $U$ であり、その形成が任意の基底変換と可換するものがあると便利である。

\medskip\noindent
$f$ が有限表示ならば、一つの選択肢は、Smoothing Ring Maps, Definition
\ref{smoothing-definition-strictly-standard} の意味でアフィン局所的に
「strictly standard」である関数によって切り出される閉部分スキーム $Z$ を考えることである。
Smoothing Ring Maps, Lemma
\ref{smoothing-lemma-strictly-standard-base-change} から次が従う。
$f' : X' \to S'$ が $f$ の、射 $S' \to S$ による基底変換ならば
$Z' \subset S' \times_S Z$ である。ここで $Z'$ は、アフィン局所的に
strictly standard である関数によって $X'$ から切り出される閉部分スキームである。
しかし等号は明らかでない。strictly standard な元という概念は
Popescu の定理の章で有用であった。これらの元により定義される閉部分スキームは、
我々の知る限り文献では用いられていない\footnote{$f$ が局所完全交叉射
（More on Morphisms, Definition \ref{more-morphisms-definition-lci}）ならば、
局所的に strictly standard な元で切り出される閉部分スキームが
考えるべき正しい対象である。}。

\medskip\noindent
$f$ が平坦かつ有限表示で、その $f$ のファイバーがすべて次元 $d$ の等次元ならば、
第 $d$ Fitting イデアル（$\Omega_{X/S}$ のもの）を用いると、よい閉部分スキームが得られる。
任意の有限型射について、$X$ の閉部分スキームであって
$\Omega_{X/S}$ の Fitting イデアルで定義されるものは、$X$ の層別化を
$\Omega_{X/S}$ の階数によって定め、
その形成は基底変換と可換する。これは役立つことがある。
これはファイバーの埋め込み次元と関係している。
Varieties, Section \ref{varieties-section-embedding-dimension} を参照せよ。

\begin{lemma}
\label{divisors-lemma-base-change-and-fitting-ideal-omega}
$f : X \to S$ を局所有限型のスキームの射とする。
$X = Z_{-1} \supset Z_0 \supset Z_1 \supset \ldots$ を
$\Omega_{X/S}$ の Fitting イデアルで定義される閉部分スキームとする。
このとき $Z_i$ の形成は任意の基底変換と可換する。
\end{lemma}

\begin{proof}
$\Omega_{X/S}$ は有限型準連接 $\mathcal{O}_X$-加群である
（Morphisms, Lemma \ref{morphisms-lemma-finite-type-differentials}）。
従って Fitting イデアルは定義される。
$f' : X' \to S'$ が $f$ の、$g : S' \to S$ による基底変換ならば
$\Omega_{X'/S'} = (g')^*\Omega_{X/S}$ である。ここで $g' : X' \to X$ は射影である
（Morphisms, Lemma \ref{morphisms-lemma-base-change-differentials}）。
従って
$(g')^{-1}\text{Fit}_i(\Omega_{X/S}) \cdot \mathcal{O}_{X'} =
\text{Fit}_i(\Omega_{X'/S'})$
である。これは、スキーム論的に
$$
Z'_i = (g')^{-1}(Z_i) = Z_i \times_X X'
$$
であることを意味し、これが補題の主張の意味である。
\end{proof}

\noindent
第 $0$ Fitting イデアル（$\Omega$ のもの）は、射の「分岐点集合」を切り出す。

\begin{lemma}
\label{divisors-lemma-zero-fitting-ideal-omega-unramified}
$f : X \to S$ を局所有限型のスキームの射とする。
閉部分スキーム $Z \subset X$ であって、第 $0$ Fitting イデアル
（$\Omega_{X/S}$ のもの）により切り出されるものは、
ちょうど $f$ が非分岐でない点の集合である。
\end{lemma}

\begin{proof}
Lemma \ref{divisors-lemma-on-subscheme-cut-out-by-Fit-0} により、$Z$ の補集合は
$\Omega_{X/S}$ が零となる点の集合にちょうど等しい。これは Morphisms, Lemma
\ref{morphisms-lemma-unramified-omega-zero} により、
$f$ が非分岐である点の集合にちょうど等しい。
\end{proof}

\begin{lemma}
\label{divisors-lemma-d-fitting-ideal-omega-smooth}
$f : X \to S$ をスキームの射とし、$d \geq 0$ を整数とする。次を仮定する。
\begin{enumerate}
\item $f$ は平坦である。
\item $f$ は局所有限表示である。
\item $f$ の空でない各ファイバーは次元 $d$ の等次元である。
\end{enumerate}
閉部分スキーム $Z \subset X$ を、第 $d$ Fitting イデアル
（$\Omega_{X/S}$ のもの）により切り出されるものとする。
このとき $Z$ は、ちょうど $f$ が滑らかでない点の集合である。
\end{lemma}

\begin{proof}
Lemma \ref{divisors-lemma-locally-free-rank-r-pullback} により、$Z$ の補集合は
$\Omega_{X/S}$ が高々 $d$ 個の元で生成される点の集合にちょうど等しい。
従って補題は Morphisms, Lemma \ref{morphisms-lemma-smooth-at-point} から従う。
\end{proof}






\section{捩れなし加群}
\label{divisors-section-torsion-free}

\noindent
この節は準連接加群に対する More on Algebra, Section
\ref{more-algebra-section-torsion-flat} の類似である。

\begin{lemma}
\label{divisors-lemma-torsion-sections}
$X$ を生成点 $\eta$ をもつ整スキームとする。$\mathcal{F}$ を準連接
$\mathcal{O}_X$-加群とする。$U \subset X$ を空でない開集合とし、
$s \in \mathcal{F}(U)$ とする。次は同値である。
\begin{enumerate}
\item ある $x \in U$ に対し、$s$ の $\mathcal{F}_x$ における像は捩れ元である。
\item すべての $x \in U$ に対し、$s$ の $\mathcal{F}_x$ における像は捩れ元である。
\item $s$ の $\mathcal{F}_\eta$ における像は零である。
\item $s$ の $j_*\mathcal{F}_\eta$ における像は零である。ここで
$j : \eta \to X$ は包含射である。
\end{enumerate}
\end{lemma}

\begin{proof}
省略する。
\end{proof}

\begin{definition}
\label{divisors-definition-torsion}
$X$ を整スキームとし、$\mathcal{F}$ を準連接
$\mathcal{O}_X$-加群とする。
\begin{enumerate}
\item $\mathcal{F}$ の局所切断が Lemma
\ref{divisors-lemma-torsion-sections} の同値な条件を満たすとき、それを {\it 捩れ切断}という。
\item $\mathcal{F}$ を {\it 捩れなし}というのは、$\mathcal{F}$ の
すべての捩れ切断が $0$ であるときである。
\end{enumerate}
\end{definition}

\noindent
次は、これを通常の代数的概念と比較するための当然の補題である。

\begin{lemma}
\label{divisors-lemma-check-torsion-on-affines}
$X$ を整スキームとし、$\mathcal{F}$ を準連接
$\mathcal{O}_X$-加群とする。次は同値である。
\begin{enumerate}
\item $\mathcal{F}$ は捩れなしである。
\item アフィン開集合 $U \subset X$ に対し、$\mathcal{F}(U)$ は
捩れなし $\mathcal{O}(U)$-加群である。
\end{enumerate}
\end{lemma}

\begin{proof}
省略する。
\end{proof}

\begin{lemma}
\label{divisors-lemma-torsion}
$X$ を整スキームとし、$\mathcal{F}$ を準連接
$\mathcal{O}_X$-加群とする。$\mathcal{F}$ の捩れ切断は準連接
$\mathcal{O}_X$-部分加群
$\mathcal{F}_{tors} \subset \mathcal{F}$
をなす。商加群 $\mathcal{F}/\mathcal{F}_{tors}$ は捩れなしである。
\end{lemma}

\begin{proof}
省略する。代数的類似については More on Algebra, Lemma
\ref{more-algebra-lemma-torsion} を参照せよ。
\end{proof}

\begin{lemma}
\label{divisors-lemma-flat-torsion-free}
$X$ を整スキームとする。任意の平坦な準連接 $\mathcal{O}_X$-加群は捩れなしである。
\end{lemma}

\begin{proof}
省略する。More on Algebra, Lemma
\ref{more-algebra-lemma-flat-torsion-free} を参照せよ。
\end{proof}

\begin{lemma}
\label{divisors-lemma-flat-pullback-torsion}
$f : X \to Y$ を整スキームの平坦射とする。
$\mathcal{G}$ を捩れなし準連接 $\mathcal{O}_Y$-加群とする。このとき
$f^*\mathcal{G}$ は捩れなし $\mathcal{O}_X$-加群である。
\end{lemma}

\begin{proof}
省略する。代数的類似については More on Algebra, Lemma
\ref{more-algebra-lemma-flat-pullback-torsion} を参照せよ。
\end{proof}

\begin{lemma}
\label{divisors-lemma-flat-over-integral-integral-fibre}
$f : X \to Y$ をスキームの平坦射とする。$Y$ が整であり、
$f$ の生成ファイバーが整ならば、$X$ は整である。
\end{lemma}

\begin{proof}
代数的類似は次である。$A$ を分数体 $K$ をもつ整域とし、
$B$ を平坦 $A$-代数で、$B \otimes_A K$ が整域となるものとする。
このとき $B$ は整域である。実際、More on Algebra, Lemma
\ref{more-algebra-lemma-flat-torsion-free} により $B$ は捩れなしであり、
従って $B \subset B \otimes_A K$ である。
\end{proof}

\begin{lemma}
\label{divisors-lemma-check-torsion}
$X$ を整スキームとし、$\mathcal{F}$ を準連接
$\mathcal{O}_X$-加群とする。このとき $\mathcal{F}$ が捩れなしであるための必要十分条件は、
$\mathcal{F}_x$ が捩れなし $\mathcal{O}_{X, x}$-加群であることが、
すべての $x \in X$ に対して成り立つことである。
\end{lemma}

\begin{proof}
省略する。More on Algebra, Lemma
\ref{more-algebra-lemma-check-torsion} を参照せよ。
\end{proof}

\begin{lemma}
\label{divisors-lemma-extension-torsion-free}
$X$ を整スキームとする。短完全列
$0 \to \mathcal{F} \to \mathcal{F}' \to \mathcal{F}'' \to 0$
（準連接 $\mathcal{O}_X$-加群のもの）
を取る。$\mathcal{F}$ と $\mathcal{F}''$ が捩れなしならば、
$\mathcal{F}'$ も捩れなしである。
\end{lemma}

\begin{proof}
省略する。代数的類似については More on Algebra, Lemma
\ref{more-algebra-lemma-extension-torsion-free} を参照せよ。
\end{proof}

\begin{lemma}
\label{divisors-lemma-torsion-free-finite-noetherian-domain}
$X$ を生成点 $\eta$ をもつ局所 Noether 整スキームとする。
$\mathcal{F}$ を零でない連接 $\mathcal{O}_X$-加群とする。次は同値である。
\begin{enumerate}
\item $\mathcal{F}$ は捩れなしである。
\item $\eta$ は $\mathcal{F}$ の唯一の随伴素イデアルである。
\item $\eta$ は $\mathcal{F}$ の台に含まれ、$\mathcal{F}$ は性質 $(S_1)$ をもつ。
\item $\eta$ は $\mathcal{F}$ の台に含まれ、$\mathcal{F}$ は
埋め込まれた随伴素イデアルをもたない。
\end{enumerate}
\end{lemma}

\begin{proof}
これは More on Algebra, Lemma
\ref{more-algebra-lemma-torsion-free-finite-noetherian-domain} を
スキームの言葉に移したものである。その翻訳は省略する。
\end{proof}

\begin{lemma}
\label{divisors-lemma-torsion-free-over-regular-dim-1}
$X$ を次元 $\leq 1$ の整正則スキームとする。
$\mathcal{F}$ を連接 $\mathcal{O}_X$-加群とする。次は同値である。
\begin{enumerate}
\item $\mathcal{F}$ は捩れなしである。
\item $\mathcal{F}$ は有限局所自由である。
\end{enumerate}
\end{lemma}

\begin{proof}
有限局所自由加群が捩れなしであることは明らかである。
逆向きには、$\mathcal{F}$ が捩れなしならば、$\mathcal{F}_x$ が自由
$\mathcal{O}_{X, x}$-加群であることを、すべての $x \in X$ に対して示す。Algebra, Lemma
\ref{algebra-lemma-finite-projective} と $\mathcal{F}$ が連接であることにより、
これで十分である。$\dim(\mathcal{O}_{X, x}) = 0$ ならば
$\mathcal{O}_{X, x}$ は体なので、主張は明らかである。
$\dim(\mathcal{O}_{X, x}) = 1$ ならば $\mathcal{O}_{X, x}$ は離散付値環である
（Algebra, Lemma \ref{algebra-lemma-characterize-dvr}）。
また $\mathcal{F}_x$ は捩れなしである。従って More on Algebra, Lemma
\ref{more-algebra-lemma-dedekind-torsion-free-flat} により
$\mathcal{F}_x$ は自由である。
\end{proof}

\begin{lemma}
\label{divisors-lemma-hom-into-torsion-free}
$X$ を整スキームとし、$\mathcal{F}$, $\mathcal{G}$ を準連接
$\mathcal{O}_X$-加群とする。$\mathcal{G}$ が捩れなしで
$\mathcal{F}$ が有限表示ならば、
$\SheafHom_{\mathcal{O}_X}(\mathcal{F}, \mathcal{G})$
は捩れなしである。
\end{lemma}

\begin{proof}
$\SheafHom_{\mathcal{O}_X}(\mathcal{F}, \mathcal{G})$ は
Schemes, Section \ref{schemes-section-quasi-coherent} により準連接なので、
この主張には意味がある。主張が正しいことについては More on Algebra, Lemma
\ref{more-algebra-lemma-hom-into-torsion-free} を参照せよ。
いくつかの詳細は省略する。
\end{proof}

\begin{lemma}
\label{divisors-lemma-isom-depth-2-torsion-free}
$X$ を整局所 Noether スキームとする。$\varphi : \mathcal{F} \to \mathcal{G}$
を準連接 $\mathcal{O}_X$-加群の写像とする。
$\mathcal{F}$ は連接、$\mathcal{G}$ は捩れなしであり、
各 $x \in X$ に対して次のいずれかが起こると仮定する。
\begin{enumerate}
\item $\mathcal{F}_x \to \mathcal{G}_x$ は同型である。
\item $\text{depth}(\mathcal{F}_x) \geq 2$ である。
\end{enumerate}
このとき $\varphi$ は同型である。
\end{lemma}

\begin{proof}
これは More on Algebra, Lemma
\ref{more-algebra-lemma-isom-depth-2-torsion-free} を
スキームの言葉に移したものである。
\end{proof}







\section{加群の階数}
\label{divisors-section-rank}

\noindent
この節は準連接加群に対する More on Algebra, Section
\ref{more-algebra-section-rank} の類似である。
$X$ が生成点 $\eta$ をもつ整スキームならば、局所環
$\mathcal{O}_{X, \eta}$ は剰余体 $\kappa(\eta)$ に等しいことを思い出そう。
Morphisms, Lemma
\ref{morphisms-lemma-integral-scheme-rational-functions} を参照せよ。

\begin{definition}
\label{divisors-definition-rank}
$X$ を生成点 $\eta$ をもつ整スキームとする。
$\mathcal{F}$ を準連接 $\mathcal{O}_X$-加群とする。
$\mathcal{F}$ の {\it 階数}とは $\dim_{\kappa(\eta)} \mathcal{F}_\eta$ である。
\end{definition}

\noindent
これは、両方の定義を適用できる場合、局所自由加群の階数の定義と矛盾しない。
スキーム $X$ が既約でない場合や、生成点における局所環
$\mathcal{O}_{X, \eta}$ が体でない場合に、準連接加群の階数を定義するのは
よい考えではないように思われる\footnote{$X$ が既約で、
$\mathcal{F}_\eta$ が $\mathcal{O}_{X, \eta}$ 上自由ならば、
この自由加群の階数を用いるのが自然に見える。文献によっては
$\mathcal{F}_\eta$ の長さを用いて階数を定義するので注意せよ。}。

\begin{lemma}
\label{divisors-lemma-check-rank-on-affines}
$X$ を整スキームとし、$\mathcal{F}$ を準連接
$\mathcal{O}_X$-加群とする。$r$ を基数とする。次は同値である。
\begin{enumerate}
\item $\mathcal{F}$ の階数は $r$ である。
\item すべての空でないアフィン開集合 $U \subset X$ に対し、
$\mathcal{F}(U)$ は $\mathcal{O}_X(U)$-加群で、その階数は $r$ である。
\item ある空でないアフィン開集合 $U \subset X$ に対し、
$\mathcal{F}(U)$ は $\mathcal{O}_X(U)$-加群で、その階数は $r$ である。
\end{enumerate}
\end{lemma}

\begin{proof}
省略する。
\end{proof}

\begin{lemma}
\label{divisors-lemma-dominant-pullback-rank}
$f : X \to Y$ を整スキームの支配的射とする。
$\mathcal{G}$ を準連接 $\mathcal{O}_Y$-加群とする。このとき
$\mathcal{G}$ の $Y$ 上での階数は、$f^*\mathcal{G}$ の $X$ 上での階数に等しい。
\end{lemma}

\begin{proof}
Lemma \ref{divisors-lemma-check-rank-on-affines},
Morphisms, Lemma \ref{morphisms-lemma-dominant-between-integral},
および More on Algebra, Lemma
\ref{more-algebra-lemma-dominant-pullback-rank} を用いる。
\end{proof}

\begin{lemma}
\label{divisors-lemma-rank-additive}
$X$ を整スキームとする。短完全列
$0 \to \mathcal{F} \to \mathcal{F}' \to \mathcal{F}'' \to 0$
（準連接 $\mathcal{O}_X$-加群のもの）を取る。このとき
$\mathcal{F}'$ の階数は $\mathcal{F}$ と $\mathcal{F}''$ の階数の和である。
\end{lemma}

\begin{proof}
省略する。More on Algebra, Lemma
\ref{more-algebra-lemma-rank-additive} を参照せよ。
\end{proof}

\begin{lemma}
\label{divisors-lemma-rank-tensor}
$X$ を整スキームとする。$\mathcal{F}$ と $\mathcal{G}$ を準連接
$\mathcal{O}_X$-加群とする。
\begin{enumerate}
\item $\mathcal{F} \otimes_{\mathcal{O}_X} \mathcal{G}$ の階数は、
$\mathcal{F}$ と $\mathcal{G}$ の階数の積である。
\item $\mathcal{F}$ が有限表示ならば、
$\SheafHom_{\mathcal{O}_X}(\mathcal{F}, \mathcal{G})$ の階数は、
$\mathcal{F}$ と $\mathcal{G}$ の階数の積である。
\end{enumerate}
\end{lemma}

\begin{proof}
省略する。(2) に意味があるのは
$\SheafHom_{\mathcal{O}_X}(\mathcal{F}, \mathcal{G})$ が
Schemes, Section \ref{schemes-section-quasi-coherent} の項目 (10) により
準連接だからである。この補題の代数版は More on Algebra, Lemma
\ref{more-algebra-lemma-rank-tensor} である。
\end{proof}

\begin{lemma}
\label{divisors-lemma-finite-module-rank}
$X$ を整スキームとし、$\mathcal{F}$ を有限型準連接
$\mathcal{O}_X$-加群とする。このとき
\begin{enumerate}
\item $\mathcal{F}$ は有限な階数 $r \geq 0$ をもつ。
\item 空でない開集合 $U \subset X$ であって、
$\mathcal{F}|_U$ が階数 $r$ の自由加群となるものが存在する。
\end{enumerate}
\end{lemma}

\begin{proof}
代数版については More on Algebra, Lemma
\ref{more-algebra-lemma-finite-module-rank} を参照せよ。
\end{proof}










\section{反射的加群}
\label{divisors-section-reflexive}

\noindent
本節では、局所 Noether スキーム上の連接加群について、
More on Algebra, Section \ref{more-algebra-section-reflexive}
の類似を述べる。連接加群を扱う理由は、
$\SheafHom_{\mathcal{O}_X}(\mathcal{F}, \mathcal{G})$ が任意の連接
$\mathcal{O}_X$-加群の組 $\mathcal{F}, \mathcal{G}$ に対して連接だからである。
Modules, Lemma \ref{modules-lemma-internal-hom-locally-kernel-direct-sum}
を参照せよ。

\begin{definition}
\label{divisors-definition-reflexive}
$X$ を整局所 Noether スキームとし、$\mathcal{F}$ を連接
$\mathcal{O}_X$-加群とする。$\mathcal{F}$ の {\it 反射包} とは
$\mathcal{O}_X$-加群
$$
\mathcal{F}^{**} = \SheafHom_{\mathcal{O}_X}(
\SheafHom_{\mathcal{O}_X}(\mathcal{F}, \mathcal{O}_X), \mathcal{O}_X)
$$
のことである。$\mathcal{F}$ が {\it 反射的} であるとは、自然な写像
$j : \mathcal{F} \longrightarrow \mathcal{F}^{**}$ が同型であることをいう。
\end{definition}

\noindent
Lemma \ref{divisors-lemma-dual-reflexive} より、反射包は反射的
$\mathcal{O}_X$-加群である。より一般の状況でも同じ定義によって
反射的加群を定義できるが、これはあまり有用ではないように思われる。
以下は、この定義と通常の代数的概念とを比較する標準的な補題である。

\begin{lemma}
\label{divisors-lemma-check-reflexive-on-affines}
$X$ を整局所 Noether スキームとし、$\mathcal{F}$ を連接
$\mathcal{O}_X$-加群とする。次は同値である。
\begin{enumerate}
\item $\mathcal{F}$ は反射的である。
\item アフィン開集合 $U \subset X$ に対して $\mathcal{F}(U)$ は
反射的 $\mathcal{O}(U)$-加群である。
\end{enumerate}
\end{lemma}

\begin{proof}
省略する。
\end{proof}

\begin{remark}
\label{divisors-remark-different-reflexive}
$X$ が体上有限型のスキームであるとき、反射的加群について異なる概念が
用いられることがある（例えば
\cite[bottom of page 5 and Definition 1.1.9]{HL} を参照せよ）。
この別の概念では、$R\SheafHom$ を双対化複体
$\omega_X^\bullet$ に取り、$\mathcal{O}_X$ への Hom の代わりに用いる。$X$ が Gorenstein
でない場合には両者が異なりうるので、これは別の名称をもつべきであろう。
例えば $X = \Spec(k[t^3, t^4, t^5])$ とすると、計算により双対化層
$\omega_X$ は本節の意味では反射的でないことが分かる。しかし、
$\omega_X \to \SheafHom(\SheafHom(\omega_X, \omega_X), \omega_X)$
が同型なので、もう一つの意味では反射的である。
\end{remark}

\begin{lemma}
\label{divisors-lemma-reflexive-torsion-free}
$X$ を整局所 Noether スキームとし、$\mathcal{F}$ を連接
$\mathcal{O}_X$-加群とする。
\begin{enumerate}
\item $\mathcal{F}$ が反射的ならば、$\mathcal{F}$ は捩れなしである。
\item 写像 $j : \mathcal{F} \longrightarrow \mathcal{F}^{**}$ が単射であることと、
$\mathcal{F}$ が捩れなしであることは同値である。
\end{enumerate}
\end{lemma}

\begin{proof}
省略する。More on Algebra, Lemma
\ref{more-algebra-lemma-reflexive-torsion-free} を参照せよ。
\end{proof}

\begin{lemma}
\label{divisors-lemma-check-reflexive}
$X$ を整局所 Noether スキームとし、$\mathcal{F}$ を連接
$\mathcal{O}_X$-加群とする。次は同値である。
\begin{enumerate}
\item $\mathcal{F}$ は反射的である。
\item $\mathcal{F}_x$ は反射的 $\mathcal{O}_{X, x}$-加群である
（すべての $x \in X$ について）。
\item $\mathcal{F}_x$ は反射的 $\mathcal{O}_{X, x}$-加群である
（すべての閉点 $x \in X$ について）。
\end{enumerate}
\end{lemma}

\begin{proof}
Modules, Lemma \ref{modules-lemma-stalk-internal-hom} により、
(1) と (2) は同値である。$X$ の各点は閉点へ特殊化するので
（Properties, Lemma \ref{properties-lemma-locally-Noetherian-closed-point}）、
(2) と (3) も同値である。
\end{proof}

\begin{lemma}
\label{divisors-lemma-flat-pullback-reflexive}
$f : X \to Y$ を整局所 Noether スキームの間の平坦射とし、
$\mathcal{G}$ を連接反射的 $\mathcal{O}_Y$-加群とする。このとき
$f^*\mathcal{G}$ は連接反射的 $\mathcal{O}_X$-加群である。
\end{lemma}

\begin{proof}
省略する。代数版については More on Algebra, Lemma
\ref{more-algebra-lemma-flat-pullback-reflexive} を参照せよ。
\end{proof}

\begin{lemma}
\label{divisors-lemma-sequence-reflexive}
$X$ を整局所 Noether スキームとする。
$0 \to \mathcal{F} \to \mathcal{F}' \to \mathcal{F}''$ を連接
$\mathcal{O}_X$-加群の完全列とする。
$\mathcal{F}'$ が反射的であり、$\mathcal{F}''$ が捩れなしならば、
$\mathcal{F}$ は反射的である。
\end{lemma}

\begin{proof}
省略する。More on Algebra, Lemma \ref{more-algebra-lemma-sequence-reflexive}
を参照せよ。
\end{proof}

\begin{lemma}
\label{divisors-lemma-dual-reflexive}
$X$ を整局所 Noether スキームとし、$\mathcal{F}$, $\mathcal{G}$ を
連接 $\mathcal{O}_X$-加群とする。$\mathcal{G}$ が反射的ならば、
$\SheafHom_{\mathcal{O}_X}(\mathcal{F}, \mathcal{G})$ は反射的である。
\end{lemma}

\begin{proof}
Cohomology of Schemes, Lemma
\ref{coherent-lemma-tensor-hom-coherent} により
$\SheafHom_{\mathcal{O}_X}(\mathcal{F}, \mathcal{G})$ は連接なので、
主張は意味をもつ。主張の証明については More on Algebra, Lemma
\ref{more-algebra-lemma-dual-reflexive} を参照せよ。細部は省略する。
\end{proof}

\begin{remark}
\label{divisors-remark-tensor}
$X$ を整局所 Noether スキームとする。Lemma
\ref{divisors-lemma-dual-reflexive} により、反射包 $\mathcal{F}^{**}$ は連接反射的
$\mathcal{O}_X$-加群である。連接反射的加群の圏を $\mathcal{C}$ とする。
ここで加群は $\mathcal{O}_X$-加群である。
反射包を取ることは、包含函手
$\mathcal{C} \to \textit{Coh}(\mathcal{O}_X)$ の左随伴を与える。
$\mathcal{C}$ は核と余核をもつ加法圏であることに注意せよ。実際、
$\varphi : \mathcal{F} \to \mathcal{G}$ が $\mathcal{C}$ における射ならば、
通常の核 $\Ker(\varphi)$ は反射的であり
（Lemma \ref{divisors-lemma-sequence-reflexive}）、通常の余核の反射包
$\Coker(\varphi)^{**}$ が $\mathcal{C}$ における余核である。
さらに $\mathcal{C}$ はテンソル積
$$
\mathcal{F} \otimes_\mathcal{C} \mathcal{G} =
(\mathcal{F} \otimes_{\mathcal{O}_X} \mathcal{G})^{**}
$$
を受け継ぎ、これは結合的かつ対称的である。また内部 Hom が存在する。
すなわち、$\mathcal{F}, \mathcal{G}, \mathcal{H}$ を $\mathcal{C}$ の
任意の三対象とすると、恒等式
$$
\Hom_\mathcal{C}(\mathcal{F} \otimes_\mathcal{C} \mathcal{G}, \mathcal{H}) =
\Hom_\mathcal{C}(\mathcal{F},
\SheafHom_{\mathcal{O}_X}(\mathcal{G}, \mathcal{H}))
$$
が成り立つ。Modules, Lemma \ref{modules-lemma-internal-hom} を参照せよ。
$\mathcal{C}$ の各対象 $\mathcal{F}$ は {\it 双対対象}
$\SheafHom_{\mathcal{O}_X}(\mathcal{F}, \mathcal{O}_X)$ をもつ。
$X$ にさらに条件を課さなければ、$\mathcal{F}, \mathcal{G}$ が階数 $1$ の
$\mathcal{C}$ の対象であるときにも
$$
\SheafHom_{\mathcal{O}_X}(\mathcal{F}, \mathcal{G}) \not \cong
\SheafHom_{\mathcal{O}_X}(\mathcal{F}, \mathcal{O}_X)
\otimes_\mathcal{C} \mathcal{G}
\quad\text{and}\quad
\mathcal{F} \otimes_\mathcal{C}
\SheafHom_{\mathcal{O}_X}(\mathcal{F}, \mathcal{O}_X)
\not \cong \mathcal{O}_X
$$
となることがある。例を作るには、$X = \Spec(R)$ とおき、$R$ を
More on Algebra, Example \ref{more-algebra-example-ring-not-S2} の環とする。
また $\mathcal{F}, \mathcal{G}$ を $M$ に対応する加群とすればよい。
計算は省略する。
\end{remark}

\begin{lemma}
\label{divisors-lemma-reflexive-depth-2}
$X$ を整局所 Noether スキームとし、$\mathcal{F}$ を連接
$\mathcal{O}_X$-加群とする。次は同値である。
\begin{enumerate}
\item $\mathcal{F}$ は反射的である。
\item 各 $x \in X$ に対して、次のいずれかが成り立つ。
\begin{enumerate}
\item $\mathcal{F}_x$ は反射的 $\mathcal{O}_{X, x}$-加群である。
\item $\text{depth}(\mathcal{F}_x) \geq 2$ である。
\end{enumerate}
\end{enumerate}
\end{lemma}

\begin{proof}
省略する。More on Algebra, Lemma \ref{more-algebra-lemma-reflexive-depth-2}
を参照せよ。
\end{proof}

\begin{lemma}
\label{divisors-lemma-reflexive-S2}
$X$ を整局所 Noether スキームとし、$\mathcal{F}$ を連接反射的
$\mathcal{O}_X$-加群とする。$x \in X$ とする。
\begin{enumerate}
\item $\text{depth}(\mathcal{O}_{X, x}) \geq 2$ ならば、
$\text{depth}(\mathcal{F}_x) \geq 2$ である。
\item $X$ が $(S_2)$ ならば、$\mathcal{F}$ は $(S_2)$ である。
\end{enumerate}
\end{lemma}

\begin{proof}
省略する。More on Algebra, Lemma \ref{more-algebra-lemma-reflexive-S2}
を参照せよ。
\end{proof}

\begin{lemma}
\label{divisors-lemma-reflexive-S2-extend}
$X$ を整局所 Noether スキームとし、$j : U \to X$ を補集合 $Z$ をもつ
開部分スキームとする。$\mathcal{O}_{X, z}$ の深さが $\geq 2$ であると、
すべての $z \in Z$ について仮定する。このとき
$j^*$ と $j_*$ は、連接反射的 $\mathcal{O}_X$-加群の圏と
連接反射的 $\mathcal{O}_U$-加群の圏との間の圏同値を定める。
\end{lemma}

\begin{proof}
$\mathcal{F}$ を連接反射的 $\mathcal{O}_X$-加群とする。$z \in Z$
に対して、Lemma \ref{divisors-lemma-reflexive-S2} により茎 $\mathcal{F}_z$
の深さは $\geq 2$ である。したがって Lemma
\ref{divisors-lemma-depth-2-hartog} により
$\mathcal{F} \to j_*j^*\mathcal{F}$ は同型である。逆に、
$\mathcal{G}$ を連接反射的 $\mathcal{O}_U$-加群とする。
$j_*\mathcal{G}$ が連接反射的 $\mathcal{O}_X$-加群であることを示せば十分である。
この証明では $X$ がアフィンであると仮定してよい。Properties, Lemma
\ref{properties-lemma-lift-finite-presentation} により、連接
$\mathcal{O}_X$-加群 $\mathcal{F}$ であって
$\mathcal{G} = j^*\mathcal{F}$ を満たすものが存在する。$\mathcal{F}$ をその反射包で置き換えることにより、
$\mathcal{F}$ は反射的であると仮定してよい
（上の議論、とくに Lemma \ref{divisors-lemma-dual-reflexive} を参照せよ）。
先に示したことから $j_*\mathcal{G} = j_*j^*\mathcal{F} = \mathcal{F}$
であり、これが求める結論である。
\end{proof}

\noindent
スキームが正規ならば、反射的であることは、捩れなしで $(S_2)$
であることと同値である。

\begin{lemma}
\label{divisors-lemma-reflexive-over-normal}
$X$ を整局所 Noether 正規スキームとし、$\mathcal{F}$ を連接
$\mathcal{O}_X$-加群とする。次は同値である。
\begin{enumerate}
\item $\mathcal{F}$ は反射的である。
\item $\mathcal{F}$ は捩れなしで、性質 $(S_2)$ をもつ。
\item 次を満たす開部分スキーム $j : U \to X$ が存在する。
\begin{enumerate}
\item $X \setminus U$ のすべての既約成分は余次元が $\geq 2$ である
（$X$ において）。
\item $j^*\mathcal{F}$ は有限局所自由である。
\item $\mathcal{F} = j_*j^*\mathcal{F}$ である。
\end{enumerate}
\end{enumerate}
\end{lemma}

\begin{proof}
Lemma \ref{divisors-lemma-check-reflexive-on-affines} を用いると、(1) と (2) の同値は
More on Algebra, Lemma \ref{more-algebra-lemma-reflexive-over-normal}
から従う。$U \subset X$ が (3) の条件を満たすとする。Properties, Lemma
\ref{properties-lemma-criterion-normal} により、
$\text{depth}(\mathcal{O}_{X, x}) \geq 2$ が $x \not \in U$ に対して成り立つ。有限局所自由加群は
反射的なので、Lemma \ref{divisors-lemma-reflexive-S2-extend} により
(3) から (1) が従う。

\medskip\noindent
(1) を仮定する。$U \subset X$ を、$j^*\mathcal{F} = \mathcal{F}|_U$
が有限局所自由となるような最大の開部分スキームとする。すると (3)(b) が成り立つ。
$x \in X$ を一点とする。$\mathcal{F}_x$ が自由 $\mathcal{O}_{X, x}$-加群ならば
$x \in U$ である。Modules, Lemma
\ref{modules-lemma-finite-presentation-stalk-free} を参照せよ。
$\dim(\mathcal{O}_{X, x}) \leq 1$ ならば、$\mathcal{O}_{X, x}$ は体または
離散付値環であり（Properties, Lemma
\ref{properties-lemma-criterion-normal}）、したがって $\mathcal{F}_x$ は自由である
（More on Algebra, Lemma
\ref{more-algebra-lemma-dedekind-torsion-free-flat}）。ゆえに
$x \not \in U \Rightarrow \dim(\mathcal{O}_{X, x}) \geq 2$ である。
Properties, Lemma \ref{properties-lemma-codimension-local-ring} により
(3)(a) が従う。また、先に用いた Properties, Lemma
\ref{properties-lemma-criterion-normal} により、
$\text{depth}(\mathcal{O}_{X, x}) \geq 2$ が $x \not \in U$ に対しても成り立つから、Lemma
\ref{divisors-lemma-reflexive-S2-extend} により (3)(c) が従う。
\end{proof}

\begin{lemma}
\label{divisors-lemma-describe-reflexive-hull}
$X$ を生成点 $\eta$ をもつ整局所 Noether 正規スキームとする。
$\mathcal{F}$, $\mathcal{G}$ を連接 $\mathcal{O}_X$-加群とし、
$T : \mathcal{G}_\eta \to \mathcal{F}_\eta$ を線形写像とする。このとき、
$T$ が写像 $\mathcal{G} \to \mathcal{F}^{**}$ に $\mathcal{O}_X$-加群の
写像として延長されるための必要十分条件は
次である。
\begin{itemize}
\item[(*)] すべての $x \in X$ で $\dim(\mathcal{O}_{X, x}) = 1$
を満たすものに対して
$$
T\left(\Im(\mathcal{G}_x \to \mathcal{G}_\eta)\right) \subset
\Im(\mathcal{F}_x \to \mathcal{F}_\eta).
$$
\end{itemize}
\end{lemma}

\begin{proof}
$\mathcal{F}^{**}$ は捩れなしであり、
$\mathcal{F}_\eta = \mathcal{F}^{**}_\eta$ であるから、延長が存在すれば一意である。
したがって補題を $X$ の開被覆の各要素上で証明すれば十分であり、
$X$ はアフィンであると仮定してよい。この場合、問題は次の代数の問題になる。
$R$ を分数体 $K$ をもつ Noether 正規整域とし、$M$, $N$ を有限
$R$-加群とし、$T : M \otimes_R K \to N \otimes_R K$ を $K$-線形写像とする。
$T$ が写像 $N \to M^{**}$ に延長されるのはいつか。More on Algebra, Lemma
\ref{more-algebra-lemma-describe-reflexive-hull} により、これは
$N_\mathfrak p$ が $(M/M_{tors})_\mathfrak p$ に写されることが、高さ $1$ の
各素イデアル $\mathfrak p$（$R$ の素イデアル）について成り立つことと同値である。これはまさに
補題の条件 $(*)$ である。
\end{proof}

\begin{lemma}
\label{divisors-lemma-reflexive-over-regular-dim-2}
$X$ を次元 $\leq 2$ の正則スキームとし、$\mathcal{F}$ を連接
$\mathcal{O}_X$-加群とする。次は同値である。
\begin{enumerate}
\item $\mathcal{F}$ は反射的である。
\item $\mathcal{F}$ は有限局所自由である。
\end{enumerate}
\end{lemma}

\begin{proof}
有限局所自由加群が反射的であることは明らかである。逆を示すため、
$\mathcal{F}$ が反射的ならば、$\mathcal{F}_x$ が自由
$\mathcal{O}_{X, x}$-加群であることを、すべての $x \in X$ に対して示す。
$\mathcal{F}$ が連接であることと Algebra, Lemma
\ref{algebra-lemma-finite-projective} により、これで十分である。
$\dim(\mathcal{O}_{X, x}) = 0$ ならば、$\mathcal{O}_{X, x}$ は体なので
主張は明らかである。$\dim(\mathcal{O}_{X, x}) = 1$ ならば、
$\mathcal{O}_{X, x}$ は離散付値環であり
（Algebra, Lemma \ref{algebra-lemma-characterize-dvr}）、
$\mathcal{F}_x$ は捩れなしである。したがって More on Algebra, Lemma
\ref{more-algebra-lemma-dedekind-torsion-free-flat} により $\mathcal{F}_x$ は自由である。
$\dim(\mathcal{O}_{X, x}) = 2$ ならば、$\mathcal{O}_{X, x}$ は次元 $2$ の
正則局所環である。More on Algebra, Lemma
\ref{more-algebra-lemma-reflexive-over-normal} により $\mathcal{F}_x$ の深さは
$\geq 2$ である。ゆえに Algebra, Lemma
\ref{algebra-lemma-regular-mcm-free} により $\mathcal{F}_x$ は自由である。
\end{proof}

\section{有効 Cartier 因子}
\label{divisors-section-effective-Cartier-divisors}

\noindent
ほかの種類の因子に先立って、有効 Cartier 因子の概念を定義する。

\begin{definition}
\label{divisors-definition-effective-Cartier-divisor}
$S$ をスキームとする。
\begin{enumerate}
\item $S$ の {\it 局所単項閉部分スキーム} とは、そのイデアル層が
局所的に一つの元で生成される閉部分スキームをいう。
\item $S$ 上の {\it 有効 Cartier 因子} とは、閉部分スキーム
$D \subset S$ であって、そのイデアル層
$\mathcal{I}_D \subset \mathcal{O}_S$ が可逆 $\mathcal{O}_S$-加群であるものをいう。
\end{enumerate}
\end{definition}

\noindent
したがって、有効 Cartier 因子は局所単項閉部分スキームであるが、逆は常には
成り立たない。有効 Cartier 因子は、最も強い意味で純余次元 $1$ の閉部分スキーム
である。すなわち、局所的には一つの非零因子によって切り出される。とくに疎である。

\begin{lemma}
\label{divisors-lemma-characterize-effective-Cartier-divisor}
$S$ をスキームとし、$D \subset S$ を閉部分スキームとする。次は同値である。
\begin{enumerate}
\item 部分スキーム $D$ は $S$ 上の有効 Cartier 因子である。
\item すべての $x \in D$ に対して、アフィン開近傍
$\Spec(A) = U \subset S$ が $x$ の近傍として存在し、
$U \cap D = \Spec(A/(f))$ であって $f \in A$ は非零因子である。
\end{enumerate}
\end{lemma}

\begin{proof}
(1) を仮定する。すべての $x \in D$ に対して、アフィン開近傍
$\Spec(A) = U \subset S$ が $x$ の近傍として存在し、
$\mathcal{I}_D|_U \cong \mathcal{O}_U$ となるものが存在する。言い換えると、
切断 $f \in \Gamma(U, \mathcal{I}_D)$ であって、制限
$\mathcal{I}_D|_U$ を自由に生成するものが存在する。したがって $f \in A$ であり、
乗法写像 $f : A \to A$ は単射である。また $\mathcal{I}_D$ は準連接だから、
$D \cap U = \Spec(A/(f))$ である。

\medskip\noindent
(2) を仮定し、$x \in D$ とする。仮定により、アフィン開近傍
$\Spec(A) = U \subset S$ が $x$ の近傍として存在し、
$U \cap D = \Spec(A/(f))$ であって $f \in A$ は非零因子である。このとき
$\mathcal{I}_D|_U \cong \mathcal{O}_U$ である。実際、これは
$\widetilde{(f)} \cong \widetilde{A} \cong \mathcal{O}_U$ に等しい。
もちろん、$\mathcal{I}_D$ の開部分スキーム $S \setminus D$ への制限は
$\mathcal{O}_{S \setminus D}$ と同型である。したがって $\mathcal{I}_D$ は
可逆 $\mathcal{O}_S$-加群である。
\end{proof}

\begin{lemma}
\label{divisors-lemma-complement-locally-principal-closed-subscheme}
$S$ をスキームとし、$Z \subset S$ を局所単項閉部分スキームとする。
$U = S \setminus Z$ とおく。このとき $U \to S$ はアフィン射である。
\end{lemma}

\begin{proof}
問題は $S$ 上局所的である。Morphisms, Lemmas
\ref{morphisms-lemma-characterize-affine} を参照せよ。したがって
$S = \Spec(A)$ および $Z = V(f)$ と、ある $f \in A$ に対して仮定してよい。
この場合 $U = D(f) = \Spec(A_f)$ はアフィンなので、$U \to S$ はアフィンである。
\end{proof}

\begin{lemma}
\label{divisors-lemma-complement-effective-Cartier-divisor}
$S$ をスキームとし、$D \subset S$ を有効 Cartier 因子とする。
$U = S \setminus D$ とおく。このとき $U \to S$ はアフィン射であり、$U$ は
$S$ においてスキーム論的に稠密である。
\end{lemma}

\begin{proof}
アフィン性は Lemma \ref{divisors-lemma-complement-locally-principal-closed-subscheme}
である。稠密性の問題は $S$ 上局所的である。Morphisms, Lemma
\ref{morphisms-lemma-characterize-scheme-theoretically-dense} を参照せよ。
したがって $S = \Spec(A)$ と仮定してよく、Lemma
\ref{divisors-lemma-characterize-effective-Cartier-divisor} により $D$ は非零因子
$f \in A$ に対応するとしてよい。すると $A \subset A_f$ であり、これから
$U \subset S$ がスキーム論的に稠密であることが従う。Morphisms, Example
\ref{morphisms-example-scheme-theoretic-closure} を参照せよ。
\end{proof}

\begin{lemma}
\label{divisors-lemma-effective-Cartier-makes-dimension-drop}
$S$ をスキームとし、$D \subset S$ を有効 Cartier 因子とする。
$s \in D$ とする。$\dim_s(S) < \infty$ ならば
$\dim_s(D) < \dim_s(S)$ である。
\end{lemma}

\begin{proof}
$\dim_s(S) < \infty$ を仮定する。アフィン開近傍
$U = \Spec(A) \subset S$ を $s$ の近傍として取り、$\dim(U) = \dim_s(S)$ かつ
$D = V(f)$ であって $f \in A$ は非零因子であるようにする
（Lemma \ref{divisors-lemma-characterize-effective-Cartier-divisor} を参照）。
$\dim(U)$ は環 $A$ の Krull 次元であり、$\dim(U \cap D)$ は環 $A/(f)$
の Krull 次元であることを思い出す。$f$ は $A$ のいかなる極小素イデアルにも
含まれない。したがって、$A/(f)$ の素イデアルの任意の極大鎖は、$A$ の
素イデアルの鎖とみなすと、極小素イデアルを加えることによって延長できる。
\end{proof}

\begin{definition}
\label{divisors-definition-sum-effective-Cartier-divisors}
$S$ をスキームとする。有効 Cartier 因子 $D_1$, $D_2$ が $S$ 上に与えられたとき、
$D = D_1 + D_2$ を、準連接イデアル層
$S$ の閉部分スキームであって
$\mathcal{I}_{D_1}\mathcal{I}_{D_2} \subset \mathcal{O}_S$ に対応するものとして定める。
これを {\it 有効 Cartier 因子 $D_1$ と $D_2$ の和}
と呼ぶ。
\end{definition}

\noindent
和 $\sum n_iD_i$ は、有限個の有効 Cartier 因子 $D_i$ が $X$ 上にあり、
$n_i$ が非負整数であるときにも定義できることは明らかである。

\begin{lemma}
\label{divisors-lemma-sum-effective-Cartier-divisors}
二つの有効 Cartier 因子の和は有効 Cartier 因子である。
\end{lemma}

\begin{proof}
省略する。局所的には $f_1, f_2 \in A$ が非零因子であり、そのとき
$f_1f_2 \in A$ も非零因子である。
\end{proof}

\begin{lemma}
\label{divisors-lemma-difference-effective-Cartier-divisors}
$X$ をスキームとし、$D, D'$ を $X$ 上の二つの有効 Cartier 因子とする。
$D \subset D'$ が $X$ の閉部分スキームとして成り立つならば、有効 Cartier 因子
$D''$ が存在して $D' = D + D''$ となる。
\end{lemma}

\begin{proof}
省略する。
\end{proof}

\begin{lemma}
\label{divisors-lemma-sum-closed-subschemes-effective-Cartier}
$X$ をスキームとし、$Z, Y$ を $X$ の二つの閉部分スキームであって、
イデアル層 $\mathcal{I}$ および $\mathcal{J}$ をもつものとする。
$\mathcal{I}\mathcal{J}$ が有効 Cartier 因子 $D \subset X$ を定めるならば、
$Z$ と $Y$ は有効 Cartier 因子であり、$D = Z + Y$ である。
\end{lemma}

\begin{proof}
Lemma \ref{divisors-lemma-characterize-effective-Cartier-divisor} を適用すると、次の代数的状況を得る。
$A$ は環、$I, J \subset A$ はイデアル、$f \in A$ は非零因子であり、
$IJ = (f)$ である。したがって結果は Algebra, Lemma
\ref{algebra-lemma-product-ideals-principal} から従う。
\end{proof}

\begin{lemma}
\label{divisors-lemma-sum-effective-Cartier-divisors-union}
$X$ をスキームとし、$D, D' \subset X$ を有効 Cartier 因子とする。
スキーム論的交叉 $D \cap D'$ が $D'$ 上の有効 Cartier 因子であると仮定する。
このとき $D + D'$ は $D$ と $D'$ のスキーム論的和である。
\end{lemma}

\begin{proof}
スキーム論的交叉と和の定義については Morphisms, Definition
\ref{morphisms-definition-scheme-theoretic-intersection-union} を参照せよ。
補題を証明するため局所的に議論すると
（Lemma \ref{divisors-lemma-characterize-effective-Cartier-divisor} を用いる）、
次の代数問題を得る。環 $A$ と非零因子 $f_1, f_2 \in A$ が与えられ、
$f_1$ の $A/f_2A$ における像も非零因子であるとき、
$f_1A \cap f_2A = f_1f_2A$ を示せ。直接的な議論は省略する。
\end{proof}

\noindent
Schemes, Definition \ref{schemes-definition-inverse-image-closed-subscheme}
において、任意のスキームの射による閉部分スキームの逆像を定義したことを思い出す。

\begin{lemma}
\label{divisors-lemma-pullback-locally-principal}
$f : S' \to S$ をスキームの射とし、$Z \subset S$ を局所単項閉部分スキームとする。
このとき逆像 $f^{-1}(Z)$ は $S'$ の局所単項閉部分スキームである。
\end{lemma}

\begin{proof}
省略する。
\end{proof}

\begin{definition}
\label{divisors-definition-pullback-effective-Cartier-divisor}
$f : S' \to S$ をスキームの射とし、$D \subset S$ を有効 Cartier 因子とする。
{\it $D$ の $f$ による引き戻しが定義される} とは、閉部分スキーム
$f^{-1}(D) \subset S'$ が有効 Cartier 因子であることをいう。この場合、それを
$f^*D$ または $f^{-1}(D)$ と表し、{\it 有効 Cartier 因子の引き戻し} と呼ぶ。
\end{definition}

\noindent
$f^{-1}(D)$ が有効 Cartier 因子であるという条件は、実際にはしばしば満たされる。
一例となる補題を挙げる。

\begin{lemma}
\label{divisors-lemma-pullback-effective-Cartier-defined}
$f : X \to Y$ をスキームの射とし、$D \subset Y$ を有効 Cartier 因子とする。
$D$ の $f$ による引き戻しは、次の各場合に定義される。
\begin{enumerate}
\item $f(x) \not \in D$ が任意の弱随伴点 $x$（$X$ の点）に対して成り立つ。
\item $X$, $Y$ は整スキームであり、$f$ は支配的である。
\item $X$ は被約であり、$f(\xi) \not \in D$ が任意の生成点 $\xi$、すなわち
$X$ の任意の既約成分の生成点に対して成り立つ。
\item $X$ は局所 Noether であり、$f(x) \not \in D$ が任意の随伴点 $x$、
すなわち $X$ の随伴点に対して成り立つ。
\item $X$ は局所 Noether で埋め込まれた点をもたず、$f(\xi) \not \in D$ が
任意の生成点 $\xi$、すなわち $X$ の既約成分の生成点に対して成り立つ。
\item $f$ は平坦である。
\item 必要に応じてここにさらに追加する。
\end{enumerate}
\end{lemma}

\begin{proof}
問題は $X$ 上局所的である。したがって、$X = \Spec(A)$,
$Y = \Spec(R)$ であり、$f$ は $\varphi : R \to A$ によって与えられ、
$D = \Spec(R/(t))$ であって $t \in R$ は非零因子である場合に帰着する。
各場合の目標は、$\varphi(t) \in A$ が非零因子であることを示すことである。

\medskip\noindent
場合 (1) は Algebra, Lemma \ref{algebra-lemma-weakly-ass-zero-divisors}
から従う。場合 (4) は Lemma \ref{divisors-lemma-ass-weakly-ass} により (1) の特殊な場合である。
場合 (5) は (4) と定義から従う。場合 (3) は Lemma
\ref{divisors-lemma-weakass-reduced} により (1) の特殊な場合である。場合 (2) は (3) の
特殊な場合である。$R \to A$ が平坦ならば、$t : R \to R$ が単射であることから
$t : A \to A$ も単射である。これで (6) が証明された。
\end{proof}

\begin{lemma}
\label{divisors-lemma-pullback-effective-Cartier-divisors-additive}
$f : S' \to S$ をスキームの射とし、$D_1$, $D_2$ を $S$ 上の有効 Cartier 因子とする。
$D_1$ と $D_2$ の引き戻しが定義されるならば、$D = D_1 + D_2$ の引き戻しも
定義され、$f^*D = f^*D_1 + f^*D_2$ である。
\end{lemma}

\begin{proof}
省略する。
\end{proof}

\section{有効 Cartier 因子と可逆層}
\label{divisors-section-effective-Cartier-invertible}

\noindent
有効 Cartier 因子は可逆なイデアル層をもつので
（Definition \ref{divisors-definition-effective-Cartier-divisor}）、次の定義には意味がある。

\begin{definition}
\label{divisors-definition-invertible-sheaf-effective-Cartier-divisor}
$S$ をスキームとし、$D \subset S$ をイデアル層 $\mathcal{I}_D$ をもつ
有効 Cartier 因子とする。
\begin{enumerate}
\item {\it 可逆層 $\mathcal{O}_S(D)$（$D$ に付随するもの）} を
$$
\mathcal{O}_S(D) =
\SheafHom_{\mathcal{O}_S}(\mathcal{I}_D, \mathcal{O}_S) =
\mathcal{I}_D^{\otimes -1}.
$$
により定義する。
\item {\it 標準切断}とは、通常 $1$ または $1_D$ と表される
$\mathcal{O}_S(D)$ の大域切断であって、包含写像
$\mathcal{I}_D \to \mathcal{O}_S$ に対応するものである。
\item
$\mathcal{O}_S(-D) = \mathcal{O}_S(D)^{\otimes -1} = \mathcal{I}_D$
と書く。
\item 第二の有効 Cartier 因子 $D' \subset S$ が与えられたとき、
$\mathcal{O}_S(D - D') =
\mathcal{O}_S(D) \otimes_{\mathcal{O}_S} \mathcal{O}_S(-D')$
と定義する。
\end{enumerate}
\end{definition}

\noindent
いくつか補足する。下で見るように、対応 $D \mapsto \mathcal{O}_S(D)$ は
有効 Cartier 因子の加法（Definition
\ref{divisors-definition-sum-effective-Cartier-divisors}）を $S$ の Picard 群における
加法へ移す（Lemma \ref{divisors-lemma-invertible-sheaf-sum-effective-Cartier-divisors}）。
しかし、上の定義における式 $D - D'$ 自体には幾何学的意味はない。
より正確には、$S$ 上の有効 Cartier 因子全体を、空の有効 Cartier 因子を零元とする
可換モノイド $\text{EffCart}(S)$ とみなせる。このとき対応
$(D, D') \mapsto \mathcal{O}_S(D - D')$ は群準同型
$$
\text{EffCart}(S)^{gp} \longrightarrow \Pic(S)
$$
を定める。ここで左辺は $\text{EffCart}(S)$ の群完備化である。言い換えると、
$\mathcal{O}_S(D - D')$ と書くとき、$D - D'$ を
$\text{EffCart}(S)^{gp}$ の元とみなしてよい。

\begin{lemma}
\label{divisors-lemma-conormal-effective-Cartier-divisor}
$S$ をスキームとし、$D \subset S$ を有効 Cartier 因子とする。このとき余法線層は
$\mathcal{C}_{D/S} = \mathcal{I}_D|_D =
\mathcal{O}_S(-D)|_D$ であり、法線層は
$\mathcal{N}_{D/S} = \mathcal{O}_S(D)|_D$ である。
\end{lemma}

\begin{proof}
これは Morphisms, Lemma \ref{morphisms-lemma-affine-conormal} から従う。
\end{proof}

\begin{lemma}
\label{divisors-lemma-ses-add-divisor}
$X$ をスキームとし、$D, D' \subset X$ を $D \subset D'$ を満たす有効 Cartier
因子とする。また $C \subset X$ を、
$D' = D + C$ を満たす有効 Cartier 因子とする（これは Lemma
\ref{divisors-lemma-difference-effective-Cartier-divisors} により存在する）。このとき
$\mathcal{O}_X$-加群の短完全列
$$
0 \to \mathcal{O}_X(-D)|_C \to \mathcal{O}_{D'} \to \mathcal{O}_D \to 0
$$
が存在する。
\end{lemma}

\begin{proof}
補題の主張と証明では Morphisms, Lemma
\ref{morphisms-lemma-i-star-equivalence} の圏同値を用いて、$X$ の閉部分スキーム上の
準連接加群を $X$ 上の準連接加群とみなす。$\mathcal{I}$ を $D$ の $D'$ における
イデアル層とし、$\mathcal{I}_D = \mathcal{O}_X(-D)$ を $D$ の $X$ における
イデアル層とする。このとき標準的な全射
$\mathcal{I}_D \to \mathcal{I}$ が存在する。したがって、この写像の核が
$\mathcal{I}_C\mathcal{I}_D$ に等しいことを示せば十分である
（記号の意味は明らかであろう）。これは局所的に確認できるので、
$X = \Spec(A)$, $D = \Spec(A/fA)$, $C = \Spec(A/gA)$ であり、
$f, g \in A$ は非零因子であると仮定してよい
（Lemma \ref{divisors-lemma-characterize-effective-Cartier-divisor}）。このとき
$\mathcal{I}_D$ は $fA$ に対応し、非零因子 $fg \in A$ は $D'$ のイデアルを生成し、
$\mathcal{I}$ は $I = fA/fgA$ に対応する。したがって結果は明らかである。
\end{proof}

\begin{lemma}
\label{divisors-lemma-invertible-sheaf-sum-effective-Cartier-divisors}
$S$ をスキームとし、$D_1$, $D_2$ を $S$ 上の有効 Cartier 因子とする。
$D = D_1 + D_2$ とおく。このとき一意な同型
$$
\mathcal{O}_S(D_1) \otimes_{\mathcal{O}_S} \mathcal{O}_S(D_2)
\longrightarrow
\mathcal{O}_S(D)
$$
であって、$1_{D_1} \otimes 1_{D_2}$ を $1_D$ に写すものが存在する。
\end{lemma}

\begin{proof}
省略する。
\end{proof}

\begin{lemma}
\label{divisors-lemma-pullback-effective-Cartier-divisors}
$f : S' \to S$ をスキームの射とし、$D$ を $S$ 上の有効 Cartier 因子とする。
$D$ の引き戻しが定義されるならば
$f^*\mathcal{O}_S(D) = \mathcal{O}_{S'}(f^*D)$
であり、標準切断 $1_D$ の引き戻しは標準切断 $1_{f^*D}$ である。
\end{lemma}

\begin{proof}
省略する。
\end{proof}

\begin{definition}
\label{divisors-definition-regular-section}
$(X, \mathcal{O}_X)$ を局所環付き空間とし、$\mathcal{L}$ を $X$ 上の可逆層とする。
大域切断 $s \in \Gamma(X, \mathcal{L})$ が {\it 正則切断} であるとは、写像
$\mathcal{O}_X \to \mathcal{L}$, $f \mapsto fs$ が単射であることをいう。
\end{definition}

\begin{lemma}
\label{divisors-lemma-regular-section-structure-sheaf}
$X$ を局所環付き空間とし、$f \in \Gamma(X, \mathcal{O}_X)$ とする。
次は同値である。
\begin{enumerate}
\item $f$ は正則切断である。
\item 任意の $x \in X$ に対して、像 $f \in \mathcal{O}_{X, x}$ は非零因子である。
\end{enumerate}
$X$ がスキームならば、これらは次とも同値である。
\begin{enumerate}
\item[(3)] 任意のアフィン開集合 $\Spec(A) = U \subset X$ に対して、像
$f \in A$ は非零因子である。
\item[(4)] アフィン開被覆 $X = \bigcup \Spec(A_i)$ であって、像 $f$ が
$A_i$ において非零因子となるものが、すべての $i$ に対して存在する。
\end{enumerate}
\end{lemma}

\begin{proof}
省略する。
\end{proof}

\noindent
$s$ を可逆 $\mathcal{O}_X$-加群 $\mathcal{L}$ の大域切断とすると、これは
$\mathcal{O}_X$-加群写像 $s : \mathcal{O}_X \to \mathcal{L}$ とみなせることに注意する。
したがって、その双対は写像
$s : \mathcal{L}^{\otimes -1} \to \mathcal{O}_X$ である。
（双対可逆層の定義については Modules, Definition
\ref{modules-definition-powers} を参照せよ。）

\begin{definition}
\label{divisors-definition-zero-scheme-s}
$X$ をスキームとし、$\mathcal{L}$ を可逆層とする。
$s \in \Gamma(X, \mathcal{L})$ を大域切断とする。$s$ の {\it 零点スキーム} とは、
閉部分スキーム $Z(s) \subset X$ であって、準連接イデアル層
$\mathcal{I} \subset \mathcal{O}_X$、すなわち写像
$s : \mathcal{L}^{\otimes -1} \to \mathcal{O}_X$ の像によって定義されるものをいう。
\end{definition}

\begin{lemma}
\label{divisors-lemma-zero-scheme}
$X$ をスキームとし、$\mathcal{L}$ を可逆層とし、
$s \in \Gamma(X, \mathcal{L})$ とする。
\begin{enumerate}
\item 閉埋め込み $i : Z \to X$ であって
$i^*s \in \Gamma(Z, i^*\mathcal{L})$ が零となるもの全体を包含関係で順序付ける。
零点スキーム $Z(s)$ はこの順序集合の最大元である。
\item 任意のスキームの射 $f : Y \to X$ に対して、
$f^*s = 0$ が $\Gamma(Y, f^*\mathcal{L})$ において成り立つことと、
$f$ が $Z(s)$ を経由することは同値である。
\item 零点スキーム $Z(s)$ は局所単項閉部分スキームである。
\item 零点スキーム $Z(s)$ が有効 Cartier 因子であることと、$s$ が
$\mathcal{L}$ の正則切断であることは同値である。
\end{enumerate}
\end{lemma}

\begin{proof}
省略する。
\end{proof}

\begin{lemma}
\label{divisors-lemma-characterize-OD}
\begin{slogan}
スキーム上の有効 Cartier 因子は、指定された正則大域切断をもつ可逆層と同じものである。
\end{slogan}
$X$ をスキームとする。
\begin{enumerate}
\item $D \subset X$ が有効 Cartier 因子ならば、標準切断 $1_D$、すなわち
$\mathcal{O}_X(D)$ の標準切断は正則である。
\item 逆に、$s$ が可逆層 $\mathcal{L}$ の正則切断ならば、一意な有効 Cartier 因子
$D = Z(s) \subset X$ と、一意な同型 $\mathcal{O}_X(D) \to \mathcal{L}$ が存在し、
この同型は $1_D$ を $s$ に写す。
\end{enumerate}
構成 $D \mapsto (\mathcal{O}_X(D), 1_D)$ と
$(\mathcal{L}, s) \mapsto Z(s)$ は互いに逆な写像
$$
\left\{
\begin{matrix}
\text{effective Cartier divisors on }X
\end{matrix}
\right\}
\leftrightarrow
\left\{
\begin{matrix}
\text{isomorphism classes of pairs }(\mathcal{L}, s)\\
\text{consisting of an invertible }
\mathcal{O}_X\text{-module}\\
\mathcal{L}\text{ and a regular global section }s
\end{matrix}
\right\}
$$
を与える。
\end{lemma}

\begin{proof}
省略する。
\end{proof}

\begin{remark}
\label{divisors-remark-ses-regular-section}
$X$ をスキーム、$\mathcal{L}$ を可逆 $\mathcal{O}_X$-加群、$s$ を
$\mathcal{L}$ の正則切断とする。このとき零点スキーム $D = Z(s)$ は $X$ 上の
有効 Cartier 因子であり、短完全列
$$
0 \to \mathcal{O}_X \to \mathcal{L} \to i_*(\mathcal{L}|_D) \to 0
\quad\text{and}\quad
0 \to \mathcal{L}^{\otimes -1} \to \mathcal{O}_X \to i_*\mathcal{O}_D \to 0.
$$
が存在する。有効 Cartier 因子 $D \subset X$ が与えられたとき、Lemmas
\ref{divisors-lemma-characterize-OD} および
\ref{divisors-lemma-conormal-effective-Cartier-divisor} を用いると
$$
0 \to \mathcal{O}_X \to \mathcal{O}_X(D) \to i_*(\mathcal{N}_{D/X}) \to 0
\quad\text{and}\quad
0 \to \mathcal{O}_X(-D) \to \mathcal{O}_X \to i_*(\mathcal{O}_D) \to 0
$$
を得る。
\end{remark}

\section{Noether スキーム上の有効 Cartier 因子}
\label{divisors-section-Noetherian-effective-Cartier}

\noindent
局所 Noether の状況では、有効 Cartier 因子と正則切断に関する議論の多くが
いくらか簡単になる。

\begin{lemma}
\label{divisors-lemma-regular-section-associated-points}
$X$ を局所 Noether スキームとし、$\mathcal{L}$ を可逆
$\mathcal{O}_X$-加群とする。$s \in \Gamma(X, \mathcal{L})$ とする。このとき
$s$ が正則切断であることと、$s$ が $X$ の随伴点で消えないことは同値である。
\end{lemma}

\begin{proof}
省略する。ヒント：アフィンの場合で $\mathcal{L}$ が自明な場合に帰着し、
Lemma \ref{divisors-lemma-regular-section-structure-sheaf} と Algebra, Lemma
\ref{algebra-lemma-ass-zero-divisors} を用いる。
\end{proof}

\begin{lemma}
\label{divisors-lemma-effective-Cartier-in-points}
$X$ を局所 Noether スキームとし、$D \subset X$ を準連接イデアル層
$\mathcal{I} \subset \mathcal{O}_X$ に対応する閉部分スキームとする。
\begin{enumerate}
\item すべての $x \in D$ に対してイデアル
$\mathcal{I}_x \subset \mathcal{O}_{X, x}$ が一つの元で生成できるならば、
$D$ は局所単項である。
\item すべての $x \in D$ に対してイデアル
$\mathcal{I}_x \subset \mathcal{O}_{X, x}$ が一つの非零因子で生成できるならば、
$D$ は有効 Cartier 因子である。
\end{enumerate}
\end{lemma}

\begin{proof}
$\Spec(A)$ を一点 $x \in D$ のアフィン近傍とし、$\mathfrak p \subset A$ を
$x$ に対応する素イデアルとする。$I \subset A$ を、$D$ の $\Spec(A)$ 上の
跡を定義するイデアルとする。$A$ は Noether 環なので（$X$ が局所 Noether
だからである）、イデアル $I$ は有限個の元で生成される。これを
$I = (f_1, \ldots, f_r)$ とする。(1) の仮定の下では
$I_\mathfrak p = (f)$ であり、ここで $f \in A_\mathfrak p$ である。したがって
$f_i = g_i f$ であり、ここで $g_i \in A_\mathfrak p$ である。さらに
$g_i = a_i/h_i$ および $f = f'/h$ と書くことができ、ここで
$a_i, h_i, f', h \in A$ かつ $h_i, h \not \in \mathfrak p$ である。このとき
$I_{h_1 \ldots h_r h} \subset A_{h_1 \ldots h_r h}$ は $f'$ で生成されるため
単項である。これで (1) が証明された。(2) については $I = (f)$ と仮定してよい。
仮定により、$f$ の $A_\mathfrak p$ における像は非零因子である。すると
Algebra, Lemma \ref{algebra-lemma-regular-sequence-in-neighbourhood} により、
$f$ は $x$ の近傍上で非零因子である。これで (2) が証明された。
\end{proof}

\begin{lemma}
\label{divisors-lemma-effective-Cartier-codimension-1}
$X$ を局所 Noether スキームとする。
\begin{enumerate}
\item $D \subset X$ を局所単項閉部分スキームとし、$\xi \in D$ を $D$ の
既約成分の生成点とする。このとき $\dim(\mathcal{O}_{X, \xi}) \leq 1$ である。
\item $D \subset X$ を有効 Cartier 因子とし、$\xi \in D$ を $D$ の
既約成分の生成点とする。このとき $\dim(\mathcal{O}_{X, \xi}) = 1$ である。
\end{enumerate}
\end{lemma}

\begin{proof}
(1) を証明する。仮定により、$X = \Spec(A)$ および
$D = \Spec(A/(f))$ と仮定してよい。ここで $A$ は Noether 環、$f \in A$ である。
$\xi$ が素イデアル $\mathfrak p \subset A$ に対応するとする。
$\xi$ が $D$ の既約成分の生成点であるという仮定は、$\mathfrak p$ が
$(f)$ 上極小であることを意味する。したがって Algebra, Lemma
\ref{algebra-lemma-minimal-over-1} により $\dim(A_\mathfrak p) \leq 1$ である。

\medskip\noindent
(2) を証明する。(1) により $\dim(\mathcal{O}_{X, \xi}) \leq 1$ である。
一方、Lemma \ref{divisors-lemma-characterize-effective-Cartier-divisor} により局所方程式
$f$ は $A_\mathfrak p$ における非零因子である。したがって次元は少なくとも
$1$ である（初等的な Algebra, Lemma
\ref{algebra-lemma-Zariski-topology} により、$A_\mathfrak p$ には $f$ を含まない
素イデアルが存在しなければならない）。
\end{proof}

\begin{lemma}
\label{divisors-lemma-integral-effective-Cartier-divisor-dvr}
$X$ を Noether スキームとし、$D \subset X$ を整閉部分スキームであって、同時に
有効 Cartier 因子であるものとする。このとき $X$ の局所環で $D$ の生成点におけるものは
離散付値環である。
\end{lemma}

\begin{proof}
Lemma \ref{divisors-lemma-characterize-effective-Cartier-divisor} により、
$X = \Spec(A)$ および $D = \Spec(A/(f))$ と仮定してよい。ここで $A$ は Noether
環であり、$f \in A$ は非零因子である。$D$ が整であるという仮定は、$(f)$ が
素イデアルであることを意味する。したがって生成点における $X$ の局所環は
$A_{(f)}$ である。これは極大イデアルが非零因子で生成される Noether 局所環である。
ゆえに Algebra, Lemma \ref{algebra-lemma-characterize-dvr} により離散付値環である。
\end{proof}

\begin{lemma}
\label{divisors-lemma-effective-Cartier-divisor-Sk}
$X$ を局所 Noether スキームとし、$D \subset X$ を有効 Cartier 因子とする。
$X$ が $(S_k)$ ならば、$D$ は $(S_{k - 1})$ である。
\end{lemma}

\begin{proof}
$x \in D$ とする。このとき
$\mathcal{O}_{D, x} = \mathcal{O}_{X, x}/(f)$ であり、ここで
$f \in \mathcal{O}_{X, x}$ は非零因子である。仮定により
$\text{depth}(\mathcal{O}_{X, x}) \geq \min(\dim(\mathcal{O}_{X, x}), k)$
である。Algebra, Lemma \ref{algebra-lemma-depth-drops-by-one} により
$\text{depth}(\mathcal{O}_{D, x}) = \text{depth}(\mathcal{O}_{X, x}) - 1$
であり、Algebra, Lemma \ref{algebra-lemma-one-equation} により
$\dim(\mathcal{O}_{D, x}) = \dim(\mathcal{O}_{X, x}) - 1$ である。したがって
$\text{depth}(\mathcal{O}_{D, x}) \geq \min(\dim(\mathcal{O}_{D, x}), k - 1)$
となり、主張を得る。
\end{proof}

\begin{lemma}
\label{divisors-lemma-normal-effective-Cartier-divisor-S1}
$X$ を局所 Noether 正規スキームとし、$D \subset X$ を有効 Cartier 因子とする。
このとき $D$ は $(S_1)$ である。
\end{lemma}

\begin{proof}
Properties, Lemma \ref{properties-lemma-criterion-normal} により $X$ は $(S_2)$
である。したがって Lemma \ref{divisors-lemma-effective-Cartier-divisor-Sk} から結論を得る。
\end{proof}

\begin{lemma}
\label{divisors-lemma-weil-divisor-is-cartier-UFD}
$X$ を局所 Noether スキームとし、$D \subset X$ を整閉部分スキームとする。
次を仮定する。
\begin{enumerate}
\item $D$ の余次元は $1$ である（$X$ において）。
\item $\mathcal{O}_{X, x}$ はすべての $x \in D$ に対して UFD である。
\end{enumerate}
このとき $D$ は有効 Cartier 因子である。
\end{lemma}

\begin{proof}
$x \in D$ とし、$A = \mathcal{O}_{X, x}$ とおく。$\mathfrak p \subset A$ を
$D$ の生成点に対応する素イデアルとする。仮定 (1) により $A_\mathfrak p$ は次元
$1$ である。したがって $\mathfrak p$ は高さ $1$ の素イデアルである。$A$ は UFD
なので $\mathfrak p = (f)$ となり、ここで $f \in A$ である。当然 $f$ は
非零因子であり、Lemma \ref{divisors-lemma-effective-Cartier-in-points} により結論を得る。
\end{proof}

\begin{lemma}
\label{divisors-lemma-codim-1-part}
$X$ を Noether スキームとし、$Z \subset X$ を閉部分スキームとする。次が存在すると仮定する。
\begin{enumerate}
\item 整有効 Cartier 因子の族 $D_i \subset X$, $i \in I$、
\item 閉部分集合 $Z' \subset X$ であって、そのすべての既約成分の余次元が
$\geq 2$ であるもの（$X$ において）
（Topology, Definition \ref{topology-definition-codimension}）。
\end{enumerate}
さらに集合論的に $Z \subset Z' \cup \bigcup_{i \in I} D_i$ とする。このとき
$a_i \geq 0$ かつ $a_i = 0$ がほとんどすべての $i \in I$ に対して成り立つ整数が存在し、
有効 Cartier 因子
$$
D = \sum a_i D_i
$$
は $Z$ に含まれ、包含射 $D \to Z$ は余次元 $2$ の外で同型である
（$X$ において）。
すなわち、開集合 $U \subset Z$ であって $D \cap U \to Z \cap U$ が同型となり、
$Z \setminus U$ のすべての既約成分の余次元が $\geq 2$ となるものが存在する
（$X$ において）。$Z$ が $X$ において疎であるとき、
$D_i$, $i \in I$ および $Z'$ の存在は、$\mathcal{O}_{X, x}$ がすべての
$x \in Z$ に対して UFD であるか、または $X$ が正則ならば保証される。
\end{lemma}

\begin{proof}
$\xi_i \in D_i$ を生成点とし、$\mathcal{O}_i = \mathcal{O}_{X, \xi_i}$ を
局所環とする。これは Lemma \ref{divisors-lemma-integral-effective-Cartier-divisor-dvr}
により離散付値環である。$a_i \geq 0$ を
$\mathcal{I}_{Z, \xi_i} \subset \mathcal{O}_i$ の元の付値の最小値とする。
もちろん $a_i > 0$ となりうるのは $D_i$ が $Z$ の既約成分である場合だけであり、
したがって $a_i > 0$ となる $i \in I$ は有限個だけである。有効 Cartier 因子
$D = \sum a_i D_i$ が求める性質をもつと主張する。

\medskip\noindent
$x \in X$ とし、$A = \mathcal{O}_{X, x}$ とおく。$D_1, \ldots, D_n$ を、因子
$D_i$ のうち $x \in D_i$ かつ $a_i > 0$ を満たす相異なるものとする。
$1 \leq i \leq n$ に対して $f_i \in A$ を $D_i$ の局所方程式とする。
このとき $f_i$ は $A$ の素元であり、$\mathcal{O}_i = A_{(f_i)}$ である。
$I = \mathcal{I}_{Z, x} \subset A$ を $Z$ のイデアル層の茎とする。
$a_i$ の選び方により $I A_{(f_i)} = f_i^{a_i}A_{(f_i)}$ である。次を主張する。
$I \subset (\prod_{i = 1, \ldots, n} f_i^{a_i})$。

\medskip\noindent
主張を証明する。局所化写像
$\varphi : A/(f_i) \to A_{(f_i)}/f_iA_{(f_i)}$ は単射である。実際、素イデアル
$(f_i)$ は極大イデアル $f_iA_{(f_i)}$ の逆像である。$n$ に関する帰納法により
$\varphi_n : A/(f_i^n)\to A_{(f_i)}/f_i^nA_{(f_i)}$ も単射である。
$\varphi_{a_i}(I) = 0$ だから $I \subset (f_i^{a_i})$ である。したがって任意の
$x \in I$ に対して $x = f_1^{a_1}x_1$ と書け、ここで $x_1 \in A$ である。
$D_1, \ldots, D_n$ は相異なるので、$f_i$ は $A_{(f_j)}$ における単元である
（$i \not = j$ のとき）。$x$ と $x_1$ を $A_{(f_i)}$ において
$n \geq i > 1$ に対して比較すると、なお $x_1 \in (f_i^{a_i})$ である。
この過程を繰り返すと、帰納的に
$x_i = f_{i + 1}^{a_{i + 1}}x_{i + 1}$ と書ける（任意の $n > i \geq 1$ に対して）。
結局、$x \in (\prod_{i = 1, \ldots n} f_i^{a_i})$ が任意の $x \in I$ に対して
成り立ち、主張が証明された。

\medskip\noindent
この主張から $\mathcal{I}_Z \subset \mathcal{I}_D$、すなわち
$D \subset Z$ が従う。さらに $D$ と $Z$ は $\xi_i$ において一致するので、
$D \to Z$ は余次元 $2$ の外で同型である（$X$ 上で）。

\medskip\noindent
最後の諸主張を見るには次のように論じる。正則局所環は UFD である
（More on Algebra, Lemma \ref{more-algebra-lemma-regular-local-UFD}）。したがって
UFD の場合を論じれば十分である。その場合、$D_i$ を $Z$ の既約成分であって
余次元が $1$ のものとする（$X$ において）。Lemma
\ref{divisors-lemma-weil-divisor-is-cartier-UFD} により各 $D_i$ は有効 Cartier 因子である。
\end{proof}

\begin{lemma}
\label{divisors-lemma-codimension-1-is-effective-Cartier}
$Z \subset X$ を Noether スキームの閉部分スキームとする。次を仮定する。
\begin{enumerate}
\item $Z$ は埋め込まれた点をもたない。
\item $Z$ のすべての既約成分は余次元 $1$ をもつ（$X$ において）。
\item 各局所環 $\mathcal{O}_{X, x}$ は $x \in Z$ に対して UFD であるか、
または $X$ は正則である。
\end{enumerate}
このとき $Z$ は有効 Cartier 因子である。
\end{lemma}

\begin{proof}
$D = \sum a_i D_i$ を Lemma \ref{divisors-lemma-codim-1-part} のものとする。ここで
$D_i \subset Z$ は $Z$ の既約成分である。$D \to Z$ が同型でなければ、
$\mathcal{O}_Z \to \mathcal{O}_D$ は余次元 $\geq 2$ に台をもつ非零な核をもつ。
これは $Z$ が埋め込まれた点をもつことを意味し、仮定 (1) に反する。
したがって望み通り $D \cong Z$ である。
\end{proof}

\begin{lemma}
\label{divisors-lemma-UFD-one-equation-CM}
$R$ を Noether UFD とし、$I \subset R$ をイデアルとする。$R/I$ は埋め込まれた
素イデアルをもたず、$I$ 上の各極小素イデアルの高さが $1$ であると仮定する。
このとき $I = (f)$ であり、ここで $f \in R$ である。
\end{lemma}

\begin{proof}
Lemma \ref{divisors-lemma-codimension-1-is-effective-Cartier} により、イデアル層
$\tilde I$ は $\Spec(R)$ 上で可逆である。More on Algebra, Lemma
\ref{more-algebra-lemma-UFD-Pic-trivial} により、一つの元で生成される。
\end{proof}

\begin{lemma}
\label{divisors-lemma-effective-Cartier-divisor-is-a-sum}
$X$ を Noether スキームとし、$D \subset X$ を有効 Cartier 因子とする。
整有効 Cartier 因子 $D_i \subset X$ であって、集合論的に
$D \subset \bigcup D_i$ となるものが存在すると仮定する。このとき
$D = \sum a_i D_i$ であり、ここで $a_i \geq 0$ である。$D_i$ の存在は、
$\mathcal{O}_{X, x}$ がすべての $x \in D$ に対して UFD であるか、または
$X$ が正則ならば保証される。
\end{lemma}

\begin{proof}
Lemma \ref{divisors-lemma-codim-1-part} のように $a_i$ を選び、
$D' = \sum a_i D_i$ とおく。このとき $D' \to D$ は有効 Cartier 因子の包含であり、
余次元 $2$ の外で同型である（$X$ 上で）。$x \in X$ を取り、
$A = \mathcal{O}_{X, x}$ とおく。$f, f' \in A$ を、$D, D'$ のイデアルを
$A$ において生成する非零因子とする。このとき $f = gf'$ であり、ここで
$g \in A$ である。さらに、任意の素イデアル $\mathfrak p$ で、高さ
$\leq 1$ のもの（$A$ の素イデアル）に対して、$g$ は $A_\mathfrak p$ の単元に写る。
これは $g$ が単元であることを意味する。実際、$(g)$ 上の極小素イデアルは高さ
$1$ である（Algebra, Lemma \ref{algebra-lemma-minimal-over-1}）。
\end{proof}

\begin{lemma}
\label{divisors-lemma-quasi-projective-Noetherian-pic-effective-Cartier}
\begin{slogan}
射影スキーム上では、すべての直線束が正則有理型切断をもつ。
\end{slogan}
$X$ を豊富な可逆層をもつ Noether スキームとする。このとき、すべての可逆
$\mathcal{O}_X$-加群は
$$
\mathcal{O}_X(D - D') =
\mathcal{O}_X(D) \otimes_{\mathcal{O}_X} \mathcal{O}_X(D')^{\otimes -1}
$$
と同型である。ここで $D, D'$ は $X$ の有効 Cartier 因子である。さらに、
有限部分集合 $E \subset X$ が与えられたとき、$D, D'$ を
$E \cap D = \emptyset$ および $E \cap D' = \emptyset$ となるように
選べる。$X$ が準アフィンならば $D' = \emptyset$ と選べる。
\end{lemma}

\begin{proof}
$x_1, \ldots, x_n$ を $X$ の随伴点とする
（Lemma \ref{divisors-lemma-finite-ass}）。

\medskip\noindent
$X$ が準アフィンで、$\mathcal{N}$ が任意の可逆 $\mathcal{O}_X$-加群ならば、
切断 $t$、すなわち $\mathcal{N}$ の切断であって、$E \cup \{x_1, \ldots, x_n\}$ のどの点でも
消えないものを取れる。Properties, Lemma
\ref{properties-lemma-quasi-affine-invertible-nonvanishing-section} を参照せよ。
すると Lemma \ref{divisors-lemma-regular-section-associated-points} により $t$ は
$\mathcal{N}$ の正則切断である。したがって $\mathcal{N} \cong \mathcal{O}_X(D)$
であり、ここで $D = Z(t)$ は $t$ に対応する有効 Cartier 因子である。Lemma
\ref{divisors-lemma-characterize-OD} を参照せよ。構成により $E \cap D = \emptyset$ なので、
この場合の証明は完了する。

\medskip\noindent
一般の場合に戻り、$\mathcal{L}$ を $X$ 上の豊富な可逆層とする。ある
$m > 0$ と切断 $s \in \Gamma(X, \mathcal{L}^{\otimes m})$ が存在して、$X_s$ は
アフィンであり、$E \cup \{x_1, \ldots, x_n\} \subset X_s$ となる
（Properties, Lemma \ref{properties-lemma-ample-finite-set-in-principal-affine}）。

\medskip\noindent
$\mathcal{N}$ を任意の可逆 $\mathcal{O}_X$-加群とする。準アフィンの場合により、切断
$t \in \mathcal{N}(X_s)$ であって、$E \cup \{x_1, \ldots, x_n\}$ のどの点でも
消えないものを取れる。Properties, Lemma
\ref{properties-lemma-invert-s-sections} により、ある $e \geq 0$ に対して切断
$s^e|_{X_s} t$ は大域切断 $\tau$、すなわち
$\mathcal{L}^{\otimes me} \otimes \mathcal{N}$ の大域切断に延長される。したがって
$\mathcal{L}^{\otimes me} \otimes \mathcal{N}$ と
$\mathcal{L}^{\otimes me}$ はともに可逆層であり、
$E \cup \{x_1, \ldots, x_n\}$ のどの点でも消えない大域切断をもつ。ゆえに Lemma
\ref{divisors-lemma-regular-section-associated-points} により、これらは正則切断である。
したがって $\mathcal{L}^{\otimes me} \otimes \mathcal{N} \cong \mathcal{O}_X(D)$ および
$\mathcal{L}^{\otimes me} \cong \mathcal{O}_X(D')$ であり、ここで $D$ と $D'$ は
有効 Cartier 因子である。Lemma
\ref{divisors-lemma-characterize-OD} を参照せよ。構成により $E \cap D = \emptyset$ および
$E \cap D' = \emptyset$ であり、証明が完了する。
\end{proof}

\begin{lemma}
\label{divisors-lemma-wedge-product-ses}
$X$ を次元 $2$ の整正則スキームとし、$i : D \to X$ を有効 Cartier 因子の
埋め込みとする。完全列 $\mathcal{F} \to \mathcal{F}' \to i_*\mathcal{G} \to 0$ を
連接 $\mathcal{O}_X$-加群の列として考える。次を仮定する。
\begin{enumerate}
\item $\mathcal{F}, \mathcal{F}'$ は階数 $r$ の局所自由加群である
（$X$ の空でない開集合上で）。
\item $D$ は整スキームである。
\item $\mathcal{G}$ は有限局所自由 $\mathcal{O}_D$-加群であり、階数は $s$ である。
\end{enumerate}
このとき $\mathcal{L} = (\wedge^r\mathcal{F})^{**}$ および
$\mathcal{L}' = (\wedge^r \mathcal{F}')^{**}$ は可逆 $\mathcal{O}_X$-加群であり、
$\mathcal{L}' \cong \mathcal{L}(k D)$ である。ここで
$k \in \{0, \ldots, \min(s, r)\}$ である。
\end{lemma}

\begin{proof}
最初の主張は Lemma \ref{divisors-lemma-reflexive-over-regular-dim-2} から従う。実際、仮定
(1) により $\mathcal{L}$ と $\mathcal{L}'$ は階数 $1$ である。
$\wedge^r$ と二重双対を取る操作は函手なので、標準的な写像
$\sigma : \mathcal{L} \to \mathcal{L}'$ を得る。これは (1) の空でない開集合上で
同型であり、したがって非零である。証明を終えるには、$\sigma$ を
$\mathcal{L}' \otimes \mathcal{L}^{\otimes -1}$ の大域切断とみたとき、それが
$X$ の任意の余次元の点では消えず、ただし $D$ の生成点では例外で、そこでの
消滅次数が高々 $\min(s, r)$ であることを見れば十分である。

\medskip\noindent
代数に移すと、次の問題を得る。$(A, \mathfrak m, \kappa)$ を分数体 $K$ をもつ
離散付値環とする。$M \to M' \to N \to 0$ を有限 $A$-加群の完全列であって、
$\dim_K(M \otimes K) = \dim_K(M' \otimes K) = r$ かつ
$N \cong \kappa^{\oplus s}$ となるものとする。誘導される写像
$L = \wedge^r(M)^{**} \to L' = \wedge^r(M')^{**}$ の消滅次数が高々
$\min(s, r)$ であることを示せ。$A$ 上の加群の構造定理を用いる。More on Algebra,
Lemma \ref{more-algebra-lemma-generalized-valuation-ring-modules} または
\ref{more-algebra-lemma-modules-PID} を参照せよ。有限 $A$-加群をその捩れ部分加群で
割っても二重双対は変わらない。したがって $M$ を $M/M_{tors}$ で、$M'$ を
$M'/\Im(M_{tors} \to M')$ で置き換え、$M$ は捩れなしであると仮定してよい。
すると $M \to M'$ は単射であり、$M'_{tors} \to N$ も単射である。ゆえに
$M'$ を $M'/M'_{tors}$ で、$N$ を $N/M'_{tors}$ で置き換えてよい。
こうして $M$ と $M'$ が階数 $r$ の自由加群であり、
$N \cong \kappa^{\oplus s}$ である場合に帰着する。この場合 $\sigma$ は
$M \to M'$ の行列式であり、その消滅次数は $s$ である。例えば Algebra, Lemma
\ref{algebra-lemma-order-vanishing-determinant} を参照せよ。
\end{proof}

\section{アフィン開集合の補集合}
\label{divisors-section-complement-affine-open}

\noindent
本節では、多様体のアフィン開集合の補集合が純余次元 $1$ をもつという結果を論じる。

\begin{lemma}
\label{divisors-lemma-affine-punctured-spec}
$(A, \mathfrak m)$ を Noether 局所環とする。穿孔スペクトル
$U = \Spec(A) \setminus \{\mathfrak m\}$、すなわち $A$ のスペクトルから閉点を除いた部分が
アフィンであるための必要十分条件は
$\dim(A) \leq 1$ である。
\end{lemma}

\begin{proof}
$\dim(A) = 0$ ならば $U$ は空であり、したがってアフィンである
（$0$ 環のスペクトルに等しい）。$\dim(A) = 1$ ならば、元 $f \in \mathfrak m$ で、
$A$ の有限個の極小素イデアルのいずれにも含まれないものを選べる
（Algebra, Lemmas \ref{algebra-lemma-Noetherian-irreducible-components} および
\ref{algebra-lemma-silly}）。すると $U = \Spec(A_f)$ はアフィンである。

\medskip\noindent
逆はより興味深い。ここではやや非標準的な証明を与え、標準的な議論は下の注意で論じる。
$U = \Spec(B)$ がアフィンであると仮定する。被約化してもアフィン性と次元は変わらない
ので、$A$ をその冪零元イデアルによる商に置き換え、$A$ は被約であると仮定してよい。
$Q = B/A$ を $A$-加群とみなす。$Q$ の台は $\{\mathfrak m\}$ である。実際、
$A_\mathfrak p = B_\mathfrak p$ は、$\mathfrak p$ が $A$ の非極大素イデアルである限り成り立つ。
$\dim(A) \geq 1$ と仮定してよいから、上と同様に $f \in \mathfrak m$ を選べる。この元は
$A$ の極小イデアルのいずれにも含まれない。$A$ は被約なので、これは $f$ が $A$ の
非零因子であることを意味する。同じ議論により $f : B \to B$ は単射である。実際、
$\Spec(B)$ は $\Spec(A)$ の開集合である。Algebra, Lemma
\ref{algebra-lemma-one-equation} により
$\dim(A/fA) = \dim(A) - 1$ であることに注意する。
$f$ による乗法を短完全列 $0 \to A \to B \to Q \to 0$ に適用して蛇の補題を用いると
$$
0 \to Q[f] \to A/fA \to B/fB \to Q/fQ \to 0
$$
を得る。ここで $Q[f] = \Ker(f : Q \to Q)$ である。これは $Q[f]$ が有限
$A$-加群であることを意味する。$Q[f]$ の台は $\{\mathfrak m\}$ なので
$l = \text{length}_A(Q[f]) < \infty$ である
（Algebra, Lemma \ref{algebra-lemma-support-point}）。
$l_n = \text{length}_A(Q[f^n])$ とおく。完全列
$$
0 \to Q[f^n] \to Q[f^{n + 1}] \xrightarrow{f^n} Q[f]
$$
から、帰納的に $l_n < \infty$ および $l_n \leq l_{n + 1}$ が分かる。完全列
$$
0 \to Q[f] \to Q[f^{n + 1}] \xrightarrow{f} Q[f^n] \to Q/fQ
$$
を考えると、$Q[f^n]$ の $Q/fQ$ における像の長さは
$l_n - l_{n + 1} + l \leq l$ である。$Q = \bigcup Q[f^n]$ なので、
$Q/fQ$ の長さは高々 $l$、すなわち有界である。したがって $Q/fQ$ は有限
$A$-加群である。ゆえに $A/fA \to B/fB$ は有限環写像であり、とくにスペクトル上の
閉写像を誘導する（Algebra, Lemmas \ref{algebra-lemma-integral-going-up} および
\ref{algebra-lemma-going-up-closed}）。一方、$\Spec(B/fB)$ は
$\Spec(A/fA)$ の穿孔スペクトルである。これは $\Spec(B/fB) = \emptyset$ でない限り
矛盾であり、したがって望み通り $\dim(A/fA) = 0$ である。
\end{proof}

\begin{remark}
\label{divisors-remark-affine-punctured-spectrum-standard-proof}
$(A, \mathfrak m)$ を次元 $\geq 2$ の Noether 局所正規整域とし、$U$ を $A$ の
穿孔スペクトルとすると、$\Gamma(U, \mathcal{O}_U) = A$ である。この Hartogs
の定理の代数版は、Algebra, Lemma
\ref{algebra-lemma-normal-domain-intersection-localizations-height-1} で見た
$A = \bigcap_{\text{height}(\mathfrak p) = 1} A_\mathfrak p$ から従う。
したがってこの場合 $U$ はアフィンではありえない（そうであれば $\mathfrak m$ が
$U$ の点であることを強制する）。これはしばしば Lemma
\ref{divisors-lemma-affine-punctured-spec} の証明の出発点として用いられる。一般の Noether
局所環の場合をこれに帰着するには、まず完備化し（Nagata 局所環を得るため）、次に
適当な極小素イデアルに対して $A$ を $A/\mathfrak q$ で置き換え、その後正規化する。
これらの各段階は次元を変えず、矛盾を得る。完備化の段階を省くこともできるが、その場合、
一般に正規化は Noether 整域ではない。しかし、それでも同じ次元の Krull 整域であり
（これは Krull--Akizuki を用いて証明される）、同じ議論を適用できる。
\end{remark}

\begin{remark}
\label{divisors-remark-affine-puctured-spectrum-general}
穿孔スペクトルがアフィンとなる非 Noether 局所環 $(A, \mathfrak m)$ をどのように
特徴付けるかは明らかでない。そのような環は、有限生成イデアル $I$ で
$\mathfrak m = \sqrt{I}$ を満たすものをもつ。もちろん、$I$ を $1$ 個の元で生成されるものに取れれば、
$A$ はアフィンな穿孔スペクトルをもつ。これにより多数の非 Noether の例が得られる。
逆に、Lemma \ref{divisors-lemma-affine-punctured-spec} の証明の議論から、そのような環は
$f \in \mathfrak m$ で $H^0_I(A/fA) = 0$ を満たす非零因子をもつことができない
（したがって $A$ は長さ $2$ の正則列をもてない）。さらに、局所環の局所準同型
の終域となる任意の環 $A'$ についても同じことが成り立つ。ただし、その準同型は
$A \to A'$ であり、$\mathfrak m_{A'} = \sqrt{\mathfrak mA'}$ を満たすものとする。
\end{remark}

\begin{lemma}
\label{divisors-lemma-complement-affine-open-immersion}
\begin{reference}
\cite[EGA IV, Corollaire 21.12.7]{EGA4}
\end{reference}
$X$ を局所 Noether スキームとし、$U \subset X$ を、包含射 $U \to X$ がアフィンとなる
開部分スキームとする。生成点 $\xi$ を任意に取る。ただし、これは $X \setminus U$ の既約成分の
生成点とする。このとき、
局所環 $\mathcal{O}_{X, \xi}$ の次元は $\leq 1$ である。$U$ が稠密であるか、または
$\xi$ が $U$ の閉包に含まれるならば、$\dim(\mathcal{O}_{X, \xi}) = 1$ である。
\end{lemma}

\begin{proof}
$\xi$ は $X \setminus U$ の生成点なので、
$$
U_\xi = U \times_X \Spec(\mathcal{O}_{X, \xi}) \subset
\Spec(\mathcal{O}_{X, \xi})
$$
は $\mathcal{O}_{X, \xi}$ の穿孔スペクトルである
（ヒント：Schemes, Lemma \ref{schemes-lemma-specialize-points} を用いる）。
$U \to X$ はアフィンなので、$U_\xi \to \Spec(\mathcal{O}_{X, \xi})$ もアフィンであり
（Morphisms, Lemma \ref{morphisms-lemma-base-change-affine}）、したがって $U_\xi$ は
アフィンである。ゆえに Lemma \ref{divisors-lemma-affine-punctured-spec} により
$\dim(\mathcal{O}_{X, \xi}) \leq 1$ である。$\xi \in \overline{U}$ ならば、
特殊化 $\eta \to \xi$ で $\eta \in U$ を満たすものが存在する。実際、$\eta$ を
$\overline{U}$ の既約成分で $\xi$ を含むものの生成点に取ればよい。$\overline{U}$ は局所
Noether であり、したがって局所的に有限個の既約成分しかもたないので、
$\eta \in U$ である。このとき $\eta \in \Spec(\mathcal{O}_{X, \xi})$ であり、
次元は $0$ ではありえない。
\end{proof}

\begin{lemma}
\label{divisors-lemma-complement-affine-open}
$X$ を分離的局所 Noether スキームとし、$U \subset X$ をアフィン開集合とする。
生成点 $\xi$ を任意に取る。ただし、これは $X \setminus U$ の既約成分の生成点とする。このとき局所環
$\mathcal{O}_{X, \xi}$ の次元は $\leq 1$ である。$U$ が稠密であるか、または
$\xi$ が $U$ の閉包に含まれるならば、$\dim(\mathcal{O}_{X, \xi}) = 1$ である。
\end{lemma}

\begin{proof}
これは Lemma \ref{divisors-lemma-complement-affine-open-immersion} から従う。実際、Morphisms,
Lemma \ref{morphisms-lemma-affine-permanence} により射 $U \to X$ はアフィンである。
\end{proof}

\noindent
次の補題は、有効 Cartier 因子を構成するために使えることがある。

\begin{lemma}
\label{divisors-lemma-complement-open-affine-effective-cartier-divisor}
$X$ を Noether 分離スキームとし、$U \subset X$ を稠密アフィン開集合とする。
$\mathcal{O}_{X, x}$ がすべての $x \in X \setminus U$ に対して UFD ならば、
有効 Cartier 因子 $D \subset X$ で $U = X \setminus D$ を満たすものが存在する。
\end{lemma}

\begin{proof}
$X$ は Noether なので、補集合 $X \setminus U$ は有限個の既約成分
$D_1, \ldots, D_r$ をもつ（$X \setminus U$ 上の被約な誘導部分スキーム構造に
Properties, Lemma \ref{properties-lemma-Noetherian-irreducible-components} を適用する）。
各 $D_i \subset X$ は Lemma \ref{divisors-lemma-complement-affine-open} により余次元 $1$
である（Properties, Lemma \ref{properties-lemma-codimension-local-ring} も参照）。
したがって Lemma \ref{divisors-lemma-weil-divisor-is-cartier-UFD} により $D_i$ は有効 Cartier
因子である。ゆえに $D = D_1 + \ldots + D_r$ と取れる。
\end{proof}

\begin{lemma}
\label{divisors-lemma-complement-open-affine-effective-cartier-divisor-bis}
$X$ をアフィン対角射をもつ Noether スキームとし、$U \subset X$ を稠密アフィン
開集合とする。$\mathcal{O}_{X, x}$ がすべての $x \in X \setminus U$ に対して UFD
ならば、有効 Cartier 因子 $D \subset X$ で $U = X \setminus D$ を満たすものが存在する。
\end{lemma}

\begin{proof}
$X$ は Noether なので、補集合 $X \setminus U$ は有限個の既約成分
$D_1, \ldots, D_r$ をもつ（$X \setminus U$ 上の被約な誘導部分スキーム構造に
Properties, Lemma \ref{properties-lemma-Noetherian-irreducible-components} を適用する）。
$D_i$ を $X$ の被約閉部分スキームとみなす。$X = \bigcup_{j \in J} X_j$ を $X$ の
アフィン開被覆とし、添字 $j$ を $J$ から任意に取り、$U_j = U \cap X_j$ とおく。
$X$ はアフィン対角射をもつので、スキーム
$$
U_j = X \times_{(X \times X)} (U \times X_j)
$$
はアフィンである。したがって $X_j$ は分離的なので、Lemma
\ref{divisors-lemma-complement-open-affine-effective-cartier-divisor} とその証明により、
すべての $j \in J$ および $1 \leq i \leq r$ に対して交叉
$D_i \cap X_j$ は空であるか、$X_j$ における有効 Cartier 因子である。
ゆえに $D_i \subset X$ は有効 Cartier 因子である（これは局所的な性質である）。
したがって $D = D_1 + \ldots + D_r$ と取れる。
\end{proof}

\begin{lemma}
\label{divisors-lemma-regular-ample-family}
$X$ をアフィン対角射をもつ準コンパクト正則スキームとする。このとき $X$ は可逆加群の
豊富な族をもつ（Morphisms, Definition
\ref{morphisms-definition-family-ample-invertible-modules}）。
\end{lemma}

\begin{proof}
$X$ は有限個の整スキームの非交和であることに注意する
（Properties, Lemmas \ref{properties-lemma-regular-normal} および
\ref{properties-lemma-normal-Noetherian}）。したがって $X$ は整かつ Noether で、
正則かつアフィン対角射をもつと仮定してよい。$x \in X$ とする。アフィン開集合
$U \subset X$ で $x$ を含むものを選ぶ。$X$ は整なので、$U$ は $X$ において稠密である。
More on Algebra, Lemma \ref{more-algebra-lemma-regular-local-UFD} により $X$ の局所環は
UFD である。したがって Lemma
\ref{divisors-lemma-complement-open-affine-effective-cartier-divisor-bis} により、有効 Cartier 因子
$D \subset X$ で、その補集合が $U$ となるものを取れる。このとき標準切断
$s = 1_D$ は、$\mathcal{L} = \mathcal{O}_X(D)$ の切断であり、
（Definition \ref{divisors-definition-invertible-sheaf-effective-Cartier-divisor} を参照）、
$D$ に沿ってちょうど消えるので $U = X_s$ である。したがって Morphisms, Definition
\ref{morphisms-definition-family-ample-invertible-modules} の両条件が成り立ち、証明が完了する。
\end{proof}

\section{ノルム}
\label{divisors-section-norms}

\noindent
$\pi : X \to Y$ をスキームの有限射とし、$d \geq 1$ を整数とする。乗法的写像
$$
\text{Norm}_\pi : \pi_*\mathcal{O}_X \to \mathcal{O}_Y
$$
であって次を満たすものが存在するとき、{\it 次数 $d$ の $\pi$ に対する
ノルム}\footnote{これは標準的でない用語法である。}が存在するという。
\begin{enumerate}
\item 合成
$\mathcal{O}_Y \xrightarrow{\pi^\sharp} \pi_*\mathcal{O}_X
\xrightarrow{\text{Norm}_\pi} \mathcal{O}_Y$ は $g \mapsto g^d$ に等しい。
\item 開集合 $V \subset Y$ に対し、$f \in \mathcal{O}_X(\pi^{-1}V)$ が
$x \in \pi^{-1}(V)$ において零ならば、$\text{Norm}_\pi(f)$ は
$\pi(x)$ において零である。
\end{enumerate}
条件 (1) から $\pi$ が全射であることに注意する。$\text{Norm}_\pi$ は乗法的なので
単元を単元に送る。したがって、$y \in Y$ を与え、$f$ が $X$ 上の正則函数で、
$x \in X$ のうち $\pi(x) = y$ を満たす任意の点で定義され、かつ消えないならば、
$\text{Norm}_\pi(f)$ は定義されており、$y$ において消えない。これは (2) を仮定
しなくても成り立つ。実際、本節の構成には条件 (1) だけを用い、特定の消滅性質
（とくに Lemma \ref{divisors-lemma-norm-ample} の証明で用いる）にのみ性質 (2) が必要となる。

\begin{lemma}
\label{divisors-lemma-finite-trivialize-invertible-upstairs}
$\pi : X \to Y$ をスキームの有限射とする。$\mathcal{L}$ を可逆
$\mathcal{O}_X$-加群とし、$y \in Y$ とする。このとき開近傍
$V \subset Y$ で、$y$ を含み、$\mathcal{L}|_{\pi^{-1}(V)}$ が自明となるものが存在する。
\end{lemma}

\begin{proof}
$Y$、したがって $X$ がアフィンであると仮定してよい。$\pi$ は有限なので、ファイバー
$\pi^{-1}(\{y\})$ は $y$ 上で有限である。$X$ はアフィンだから、
$s \in \Gamma(X, \mathcal{L})$ を $\pi^{-1}(\{y\})$ のどの点でも消えないように
選べる。これは Properties, Lemma
\ref{properties-lemma-quasi-affine-invertible-nonvanishing-section}
から従うが、直接の議論も与える。すなわち、閉点の有限集合 $E \subset X$ を、
各 $x \in \pi^{-1}(\{y\})$ が $E$ のある点へ特殊化するように選べる。
$x \in E$ に対し、閉埋入を $i_x : x \to X$ と書く。このとき
$\mathcal{L} \to \bigoplus_{x \in E} i_{x, *}i_x^*\mathcal{L}$
は準連接 $\mathcal{O}_X$-加群の全射である。したがって写像
$$
\Gamma(X, \mathcal{L}) \to
\bigoplus\nolimits_{x \in E} \mathcal{L}_x/\mathfrak m_x\mathcal{L}_x
$$
は全射である（準連接 $\mathcal{O}_X$-加群の圏では大域切断を取る函数が完全だから
である。Schemes, Lemma \ref{schemes-lemma-equivalence-quasi-coherent} を参照）。
ゆえに、$s \in \Gamma(X, \mathcal{L})$ を、$E$ の点へ特殊化するどの点でも
消えないように選べる。このとき $X_s \subset X$ は $\pi^{-1}(\{y\})$ の開近傍である。
$\pi$ は有限、したがって閉であるから、開近傍 $V \subset Y$ で $y$ を含み、その逆像が
$X_s$ に含まれるものが存在し、これが求めるものである。
\end{proof}

\begin{lemma}
\label{divisors-lemma-norm-invertible}
$\pi : X \to Y$ をスキームの有限射とする。次数 $d$ の $\pi$ に対するノルムが存在する
ならば、アーベル群の準同型
$$
\text{Norm}_\pi : \Pic(X) \to \Pic(Y)
$$
で、等式 $\text{Norm}_\pi(\pi^*\mathcal{N}) \cong \mathcal{N}^{\otimes d}$ が、任意の可逆
$\mathcal{O}_Y$-加群 $\mathcal{N}$ に対して成り立つものが存在する。
\end{lemma}

\begin{proof}
可逆 $\mathcal{O}_X$-加群の同型類と
$H^1(X, \mathcal{O}_X^*)$ の元との対応（Cohomology, Lemma
\ref{cohomology-lemma-h1-invertible}）を、以下では断りなく用いる。可逆
$\mathcal{O}_X$-加群 $\mathcal{L}$ のノルムの取り方を説明する。Lemma
\ref{divisors-lemma-finite-trivialize-invertible-upstairs} により、開被覆
$Y = \bigcup V_j$ で $\mathcal{L}|_{\pi^{-1}V_j}$ が自明となるものが存在する。
生成切断 $s_j \in \mathcal{L}(\pi^{-1}V_j)$ を各 $j$ について選ぶ。重なり
$\pi^{-1}V_j \cap \pi^{-1}V_{j'}$ 上では
$$
s_j = u_{jj'} s_{j'}
$$
と書ける。ここで
$u_{jj'} \in \mathcal{O}^*_X(\pi^{-1}V_j \cap \pi^{-1}V_{j'})$ は一意である。
したがって元
$$
v_{jj'} = \text{Norm}_\pi(u_{jj'}) \in \mathcal{O}_Y^*(V_j \cap V_{j'})
$$
を考えられる。これらの元はコサイクル条件を満たす（$u_{jj'}$ がそうであり、
$\text{Norm}_\pi$ が乗法的だからである）。ゆえに可逆 $\mathcal{O}_Y$-加群を定める。
これが良定義であること、Picard 群上で加法的であること、および等式
$\text{Norm}_\pi(\pi^*\mathcal{N}) \cong \mathcal{N}^{\otimes d}$ が任意の可逆
$\mathcal{O}_Y$-加群 $\mathcal{N}$ に対して成り立つことの確認は省略する。
\end{proof}

\begin{lemma}
\label{divisors-lemma-norm-map-invertible}
$\pi : X \to Y$ をスキームの有限射とする。次数 $d$ の $\pi$ に対するノルムが存在すると
仮定する。可逆 $\mathcal{O}_X$-加群の間の任意の写像
$\varphi : \mathcal{L} \to \mathcal{L}'$ が $\mathcal{O}_X$-線形ならば、$\mathcal{O}_Y$-線形写像
$$
\text{Norm}_\pi(\varphi) :
\text{Norm}_\pi(\mathcal{L})
\longrightarrow
\text{Norm}_\pi(\mathcal{L}')
$$
が存在する。ここで $\text{Norm}_\pi(\mathcal{L})$ と
$\text{Norm}_\pi(\mathcal{L}')$ は Lemma \ref{divisors-lemma-norm-invertible} のものとする。
さらに、$y \in Y$ に対して次は同値である。
\begin{enumerate}
\item $\varphi$ はある点 $x \in X$ において零であり、その点は $\pi(x) = y$ を満たす。
\item $\text{Norm}_\pi(\varphi)$ は $y$ において零である。
\end{enumerate}
\end{lemma}

\begin{proof}
開被覆 $Y = \bigcup V_j$ を、$\mathcal{L}$ と $\mathcal{L}'$ が各開集合
$\pi^{-1}V_j$ 上で自明となるように選ぶ。これは Lemma
\ref{divisors-lemma-finite-trivialize-invertible-upstairs} により可能である。
生成切断 $s_j$ と $s'_j$ を、それぞれ $\mathcal{L}$ と $\mathcal{L}'$ について、開集合
$\pi^{-1}V_j$ 上で選ぶ。このとき $\varphi(s_j) = f_js'_j$ であり、ある
$f_j \in \mathcal{O}_X(\pi^{-1}V_j)$ が存在する。$\text{Norm}_\pi(\varphi)$ を、
$\text{Norm}_\pi(f_j)$ による $V_j$ 上の乗法として定める。Lemma
\ref{divisors-lemma-norm-invertible} の証明で $\text{Norm}_\pi(\mathcal{L})$ と
$\text{Norm}_\pi(\mathcal{L}')$ の構成に用いたコサイクルについて簡単な計算をすれば、
これが補題に述べた写像を定めることが分かる。最後の主張は、次数 $d$ のノルム写像の
定義における条件 (2) から従う。細部は省略する。
\end{proof}

\begin{lemma}
\label{divisors-lemma-norm-ample}
$\pi : X \to Y$ をスキームの有限射とする。$X$ が豊富な可逆層をもち、次数 $d$ の
$\pi$ に対するノルムが存在すると仮定する。このとき $Y$ は豊富な可逆層をもつ。
\end{lemma}

\begin{proof}
$\mathcal{L}$ を、仮定で与えられた $X$ 上の豊富な可逆層とする。
$\mathcal{N} = \text{Norm}_\pi(\mathcal{L})$ が $Y$ 上で豊富であることを示す。

\medskip\noindent
$X$ は準コンパクトであり（Properties, Definition
\ref{properties-definition-ample}）、$X \to Y$ は全射である
（$\text{Norm}_\pi$ の存在による）から、$Y$ は準コンパクトである。
点 $y \in Y$ を取る。証明を終えるため、切断 $t$ で、$\mathcal{N}$ のある正の
テンソル冪に属し、$y$ において消えず、かつ $Y_t$ がアフィンとなるものが存在することを示す。
そのため、アフィン開近傍 $V \subset Y$ で $y$ を含むものを選ぶ。$n \gg 0$ と切断
$s \in \Gamma(X, \mathcal{L}^{\otimes n})$ を、
$$
\pi^{-1}(\{y\}) \subset X_s \subset \pi^{-1}V
$$
となるように選ぶ。これは Properties, Lemma
\ref{properties-lemma-ample-finite-set-in-principal-affine} による。このとき
$t = \text{Norm}_\pi(s)$ は $\mathcal{N}^{\otimes n}$ の切断で、$x$ において消えず、
$Y_t \subset V$ を満たす。Lemma \ref{divisors-lemma-norm-map-invertible} を参照。すると
$Y_t$ は Properties, Lemma \ref{properties-lemma-affine-cap-s-open} によりアフィンである。
\end{proof}

\begin{lemma}
\label{divisors-lemma-norm-quasi-affine}
$\pi : X \to Y$ をスキームの有限射とする。$X$ が準アフィンで、次数 $d$ の
$\pi$ に対するノルムが存在すると仮定する。このとき $Y$ は準アフィンである。
\end{lemma}

\begin{proof}
Properties, Lemma \ref{properties-lemma-quasi-affine-O-ample} により
$\mathcal{O}_X$ は $X$ 上の豊富な可逆層である。Lemma \ref{divisors-lemma-norm-ample} の証明から
$\text{Norm}_\pi(\mathcal{O}_X) = \mathcal{O}_Y$ は豊富な可逆
$\mathcal{O}_Y$-加群である。したがって Properties, Lemma
\ref{properties-lemma-quasi-affine-O-ample} により $Y$ は準アフィンである。
\end{proof}

\begin{lemma}
\label{divisors-lemma-finite-locally-free-has-norm}
$\pi : X \to Y$ を次数 $d \geq 1$ の有限局所自由射とする。このとき、任意の基底変換と
可換する次数 $d$ の標準的なノルムが存在する。
\end{lemma}

\begin{proof}
$V \subset Y$ を、$(\pi_*\mathcal{O}_X)|_V$ が階数 $d$ の有限自由となるアフィン開集合
とする。基底を選ぶと同型
$$
\mathcal{O}_V^{\oplus d} \cong (\pi_*\mathcal{O}_X)|_V
$$
を得る。任意の $f \in \pi_*\mathcal{O}_X(V) = \mathcal{O}_X(\pi^{-1}(V))$ に対し、
$f$ による乗法は、表示された自由ベクトル束の $\mathcal{O}_V$-線形自己準同型 $m_f$
を定める。したがって $d \times d$ 行列
$M_f \in \text{Mat}(d \times d, \mathcal{O}_Y(V))$ を得て、
$$
\text{Norm}_\pi(f) = \det(M_f)
$$
とおける。行列の行列式は選んだ基底に依存しないので、これは良定義であり、この構成が
貼り合わさって所望の大域写像を与えることも分かる。基底変換との両立性は構成から直ちに従う。

\medskip\noindent
性質 (1) は、$d \times d$ 対角行列で成分が $g, g, \ldots, g$ であるものの行列式が
$g^d$ であることから従う。性質 (2) を見るため、基底変換を行い、$Y$ が体 $k$ の
スペクトルであると仮定してよい。このとき $X = \Spec(A)$ であり、$A$ は
$k$-代数で、$\dim_k(A) = d$ を満たす。$x \in X$ が存在し、$f \in A$ が $x$ で
消えるならば、体への写像 $A \to \kappa$ で $f$ が $\kappa$ において零へ写るものが
存在する。このとき $f : A \to A$ は全射ではありえず、したがって
$\det(f : A \to A) = 0$ となり、求める主張を得る。
\end{proof}

\begin{lemma}
\label{divisors-lemma-norm-in-normal-case}
$\pi : X \to Y$ を有限全射とし、$X$ と $Y$ は整で、$Y$ は正規であるとする。
このとき次数 $[R(X) : R(Y)]$ の $\pi$ に対するノルムが存在する。
\end{lemma}

\begin{proof}
$\Spec(B) \subset Y$ をアフィン開部分集合とし、その逆像を
$\Spec(A) \subset X$ とする。このとき $A$ と $B$ は整域である。
$K$ を $A$ の分数体とし、$L$ を $B$ の分数体とする。図式は次である。
$$
\xymatrix{
L \ar[r] & K \\
B \ar[u] \ar[r] & A \ar[u]
}
$$
$K/L$ は有限拡大なので、ノルム写像
$\text{Norm}_{K/L} : K^* \to L^*$ が存在し、その次数は $d = [K : L]$ である。これは Lemma
\ref{divisors-lemma-finite-locally-free-has-norm} の証明と同様に、$f \in K$ を
$\det_L(f : K \to K)$ へ写すことで与えられる。$f : K \to K$ の特性多項式は
$f$ の $L$ 上の最小多項式の冪であることに注意する。とくに
$\text{Norm}_{K/L}(f)$ は、$f$ の $L$ 上の最小多項式の定数項の冪である。
したがって Algebra, Lemma
\ref{algebra-lemma-minimal-polynomial-normal-domain} により
$\text{Norm}_{K/L}$ は $A$ を $B$ の中へ写す。これによりアフィン開上の切断について
両立する写像系が定まり、したがって大域ノルム写像
$\text{Norm}_\pi$ が定まる。その次数は $d$ である。

\medskip\noindent
性質 (1) は構成から直ちに従う。性質 (2) を見るため、$f \in A$ を取り、これは素イデアル
$\mathfrak p \subset A$ に含まれるとする。多項式
$f^m + b_1 f^{m - 1} + \ldots + b_m$ を、$f$ の $L$ 上の最小多項式とする。Algebra, Lemma
\ref{algebra-lemma-minimal-polynomial-normal-domain} により $b_i \in B$ である。
したがって $b_0 \in B \cap \mathfrak p$ である。上で見たように
$\text{Norm}_{K/L}(f) = b_0^{d/m}$ だから、ノルムは $\mathfrak p$ の像の点で
消えることが分かる。
\end{proof}

\begin{lemma}
\label{divisors-lemma-Frobenius-gives-norm-for-reduction}
$X$ を Noether スキームとする。$p$ を $p\mathcal{O}_X = 0$ を満たす素数とする。
このとき、ある $e > 0$ に対して、次数 $p^e$ の $X_{red} \to X$ に対するノルムが
存在する。ここで $X_{red}$ は $X$ の被約化である。
\end{lemma}

\begin{proof}
$A$ を $pA = 0$ を満たす Noether 環とする。$I \subset A$ を冪零元のイデアルとする。
このとき $I^n = 0$ となる $n$ が存在する（Algebra, Lemma
\ref{algebra-lemma-Noetherian-power}）。$e$ を $p^e \geq n$ となるように選ぶ。このとき
$$
A/I \longrightarrow A,\quad
f \bmod I \longmapsto f^{p^e}
$$
は良定義である。これにより次数 $p^e$ の $\Spec(A/I) \to \Spec(A)$ に対するノルムを
得る。ここで $X$ がアフィンスキーム $\Spec(A_i)$ をいくつか貼り合わせて得られる
ならば、ある $e \gg 0$ に対してこれらの写像は貼り合わさり、
$X_{red} \to X$ に対するノルム写像を与える。細部は省略する。
\end{proof}

\begin{proposition}
\label{divisors-proposition-push-down-ample}
$\pi : X \to Y$ をスキームの有限全射とする。$X$ が豊富な可逆
$\mathcal{O}_X$-加群をもつと仮定する。次のいずれかが成り立つとする。
\begin{enumerate}
\item $\pi$ は有限局所自由である。
\item $Y$ は整正規スキームである。
\item $Y$ は Noether で、$p\mathcal{O}_Y = 0$ かつ $X = Y_{red}$ である。
\end{enumerate}
このとき $Y$ は豊富な可逆 $\mathcal{O}_Y$-加群をもつ。
\end{proposition}

\begin{proof}
場合 (1) は Lemmas \ref{divisors-lemma-finite-locally-free-has-norm} と
\ref{divisors-lemma-norm-ample} を組み合わせれば従う。場合 (3) は Lemmas
\ref{divisors-lemma-Frobenius-gives-norm-for-reduction} と
\ref{divisors-lemma-norm-ample} を組み合わせれば従う。場合 (2) では、まず $X$ を、$X$ の
既約成分で $Y$ を支配するものに置き換える（これを $X$ の被約閉部分スキームとみなす）。
その後 Lemma \ref{divisors-lemma-norm-in-normal-case} を適用できる。
\end{proof}

\begin{lemma}
\label{divisors-lemma-push-down-quasi-affine}
$\pi : X \to Y$ をスキームの有限全射とする。$X$ が準アフィンであると仮定する。
次のいずれかが成り立つとする。
\begin{enumerate}
\item $\pi$ は有限局所自由である。
\item $Y$ は整正規スキームである。
\end{enumerate}
このとき $Y$ は準アフィンである。
\end{lemma}

\begin{proof}
場合 (1) は Lemmas \ref{divisors-lemma-finite-locally-free-has-norm} と
\ref{divisors-lemma-norm-quasi-affine} を組み合わせれば従う。場合 (2) では、まず $X$ を、$X$ の
既約成分で $Y$ を支配するものに置き換える（これを $X$ の被約閉部分スキームとみなす）。
その後 Lemma \ref{divisors-lemma-norm-in-normal-case} を適用できる。
\end{proof}

\section{相対有効 Cartier 因子}
\label{divisors-section-effective-Cartier-morphisms}

\noindent
次の補題は、基底上平坦な有効 Cartier 因子が、実際に基底上の
「有効 Cartier 因子の族」であることを示す。たとえば、任意のファイバーへの制限は
有効 Cartier 因子である。

\begin{lemma}
\label{divisors-lemma-relative-Cartier}
$f : X \to S$ をスキームの射とし、$D \subset X$ を閉部分スキームとする。
次を仮定する。
\begin{enumerate}
\item $D$ は有効 Cartier 因子である。
\item $D \to S$ は平坦射である。
\end{enumerate}
このとき、スキームの任意の射 $g : S' \to S$ に対して、引き戻し
$(g')^{-1}D$ は $X' = S' \times_S X$ 上の有効 Cartier 因子である。ここで
$g' : X' \to X$ は射影である。
\end{lemma}

\begin{proof}
Lemma \ref{divisors-lemma-characterize-effective-Cartier-divisor} を用いて、これを次のように代数へ
翻訳する。$A \to B$ を環写像とし、$h \in B$ とする。$h$ が非零因子であり、
$B/hB$ が $A$ 上平坦であると仮定する。このとき
$$
0 \to B \xrightarrow{h} B \to B/hB \to 0
$$
は $A$-加群の短完全列であり、$B/hB$ は $A$ 上平坦である。Algebra, Lemma
\ref{algebra-lemma-flat-tor-zero} により、この列は任意の加群との $A$ 上の
テンソル積を取っても完全である。とくに、任意の $A$-代数 $A'$ について完全である。
\end{proof}

\noindent
この補題が次の定義の動機である。

\begin{definition}
\label{divisors-definition-relative-effective-Cartier-divisor}
$f : X \to S$ をスキームの射とする。$X/S$ 上の{\it 相対有効 Cartier 因子}とは、
有効 Cartier 因子 $D \subset X$ であって、$D \to S$ がスキームの平坦射となるものをいう。
\end{definition}

\noindent
これは標準的でない用語法かもしれないので注意する。とくに
\cite[IV, Section 21.15]{EGA} では、$X \to S$ が平坦かつ局所有限表示である場合に
限って相対因子の概念を論じている。本書の定義は少し一般的である。しかし多くの場合
（常にではない）、$x \in D$ ならば $X \to S$ は $x$ において平坦であることが分かる。

\begin{lemma}
\label{divisors-lemma-sum-relative-effective-Cartier-divisor}
$f : X \to S$ をスキームの射とする。$D_1, D_2 \subset X$ が $X/S$ 上の相対有効
Cartier 因子ならば、$D_1 + D_2$ もそうである
（Definition \ref{divisors-definition-sum-effective-Cartier-divisors}）。
\end{lemma}

\begin{proof}
これは次の代数的事実に翻訳される。$A \to B$ を環写像とし、
$h_1, h_2 \in B$ とする。$h_i$ が非零因子で、$B/h_iB$ が $A$ 上平坦であると仮定する。
このとき $h_1h_2$ は非零因子で、$B/h_1h_2B$ は $A$ 上平坦である。その理由は、短完全列
$$
0 \to B/h_1B \to B/h_1h_2B \to B/h_2B \to 0
$$
が存在し、最初の矢印が $h_2$ による乗法で与えられるからである。外側の二項は
$A$ 上平坦なので、中央の項も平坦である。Algebra, Lemma
\ref{algebra-lemma-flat-ses} を参照。
\end{proof}

\begin{lemma}
\label{divisors-lemma-difference-relative-effective-Cartier-divisor}
$f : X \to S$ をスキームの射とする。$D_1, D_2 \subset X$ が $X/S$ 上の相対有効
Cartier 因子で、閉部分スキームとして $D_1 \subset D_2$ であるとする。このとき
有効 Cartier 因子 $D$ で $D_2 = D_1 + D$ を満たすもの
（Lemma \ref{divisors-lemma-difference-effective-Cartier-divisors}）は
$X/S$ 上の相対有効 Cartier 因子である。
\end{lemma}

\begin{proof}
これは次の代数的事実に翻訳される。$A \to B$ を環写像とし、
$h_1, h_2 \in B$ とする。$h_i$ が非零因子で、$B/h_iB$ が $A$ 上平坦であり、
$(h_2) \subset (h_1)$ であると仮定する。このとき $h_2 = h h_1$ と書ける。ここで
$h \in B$ は非零因子である。短完全列
$$
0 \to B/hB \to B/h_2B \to B/h_1B \to 0
$$
を得る。最初の矢印は $h_1$ による乗法で与えられる。右側の二項は $A$ 上平坦なので、
中央の項も平坦である。Algebra, Lemma \ref{algebra-lemma-flat-ses} を参照。
\end{proof}

\begin{lemma}
\label{divisors-lemma-flat-at-x}
$f : X \to S$ をスキームの射とし、$D \subset X$ を $X/S$ 上の相対有効 Cartier
因子とする。$x \in D$ かつ $\mathcal{O}_{X, x}$ が Noether 環ならば、$f$ は
$x$ において平坦である。
\end{lemma}

\begin{proof}
$A = \mathcal{O}_{S, f(x)}$ および $B = \mathcal{O}_{X, x}$ とおく。
$h \in B$ を $D$ のイデアルを生成する元とする。このとき $h$ は $B$ の非零因子で、
$B/hB$ は平坦な局所 $A$-代数である。$I \subset A$ を有限生成イデアルとする。
可換図式
$$
\xymatrix{
0 \ar[r] &
B \ar[r]_h &
B \ar[r] &
B/hB \ar[r] & 0 \\
0 \ar[r] &
B \otimes_A I \ar[r]^h \ar[u] &
B \otimes_A I \ar[r] \ar[u] &
B/hB \otimes_A I \ar[r] \ar[u] & 0
}
$$
を考える。$B/hB$ は $A$ 上平坦なので、下の列は短完全である。Algebra, Lemma
\ref{algebra-lemma-flat-tor-zero} を参照。$B/hB$ は $A$ 上平坦なので、右端の
垂直矢印は単射である。Algebra, Lemma \ref{algebra-lemma-flat} を参照。
したがって、$h$ による乗法は中央の垂直矢印の核 $K$ 上で全射である。
中山の補題（Algebra, Lemma \ref{algebra-lemma-NAK}）により $K= 0$ と分かる。
したがって $B$ は $A$ 上平坦である。Algebra, Lemma \ref{algebra-lemma-flat} を参照。
\end{proof}

\noindent
次の補題は、平坦性の軌跡が開であることの代数版に依拠する。スキーム論的な版は
More on Morphisms, Section \ref{more-morphisms-section-open-flat} にある。

\begin{lemma}
\label{divisors-lemma-flat-relative-Cartier-divisor}
$f : X \to S$ をスキームの射とし、$D \subset X$ を相対有効 Cartier 因子とする。
$f$ が局所有限表示ならば、開部分スキーム $U \subset X$ で、$D \subset U$ かつ
$f|_U : U \to S$ が平坦となるものが存在する。
\end{lemma}

\begin{proof}
$x \in D$ を取る。開近傍 $U \subset X$ で $x$ を含み、$f|_U$ が平坦となるものを見つければ
十分である。したがって、補題は $X = \Spec(B)$ と $S = \Spec(A)$ がアフィンで、
$D$ が非零因子 $h \in B$ により与えられる場合に帰着する。仮定により $B$ は有限表示
$A$-代数で、$B/hB$ は平坦な $A$-代数である。絶対 Noether 近似を用いる。

\medskip\noindent
$B = A[x_1, \ldots, x_n]/(g_1, \ldots, g_m)$ と書く。$h$ が
$h' \in A[x_1, \ldots, x_n]$ の像であると仮定する。有限型
$\mathbf{Z}$-部分代数 $A_0 \subset A$ を、多項式
$h', g_1, \ldots, g_m$ のすべての係数が $A_0$ に属するように選ぶ。このとき
$B_0 = A_0[x_1, \ldots, x_n]/(g_1, \ldots, g_m)$ とおき、$h_0$ を $h'$ の
$B_0$ における像とおける。このとき $B = B_0 \otimes_{A_0} A$ かつ
$B/hB = B_0/h_0B_0 \otimes_{A_0} A$ である。Algebra, Lemma
\ref{algebra-lemma-flat-finite-presentation-limit-flat} により、$A_0$ を拡大した後、
$B_0/h_0B_0$ が $A_0$ 上平坦であると仮定してよい。
$K_0 = \Ker(h_0 : B_0 \to B_0)$ とおく。$B_0$ は $\mathbf{Z}$ 上有限型なので、
$K_0$ は有限生成イデアルである。$A_1 \subset A$ を有限型
$\mathbf{Z}$-部分代数で $A_0$ を含むものとし、対応する対象 $B_1$、$h_1$、$K_1$ を
$A_1$ 上のものとする。
More on Algebra, Lemma \ref{more-algebra-lemma-base-change-H1-regular} により、写像
$K_0 \otimes_{A_0} A_1 \to K_1$ は全射である。一方、仮定により
$h : B \to B$ の核は零である。したがって $K_0$ の各元は $K_1$ において零へ写る。
これは十分大きい部分環 $A_1 \subset A$ に対して成り立つ。$K_0$ は有限生成なので、
$K_1 = 0$ となるように適当な $A_1$ を選べる。

\medskip\noindent
$f_1 : X_1 \to S_1$ を、$\Spec$ を環写像 $A_1 \to B_1$ に適用して得られる射とおく。
$D_1 = \Spec(B_1/h_1B_1)$ とおく。$B = B_1 \otimes_{A_1} A$、すなわち
$X = X_1 \times_{S_1} S$ なので、$X_1 \to S_1$ と相対有効 Cartier 因子 $D_1$
について補題を証明すれば十分である。Morphisms, Lemma
\ref{morphisms-lemma-base-change-module-flat} を参照。したがって $A$ が Noether
環である場合に帰着した。この場合、環写像 $A \to B$ は、素イデアル
$\mathfrak q$ が $V(h)$ の点であれば、その点で平坦であることが Lemma \ref{divisors-lemma-flat-at-x} から分かる。
この場合に平坦性の軌跡が開であること（Algebra, Theorem
\ref{algebra-theorem-openness-flatness}）と組み合わせれば証明が終わる。
\end{proof}

\noindent
有限性の仮定をまったく必要としない次の補題もある（着想は Michael Artin によるもの
らしい。\cite{Nobile} を参照）。

\begin{lemma}
\label{divisors-lemma-michael-artin}
$f : X \to S$ をスキームの射とする。$D \subset X$ を $X/S$ 上の相対有効 Cartier
因子とする。$f$ が $X \setminus D$ のすべての点で平坦ならば、$f$ は平坦である。
\end{lemma}

\begin{proof}
これは次の代数的事実に翻訳される。$A \to B$ を環写像とし、$h \in B$ とする。
$h$ が非零因子で、$B/hB$ が $A$ 上平坦であり、局所化 $B_h$ が $A$ 上平坦であると
仮定する。このとき $B$ は $A$ 上平坦である。その理由は短完全列
$$
0 \to B \to B_h \to \colim_n (1/h^n)B/B \to 0
$$
があり、第二項と第三項が $A$ 上平坦であるため、$B$ も $A$ 上平坦となるからである
（Algebra, Lemma \ref{algebra-lemma-flat-ses} を参照）。平坦加群のフィルター余極限は
平坦であること（Algebra, Lemma \ref{algebra-lemma-colimit-flat}）と、各
$n$ に関する帰納法により、各 $(1/h^n)B/B \cong B/h^nB$ が $A$ 上平坦であることに
注意する。後者は短完全列
$$
0 \to B/h^{n - 1}B \xrightarrow{h} B/h^nB \to B/hB \to 0
$$
に入ることから従う。細部は省略する。
\end{proof}

\begin{example}
\label{divisors-example-relative-cartier-ambient-space-not-flat}
相対有効 Cartier 因子 $D$ で、その近傍において周囲のスキームが平坦でない例を与える。
この $D$ の具体例として、$A = k[t]$ とし、
$$
B = k[t, x, y, x^{-1}y, x^{-2}y, \ldots]/(ty, tx^{-1}y, tx^{-2}y, \ldots)
$$
とおく。このとき $B$ は $A$ 上平坦でないが、$B/xB \cong A$ は $A$ 上平坦である。
さらに $x$ は非零因子なので、$\Spec(B)$ の中で $\Spec(A)$ 上の相対有効 Cartier
因子を定める。
\end{example}

\noindent
周囲のスキームが基底上平坦かつ局所有限表示ならば、相対有効 Cartier 因子をその
ファイバーにより特徴付けられる。この補題の少し異なる扱いについては
More on Morphisms, Lemma \ref{more-morphisms-lemma-slice-given-element} も参照。

\begin{lemma}
\label{divisors-lemma-fibre-Cartier}
$\varphi : X \to S$ を平坦かつ局所有限表示な射とする。$Z \subset X$ を閉部分スキーム
とし、$x \in Z$ の像を $s \in S$ とする。
\begin{enumerate}
\item $Z_s \subset X_s$ が $x$ の近傍で Cartier 因子ならば、開集合 $U \subset X$ と
相対有効 Cartier 因子 $D \subset U$ で、$Z \cap U \subset D$ かつ
$Z_s \cap U = D_s$ を満たすものが存在する。
\item $Z_s \subset X_s$ が $x$ の近傍で Cartier 因子で、射 $Z \to X$ が有限表示、
$Z \to S$ が $x$ において平坦ならば、$U$ と $D$ を $Z \cap U = D$ となるように選べる。
\item $Z_s \subset X_s$ が $x$ の近傍で Cartier 因子で、$Z$ が $X$ の局所単項
閉部分スキームとして $x$ の近傍で与えられるならば、$U$ と $D$ を
$Z \cap U = D$ となるように選べる。
\end{enumerate}
とくに、$Z \to S$ が局所有限表示かつ平坦で、すべてのファイバー
$Z_s \subset X_s$ が有効 Cartier 因子ならば、$Z$ は相対有効 Cartier 因子である。
同様に、$Z$ が $X$ の局所単項閉部分スキームで、すべてのファイバー
$Z_s \subset X_s$ が有効 Cartier 因子ならば、$Z$ は相対有効 Cartier 因子である。
\end{lemma}

\begin{proof}
$\Spec(A)$ を $s$ のアフィン開近傍、$\Spec(B)$ を $x$ のアフィン開近傍として、
$\varphi(\Spec(B)) \subset \Spec(A)$ となるように選ぶ。素イデアル
$\mathfrak p \subset A$ に $s$ が対応するとする。素イデアル
$\mathfrak q \subset B$ に $x$ が対応するとし、イデアル $I \subset B$ を
$Z$ に対応するものとする。
補題の最初の仮定により、$A \to B$ は平坦かつ有限表示である。(1) の仮定は、
$\Spec(B)$ を縮小した後、$I(B \otimes_A \kappa(\mathfrak p))$ が
$B \otimes_A \kappa(\mathfrak p)$ の非零因子である一つの元により生成されると仮定
できることを意味する。$f \in I$ がこの生成元へ写るとする。ある元 $g \in B$ で
$g \not \in \mathfrak q$ を満たすものを可逆にした後、閉部分スキーム
$D = V(f) \subset \Spec(B_g)$ が相対有効 Cartier 因子であると主張する。

\medskip\noindent
Algebra, Lemma \ref{algebra-lemma-flat-finite-presentation-limit-flat} により、
Noether 環の平坦有限型環写像 $A_0 \to B_0$、元 $f_0 \in B_0$、環写像
$A_0 \to A$、および同型 $A \otimes_{A_0} B_0 \cong B$ を見つけられる。
$\mathfrak p_0 = A_0 \cap \mathfrak p$ ならば
$$
B \otimes_A \kappa(\mathfrak p) =
\left(B_0 \otimes_{A_0} \kappa(\mathfrak p_0)\right)
\otimes_{\kappa(\mathfrak p_0))} \kappa(\mathfrak p)
$$
である。したがって $f_0$ は $B_0 \otimes_{A_0} \kappa(\mathfrak p_0)$ の非零因子である。
Algebra, Lemma \ref{algebra-lemma-grothendieck} により、$f_0$ は
$(B_0)_{\mathfrak q_0}$ の非零因子である。ここで
$\mathfrak q_0 = B_0 \cap \mathfrak q$ であり、さらに
$(B_0/f_0B_0)_{\mathfrak q_0}$ は $A_0$ 上平坦である。したがって Algebra, Lemma
\ref{algebra-lemma-regular-sequence-in-neighbourhood} および Algebra, Theorem
\ref{algebra-theorem-openness-flatness} により、
$g_0 \in B_0$ かつ $g_0 \not \in \mathfrak q_0$ となるものが存在し、
$f_0$ は $(B_0)_{g_0}$ の非零因子で、
$(B_0/f_0B_0)_{g_0}$ は $A_0$ 上平坦となる。ゆえに
$D_0 = V(f_0) \subset \Spec((B_0)_{g_0})$ は相対有効 Cartier 因子である。
この性質は基底変換で保たれることが Lemma \ref{divisors-lemma-relative-Cartier} から分かるので、
上の主張を得る。このとき $g$ は $g_0$ の $B$ における像とする。

\medskip\noindent
ここまでで (1) を証明した。(2) を見るため閉埋入 $Z \to D$ を考える。全射環写像
$u : \mathcal{O}_{D, x} \to \mathcal{O}_{Z, x}$ は、平坦な局所
$\mathcal{O}_{S, s}$-代数の間の写像である。これらは本質的有限表示であり、
$\mathfrak m_s$ で割った後に同型となる。したがってこの写像は同型である。Algebra,
Lemma \ref{algebra-lemma-mod-injective-general} を参照。ゆえに $Z \to D$ は
$x$ の近傍で同型である。Algebra, Lemma \ref{algebra-lemma-local-isomorphism} を参照。
(3) を見るため、必要なら $U$ を縮小し、$D$ のイデアルが一つの非零因子 $f$ により
生成され、$Z$ のイデアルが元 $g$ により生成されると仮定してよい。このとき
$f = gh$ である。しかし $g|_{U_s}$ と $f|_{U_s}$ は $x$ の近傍で同じ有効 Cartier
因子を切り出す。したがって $h|_{X_s}$ は $\mathcal{O}_{X_s, x}$ の単元であり、
$h$ は $\mathcal{O}_{X, x}$ の単元であり、したがって $h$ は $x$ のある開近傍で
単元である。すなわち $Z \cap U = D$ となるように $U$ を縮小できる。

\medskip\noindent
補題の最後の主張は、(2) と (3)、および $Z \to S$ が局所有限表示であることと
$Z \to X$ が有限表示であることが同値であるという事実を組み合わせれば直ちに従う。
Morphisms, Lemmas \ref{morphisms-lemma-composition-finite-presentation} および
\ref{morphisms-lemma-finite-presentation-permanence} を参照。
\end{proof}

\section{埋め込みの法錐}
\label{divisors-section-normal-cone}

\noindent
$i : Z \to X$ を閉埋め込みとし、$\mathcal{I} \subset \mathcal{O}_X$ を対応する
準連接イデアル層とする。次数付き $\mathcal{O}_X$-代数の準連接層
$\bigoplus_{n \geq 0} \mathcal{I}^n/\mathcal{I}^{n + 1}$ を考える。各層
$\mathcal{I}^n/\mathcal{I}^{n + 1}$ は $\mathcal{I}$ により零化されるので、この
次数付き代数は Morphisms, Lemma \ref{morphisms-lemma-i-star-equivalence} により、
次数付き $\mathcal{O}_Z$-代数の準連接層に対応する。この次数付き
$\mathcal{O}_Z$-代数の準連接層を{\it $Z$ の $X$ における余法代数}と呼び、
Morphisms, Section \ref{morphisms-section-closed-immersions-quasi-coherent} で述べた
記法の濫用により、しばしば単に
$\bigoplus_{n \geq 0} \mathcal{I}^n/\mathcal{I}^{n + 1}$ と書く。

\medskip\noindent
$f : Z \to X$ を埋め込みとする。$f$ の余法代数を、閉埋め込み
$i : Z \to X \setminus \partial Z$ の余法代数として定義する。ここで
$\partial Z = \overline{Z} \setminus Z$ である。これはしばしば
$\bigoplus_{n \geq 0} \mathcal{I}^n/\mathcal{I}^{n + 1}$ と書かれる。ここで
$\mathcal{I}$ は閉埋め込み $i : Z \to X \setminus \partial Z$ のイデアル層である。

\begin{definition}
\label{divisors-definition-conormal-sheaf}
$f : Z \to X$ を埋め込みとする。{\it 余法代数
$\mathcal{C}_{Z/X, *}$（$Z$ の $X$ におけるもの）}、または{\it $f$ の余法代数}とは、
上で述べた次数付き $\mathcal{O}_Z$-代数の準連接層
$\bigoplus_{n \geq 0} \mathcal{I}^n/\mathcal{I}^{n + 1}$ をいう。
\end{definition}

\noindent
したがって $\mathcal{C}_{Z/X, 1} = \mathcal{C}_{Z/X}$ は埋め込みの余法層である。
また $\mathcal{C}_{Z/X, 0} = \mathcal{O}_Z$ であり、
$\mathcal{C}_{Z/X, n}$ は次の性質で特徴付けられる準連接
$\mathcal{O}_Z$-加群である。
\begin{equation}
\label{divisors-equation-conormal-in-degree-n}
i_*\mathcal{C}_{Z/X, n} = \mathcal{I}^n/\mathcal{I}^{n + 1}
\end{equation}
ここで $i : Z \to X \setminus \partial Z$ であり、$\mathcal{I}$ は上と同じく
$i$ のイデアル層である。最後に、次数付き準連接 $\mathcal{O}_Z$-代数の標準全射
\begin{equation}
\label{divisors-equation-conormal-algebra-quotient}
\text{Sym}^*(\mathcal{C}_{Z/X}) \longrightarrow \mathcal{C}_{Z/X, *}
\end{equation}
が存在し、次数 $0$ および $1$ では同型であることに注意する。

\begin{lemma}
\label{divisors-lemma-affine-conormal-sheaf}
$i : Z \to X$ を埋め込みとする。$i$ の余法代数は次の性質をもつ。
\begin{enumerate}
\item $U \subset X$ を、$i(Z)$ が $U$ の閉部分集合となる任意の開集合とする。
$\mathcal{I} \subset \mathcal{O}_U$ を、閉部分スキーム $i(Z) \subset U$ に対応する
イデアル層とする。このとき
$$
\mathcal{C}_{Z/X, *} =
i^*\left(\bigoplus\nolimits_{n \geq 0} \mathcal{I}^n\right) =
i^{-1}\left(
\bigoplus\nolimits_{n \geq 0} \mathcal{I}^n/\mathcal{I}^{n + 1}
\right)
$$
\item 任意のアフィン開集合 $\Spec(R) = U \subset X$ で
$Z \cap U = \Spec(R/I)$ を満たすものに対して、標準同型
$\Gamma(Z \cap U, \mathcal{C}_{Z/X, *}) = \bigoplus_{n \geq 0} I^n/I^{n + 1}$
が存在する。
\end{enumerate}
\end{lemma}

\begin{proof}
大部分は定義から明らかである。環 $R$ とイデアル $I$ を与え、後者が $R$ のイデアルならば、
$I^n/I^{n + 1} = I^n \otimes_R R/I$ であることに注意する。細部は省略する。
\end{proof}

\begin{lemma}
\label{divisors-lemma-conormal-algebra-functorial}
スキームの圏における可換図式
$$
\xymatrix{
Z \ar[r]_i \ar[d]_f & X \ar[d]^g \\
Z' \ar[r]^{i'} & X'
}
$$
を考える。$i$ と $i'$ は埋め込みであると仮定する。このとき次数付き
$\mathcal{O}_Z$-代数の標準写像
$$
f^*\mathcal{C}_{Z'/X', *}
\longrightarrow
\mathcal{C}_{Z/X, *}
$$
が存在し、次の性質で特徴付けられる。アフィン開集合の任意の組
$(\Spec(R) = U \subset X, \Spec(R') = U' \subset X')$ で
$g(U) \subset U'$ を満たし、かつ
$Z \cap U = \Spec(R/I)$ および $Z' \cap U' = \Spec(R'/I')$ となるものに対し、
誘導写像
$$
\Gamma(Z' \cap U', \mathcal{C}_{Z'/X', *}) =
\bigoplus\nolimits (I')^n/(I')^{n + 1}
\longrightarrow
\bigoplus\nolimits_{n \geq 0} I^n/I^{n + 1} =
\Gamma(Z \cap U, \mathcal{C}_{Z/X, *})
$$
は環写像 $f^\sharp : R' \to R$ から誘導される写像である。この環写像は
$f^\sharp(I') \subset I$ を満たす。
\end{lemma}

\begin{proof}
$\partial Z' = \overline{Z'} \setminus Z'$ および
$\partial Z = \overline{Z} \setminus Z$ とおく。これらはそれぞれ $X'$ と $X$ の
閉部分集合である。$X'$ を $X' \setminus \partial Z'$ で、$X$ を
$X \setminus \big(g^{-1}(\partial Z') \cup \partial Z\big)$ で置き換えると、
$i$ と $i'$ は閉埋め込みであると仮定してよい。

\medskip\noindent
$g \circ i$ が $i'$ を経由して分解することから、$g^*\mathcal{I}'$ は
$\mathcal{I}$ の中へ写る。これは標準写像 $g^*\mathcal{I}' \to \mathcal{O}_X$ のもとでの
主張である。Schemes, Lemmas
\ref{schemes-lemma-characterize-closed-subspace} および
\ref{schemes-lemma-restrict-map-to-closed} を参照。したがって準連接層の誘導写像
$g^*((\mathcal{I}')^n/(\mathcal{I}')^{n + 1}) \to
\mathcal{I}^n/\mathcal{I}^{n + 1}$ を得る。$i$ で引き戻すと
$i^*g^*((\mathcal{I}')^n/(\mathcal{I}')^{n + 1}) \to
i^*(\mathcal{I}^n/\mathcal{I}^{n + 1})$ を得る。
$i^*(\mathcal{I}^n/\mathcal{I}^{n + 1}) = \mathcal{C}_{Z/X, n}$ であることに注意する。
一方、
$i^*g^*((\mathcal{I}')^n/(\mathcal{I}')^{n + 1}) =
f^*(i')^*((\mathcal{I}')^n/(\mathcal{I}')^{n + 1}) =
f^*\mathcal{C}_{Z'/X', n}$ である。これが所望の写像を与える。

\medskip\noindent
この写像が局所的に与えられた写像
$(I')^n/(I')^{n + 1} \to I^n/I^{n + 1}$ で記述されることの確認は、定義を展開
するだけであり、省略する。もう一つ注意しておくと、任意の $x \in i(Z)$ に対して、
アフィン開近傍 $U$ と $U'$ で、$f(U) \subset U'$ を満たし、
$Z \cap U$ および $U' \cap Z'$ が閉であり、かつ $x \in U$ となるものが実際に
存在する。証明は省略する。したがって補題の条件は確かにこの写像を特徴付け
（また、この条件を用いて写像を定義することもできた）。
\end{proof}

\begin{lemma}
\label{divisors-lemma-conormal-algebra-functorial-flat}
スキームの圏におけるファイバー積図式
$$
\xymatrix{
Z \ar[r]_i \ar[d]_f & X \ar[d]^g \\
Z' \ar[r]^{i'} & X'
}
$$
を考え、$i$ と $i'$ は埋め込みであるとする。このとき Lemma
\ref{divisors-lemma-conormal-algebra-functorial} の標準写像
$f^*\mathcal{C}_{Z'/X', *} \to \mathcal{C}_{Z/X, *}$ は全射である。
$g$ が平坦ならば、これは同型である。
\end{lemma}

\begin{proof}
$R' \to R$ を環写像とし、$I' \subset R'$ をイデアルとする。$I = I'R$ とおく。
このとき $(I')^n/(I')^{n + 1} \otimes_{R'} R \to I^n/I^{n + 1}$ は全射である。
$R' \to R$ が平坦ならば、$I^n = (I')^n \otimes_{R'} R$ であり、この写像は同型である。
\end{proof}

\begin{definition}
\label{divisors-definition-normal-cone}
$i : Z \to X$ をスキームの埋め込みとする。{\it 法錐 $C_ZX$（$Z$ の $X$ におけるもの）}とは
$$
C_ZX = \underline{\Spec}_Z(\mathcal{C}_{Z/X, *})
$$
をいう。Constructions, Definitions \ref{constructions-definition-cone} および
\ref{constructions-definition-abstract-cone} を参照。$Z$ の $X$ における
{\it 法束}とは、ベクトル束
$$
N_ZX = \underline{\Spec}_Z(\text{Sym}(\mathcal{C}_{Z/X}))
$$
をいう。Constructions, Definitions \ref{constructions-definition-vector-bundle} および
\ref{constructions-definition-abstract-vector-bundle} を参照。
\end{definition}

\noindent
したがって $C_ZX \to Z$ は $Z$ 上の錐であり、$N_ZX \to Z$ は $Z$ 上のベクトル束
である（本書の用語法では、これは余法層が有限局所自由層であることを意味しない）。
さらに、次数付き代数の標準全射 (\ref{divisors-equation-conormal-algebra-quotient}) は
$Z$ 上の錐の標準閉埋め込み
\begin{equation}
\label{divisors-equation-normal-cone-in-normal-bundle}
C_ZX \longrightarrow N_ZX
\end{equation}
を定める。

\section{正則イデアル層}
\label{divisors-section-regular-ideal-sheaves}

\noindent
本節では、有効 Cartier 因子の概念をより高い余次元へ一般化する。列
$f_1, \ldots, f_r$ が環 $R$ の元からなるとする。これが{\it 正則列}であるとは、
各 $i = 1, \ldots, r$ に対して
$f_i$ が $R/(f_1, \ldots, f_{i - 1})$ 上の非零因子であり、かつ
$R/(f_1, \ldots, f_r) \not = 0$ となることだった。Algebra, Definition
\ref{algebra-definition-regular-sequence} を参照。これに近い三つの弱い条件を課すことも
できる。第一は、$f_1, \ldots, f_r$ が{\it Koszul 正則列}、すなわち
$H_i(K_\bullet(f_1, \ldots, f_r)) = 0$ が $i > 0$ に対して成り立つと仮定することである。
More on Algebra, Definition
\ref{more-algebra-definition-koszul-regular-sequence} を参照。列が
{\it $H_1$-正則列}であるとは、$H_1(K_\bullet(f_1, \ldots, f_r)) = 0$ となることをいう。
もう一つ課せる条件は、$J = (f_1, \ldots, f_r)$ とおいたとき、写像
$$
R/J[T_1, \ldots, T_r]
\longrightarrow
\bigoplus\nolimits_{n \geq 0}
J^n/J^{n + 1}
$$
が同型となることである。この写像は $T_i$ を $f_i \bmod J^2$ へ送る。この場合、
$f_1, \ldots, f_r$ を{\it 準正則列}という。Algebra, Definition
\ref{algebra-definition-quasi-regular-sequence} を参照。$R$-加群 $M$ を与えたときには、
$M$-正則列および $M$-準正則列という概念もある。

\medskip\noindent
これは次のように環付き空間の場合へ一般化できる。$X$ を環付き空間とし、
$f_1, \ldots, f_r \in \Gamma(X, \mathcal{O}_X)$ とする。$f_1, \ldots, f_r$ が
{\it 正則列}であるとは、各 $i = 1, \ldots, r$ に対して写像
\begin{equation}
\label{divisors-equation-map-regular}
f_i :
\mathcal{O}_X/(f_1, \ldots, f_{i - 1})
\longrightarrow
\mathcal{O}_X/(f_1, \ldots, f_{i - 1})
\end{equation}
が層の単射となることをいう。$f_1, \ldots, f_r$ が{\it Koszul 正則列}であるとは、
Koszul 複体
\begin{equation}
\label{divisors-equation-koszul}
K_\bullet(\mathcal{O}_X, f_\bullet),
\end{equation}
が次数 $> 0$ で非輪状となることをいう。Modules, Definition
\ref{modules-definition-koszul-complex} を参照。$f_1, \ldots, f_r$ が
{\it $H_1$-正則列}であるとは、Koszul 複体
$K_\bullet(\mathcal{O}_X, f_\bullet)$ が次数 $1$ で完全となることをいう。最後に、
$f_1, \ldots, f_r$ が{\it 準正則列}であるとは、写像
\begin{equation}
\label{divisors-equation-map-quasi-regular}
\mathcal{O}_X/\mathcal{J}[T_1, \ldots, T_r]
\longrightarrow
\bigoplus\nolimits_{d \geq 0} \mathcal{J}^d/\mathcal{J}^{d + 1}
\end{equation}
が層の同型となることをいう。ここで $\mathcal{J} \subset \mathcal{O}_X$ は
$f_1, \ldots, f_r$ により生成されるイデアル層である。（$\mathcal{F}$-正則列および
$\mathcal{F}$-準正則列という概念もある。これは与えられた
$\mathcal{O}_X$-加群 $\mathcal{F}$ に対するもので、必要になればここで導入する。）

\begin{lemma}
\label{divisors-lemma-types-regular-sequences-implications}
$X$ を環付き空間とし、
$f_1, \ldots, f_r \in \Gamma(X, \mathcal{O}_X)$ とする。このとき次の含意が成り立つ。
$f_1, \ldots, f_r$ が正則列 $\Rightarrow$
$f_1, \ldots, f_r$ が Koszul 正則列 $\Rightarrow$
$f_1, \ldots, f_r$ が $H_1$-正則列 $\Rightarrow$
$f_1, \ldots, f_r$ が準正則列。
\end{lemma}

\begin{proof}
完全性は茎で確認できるので、列 $f_1, \ldots, f_r$ が正則列であるための必要十分条件は、
写像
$$
f_i :
\mathcal{O}_{X, x}/(f_1, \ldots, f_{i - 1})
\longrightarrow
\mathcal{O}_{X, x}/(f_1, \ldots, f_{i - 1})
$$
がすべての $x \in X$ に対して単射となることである。他の種類の正則性も茎ごとに
確認できる（細部は省略する）。したがって含意は More on Algebra, Lemmas
\ref{more-algebra-lemma-regular-koszul-regular},
\ref{more-algebra-lemma-koszul-regular-H1-regular}, および
\ref{more-algebra-lemma-H1-regular-quasi-regular} から従う。
\end{proof}

\begin{definition}
\label{divisors-definition-regular-ideal-sheaf}
\begin{reference}
Koszul 正則イデアル層の概念は \cite[Expose VII, Definition 1.4]{SGA6} で導入され、
そこでは正則イデアル層と呼ばれた。
\end{reference}
$X$ を環付き空間とし、$\mathcal{J} \subset \mathcal{O}_X$ をイデアル層とする。
\begin{enumerate}
\item $\mathcal{J}$ が{\it 正則}であるとは、すべての
$x \in \text{Supp}(\mathcal{O}_X/\mathcal{J})$ に対して、開近傍
$x \in U \subset X$ と正則列 $f_1, \ldots, f_r \in \mathcal{O}_X(U)$ が存在し、
$\mathcal{J}|_U$ が $f_1, \ldots, f_r$ により生成されることをいう。
\item $\mathcal{J}$ が{\it Koszul 正則}であるとは、すべての
$x \in \text{Supp}(\mathcal{O}_X/\mathcal{J})$ に対して、開近傍
$x \in U \subset X$ と Koszul 正則列 $f_1, \ldots, f_r \in \mathcal{O}_X(U)$ が存在し、
$\mathcal{J}|_U$ が $f_1, \ldots, f_r$ により生成されることをいう。
\item $\mathcal{J}$ が{\it $H_1$-正則}であるとは、すべての
$x \in \text{Supp}(\mathcal{O}_X/\mathcal{J})$ に対して、開近傍
$x \in U \subset X$ と $H_1$-正則列 $f_1, \ldots, f_r \in \mathcal{O}_X(U)$ が存在し、
$\mathcal{J}|_U$ が $f_1, \ldots, f_r$ により生成されることをいう。
\item $\mathcal{J}$ が{\it 準正則}であるとは、すべての
$x \in \text{Supp}(\mathcal{O}_X/\mathcal{J})$ に対して、開近傍
$x \in U \subset X$ と準正則列 $f_1, \ldots, f_r \in \mathcal{O}_X(U)$ が存在し、
$\mathcal{J}|_U$ が $f_1, \ldots, f_r$ により生成されることをいう。
\end{enumerate}
\end{definition}

\noindent
この概念の多くの性質は、環における正則列と準正則列についての対応する概念から
直ちに従う。

\begin{lemma}
\label{divisors-lemma-regular-quasi-regular-scheme}
$X$ を環付き空間とし、$\mathcal{J}$ をイデアル層とする。このとき次の含意が成り立つ。
$\mathcal{J}$ が正則 $\Rightarrow$
$\mathcal{J}$ が Koszul 正則 $\Rightarrow$
$\mathcal{J}$ が $H_1$-正則 $\Rightarrow$
$\mathcal{J}$ が準正則。
\end{lemma}

\begin{proof}
この補題は直ちに Lemma \ref{divisors-lemma-types-regular-sequences-implications} に帰着する。
\end{proof}

\begin{lemma}
\label{divisors-lemma-quasi-regular-ideal}
$X$ を局所環付き空間とし、$\mathcal{J} \subset \mathcal{O}_X$ をイデアル層とする。
このとき $\mathcal{J}$ が準正則であるための必要十分条件は、次の条件が成り立つことである。
\begin{enumerate}
\item $\mathcal{J}$ は有限型 $\mathcal{O}_X$-加群である。
\item $\mathcal{J}/\mathcal{J}^2$ は有限局所自由
$\mathcal{O}_X/\mathcal{J}$-加群である。
\item 標準写像
$$
\text{Sym}^n_{\mathcal{O}_X/\mathcal{J}}(\mathcal{J}/\mathcal{J}^2)
\longrightarrow
\mathcal{J}^n/\mathcal{J}^{n + 1}
$$
はすべての $n \geq 0$ に対して同型である。
\end{enumerate}
\end{lemma}

\begin{proof}
$U \subset X$ を、$\mathcal{J}|_U$ が準正則列
$f_1, \ldots, f_r \in \mathcal{O}_X(U)$ により生成されるような開集合とすると、
$\mathcal{J}|_U$ は有限型で、$\mathcal{J}|_U/\mathcal{J}^2|_U$ は
$f_1, \ldots, f_r$ を基底とする自由加群であり、(3) の写像は同型である。
これは、それらが (\ref{divisors-equation-map-quasi-regular}) の次数 $n$ の部分を座標によらず
表したものだからである。したがって、準正則ならば条件 (1)、(2)、(3) が成り立つ。

\medskip\noindent
逆に、(1)、(2)、(3) が成り立つと仮定する。点
$x \in \text{Supp}(\mathcal{O}_X/\mathcal{J})$ を取る。このとき近傍
$U \subset X$ で $x$ を含み、$\mathcal{J}|_U/\mathcal{J}^2|_U$ が階数
$r$ の自由 $\mathcal{O}_U/\mathcal{J}|_U$-加群となるものが存在する。
$U$ を必要なら縮小し、基底へ写る元
$f_1, \ldots, f_r \in \mathcal{J}(U)$ が存在すると仮定してよい。この基底とは、
$\mathcal{J}|_U/\mathcal{J}^2|_U$ を $\mathcal{O}_U/\mathcal{J}|_U$-加群とみなした
ときの基底である。とくに、$f_1, \ldots, f_r$ の
$\mathcal{J}_x/\mathcal{J}^2_x$ における像はこれを生成する。したがって中山の補題
（Algebra, Lemma \ref{algebra-lemma-NAK}）により、$f_1, \ldots, f_r$ は茎
$\mathcal{J}_x$ を生成する。$\mathcal{J}$ は有限型なので、Modules, Lemma
\ref{modules-lemma-finite-type-surjective-on-stalk} により、$U$ を縮小した後、
$f_1, \ldots, f_r$ が $\mathcal{J}$ を生成すると仮定してよい。最後に (3) と同型
$\mathcal{J}|_U/\mathcal{J}^2|_U = \bigoplus \mathcal{O}_U/\mathcal{J}|_U f_i$
から、$f_1, \ldots, f_r \in \mathcal{O}_X(U)$ が準正則列であることは明らかである。
\end{proof}

\begin{lemma}
\label{divisors-lemma-generate-regular-ideal}
$(X, \mathcal{O}_X)$ を局所環付き空間とし、
$\mathcal{J} \subset \mathcal{O}_X$ をイデアル層とする。$x \in X$ とし、
$f_1, \ldots, f_r \in \mathcal{J}_x$ の像が $\kappa(x)$-ベクトル空間
$\mathcal{J}_x/\mathfrak m_x\mathcal{J}_x$ の基底を与えるとする。
\begin{enumerate}
\item $\mathcal{J}$ が準正則ならば、開近傍が存在し、その上で
$f_1, \ldots, f_r \in \mathcal{O}_X(U)$ は $\mathcal{J}|_U$ を生成する準正則列となる。
\item $\mathcal{J}$ が $H_1$-正則ならば、開近傍が存在し、その上で
$f_1, \ldots, f_r \in \mathcal{O}_X(U)$ は $H_1$-正則列となり、
$\mathcal{J}|_U$ を生成する。
\item $\mathcal{J}$ が Koszul 正則ならば、開近傍が存在し、その上で
$f_1, \ldots, f_r \in \mathcal{O}_X(U)$ は $\mathcal{J}|_U$ を生成する
Koszul 正則列となる。
\end{enumerate}
\end{lemma}

\begin{proof}
まず $\mathcal{J}$ が準正則であると仮定する。開近傍 $U \subset X$ で $x$ を含む
ものと、準正則列 $g_1, \ldots, g_s \in \mathcal{O}_X(U)$ で
$\mathcal{J}|_U$ を生成するものを選べる。これは
$\mathcal{J}/\mathcal{J}^2$ が階数 $s$ の自由 $\mathcal{O}_U/\mathcal{J}|_U$-加群
であることを意味する（Lemma \ref{divisors-lemma-quasi-regular-ideal} とその証明を参照）。
したがって $r = s$ である。$U$ を縮小し、
$f_1, \ldots, f_r \in \mathcal{J}(U)$ と仮定してよい。ゆえに
$$
f_i = \sum a_{ij} g_j
$$
と書ける。ここで $a_{ij} \in \mathcal{O}_X(U)$ である。仮定により行列
$A = (a_{ij})$ は $\kappa(x)$ 上の可逆行列へ写る。したがって $U$ をもう一度縮小し、
$(a_{ij})$ は可逆であると仮定してよい。ゆえに $f_1, \ldots, f_r$ は
$(\mathcal{J}/\mathcal{J}^2)|_U$ の基底を与え、このことから
$f_1, \ldots, f_r$ が $U$ 上の準正則列であると分かる。

\medskip\noindent
(2) と (3) の仮定は (1) の仮定より強いので、これらを証明する際には、すでに
$f_1, \ldots, f_r \in \mathcal{J}(U)$ かつ $f_i = \sum a_{ij}g_j$ であり、
上と同じく $(a_{ij})$ は可逆であると仮定してよい。ここで今度は
$g_1, \ldots, g_r$ が $H_1$-正則列または Koszul 正則列である。
$f_1, \ldots, f_r$ 上の Koszul 複体は、$g_1, \ldots, g_r$ 上の Koszul 複体と
行列 $(a_{ij})$ を介して同型である（More on Algebra, Lemma
\ref{more-algebra-lemma-change-basis} を参照）。したがって
$f_1, \ldots, f_r$ は所望のとおり $H_1$-正則または Koszul 正則である。
\end{proof}

\begin{lemma}
\label{divisors-lemma-regular-ideal-sheaf-quasi-coherent}
スキーム上の正則、Koszul 正則、$H_1$-正則、または準正則なイデアル層は、いずれも
有限型準連接イデアル層である。
\end{lemma}

\begin{proof}
このようなイデアル層は局所的に有限個の切断で生成されるので、主張は従う。
また、スキーム上で切断により局所的に生成される任意のイデアル層は準連接である。
Schemes, Lemma \ref{schemes-lemma-closed-subspace-scheme} を参照。
\end{proof}

\begin{lemma}
\label{divisors-lemma-regular-ideal-sheaf-scheme}
$X$ をスキームとし、$\mathcal{J}$ をイデアル層とする。このとき
$\mathcal{J}$ が正則（それぞれ Koszul 正則、$H_1$-正則、準正則）であるための必要十分条件は、
すべての $x \in \text{Supp}(\mathcal{O}_X/\mathcal{J})$ に対してアフィン開近傍
$x \in U \subset X$、$U = \Spec(A)$ が存在し、$\mathcal{J}|_U = \widetilde{I}$ かつ
$I$ が環内の正則列（それぞれ Koszul 正則列、$H_1$-正則列、準正則列）
$f_1, \ldots, f_r \in A$ により生成されることである。
\end{lemma}

\begin{proof}
仮定により、開近傍 $U$ で $x$ を含み、その上で $\mathcal{J}$ が正則列（それぞれ
Koszul 正則列、$H_1$-正則列、準正則列）
$f_1, \ldots, f_r \in \mathcal{O}_X(U)$ により生成されるものを見つけられる。
$U$ を縮小し、$U$ はアフィン、たとえば $U = \Spec(A)$ であると仮定してよい。
Lemma \ref{divisors-lemma-regular-ideal-sheaf-quasi-coherent} により $\mathcal{J}$ は準連接なので、
$\mathcal{J}|_U = \widetilde{I}$ である。ここで $I \subset A$ はあるイデアルである。
ここで、完全性を保つ圏同値
$$
\widetilde{\ } : \text{Mod}_A \longrightarrow \QCoh(\mathcal{O}_U)
$$
を用いる。たとえば、函数 $f_i$ が $\mathcal{J}$ を生成することは、$f_i$ を
$A$ の元とみたとき、それらが $I$ を生成することを意味する。
(\ref{divisors-equation-map-regular}) が単射であること（それぞれ
(\ref{divisors-equation-koszul}) が完全、(\ref{divisors-equation-koszul}) が次数 $1$ で完全、
(\ref{divisors-equation-map-quasi-regular}) が同型であること）から、対応する写像
$A/(f_1, \ldots, f_{i - 1}) \to A/(f_1, \ldots, f_{i - 1})$
（それぞれ複体 $K_\bullet(A, f_1, \ldots, f_r)$、写像
$A/I[T_1, \ldots, T_r] \to \bigoplus I^n/I^{n + 1}$）の対応する性質が従う。
したがって $f_1, \ldots, f_r \in A$ は正則列（それぞれ Koszul 正則列、
$H_1$-正則列、準正則列）であり、これは環 $A$ の列である。
\end{proof}

\begin{lemma}
\label{divisors-lemma-Noetherian-scheme-regular-ideal}
$X$ を局所 Noether スキームとし、$\mathcal{J} \subset \mathcal{O}_X$ を準連接
イデアル層とする。$x$ を取り、これは $\mathcal{O}_X/\mathcal{J}$ の台の点とする。
次は同値である。
\begin{enumerate}
\item $\mathcal{J}_x$ は $\mathcal{O}_{X, x}$ の正則列により生成される。
\item $\mathcal{J}_x$ は $\mathcal{O}_{X, x}$ の Koszul 正則列により生成される。
\item $\mathcal{J}_x$ は $H_1$-正則列、すなわち $\mathcal{O}_{X, x}$ の列により生成される。
\item $\mathcal{J}_x$ は $\mathcal{O}_{X, x}$ の準正則列により生成される。
\item アフィン近傍 $U = \Spec(A)$ で $x$ を含み、$\mathcal{J}|_U = \widetilde{I}$ かつ
$I$ が $A$ の正則列により生成されるものが存在する。
\item アフィン近傍 $U = \Spec(A)$ で $x$ を含み、$\mathcal{J}|_U = \widetilde{I}$ かつ
$I$ が $A$ の Koszul 正則列により生成されるものが存在する。
\item アフィン近傍 $U = \Spec(A)$ で $x$ を含み、$\mathcal{J}|_U = \widetilde{I}$ かつ
$I$ が $H_1$-正則列、すなわち $A$ の列により生成されるものが存在する。
\item アフィン近傍 $U = \Spec(A)$ で $x$ を含み、$\mathcal{J}|_U = \widetilde{I}$ かつ
$I$ が $A$ の準正則列により生成されるものが存在する。
\item 近傍 $U$ で $x$ を含み、$\mathcal{J}|_U$ が正則となるものが存在する。
\item 近傍 $U$ で $x$ を含み、$\mathcal{J}|_U$ が Koszul 正則となるものが存在する。
\item 近傍 $U$ で $x$ を含み、$\mathcal{J}|_U$ が $H_1$-正則となるものが存在する。
\item 近傍 $U$ で $x$ を含み、$\mathcal{J}|_U$ が準正則となるものが存在する。
\end{enumerate}
とくに、局所 Noether スキーム上では、正則、Koszul 正則、$H_1$-正則、準正則の
各イデアル層の概念はすべて一致する。
\end{lemma}

\begin{proof}
Lemma \ref{divisors-lemma-regular-ideal-sheaf-scheme} から、
(5) $\Leftrightarrow$ (9)、(6) $\Leftrightarrow$ (10)、
(7) $\Leftrightarrow$ (11)、(8) $\Leftrightarrow$ (12) が従う。
(5) $\Rightarrow$ (1)、(6) $\Rightarrow$ (2)、
(7) $\Rightarrow$ (3)、(8) $\Rightarrow$ (4) は明らかである。
(1) $\Rightarrow$ (5) は Algebra, Lemma
\ref{algebra-lemma-regular-sequence-in-neighbourhood} による。
(9) $\Rightarrow$ (10) $\Rightarrow$ (11) $\Rightarrow$ (12) は
Lemma \ref{divisors-lemma-regular-quasi-regular-scheme} による。最後に、
(4) $\Rightarrow$ (1) は Algebra, Lemma
\ref{algebra-lemma-quasi-regular-regular} による。以上で十二の主張はすべて同値である。
\end{proof}

\section{正則埋め込み}
\label{divisors-section-regular-immersions}

\noindent
$i : Z \to X$ をスキームの埋め込みとする。定義により、ある開部分スキーム
$U \subset X$ が存在して、$Z$ は $U$ の閉部分スキームと同一視される。
$\mathcal{I} \subset \mathcal{O}_U$ を対応する準連接イデアル層とする。
$U' \subset X$ をこの性質をもつ別の開部分スキームとし、
$\mathcal{I}' \subset \mathcal{O}_{U'}$ を対応する準連接イデアル層と書く。このとき
$\mathcal{I}|_{U \cap U'} = \mathcal{I}'|_{U \cap U'}$ である。
さらに、$\mathcal{O}_U/\mathcal{I}$ の台は $Z$ であり、これは
$U \cap U'$ に含まれ、同時に $\mathcal{O}_{U'}/\mathcal{I}'$ の台でもある。
したがって Definition \ref{divisors-definition-regular-ideal-sheaf} より、
$\mathcal{I}$ が正則イデアルであることと $\mathcal{I}'$ が正則イデアルであることは
同値である。Koszul 正則、$H_1$-正則、または準正則であることについても同様である。

\begin{definition}
\label{divisors-definition-regular-immersion}
\begin{reference}
Koszul 正則埋め込みの概念は
\cite[Expose VII, Definition 1.4]{SGA6} で導入され、そこでは
正則埋め込みと呼ばれた。
\end{reference}
$i : Z \to X$ をスキームの埋め込みとする。開部分スキーム
$U \subset X$ を選び、$i$ が $Z$ を $U$ の閉部分スキームと同一視するものとし、
$\mathcal{I} \subset \mathcal{O}_U$ を対応する準連接イデアル層と書く。
\begin{enumerate}
\item $i$ を {\it 正則埋め込み}というのは、$\mathcal{I}$ が正則であるときである。
\item $i$ を {\it Koszul 正則埋め込み}というのは、
$\mathcal{I}$ が Koszul 正則であるときである。
\item $i$ を {\it $H_1$-正則埋め込み}というのは、
$\mathcal{I}$ が $H_1$-正則であるときである。
\item $i$ を {\it 準正則埋め込み}というのは、
$\mathcal{I}$ が準正則であるときである。
\end{enumerate}
\end{definition}

\noindent
上の議論により、これは $U$ の選択によらない。これらの条件は強いものから順に
並べてあり、Lemma \ref{divisors-lemma-regular-quasi-regular-immersion} を参照。
Koszul 正則閉埋め込みは平滑位相で局所的に正則埋め込みとなる。
Lemma \ref{divisors-lemma-koszul-regular-smooth-locally-regular} を参照。
局所 Noether の場合には四つの概念はすべて一致する。
Lemma \ref{divisors-lemma-Noetherian-scheme-regular-ideal} を参照。

\begin{lemma}
\label{divisors-lemma-regular-quasi-regular-immersion}
$i : Z \to X$ をスキームの埋め込みとする。このとき次の含意が成り立つ。
$i$ が正則 $\Rightarrow$
$i$ が Koszul 正則 $\Rightarrow$
$i$ が $H_1$-正則 $\Rightarrow$
$i$ が準正則。
\end{lemma}

\begin{proof}
主張は直ちに Lemma \ref{divisors-lemma-regular-quasi-regular-scheme} に帰着する。
\end{proof}

\begin{lemma}
\label{divisors-lemma-regular-immersion-noetherian}
$i : Z \to X$ をスキームの埋め込みとし、$X$ は局所 Noether であると仮定する。
このとき
$i$ が正則 $\Leftrightarrow$
$i$ が Koszul 正則 $\Leftrightarrow$
$i$ が $H_1$-正則 $\Leftrightarrow$
$i$ が準正則。
\end{lemma}

\begin{proof}
Lemma \ref{divisors-lemma-regular-quasi-regular-immersion} および
Lemma \ref{divisors-lemma-Noetherian-scheme-regular-ideal} から直ちに従う。
\end{proof}

\begin{lemma}
\label{divisors-lemma-flat-base-change-regular-immersion}
$i : Z \to X$ を正則（それぞれ Koszul 正則、$H_1$-正則、準正則）
埋め込みとする。$X' \to X$ を平坦射とする。このとき基底変換
$i' : Z \times_X X' \to X'$ は正則（それぞれ Koszul 正則、
$H_1$-正則、準正則）埋め込みである。
\end{lemma}

\begin{proof}
Lemma \ref{divisors-lemma-regular-ideal-sheaf-scheme} を介すると、これは
Algebra, Lemmas \ref{algebra-lemma-flat-increases-depth} および
\ref{algebra-lemma-flat-base-change-quasi-regular}
および More on Algebra,
Lemma \ref{more-algebra-lemma-koszul-regular-flat-base-change}
の代数的主張に移される。
\end{proof}

\begin{lemma}
\label{divisors-lemma-quasi-regular-immersion}
$i : Z \to X$ をスキームの埋め込みとする。このとき $i$ が準正則埋め込みであるための
必要十分条件は、次の条件が成り立つことである。
\begin{enumerate}
\item $i$ は局所有限表示である。
\item 余法層 $\mathcal{C}_{Z/X}$ は有限局所自由である。
\item 写像 (\ref{divisors-equation-conormal-algebra-quotient}) は同型である。
\end{enumerate}
\end{lemma}

\begin{proof}
開埋め込みは局所有限表示である。したがって、$X$ を開部分スキーム
$U \subset X$ で置き換えてよい。ただし $i$ は $Z$ を $U$ の閉部分スキームと
同一視するものとする。すなわち、$i$ は閉埋め込みであると仮定してよい。
$\mathcal{I} \subset \mathcal{O}_X$ を対応する準連接イデアル層とする。
Morphisms, Lemma \ref{morphisms-lemma-closed-immersion-finite-presentation}
を参照して思い出すと、$\mathcal{I}$ が有限型であることと $i$ が局所有限表示であることは
同値である。したがって同値性は Lemma \ref{divisors-lemma-quasi-regular-ideal} と
定義を書き下すことから従う。
\end{proof}

\begin{lemma}
\label{divisors-lemma-transitivity-conormal-quasi-regular}
$Z \to Y \to X$ をスキームの埋め込みとする。$Z \to Y$ は
$H_1$-正則であると仮定する。このとき Morphisms, Lemma
\ref{morphisms-lemma-transitivity-conormal} の標準列
$$
0 \to i^*\mathcal{C}_{Y/X} \to
\mathcal{C}_{Z/X} \to
\mathcal{C}_{Z/Y} \to 0
$$
は完全であり、局所的に分裂する。
\end{lemma}

\begin{proof}
$\mathcal{C}_{Z/Y}$ は有限局所自由なので（Lemma
\ref{divisors-lemma-quasi-regular-immersion} および
Lemma \ref{divisors-lemma-regular-quasi-regular-scheme} を参照）、列が完全であることを
示せば十分である。Morphisms, Lemma
\ref{morphisms-lemma-transitivity-conormal} で示されたことにより、
第一の写像が単射であることを示せば十分である。
アフィン局所で考えると、これは次の問題に帰着する。環 $A$ とイデアル
$I \subset J \subset A$ が与えられ、$J/I \subset A/I$ は
$H_1$-正則列により生成されると仮定する。このとき
$I/I^2 \otimes_A A/J \to J/J^2$ は単射か。
$I/I^2 \otimes_A A/J = I/IJ$ であることに注意せよ。したがって示そうとしているのは
$I \cap J^2 = IJ$ である。これは More on Algebra, Lemma
\ref{more-algebra-lemma-conormal-sequence-H1-regular} の結果である。
\end{proof}

\noindent
準正則埋め込みの合成は準正則とは限らない。Algebra, Remark
\ref{algebra-remark-join-quasi-regular-sequences} を参照。
それ以外の種類の正則埋め込みは合成で保たれる。

\begin{lemma}
\label{divisors-lemma-composition-regular-immersion}
$i : Z \to Y$ および $j : Y \to X$ をスキームの埋め込みとする。
\begin{enumerate}
\item $i$ と $j$ が正則埋め込みならば、$j \circ i$ もそうである。
\item $i$ と $j$ が Koszul 正則埋め込みならば、$j \circ i$ もそうである。
\item $i$ と $j$ が $H_1$-正則埋め込みならば、$j \circ i$ もそうである。
\item $i$ が $H_1$-正則埋め込みで $j$ が準正則埋め込みならば、
$j \circ i$ は準正則埋め込みである。
\end{enumerate}
\end{lemma}

\begin{proof}
(1) の代数的な版は Algebra, Lemma
\ref{algebra-lemma-join-regular-sequences} である。
(2) の代数的な版は More on Algebra, Lemma
\ref{more-algebra-lemma-join-koszul-regular-sequences} である。
(3) の代数的な版は More on Algebra, Lemma
\ref{more-algebra-lemma-join-H1-regular-sequences} である。
(4) の代数的な版は More on Algebra, Lemma
\ref{more-algebra-lemma-join-quasi-regular-H1-regular} である。
\end{proof}

\begin{lemma}
\label{divisors-lemma-permanence-regular-immersion}
$i : Z \to Y$ および $j : Y \to X$ をスキームの埋め込みとする。
$j$ は局所有限表示であり、Morphisms, Lemma
\ref{morphisms-lemma-transitivity-conormal} の列
$$
0 \to i^*\mathcal{C}_{Y/X} \to
\mathcal{C}_{Z/X} \to
\mathcal{C}_{Z/Y} \to 0
$$
は完全で局所的に分裂すると仮定する。
\begin{enumerate}
\item $j \circ i$ が準正則埋め込みならば、$i$ もそうである。
\item $j \circ i$ が $H_1$-正則埋め込みならば、$i$ もそうである。
\item $j$ と $j \circ i$ がともに Koszul 正則埋め込みならば、$i$ もそうである。
\end{enumerate}
\end{lemma}

\begin{proof}
$Y$ と $X$ を縮小することにより、$i$ と $j$ は閉埋め込みであると仮定してよい。
$\mathcal{I} \subset \mathcal{O}_X$ を $Y$ のイデアル層、
$\mathcal{J} \subset \mathcal{O}_X$ を $Z$ のイデアル層と書く。
余法列は $0 \to \mathcal{I}/\mathcal{I}\mathcal{J}
\to \mathcal{J}/\mathcal{J}^2 \to
\mathcal{J}/(\mathcal{I} + \mathcal{J}^2) \to 0$ である。
$z \in Z$ とし、$y = i(z)$、$x = j(y) = j(i(z))$ と置く。
$f_1, \ldots, f_n \in \mathcal{I}_x$ を選び、
$\mathcal{I}_x/\mathfrak m_z\mathcal{I}_x$ の基底へ写るものとする。これを
$f_1, \ldots, f_n, g_1, \ldots, g_m \in \mathcal{J}_x$ へ拡張し、
$\mathcal{J}_x/\mathfrak m_z\mathcal{J}_x$ の基底へ写るものとする。
これは可能である。実際、余法層の列は $z$ の近傍で分裂すると仮定したので、
$\mathcal{I}_x/\mathfrak m_x\mathcal{I}_x \to
\mathcal{J}_x/\mathfrak m_x\mathcal{J}_x$ は単射である。

\medskip\noindent
(1) の証明。Lemma \ref{divisors-lemma-generate-regular-ideal} により、アフィン開近傍
$U$ を $x$ について取り、そこで
$f_1, \ldots, f_n, g_1, \ldots, g_m$ が $\mathcal{J}$ を生成する準正則列を
なすようにできる。$i$ は局所有限表示なので、さらに
$\mathcal{I}$ は $f_1, \ldots, f_n$ により生成されると仮定してよい
（Morphisms, Lemma
\ref{morphisms-lemma-closed-immersion-finite-presentation}、Nakayama の補題、および
Modules, Lemma \ref{modules-lemma-finite-type-surjective-on-stalk} を用いる）。
したがって Algebra, Lemma
\ref{algebra-lemma-truncate-quasi-regular} により、
$g_1, \ldots, g_m$ は $Y \cap U$ 上で $Z$ を切り出す準正則列を誘導する。

\medskip\noindent
(2) の証明。More on Algebra, Lemma
\ref{more-algebra-lemma-truncate-H1-regular} を用いることを除けば
(1) の証明とまったく同じである。

\medskip\noindent
(3) の証明。Lemma \ref{divisors-lemma-generate-regular-ideal} を（二度）適用すると、
アフィン開近傍 $U$ を $x$ について取り、そこで
$f_1, \ldots, f_n$ が $\mathcal{I}$ を生成する Koszul 正則列をなし、かつ
$f_1, \ldots, f_n, g_1, \ldots, g_m$ が $\mathcal{J}$ を生成する
Koszul 正則列をなすようにできる。
したがって More on Algebra, Lemma
\ref{more-algebra-lemma-truncate-koszul-regular} により、
$g_1, \ldots, g_m$ は $Y \cap U$ 上で $Z$ を切り出す
Koszul 正則列を誘導する。
\end{proof}

\begin{lemma}
\label{divisors-lemma-extra-permanence-regular-immersion-noetherian}
$i : Z \to Y$ および $j : Y \to X$ をスキームの埋め込みとする。
$z \in Z$ を取り、対応する点を $y \in Y$、$x \in X$ と書く。
$X$ は局所 Noether であると仮定する。次は同値である。
\begin{enumerate}
\item $i$ は $z$ の近傍で正則埋め込みであり、$j$ は
$y$ の近傍で正則埋め込みである。
\item $i$ と $j \circ i$ は $z$ の近傍で正則埋め込みである。
\item $j \circ i$ は $z$ の近傍で正則埋め込みであり、余法列
$$
0 \to i^*\mathcal{C}_{Y/X} \to
\mathcal{C}_{Z/X} \to
\mathcal{C}_{Z/Y} \to 0
$$
は $z$ の近傍で分裂完全である。
\end{enumerate}
\end{lemma}

\begin{proof}
$X$（したがって $Y$）は局所 Noether なので、四種類の正則埋め込みはすべて一致する。
さらに、ある射が正則埋め込みであるかどうかは局所環のレベルで判定できる。
Lemma \ref{divisors-lemma-Noetherian-scheme-regular-ideal} を参照。
(1) $\Rightarrow$ (2) は Lemma
\ref{divisors-lemma-composition-regular-immersion} である。
(2) $\Rightarrow$ (3) は Lemma
\ref{divisors-lemma-transitivity-conormal-quasi-regular} である。
したがって (3) が (1) を含意することを示せば十分である。

\medskip\noindent
(3) を仮定する。$A = \mathcal{O}_{X, x}$ と置く。
$I \subset A$ を全射 $\mathcal{O}_{X, x} \to \mathcal{O}_{Y, y}$ の核と書き、
$J \subset A$ を全射 $\mathcal{O}_{X, x} \to \mathcal{O}_{Z, z}$ の核と書く。
$J$ を $A$ において生成する任意の極小生成列は準正則列、
したがって正則列であることに注意せよ。Lemma
\ref{divisors-lemma-generate-regular-ideal} を参照。仮定により、余法列
$$
0 \to I/IJ \to J/J^2 \to J/(I + J^2) \to 0
$$
は $A/J$-加群の列として分裂完全である。したがって、極小生成系
$f_1, \ldots, f_n, g_1, \ldots, g_m$ を $J$ について選べて、
$f_1, \ldots, f_n \in I$ が $I$ の極小生成系であるものを選べる。
上で指摘したように、$f_1, \ldots, f_n, g_1, \ldots, g_m$ は
$A$ の正則列である。正則列の定義から直接、
$f_1, \ldots, f_n$ は $A$ の正則列であり、
$\overline{g}_1, \ldots, \overline{g}_m$ は $A/I$ の正則列である。
したがって $j$ は $y$ における正則埋め込みであり、$i$ は
$z$ における正則埋め込みである。
\end{proof}

\begin{remark}
\label{divisors-remark-not-always-extra-permanence}
Lemma \ref{divisors-lemma-extra-permanence-regular-immersion-noetherian} の状況で、
(1)、(2)、(3) は「$j \circ i$ と $j$ がそれぞれ $z$ と $y$ における正則埋め込みである」
こととは {\bf 同値ではない}。例として
$X = \mathbf{A}^1_k = \Spec(k[x])$、
$Y = \Spec(k[x]/(x^2))$、$Z = \Spec(k[x]/(x))$ がある。
\end{remark}

\begin{lemma}
\label{divisors-lemma-koszul-regular-smooth-locally-regular}
$i : Z \to X$ を Koszul 正則閉埋め込みとする。このとき、
全射平滑射 $X' \to X$ が存在し、そのもとでの基底変換
$i' : Z \times_X X' \to X'$（$i$ の基底変換）は正則埋め込みとなる。
\end{lemma}

\begin{proof}
$X$ はアフィンであり、$Z$ のイデアルは Koszul 正則列により生成されると仮定してよい。
これは $X$ を適当なアフィン開被覆の各要素で置き換えることによる
（Lemma \ref{divisors-lemma-regular-ideal-sheaf-scheme} にあるようなアフィン開集合）。
アフィンの場合は More on Algebra, Lemma
\ref{more-algebra-lemma-Koszul-regular-flat-locally-regular} である。
\end{proof}

\begin{lemma}
\label{divisors-lemma-immersion-regular-regular-immersion}
$i : Z \to X$ を埋め込みとする。$Z$ と $X$ が正則スキームならば、
$i$ は正則埋め込みである。
\end{lemma}

\begin{proof}
$z \in Z$ とする。Lemma \ref{divisors-lemma-Noetherian-scheme-regular-ideal} により、
$\mathcal{O}_{X, z} \to \mathcal{O}_{Z, z}$ の核が正則列により生成されることを
示せば十分である。これは Algebra, Lemmas
\ref{algebra-lemma-regular-quotient-regular} および
\ref{algebra-lemma-regular-ring-CM} から従う。
\end{proof}

\section{相対正則埋め込み}
\label{divisors-section-relative-regular-immersion}

\noindent
この節では正則埋め込みの基底変換に関する性質を考える。次の補題は正則埋め込み、
あるいは Koszul 埋め込みについては成り立たない。Examples, Lemma
\ref{examples-lemma-base-change-regular-sequence} を参照。

\begin{lemma}
\label{divisors-lemma-relative-regular-immersion}
$f : X \to S$ をスキームの射とし、$i : Z \subset X$ を埋め込みとする。
次を仮定する。
\begin{enumerate}
\item $i$ は $H_1$-正則（それぞれ準正則）埋め込みである。
\item $Z \to S$ は平坦射である。
\end{enumerate}
このとき、任意のスキームの射 $g : S' \to S$ に対し、基底変換
$Z' = S' \times_S Z \to X' = S' \times_S X$ は
$H_1$-正則（それぞれ準正則）埋め込みである。
\end{lemma}

\begin{proof}
定義を書き下し Lemma \ref{divisors-lemma-regular-ideal-sheaf-scheme} を用いると、
これは More on Algebra, Lemma
\ref{more-algebra-lemma-relative-regular-immersion-algebra}
の主張に移される。
\end{proof}

\noindent
この補題が次の定義の動機である。

\begin{definition}
\label{divisors-definition-relative-H1-regular-immersion}
$f : X \to S$ をスキームの射とし、$i : Z \to X$ を埋め込みとする。
\begin{enumerate}
\item $i$ を {\it 相対準正則埋め込み}というのは、$Z \to S$ が平坦であり、
$i$ が準正則埋め込みであるときである。
\item $i$ を {\it 相対 $H_1$-正則埋め込み}というのは、$Z \to S$ が平坦であり、
$i$ が $H_1$-正則埋め込みであるときである。
\end{enumerate}
\end{definition}

\noindent
これは標準的でない用語法かもしれないことを注意しておく。
Lemma \ref{divisors-lemma-relative-regular-immersion} により、相対準正則（それぞれ
$H_1$-正則）埋め込みは任意の基底変換で保たれる。相対 $H_1$-正則埋め込みは
相対準正則埋め込みである。Lemma
\ref{divisors-lemma-regular-quasi-regular-immersion} を参照。
Lemma \ref{divisors-lemma-flat-relative-H1-regular}（または Lemma
\ref{divisors-lemma-relative-regular-immersion-flat-in-neighbourhood}）も参照されたい。
そこでは、$Z \to X$ が相対 $H_1$-正則（または準正則）埋め込みであり、
周囲のスキームが $S$ 上（平坦かつ）局所有限表示ならば、$Z \to X$ は実際に
正則埋め込みであり、任意の基底変換後にも同じことが成り立つと示されている。

\begin{lemma}
\label{divisors-lemma-quasi-regular-immersion-flat-at-x}
$f : X \to S$ をスキームの射とし、$Z \to X$ を相対準正則埋め込みとする。
$x \in Z$ かつ $\mathcal{O}_{X, x}$ が Noether ならば、$f$ は $x$ で平坦である。
\end{lemma}

\begin{proof}
$f_1, \ldots, f_r \in \mathcal{O}_{X, x}$ を、$Z$ の $x$ におけるイデアルを
切り出す準正則列とする。Algebra, Lemma
\ref{algebra-lemma-quasi-regular-regular} により、
$f_1, \ldots, f_r$ は正則列である。したがって $f_r$ は
$\mathcal{O}_{X, x}/(f_1, \ldots, f_{r - 1})$ 上の非零因子であり、
その商は平坦な $\mathcal{O}_{S, f(x)}$-加群である。
Lemma \ref{divisors-lemma-flat-at-x} により、
$\mathcal{O}_{X, x}/(f_1, \ldots, f_{r - 1})$ は平坦な
$\mathcal{O}_{S, f(x)}$-加群である。帰納的に続ければ、
$\mathcal{O}_{X, x}$ は平坦な $\mathcal{O}_{S, s}$-加群であることが分かる。
\end{proof}

\begin{lemma}
\label{divisors-lemma-relative-regular-immersion-flat-in-neighbourhood}
$X \to S$ をスキームの射とし、$Z \to X$ を埋め込みとする。次を仮定する。
\begin{enumerate}
\item $X \to S$ は平坦かつ局所有限表示である。
\item $Z \to X$ は相対準正則埋め込みである。
\end{enumerate}
このとき $Z \to X$ は正則埋め込みであり、任意の基底変換後にも
同じことが成り立つ。
\end{lemma}

\begin{proof}
$x \in Z$ を取り、その像を $s \in S$ とする。これを証明するには、
$x$ の $U$ に含まれるアフィン近傍で、結果がそのアフィン開集合上で
成り立つものを見つければ十分である。したがって $X$ はアフィンであり、
ある準正則列 $f_1, \ldots, f_r \in \Gamma(X, \mathcal{O}_X)$ が存在して
$Z = V(f_1, \ldots, f_r)$ であると仮定してよい。More on Algebra, Lemma
\ref{more-algebra-lemma-relative-regular-immersion-algebra} により、列
$f_1|_{X_s}, \ldots, f_r|_{X_s}$ は
$\Gamma(X_s, \mathcal{O}_{X_s})$ の準正則列である。
$X_s$ は Noether なので、必要なら $X$ を少し縮小することにより、
$f_1|_{X_s}, \ldots, f_r|_{X_s}$ は正則列となる。Algebra, Lemmas
\ref{algebra-lemma-quasi-regular-regular} および
\ref{algebra-lemma-regular-sequence-in-neighbourhood} を参照。
Lemma \ref{divisors-lemma-fibre-Cartier} により、
$Z_1 = V(f_1) \subset X$ は相対有効 Cartier 因子となる。ここでも必要なら
$X$ を少し縮小する。同じ補題を今度は
$Z_2 = V(f_1, f_2) \subset Z_1$ に適用すると、$Z_2 \subset Z_1$ は
相対有効 Cartier 因子であることが分かる。このようにして
$Z = Z_n = V(f_1, \ldots, f_n)$ に至る。相対有効 Cartier 因子であることは
任意の基底変換で保たれるので（Lemma \ref{divisors-lemma-relative-Cartier} を参照）、
補題の最後の主張も従う。
\end{proof}

\begin{remark}
\label{divisors-remark-relative-regular-immersion-elements}
相対準正則埋め込みの余次元は、それが一定ならば、基底変換後にも変わらない。
実際、環準同型 $A \to B$ と準正則列
$f_1, \ldots, f_r \in B$ があり、$B/(f_1, \ldots, f_r)$ は
$A$ 上平坦であるとする。このとき任意の環準同型 $A \to A'$ に対し、
More on Algebra, Lemma
\ref{more-algebra-lemma-relative-regular-immersion-algebra} により
$f_1 \otimes 1, \ldots, f_r \otimes 1$ は
$B' = B \otimes_A A'$ の準正則列である（これは上の Lemma
\ref{divisors-lemma-relative-regular-immersion} の証明で用いた）。
ここで Lemma
\ref{divisors-lemma-relative-regular-immersion-flat-in-neighbourhood} の証明から、
$A \to B$ が平坦かつ局所有限表示ならば、任意の素イデアル
$\mathfrak q' \subset B'$ に対し、列
$f_1 \otimes 1, \ldots, f_r \otimes 1$ は局所環
$B'_{\mathfrak q'}$ において正則列でさえあることが分かる。
\end{remark}

\begin{lemma}
\label{divisors-lemma-flat-relative-H1-regular}
$X \to S$ をスキームの射とし、$Z \to X$ を相対 $H_1$-正則埋め込みとする。
$X \to S$ は局所有限表示であると仮定する。このとき次が成り立つ。
\begin{enumerate}
\item ある開部分スキーム $U \subset X$ が存在して、$Z \subset U$ かつ
$U \to S$ は平坦である。
\item $Z \to X$ は正則埋め込みであり、任意の基底変換後にも
同じことが成り立つ。
\end{enumerate}
\end{lemma}

\begin{proof}
$x \in Z$ を取る。(1) を証明するには、開近傍 $U \subset X$ を $x$ について取り、
$U \to S$ が平坦となるものを見つければ十分である。したがって補題は、
$X = \Spec(B)$ と $S = \Spec(A)$ がアフィンであり、$Z$ が
$H_1$-正則列 $f_1, \ldots, f_r \in B$ により定められる場合に帰着する。
仮定により $B$ は有限表示 $A$-代数であり、
$B/(f_1, \ldots, f_r)B$ は平坦な $A$-代数である。
絶対 Noether 近似を用いる。

\medskip\noindent
$B = A[x_1, \ldots, x_n]/(g_1, \ldots, g_m)$ と書く。
$f_i$ は $f_i' \in A[x_1, \ldots, x_n]$ の像であるとする。
有限型 $\mathbf{Z}$-部分代数 $A_0 \subset A$ で、すべての多項式
$f_1', \ldots, f_r', g_1, \ldots, g_m$ の係数が $A_0$ に入るものを選ぶ。
$B_0 = A_0[x_1, \ldots, x_n]/(g_1, \ldots, g_m)$ と置き、
$f_{i, 0}$ を $f_i'$ の $B_0$ における像と書く。このとき
$B = B_0 \otimes_{A_0} A$ であり、
$$
B/(f_1, \ldots, f_r) =
B_0/(f_{0, 1}, \ldots, f_{0, r}) \otimes_{A_0} A.
$$
Algebra, Lemma \ref{algebra-lemma-flat-finite-presentation-limit-flat}
により、$A_0$ を拡大した後、
$B_0/(f_{0, 1}, \ldots, f_{0, r})$ は $A_0$ 上平坦であると仮定してよい。
この時点では Koszul コホモロジー群
$H_1(K_\bullet(B_0, f_{0, 1}, \ldots, f_{0, r}))$ が零とは限らない。
一方、$B_0$ は Noether なので、これは有限生成 $B_0$-加群である。生成元を
$\xi_1, \ldots, \xi_n \in H_1(K_\bullet(B_0, f_{0, 1}, \ldots, f_{0, r}))$
とする。$A_0 \subset A_1 \subset A$ をより大きい有限型
$\mathbf{Z}$-部分代数で、$A$ の部分代数であるものとする。$f_{1, i}$ を
$f_{0, i}$ の $B_1 = B_0 \otimes_{A_0} A_1$ における像と書く。
More on Algebra, Lemma \ref{more-algebra-lemma-base-change-H1-regular}
により、写像
$$
H_1(K_\bullet(B_0, f_{0, 1}, \ldots, f_{0, r})) \otimes_{A_0} A_1
\longrightarrow
H_1(K_\bullet(B_1, f_{1, 1}, \ldots, f_{1, r}))
$$
は全射である。さらに、上のような $A_1$ のすべての選択にわたる複体
$K_\bullet(B_1, f_{1, 1}, \ldots, f_{1, r})$ の余極限は、複体
$K_\bullet(B, f_1, \ldots, f_r)$ であり、これは次数 $1$ で非輪状であることは
明らかである。したがって
$$
\colim_{A_0 \subset A_1 \subset A}
H_1(K_\bullet(B_1, f_{1, 1}, \ldots, f_{1, r}))
= 0
$$
である。Algebra, Lemma
\ref{algebra-lemma-directed-colimit-exact} を参照。
よって、$A_1$ を、$\xi_1, \ldots, \xi_n$ がすべて
$H_1(K_\bullet(B_1, f_{1, 1}, \ldots, f_{1, r}))$ の零へ写るように選べる。
言い換えると、Koszul コホモロジー群
$H_1(K_\bullet(B_1, f_{1, 1}, \ldots, f_{1, r}))$ は零である。

\medskip\noindent
アフィンスキームの射 $X_1 \to S_1$ を考える。これは
$\Spec$ を環準同型 $A_1 \to B_1$ に適用して得られる。さらに
$Z_1 = \Spec(B_1/(f_{1, 1}, \ldots, f_{1, r}))$ を考える。
$B = B_1 \otimes_{A_1} A$、すなわち
$X = X_1 \times_{S_1} S$ であり、同様に
$Z = Z_1 \times_S S_1$ であるから、今や (1) を
$X_1 \to S_1$ および相対 $H_1$-正則埋め込み $Z_1 \to X_1$
について証明すれば十分である。Morphisms, Lemma
\ref{morphisms-lemma-base-change-module-flat} を参照。
したがって、$X \to S$ が Noether スキーム間の有限型射である場合に帰着した。
この場合 Lemma \ref{divisors-lemma-quasi-regular-immersion-flat-at-x} により、
$X \to S$ は $Z$ のすべての点で平坦である。さらにこの場合、平坦性の軌跡が
開であること（Algebra, Theorem
\ref{algebra-theorem-openness-flatness} を参照）と合わせると、(1) が従う。
(2) は Lemma
\ref{divisors-lemma-relative-regular-immersion-flat-in-neighbourhood} を適用して従う。
\end{proof}

\noindent
周囲のスキームが基底上平坦かつ局所有限表示ならば、相対準正則埋め込みを
そのファイバーによって特徴づけられる。

\begin{lemma}
\label{divisors-lemma-fibre-quasi-regular}
$\varphi : X \to S$ を平坦かつ局所有限表示な射とし、
$T \subset X$ を閉部分スキームとする。$x \in T$ とし、その像を $s \in S$ とする。
\begin{enumerate}
\item $T_s \subset X_s$ が $x$ の近傍で準正則埋め込みならば、開集合
$U \subset X$ と相対準正則埋め込み $Z \subset U$ が存在して、
$Z_s = T_s \cap U_s$ かつ $T \cap U \subset Z$ となる。
\item $T_s \subset X_s$ が $x$ の近傍で準正則埋め込みであり、射
$T \to X$ が有限表示、$T \to S$ が $x$ で平坦ならば、(1) の $U$ と
$Z$ を $T \cap U = Z$ となるように選べる。
\item $T_s \subset X_s$ が $x$ の近傍で準正則埋め込みであり、$T$ が
$c$ 個の方程式により $x$ の近傍で切り出され、ここで
$c = \dim_x(X_s) - \dim_x(T_s)$ ならば、(1) の $U$ と $Z$ を
$T \cap U = Z$ となるように選べる。
\end{enumerate}
いずれの場合も $Z \to U$ は Lemma
\ref{divisors-lemma-relative-regular-immersion-flat-in-neighbourhood} により正則埋め込みである。
とくに、$T \to S$ が局所有限表示かつ平坦で、すべてのファイバー
$T_s \subset X_s$ が準正則埋め込みならば、$T \to X$ は相対準正則埋め込みである。
\end{lemma}

\begin{proof}
アフィン開近傍 $\Spec(A)$ を $s$ について、またアフィン開近傍
$\Spec(B)$ を $x$ について、$\varphi(\Spec(B)) \subset \Spec(A)$ となるように選ぶ。
素イデアル $\mathfrak p \subset A$ を $s$ に対応するものとし、素イデアル
$\mathfrak q \subset B$ を $x$ に対応するものとする。イデアル
$I \subset B$ を $T$ に対応するものとする。補題の最初の仮定により、
$A \to B$ は平坦かつ有限表示である。
(1) の仮定は、$\Spec(B)$ を縮小した後、
$I(B \otimes_A \kappa(\mathfrak p))$ が元の準正則列により生成されると
仮定できることを意味する。必要なら $B$ をある $g \in B$、
$g \not \in \mathfrak q$ で局所化した後、ある
$f_1, \ldots, f_r \in I$ が存在して、それらが
$B \otimes_A \kappa(\mathfrak p)$ の準正則列に写り、
$I(B \otimes_A \kappa(\mathfrak p))$ を生成すると仮定してよい。
Algebra, Lemmas \ref{algebra-lemma-quasi-regular-regular} および
\ref{algebra-lemma-regular-sequence-in-neighbourhood} により、
さらに局所化した後、$f_1, \ldots, f_r \in I$ は
$B \otimes_A \kappa(\mathfrak p)$ の正則列をなすと仮定してよい。
Lemma \ref{divisors-lemma-fibre-Cartier} により、
$Z_1 = V(f_1) \subset \Spec(B)$ は相対有効 Cartier 因子となる。ここでも必要なら
$B$ を局所化する。
同じ補題を今度は $Z_2 = V(f_1, f_2) \subset Z_1$ に適用すると、
$Z_2 \subset Z_1$ は相対有効 Cartier 因子であることが分かる。
このようにして $Z = Z_n = V(f_1, \ldots, f_n)$ に至る。このとき
$Z \to \Spec(B)$ は正則埋め込みで、$Z$ は $S$ 上平坦である。
とくに $Z \to \Spec(B)$ は $\Spec(A)$ 上の相対準正則埋め込みである。
これで (1) が示された。

\medskip\noindent
(2) を示すため、閉埋め込み $Z \to D$ を考える。全射環準同型
$u : \mathcal{O}_{D, x} \to \mathcal{O}_{Z, x}$ は、
本質的有限表示な平坦局所 $\mathcal{O}_{S, s}$-代数の間の写像であり、
$\mathfrak m_s$ で割ると同型となる。したがってこれは同型である。
Algebra, Lemma \ref{algebra-lemma-mod-injective-general} を参照。
ゆえに $Z \to D$ は $x$ の近傍で同型である。Algebra, Lemma
\ref{algebra-lemma-local-isomorphism} を参照。

\medskip\noindent
(3) を示すため、必要なら $U$ を縮小し、$Z$ のイデアルは正則列
$f_1, \ldots, f_r$ により生成され（上の $Z$ の構成を参照）、$T$ のイデアルは
$g_1, \ldots, g_c$ により生成されると仮定してよい。$c = r$ であると主張する。
実際、
\begin{align*}
\dim_x(X_s) & = \dim(\mathcal{O}_{X_s, x}) +
\text{trdeg}_{\kappa(s)}(\kappa(x)), \\
\dim_x(T_s) & = \dim(\mathcal{O}_{T_s, x}) +
\text{trdeg}_{\kappa(s)}(\kappa(x)), \\
\dim(\mathcal{O}_{X_s, x}) & = \dim(\mathcal{O}_{T_s, x}) + r
\end{align*}
である。最初の二つの等式は Algebra, Lemma
\ref{algebra-lemma-dimension-at-a-point-finite-type-field} により、
第二の等式は Algebra, Lemma \ref{algebra-lemma-one-equation} を
$r$ 回適用することによる。$T \subset Z$ なので
$f_i = \sum b_{ij} g_j$ と書ける。しかし $Z$ と $T$ のイデアルは、
$X_s$ の、$x$ の近傍における同じ準正則閉部分スキームを切り出す。
したがって行列 $(b_{ij}) \bmod \mathfrak m_x$ は可逆である
（いくつかの細部は省略する）。ゆえに $(b_{ij})$ は $x$ の開近傍で可逆である。
言い換えると、$T \cap U = Z$ となり、これは $U$ を縮小した後に成り立つ。

\medskip\noindent
補題の最後の主張は、(2) と、$Z \to S$ が局所有限表示であることと
$Z \to X$ が有限表示であることの同値性を合わせれば直ちに従う。
Morphisms, Lemmas \ref{morphisms-lemma-composition-finite-presentation} および
\ref{morphisms-lemma-finite-presentation-permanence} を参照。
\end{proof}

\noindent
次の補題は Morphisms, Lemma
\ref{morphisms-lemma-section-smooth-morphism} の強化である。

\begin{lemma}
\label{divisors-lemma-section-smooth-regular-immersion}
$f : X \to S$ をスキームの平滑射とし、$\sigma : S \to X$ を
$f$ の切断とする。このとき $\sigma$ は正則埋め込みである。
\end{lemma}

\begin{proof}
Schemes, Lemma \ref{schemes-lemma-semi-diagonal} により、射
$\sigma$ は埋め込みである。$X$ を $\sigma(S)$ の開近傍で置き換え、
$\sigma$ は閉埋め込みであると仮定してよい。$T = \sigma(S)$ を対応する
$X$ の閉部分スキームとする。$T \to S$ は同型なので平坦かつ有限表示である。
また平滑射は平坦かつ局所有限表示である。Morphisms, Lemmas
\ref{morphisms-lemma-smooth-flat} および
\ref{morphisms-lemma-smooth-locally-finite-presentation} を参照。
したがって Lemma \ref{divisors-lemma-fibre-quasi-regular} により、
$T_s \subset X_s$ が準正則閉部分スキームであることを示せば十分である。
これは Morphisms, Lemma
\ref{morphisms-lemma-section-smooth-morphism} から直ちに従うが、
次のように直接見ることもできる。$k$ を体、$A$ を平滑 $k$-代数とする。
$\mathfrak m \subset A$ を、剰余体が $k$ である極大イデアルとする。
このとき $\mathfrak m$ は準正則列により生成される。必要なら
$A$ を $A_g$ で置き換える。ここで $g \in A$、$g \not \in \mathfrak m$ である。
Algebra, Lemma \ref{algebra-lemma-characterize-smooth-over-field} で
$A_{\mathfrak m}$ が正則局所環であることを示したので、
$\mathfrak mA_{\mathfrak m}$ は正則列により生成される。これは実際、
$\mathfrak m$ が正則列により生成されることを含意する
（$A$ を $A_g$ で置き換えた後。ここで $g \in A$、
$g \not \in \mathfrak m$ である）。Algebra, Lemma
\ref{algebra-lemma-regular-sequence-in-neighbourhood} を参照。
\end{proof}

\noindent
次の補題にはある種の逆がある。Lemma
\ref{divisors-lemma-push-regular-immersion-thru-smooth} を参照。

\begin{lemma}
\label{divisors-lemma-lift-regular-immersion-to-smooth}
次の可換図式
$$
\xymatrix{
Y \ar[rd]_j \ar[rr]_i & & X \ar[ld] \\
& S
}
$$
をスキームの射の図式とする。$X \to S$ は平滑で、$i$、$j$ は埋め込みと仮定する。
$j$ が正則（それぞれ Koszul 正則、$H_1$-正則、準正則）埋め込みならば、
$i$ もそうである。
\end{lemma}

\begin{proof}
$i$ を合成
$$
Y \to Y \times_S X \to X
$$
と書ける。Lemma \ref{divisors-lemma-section-smooth-regular-immersion} により、
第一の矢印は正則埋め込みである。第二の矢印は $Y \to S$ の平坦基底変換なので、
正則（それぞれ Koszul 正則、$H_1$-正則、準正則）埋め込みである。Lemma
\ref{divisors-lemma-flat-base-change-regular-immersion} を参照。
Lemma \ref{divisors-lemma-composition-regular-immersion} を適用して結論を得る。
\end{proof}

\begin{lemma}
\label{divisors-lemma-immersion-lci-into-smooth-regular-immersion}
次の可換図式
$$
\xymatrix{
Y \ar[rd] \ar[rr]_i & & X \ar[ld] \\
& S
}
$$
をスキームの射の図式とする。$Y \to S$ は syntomic、$X \to S$ は平滑で、
$i$ は埋め込みであると仮定する。このとき $i$ は正則埋め込みである。
\end{lemma}

\begin{proof}
$X$ を $i(Y)$ の開近傍で置き換え、$i$ は閉埋め込みであると仮定してよい。
$T = i(Y)$ を対応する $X$ の閉部分スキームとする。$T \cong Y$ なので、
射 $T \to S$ は平坦かつ有限表示である（Morphisms, Lemmas
\ref{morphisms-lemma-syntomic-locally-finite-presentation} および
\ref{morphisms-lemma-syntomic-flat}）。
また平滑射は平坦かつ局所有限表示である（Morphisms, Lemmas
\ref{morphisms-lemma-smooth-flat} および
\ref{morphisms-lemma-smooth-locally-finite-presentation}）。
したがって Lemma \ref{divisors-lemma-fibre-quasi-regular} により、
$T_s \subset X_s$ が準正則閉部分スキームであることを示せば十分である。
$X_s$ は体上局所有限型なので Noether である（Morphisms, Lemma
\ref{morphisms-lemma-finite-type-noetherian}）。したがって
$T_s \subset X_s$ が準正則埋め込みであることは点ごとに判定できる。
Lemma \ref{divisors-lemma-Noetherian-scheme-regular-ideal} を参照。
$t \in T_s$ を取る。Morphisms, Lemma
\ref{morphisms-lemma-local-complete-intersection} により、局所環
$\mathcal{O}_{T_s, t}$ は $\kappa(s)$ 上の局所完全交叉である。
局所環 $\mathcal{O}_{X_s, t}$ は正則である。Algebra, Lemma
\ref{algebra-lemma-characterize-smooth-over-field} を参照。
Algebra, Lemma \ref{algebra-lemma-lci-local} により、全射の核
$\mathcal{O}_{X_s, t} \to \mathcal{O}_{T_s, t}$ は正則列により生成される。
これが示すべきことであった。
\end{proof}

\begin{lemma}
\label{divisors-lemma-immersion-smooth-into-smooth-regular-immersion}
次の可換図式
$$
\xymatrix{
Y \ar[rd] \ar[rr]_i & & X \ar[ld] \\
& S
}
$$
をスキームの射の図式とする。$Y \to S$ は平滑、$X \to S$ は平滑で、
$i$ は埋め込みであると仮定する。このとき $i$ は正則埋め込みである。
\end{lemma}

\begin{proof}
これは Lemma
\ref{divisors-lemma-immersion-lci-into-smooth-regular-immersion} の特別な場合である。
実際、平滑射は syntomic である。Morphisms, Lemma
\ref{morphisms-lemma-smooth-syntomic} を参照。
\end{proof}

\begin{lemma}
\label{divisors-lemma-push-regular-immersion-thru-smooth}
次の可換図式
$$
\xymatrix{
Y \ar[rd]_j \ar[rr]_i & & X \ar[ld] \\
& S
}
$$
をスキームの射の図式とする。$X \to S$ は平滑で、$i$ と $j$ は埋め込みと仮定する。
$i$ が Koszul 正則（それぞれ $H_1$-正則、準正則）埋め込みならば、
$j$ もそうである。
\end{lemma}

\begin{proof}
以下では Lemma \ref{divisors-lemma-regular-quasi-regular-immersion} を
断りなく用いる。任意の点 $y \in Y$ を取る。$x = i(y)$、$s = j(y)$ と置く。
$X$ と $S$ をそれぞれ開近傍 $U$ と $V$ で置き換える。これらは
$x$ と $s$ の近傍である。また $Y$ を $y$ の、
$i^{-1}(U) \cap j^{-1}(V)$ における開近傍で置き換えた後に
結果を証明すれば十分である。

\medskip\noindent
まず $X = \mathbf{A}^n_S$ の場合に結果を証明する。$S$ をアフィン開集合
$V$ で置き換え、$Y$ を $j^{-1}(V)$ で置き換えることにより、
$j$ は閉埋め込みで $S$ はアフィンであると仮定してよい。
$S = \Spec(A)$ と書く。このとき $j : Y \to S$ は $Y$ を閉部分スキーム
$\Spec(A/I)$ と同一視する。ここで $I \subset A$ はあるイデアルである。
写像
$i : Y = \Spec(A/I) \to
\mathbf{A}^n_S = \Spec(A[x_1, \ldots, x_n])$
は $A$-代数準同型
$i^\sharp : A[x_1, \ldots, x_n] \to A/I$
に対応する。$a_i \in A$ を選び、$i^\sharp(x_i)$ の $A/I$ における像になるものとする。
閉埋め込み $i$ のイデアルは
$$
J = (x_1 - a_1, \ldots, x_n - a_n) + IA[x_1, \ldots, x_n].
$$
であることに注意せよ。$K = (x_1 - a_1, \ldots, x_n - a_n)$ と置く。
列
$$
0 \to K/KJ \to J/J^2 \to J/(K + J^2) \to 0
$$
は分裂完全であると主張する。これを見るため、$K/K^2$ は基底
$x_i - a_i$ をもち、環 $A[x_1, \ldots, x_n]/K \cong A$ 上の自由加群であることに
注意する。したがって $K/KJ$ は
$A[x_1, \ldots, x_n]/J \cong A/I$ 上、同じ基底をもつ自由加群である。
一方、微分を取ると写像
$$
\text{d}_{A[x_1, \ldots, x_n]/A} :
J/J^2
\longrightarrow
\Omega_{A[x_1, \ldots, x_n]/A} \otimes_{A[x_1, \ldots, x_n]}
A[x_1, \ldots, x_n]/J
$$
を得る。これは生成元 $x_i - a_i$ を右辺の自由加群の基底元
$\text{d}x_i$ へ写すので、主張が従う。さらに、
$x_1 - a_1, \ldots, x_n - a_n$ は
$A[x_1, \ldots, x_n]$ の正則列であり、その商環は
$A[x_1, \ldots, x_n]/(x_1 - a_1, \ldots, x_n - a_n) \cong A$ である。
したがって閉埋め込み $i$ は次のように分解される。
$$
Y \to V(x_1 - a_1, \ldots, x_n - a_n) \to \mathbf{A}^n_S
$$
ここで合成は Koszul 正則（それぞれ $H_1$-正則、準正則）、第二の矢印は
正則埋め込みで、付随する余法列は分裂する。よって Lemma
\ref{divisors-lemma-permanence-regular-immersion} から結果が従う。

\medskip\noindent
次に、$i$ が $H_1$-正則ならば結果が成り立つことを証明する。
すなわち、証明の第一段落のように縮小して、$Y$、$X$、$S$ はアフィンであると
仮定してよい。この場合、閉埋め込み $h : X \to \mathbf{A}^n_S$ を
$S$ 上で、ある $n$ に対して選べる。Lemma
\ref{divisors-lemma-immersion-smooth-into-smooth-regular-immersion} により、
$h$ は正則埋め込みである。したがって $h \circ i$ は
$H_1$-正則埋め込みである。Lemma
\ref{divisors-lemma-composition-regular-immersion} を参照
（この手順は「準正則の場合」には使えない）。
ゆえに、上で証明した $X = \mathbf{A}^n_S$ かつ $S$ がアフィンの場合に帰着する。

\medskip\noindent
最後に $i$ は準正則であると仮定する。証明の第一段落のように縮小した後、
Morphisms, Lemma
\ref{morphisms-lemma-smooth-etale-over-affine-space} により、$f$ を
$X \to \mathbf{A}^n_S \to S$ と分解でき、第一の射
$X \to \mathbf{A}^n_S$ は \'etale である。これにより問題は
(a) $X = \mathbf{A}^n_S$ の場合と (b) $f$ が \'etale の場合に帰着する。
(a) は証明の第二段落で扱った。(b) は次の段落で扱う。

\medskip\noindent
$f$ は \'etale であると仮定する。縮小した後、$X$、$Y$、$S$ はアフィンで、
$i$ と $j$ は閉埋め込みであると仮定してよい（細部を一つ省略する）。
$S = \Spec(A)$、$X = \Spec(B)$、
$Y = \Spec(B/J) = \Spec(A/I)$ とする。さらに縮小して、$J$ は準正則列により
生成されると仮定してよい。環準同型 $A \to B$ は \'etale、したがって
形式的 \'etale である（Algebra, Lemma
\ref{algebra-lemma-formally-etale-etale}）。ゆえに Algebra, Lemma
\ref{algebra-lemma-formally-etale-lift-infinitesimal} により
$\bigoplus I^n/I^{n + 1} \cong \bigoplus J^n/J^{n + 1}$ である。
$J$ は準正則列により生成されるので、$I$ もそうである。これで証明は終わる。
\end{proof}

\section{有理型関数と有理型切断}
\label{divisors-section-meromorphic-functions}

\noindent
この節では一般的な定義といくつかの初等的な結果だけを扱う。起こりうる落とし穴については
\cite{misconceptions} を参照\footnote{危険だ、ウィル・ロビンソン！}。

\medskip\noindent
$(X, \mathcal{O}_X)$ を局所環付き空間とする。任意の開集合 $U \subset X$ に対し、
集合 $\mathcal{S}(U) \subset \mathcal{O}_X(U)$、すなわち
$\mathcal{O}_X$ の $U$ 上の正則切断全体を定義した。Definition
\ref{divisors-definition-regular-section} を参照。正則切断をより小さい開集合へ制限しても
正則である。したがって $\mathcal{S} : U \mapsto \mathcal{S}(U)$ は
$\mathcal{O}_X$ の（集合の）部分層である。$\mathcal{S} = \mathcal{S}_X$ と書いて
$X$ への依存を明示することもある。さらに、$\mathcal{S}(U)$ は環
$\mathcal{O}_X(U)$ の乗法的部分集合であり、これは各 $U$ について成り立つ。
したがって環の前層
$$
U \longmapsto \mathcal{S}(U)^{-1} \mathcal{O}_X(U),
$$
を考えられる。Modules, Lemma \ref{modules-lemma-simple-invert} を参照。

\begin{definition}
\label{divisors-definition-sheaf-meromorphic-functions}
$(X, \mathcal{O}_X)$ を局所環付き空間とする。{\it $X$ 上の有理型関数の層}とは、
上に表示した前層に付随する層 {\it $\mathcal{K}_X$} のことである。
$X$ 上の {\it 有理型関数}とは $\mathcal{K}_X$ の大域切断である。
\end{definition}

\noindent
各 $\mathcal{S}(U)$ の各元は $\mathcal{O}_X(U)$ 上の非零因子なので、
環の層の自然な写像 $\mathcal{O}_X \to \mathcal{K}_X$ は単射である。

\begin{example}
\label{divisors-example-no-change}
$A = \mathbf{C}[x, \{y_\alpha\}_{\alpha \in \mathbf{C}}]/
((x - \alpha)y_\alpha, y_\alpha y_\beta)$ とする。$A$ の任意の元は一意的に
$f(x) + \sum \lambda_\alpha y_\alpha$ と書ける。ここで
$f(x) \in \mathbf{C}[x]$ かつ $\lambda_\alpha \in \mathbf{C}$ である。
$X = \Spec(A)$ とする。この場合 $\mathcal{O}_X = \mathcal{K}_X$ である。
実際、任意のアフィン開集合 $D(f)$ 上で、環 $A_f$ の任意の非零因子は
単元である（証明は省略する）。
\end{example}

\noindent
$(X, \mathcal{O}_X)$ を局所環付き空間とする。$\mathcal{F}$ を
$\mathcal{O}_X$-加群の層とする。前層
$U \mapsto \mathcal{S}(U)^{-1}\mathcal{F}(U)$ を考える。その層化は
$\mathcal{F} \otimes_{\mathcal{O}_X} \mathcal{K}_X$ である。
Modules, Lemma \ref{modules-lemma-simple-invert-module} を参照。

\begin{definition}
\label{divisors-definition-meromorphic-section}
$X$ を局所環付き空間とし、$\mathcal{F}$ を $\mathcal{O}_X$-加群の層とする。
\begin{enumerate}
\item $\mathcal{K}_X(\mathcal{F})$ により、$\mathcal{K}_X$-加群の層であって、
前層 $U \mapsto \mathcal{S}(U)^{-1}\mathcal{F}(U)$ の層化であるものを表す。
同値な表示は
$\mathcal{K}_X(\mathcal{F}) =
\mathcal{F} \otimes_{\mathcal{O}_X} \mathcal{K}_X$ である（上を参照）。
\item {\it $\mathcal{F}$ の有理型切断}とは、
$\mathcal{K}_X(\mathcal{F})$ の大域切断である。
\end{enumerate}
\end{definition}

\noindent
とくに、任意の点 $x \in X$ に対し
$$
\mathcal{K}_X(\mathcal{F})_x
=
\mathcal{F}_x \otimes_{\mathcal{O}_{X, x}} \mathcal{K}_{X, x}
=
\mathcal{S}_x^{-1}\mathcal{F}_x
$$
である。しかし、$\mathcal{S}_x$ が局所環
$\mathcal{O}_{X, x}$ の非零因子全体とは限らないので注意が必要である。
実際、全商環への単射
$$
\mathcal{K}_{X, x} \longrightarrow Q(\mathcal{O}_{X, x})
$$
は常に存在する。これが全射であることと、$\mathcal{S}_x$ が
$\mathcal{O}_{X, x}$ の非零因子全体であることは同値である。
有理型切断の層は一般には準連接加群ではないが、準連接加群と共通する性質を
いくつかもつ。

\begin{definition}
\label{divisors-definition-pullback-meromorphic-sections}
$f : (X, \mathcal{O}_X) \to (Y, \mathcal{O}_Y)$ を局所環付き空間の射とする。
{\it $f$ に対して有理型関数の引き戻しが定義されている}とは、
任意の開集合の組 $U \subset X$、$V \subset Y$ で
$f(U) \subset V$ となるもの、および任意の切断
$s \in \Gamma(V, \mathcal{S}_Y)$ に対し、その引き戻し
$f^\sharp(s) \in \Gamma(U, \mathcal{O}_X)$ が
$\Gamma(U, \mathcal{S}_X)$ の元となることをいう。
\end{definition}

\noindent
この場合、誘導写像
$f^\sharp : f^{-1}\mathcal{K}_Y \to \mathcal{K}_X$
があり、言い換えると環付き空間の射の可換図式
$$
\xymatrix{
(X, \mathcal{K}_X) \ar[r] \ar[d]^f &
(X, \mathcal{O}_X) \ar[d]^f \\
(Y, \mathcal{K}_Y) \ar[r] &
(Y, \mathcal{O}_Y)
}
$$
を得る。$f^*(s) = f^\sharp(s)$ と書くこともある。ここで
$s \in \Gamma(Y, \mathcal{K}_Y)$ は切断である。

\begin{lemma}
\label{divisors-lemma-pullback-meromorphic-sections-defined}
$f : X \to Y$ をスキームの射とする。次の各場合には有理型関数の引き戻しが
定義されている。
\begin{enumerate}
\item $X$ のすべての弱随伴点は $Y$ の既約成分の一般点へ写る。
\item $X$、$Y$ は整であり、$f$ は支配的である。
\item $X$ は整であり、$X$ の一般点は $Y$ の既約成分の一般点へ写る。
\item $X$ は被約であり、$X$ の各既約成分の各一般点は
$Y$ の既約成分の一般点へ写る。
\item $X$ は局所 Noether であり、$X$ の任意の随伴点は
$Y$ の既約成分の一般点へ写る。
\item $X$ は局所 Noether で埋め込まれた点をもたず、$X$ の既約成分の
任意の一般点は $Y$ の既約成分の一般点へ写る。
\item $f$ は平坦である。
\end{enumerate}
\end{lemma}

\begin{proof}
問題は $X$ と $Y$ 上局所的である。したがって、
$X = \Spec(A)$、$Y = \Spec(R)$ で、$f$ が環準同型
$\varphi : R \to A$ により与えられる場合に帰着する。
Lemma \ref{divisors-lemma-regular-section-structure-sheaf} による構造層の正則切断の
特徴づけから、$R \to A$ が非零因子を非零因子へ写すことを示せばよい。
$t \in R$ を非零因子とする。

\medskip\noindent
$R \to A$ が平坦ならば、$t : R \to R$ が単射であることから
$t : A \to A$ も単射である。これで (7) が示された。

\medskip\noindent
他の場合には、$t$ は $R$ のどの極小素イデアルにも含まれないことに注意する
（環の極小素イデアルの各元は零因子だからである）。
したがって (1) の場合、$\varphi(t)$ は $A$ のどの弱随伴素イデアルにも
含まれない。よってこの場合は Algebra, Lemma
\ref{algebra-lemma-weakly-ass-zero-divisors} から従う。
(5) は Lemma \ref{divisors-lemma-ass-weakly-ass} により (1) の特別な場合である。
(6) は (5) と定義から従う。(4) は Lemma
\ref{divisors-lemma-weakass-reduced} により (1) の特別な場合である。
(2) と (3) は (4) の特別な場合である。
\end{proof}

\begin{lemma}
\label{divisors-lemma-meromorphic-weakass-finite}
$X$ を次を満たすスキームとする。
\begin{enumerate}
\item[(a)] $X$ のすべての弱随伴点は $X$ の既約成分の一般点である。
\item[(b)] 任意の準コンパクト開集合は有限個の既約成分をもつ。
\end{enumerate}
$X^0$ を $X$ の既約成分の一般点全体とする。このとき
$$
\mathcal{K}_X =
\bigoplus\nolimits_{\eta \in X^0} j_{\eta, *}\mathcal{O}_{X, \eta} =
\prod\nolimits_{\eta \in X^0} j_{\eta, *}\mathcal{O}_{X, \eta}
$$
である。ここで $j_\eta : \Spec(\mathcal{O}_{X, \eta}) \to X$ は
Schemes, Section \ref{schemes-section-points} の標準写像である。さらに、
\begin{enumerate}
\item $\mathcal{K}_X$ は準連接な $\mathcal{O}_X$-代数の層である。
\item 任意の準連接 $\mathcal{O}_X$-加群 $\mathcal{F}$ に対し、
$$
\mathcal{K}_X(\mathcal{F}) =
\bigoplus\nolimits_{\eta \in X^0} j_{\eta, *}\mathcal{F}_\eta =
\prod\nolimits_{\eta \in X^0} j_{\eta, *}\mathcal{F}_\eta
$$
という $\mathcal{F}$ の有理型切断の層は準連接である。
\item $\mathcal{S}_x \subset \mathcal{O}_{X, x}$ は、任意の $x \in X$ に対し
非零因子全体である。
\item $\mathcal{K}_{X, x}$ は $\mathcal{O}_{X, x}$ の全商環であり、
これは任意の $x \in X$ について成り立つ。
\item $\mathcal{K}_X(U)$ は $\mathcal{O}_X(U)$ の全商環に等しく、
これは任意のアフィン開集合 $U \subset X$ について成り立つ。
\item $X$ の有理関数の環（Morphisms, Definition
\ref{morphisms-definition-rational-function}）は $X$ 上の有理型関数の環であり、
式で書けば $R(X) = \Gamma(X, \mathcal{K}_X)$ である。
\end{enumerate}
\end{lemma}

\begin{proof}
たとえば茎で確認できるので、加群の層の局所有限直和は積に等しいことに注意する。
次に $\mathcal{K}_X(\mathcal{F}) =
\mathcal{F} \otimes_{\mathcal{O}_X} \mathcal{K}_X$ なので、(2) は他の主張から従う。
また、$X^0$ が有限ならば (6) は補題の最初の主張と Morphisms, Lemma
\ref{morphisms-lemma-integral-scheme-rational-functions} から従い、
一般の (6) は貼り合わせによりこれから従う（議論は省略する）。

\medskip\noindent
$j : Y = \coprod\nolimits_{\eta \in X^0} \Spec(\mathcal{O}_{X, \eta}) \to X$
を射 $j_\eta$ の余積とする。示すべきことは
$\mathcal{K}_X = j_*\mathcal{O}_Y$ である。まず
$\mathcal{K}_Y = \mathcal{O}_Y$ であることに注意する。実際、$Y$ は次元
$0$ の局所環のスペクトルの非交和であり、零次元局所環の任意の非零因子は
単元である。次に、$j$ に対して
Lemma \ref{divisors-lemma-pullback-meromorphic-sections-defined} により
有理型関数の引き戻しが定義されていることに注意する。これにより写像
$$
\mathcal{K}_X \longrightarrow j_*\mathcal{O}_Y.
$$
を得る。$\Spec(A) = U \subset X$ をアフィン開集合とする。このとき $A$ は
有限個の極小素イデアル $\mathfrak q_1, \ldots, \mathfrak q_t$ をもつ環であり、
$A$ の各弱随伴素イデアルは $\mathfrak q_i$ のいずれかである。
$Q(A) = \prod A_{\mathfrak q_i}$ を Algebra, Lemmas
\ref{algebra-lemma-total-ring-fractions-no-embedded-points}
および \ref{algebra-lemma-weakly-ass-zero-divisors} により得る。
言い換えると、前層の値
$U \mapsto \mathcal{S}(U)^{-1}\mathcal{O}_X(U)$ はすでに
$j_*\mathcal{O}_Y(U)$ と、このアフィン開集合 $U$ 上で一致する。
したがって表示した写像は同型であり、
補題の主張における最初の表示等式が示された。

\medskip\noindent
最後に (1)、(3)、(4)、(5) を証明する。(5) については、証明の途中で
$\mathcal{K}_X = j_*\mathcal{O}_Y$ を見た。$j$ は準コンパクトである。
実際、$X$ の既約成分全体が局所有限であるという仮定による。
したがって $j$ は準コンパクトかつ準分離的である（$Y$ は分離的だからである）。
Schemes, Lemma \ref{schemes-lemma-push-forward-quasi-coherent} により、
$j_*\mathcal{O}_Y$ は準連接である。これで (1) が示された。
$x \in X$ とする。$U = \Spec(A)$ という $x$ のアフィン開近傍で、そのすべての
既約成分が $x$ を通るものを選べる。このとき $A \subset A_\mathfrak p$ である。
実際、$A$ の各弱随伴素イデアルは $\mathfrak p$ に含まれるので、
$A \setminus \mathfrak p$ の各元は Algebra, Lemma
\ref{algebra-lemma-weakly-ass-zero-divisors} により非零因子である。
これより、$A_\mathfrak p$ の任意の非零因子は、（必要ならより小さい）
$x$ のアフィン開近傍上の非零因子の像であることが容易に従う。これで (3) が示された。
(4) は茎を計算することにより (3) から従う。
\end{proof}

\begin{definition}
\label{divisors-definition-regular-meromorphic-section}
$X$ を局所環付き空間とし、$\mathcal{L}$ を可逆
$\mathcal{O}_X$-加群とする。有理型切断 $s$ が $\mathcal{L}$ のものとして
{\it 正則}であるとは、誘導写像
$\mathcal{K}_X \to \mathcal{K}_X(\mathcal{L})$ が単射であることをいう。
言い換えると、$s$ は可逆 $\mathcal{K}_X$-加群
$\mathcal{K}_X(\mathcal{L})$ の正則切断である。Definition
\ref{divisors-definition-regular-section} を参照。
\end{definition}

\noindent
（正則）有理型切断を引き戻せる条件を詳しく述べる。

\begin{lemma}
\label{divisors-lemma-meromorphic-sections-pullback}
$f : X \to Y$ を局所環付き空間の射とする。$f$ に対して有理型関数の
引き戻しが定義されていると仮定する（Definition
\ref{divisors-definition-pullback-meromorphic-sections} を参照）。
\begin{enumerate}
\item $\mathcal{F}$ を $\mathcal{O}_Y$-加群の層とする。
$f^* : \Gamma(Y, \mathcal{K}_Y(\mathcal{F})) \to
\Gamma(X, \mathcal{K}_X(f^*\mathcal{F}))$
という、$\mathcal{F}$ の有理型切断に対する標準引き戻し写像がある。
\item $\mathcal{L}$ を可逆 $\mathcal{O}_Y$-加群とする。
正則有理型切断 $s$ が $\mathcal{L}$ のものならば、その引き戻し
$f^*s$ は $f^*\mathcal{L}$ の正則有理型切断である。
\end{enumerate}
\end{lemma}

\begin{proof}
省略する。
\end{proof}

\begin{lemma}
\label{divisors-lemma-regular-meromorphic-ideal-denominators}
$X$ をスキームとする。$\mathcal{L}$ を可逆 $\mathcal{O}_X$-加群とする。
$s$ を $\mathcal{L}$ の正則有理型切断とする。
$\mathcal{I} \subset \mathcal{O}_X$ を、規則
$$
\mathcal{I}(V)
=
\{f \in \mathcal{O}_X(V) \mid fs \in \mathcal{L}(V)\}.
$$
により定まるイデアル層とする。
$\mathcal{L}(V) \subset \mathcal{K}_X(\mathcal{L})(V)$ なので、この式には意味がある。
このとき $\mathcal{I}$ は準連接イデアル層であり、単射
$$
1 : \mathcal{I} \longrightarrow \mathcal{O}_X,
\quad
s : \mathcal{I} \longrightarrow \mathcal{L}
$$
をもち、それらの余核は $X$ の閉疎部分集合上に台をもつ。
\end{lemma}

\begin{proof}
問題は $X$ 上局所的である。したがって $X = \Spec(A)$ かつ
$\mathcal{L} = \mathcal{O}_X$ と仮定してよい。さらに縮小して、
$s = a/b$ かつ $a, b \in A$ は {\it ともに} $A$ の非零因子であると仮定してよい。
$I = \{x \in A \mid x(a/b) \in A\}$ と置く。

\medskip\noindent
$\mathcal{I}$ が準連接であることを示すには、
$I_f = \{x \in A_f \mid x(a/b) \in A_f\}$ が任意の
$f \in A$ に対して成り立つことを示せばよい。
$c/f^n \in A_f$ かつ $(c/f^n)(a/b) \in A_f$ ならば、
$f^mc(a/b) \in A$ となる $m$ が存在し、したがって $c/f^n \in I_f$ である。
逆に、$I_f$ が $\{x \in A_f \mid x(a/b) \in A_f\}$ に含まれることは容易である。
これで準連接性が示された。

\medskip\noindent
最後の主張を証明する。$(b) \subset I$ は明らかである。したがって
$V(I) \subset V(b)$ は、$b$ が非零因子なので疎部分集合である。
ゆえに $1$ の余核は閉疎集合上に台をもつ。
$s$ の余核についても同じ議論が使える。実際
$s(b) = (a) \subset sI \subset A$ である。
\end{proof}

\begin{definition}
\label{divisors-definition-regular-meromorphic-ideal-denominators}
$X$ をスキームとし、$\mathcal{L}$ を可逆 $\mathcal{O}_X$-加群、
$s$ を $\mathcal{L}$ の正則有理型切断とする。
Lemma \ref{divisors-lemma-regular-meromorphic-ideal-denominators} で構成した
イデアル層 $\mathcal{I}$ を {\it $s$ の分母のイデアル層}という。
\end{definition}

\section{有理型関数と有理型切断：Noether の場合}
\label{divisors-section-meromorphic-noetherian}

\noindent
局所 Noether スキームについては、有理型関数の層に関するいくつかの結果を
証明できる。しかし \cite{misconceptions} には、$\mathcal{K}_X$ が
準連接であるとは限らない Noether スキーム $X$ の例がある。

\begin{lemma}
\label{divisors-lemma-meromorphic-section-restricts-to-zero}
$X$ を準コンパクトスキームとする。$h \in \Gamma(X, \mathcal{O}_X)$ および
$f \in \Gamma(X, \mathcal{K}_X)$ を、$f$ の $X_h$ への制限が零となるものとする。
このとき $h^n f = 0$ となる $n \gg 0$ が存在する。
\end{lemma}

\begin{proof}
$X$ のアフィン開集合 $U$ による被覆で、
$f|_U = s^{-1}a$、$a \in \mathcal{O}_X(U)$、$s \in \mathcal{S}(U)$
となるものを取れる。$X$ は準コンパクトなので、この形の有限個のアフィン開集合で
被覆できる。したがって $X = \Spec(A)$ かつ $f = s^{-1}a$ の場合に補題を
証明すれば十分である。$s \in A$ は非零因子であるから、$f = a$ の場合を
証明すれば十分である。条件 $f|_{X_h} = 0$ は、$a$ が
$A_h = \mathcal{O}_X(X_h)$ で零に写ることを意味する。ここでは
$\mathcal{O}_X \subset \mathcal{K}_X$ を用いた。したがって、望んだとおり
$h^na = 0$ となる $n > 0$ が存在する。
\end{proof}

\begin{lemma}
\label{divisors-lemma-locally-Noetherian-K}
$X$ を局所 Noether スキームとする。
\begin{enumerate}
\item 任意の $x \in X$ に対して、$\mathcal{S}_x \subset \mathcal{O}_{X, x}$ は
非零因子全体であり、したがって $\mathcal{K}_{X, x}$ は
$\mathcal{O}_{X, x}$ の全商環である。
\item 任意のアフィン開集合 $U \subset X$ に対して、環
$\mathcal{K}_X(U)$ は $\mathcal{O}_X(U)$ の全商環に等しい。
\end{enumerate}
\end{lemma}

\begin{proof}
この補題を証明するには、$X$ が Noether 環 $A$ のスペクトルであると仮定してよい。
$x \in X$ が $\mathfrak p \subset A$ に対応するとする。

\medskip\noindent
(1) の証明。$\mathcal{S}_x$ が
$\mathcal{O}_{X, x} = A_\mathfrak p$ の非零因子全体に含まれることは明らかである。
逆に、$f, g \in A$、$g \not \in \mathfrak p$ とし、
$f/g$ が $A_{\mathfrak p}$ の非零因子であると仮定する。
$I = \{a \in A \mid af = 0\}$ と置く。局所化の完全性から
$I_{\mathfrak p} = 0$ である。$A$ は Noether なので $I$ は有限生成であり、
したがって $g'I = 0$ となるある $g' \in A$、$g' \not \in \mathfrak p$
が存在する。よって $f$ は $A_{g'}$、すなわち $\mathfrak p$ の
Zariski 開近傍における非零因子である。したがって $f/g$ は
$\mathcal{S}_x$ の元である。

\medskip\noindent
(2) の証明。$f \in \Gamma(X, \mathcal{K}_X)$ を有理型関数とする。
$I = \{a \in A \mid af \in A\}$ と置く。素イデアル $\mathfrak p \subset A$ で、
点 $x \in X$ に対応するものを固定する。(1) により、$f$ の
$\mathfrak p$ における茎での像を $a/b$ と書ける。ここで
$a, b \in A_{\mathfrak p}$ であり、
$b \in A_{\mathfrak p}$ は零因子でない。$b = c/d$ と書く。ただし
$c, d \in A$、$d \not \in \mathfrak p$ である。このとき $ad - cf$ は
$\mathcal{K}_X$ の切断であり、$x$ のある開近傍で消える。
$D(e)$ 上で消えるとする。ここで $e \in A$、$e \not \in \mathfrak p$ である。
Lemma \ref{divisors-lemma-meromorphic-section-restricts-to-zero} により、
$e^n(ad - cf) = 0$ となる $n \gg 0$ が存在する。
したがって $e^nc \in I$ であり、$e^nc$ は $A_{\mathfrak p}$ で非零因子に写る。
$\text{Ass}(A) = \{\mathfrak q_1, \ldots, \mathfrak q_t\}$ を $A$ の
随伴素イデアル全体とする。$IA_{\mathfrak q_i}$ を調べ、
Algebra, Lemma \ref{algebra-lemma-associated-primes-localize} を用いると、
上の議論から $I \not \subset \mathfrak q_i$ であり、これは各 $i$ について成り立つ。
Algebra, Lemma \ref{algebra-lemma-silly} により、
$z \in I$ かつ $z \not \in \bigcup \mathfrak q_i$ である元が存在する。
Algebra, Lemma \ref{algebra-lemma-ass-zero-divisors} により、
$z$ は $A$ 上の零因子でない。したがって
$f = (zf)/z$ は $A$ の全商環の元である。これで (2) が示された。
\end{proof}

\begin{lemma}
\label{divisors-lemma-quasi-coherent-K}
$X$ を埋め込まれた点をもたない局所 Noether スキームとする。
$X^0$ を $X$ の既約成分の一般点全体とする。このとき
$$
\mathcal{K}_X =
\bigoplus\nolimits_{\eta \in X^0} j_{\eta, *}\mathcal{O}_{X, \eta} =
\prod\nolimits_{\eta \in X^0} j_{\eta, *}\mathcal{O}_{X, \eta}
$$
である。ここで $j_\eta : \Spec(\mathcal{O}_{X, \eta}) \to X$ は
Schemes, Section \ref{schemes-section-points} の標準写像である。さらに、
\begin{enumerate}
\item $\mathcal{K}_X$ は準連接な $\mathcal{O}_X$-代数の層である。
\item 任意の準連接 $\mathcal{O}_X$-加群 $\mathcal{F}$ に対し、
$$
\mathcal{K}_X(\mathcal{F}) =
\bigoplus\nolimits_{\eta \in X^0} j_{\eta, *}\mathcal{F}_\eta =
\prod\nolimits_{\eta \in X^0} j_{\eta, *}\mathcal{F}_\eta
$$
という $\mathcal{F}$ の有理型切断の層は準連接である。
\item $X$ の有理関数の環は $X$ 上の有理型関数の環であり、式で書けば
$R(X) = \Gamma(X, \mathcal{K}_X)$ である。
\end{enumerate}
\end{lemma}

\begin{proof}
これは Lemma \ref{divisors-lemma-meromorphic-weakass-finite} の特別な場合である。
実際、局所 Noether の場合、弱随伴点と随伴点は
Lemma \ref{divisors-lemma-ass-weakly-ass} により同じものである。
\end{proof}

\begin{lemma}
\label{divisors-lemma-regular-meromorphic-section-exists-noetherian}
$X$ を埋め込まれた点をもたない局所 Noether スキームとする。
$\mathcal{L}$ を可逆 $\mathcal{O}_X$-加群とする。このとき
$\mathcal{L}$ は正則有理型切断をもつ。
\end{lemma}

\begin{proof}
一般点 $\eta$ を $X$ の各一般点として取るごとに、生成元 $s_\eta$ を、
自由階数 $1$ の加群
$\mathcal{L}_\eta$ の生成元として選ぶ。この加群は Artin 局所環
$\mathcal{O}_{X, \eta}$ 上のものである。Lemma \ref{divisors-lemma-quasi-coherent-K} における
$\mathcal{K}_X$ と $\mathcal{K}_X(\mathcal{L})$ の記述から直ちに、
$s = \prod s_\eta$ は $\mathcal{L}$ の正則有理型切断である。
\end{proof}

\begin{lemma}
\label{divisors-lemma-make-maps-regular-section}
次が与えられているとする。
\begin{enumerate}
\item $X$ は局所 Noether スキームである。
\item $\mathcal{L}$ は可逆 $\mathcal{O}_X$-加群である。
\item $s$ は $\mathcal{L}$ の正則有理型切断である。
\item $\mathcal{F}$ は $X$ 上の連接加群であり、埋め込まれた随伴点をもたず、
$\text{Supp}(\mathcal{F}) = X$ である。
\end{enumerate}
$\mathcal{I} \subset \mathcal{O}_X$ を $s$ の分母のイデアルとする。
$T \subset X$ を $\mathcal{O}_X/\mathcal{I}$ と
$\mathcal{L}/s(\mathcal{I})$ の台の和集合とする。これは閉疎部分集合
$T \subset X$ であることが
Lemma \ref{divisors-lemma-regular-meromorphic-ideal-denominators} から従う。
このとき標準的な単射
$$
1 : \mathcal{I}\mathcal{F} \to \mathcal{F}, \quad
s : \mathcal{I}\mathcal{F} \to \mathcal{F} \otimes_{\mathcal{O}_X}\mathcal{L}
$$
があり、その余核は $T$ 上に台をもつ。
\end{lemma}

\begin{proof}
$\mathcal{L} \cong \mathcal{O}_X$ かつ $s = a/b$ で、
$a, b \in A$ がともに非零因子であるアフィンの場合に帰着する。
帰着の証明は省略する。$\mathcal{F} = \widetilde{M}$ と書く。
$I = \{x \in A \mid x(a/b) \in A\}$ と置けば、
$\mathcal{I} = \widetilde{I}$ である（Lemma
\ref{divisors-lemma-regular-meromorphic-ideal-denominators} の証明を参照）。
$T = V(I) \cup V((a/b)I)$ であることに注意する。任意の $A$-加群 $M$ に対し、
写像 $1 : IM \to M$ を考える。これが補題の写像 $1$ を与える写像である。
一方、写像 $\sigma : IM \to M_b, x \mapsto ax/b$ を考える。
$b$ は $A$ の零因子でなく、$M$ は台 $\Spec(A)$ をもち埋め込まれた素イデアルを
もたないので、$b$ は $M$ 上でも非零因子である。したがって $M \subset M_b$ である。
$I$ の定義により、$\sigma(IM) \subset M$ が $M_b$ の部分加群として成り立つ。
よって $A$-加群写像 $s : IM \to M$ を得る（すなわち、
$s(z)/1 = \sigma(z)$ が $M_b$ で成り立つ、すべての $z \in IM$ に対する
一意的な写像）。
$a$ も（$A$ と $M$ のいずれにおいても）非零因子なので、この写像は単射である。
$M/IM$ が $I$ で零化され、$M/s(IM)$ が $(a/b)I$ で零化されることは明らかである。
したがって補題が従う。
\end{proof}





\section{有理型関数と有理型切断：被約の場合}
\label{divisors-section-meromorphic-reduced}

\noindent
被約であって局所的に有限個の既約成分しかもたないスキームでは、
有理型関数の層は準連接である。

\begin{lemma}
\label{divisors-lemma-reduced-finite-irreducible}
$X$ を被約スキームで、任意の準コンパクト開集合が有限個の既約成分をもつものとする。
$X^0$ を $X$ の既約成分の一般点全体とする。このとき
$$
\mathcal{K}_X =
\bigoplus\nolimits_{\eta \in X^0} j_{\eta, *}\kappa(\eta) =
\prod\nolimits_{\eta \in X^0} j_{\eta, *}\kappa(\eta)
$$
である。ここで $j_\eta : \Spec(\kappa(\eta)) \to X$ は
Schemes, Section \ref{schemes-section-points} の標準写像である。さらに、
\begin{enumerate}
\item $\mathcal{K}_X$ は準連接な $\mathcal{O}_X$-代数の層である。
\item 任意の準連接 $\mathcal{O}_X$-加群 $\mathcal{F}$ に対して、
$$
\mathcal{K}_X(\mathcal{F}) =
\bigoplus\nolimits_{\eta \in X^0} j_{\eta, *}\mathcal{F}_\eta =
\prod\nolimits_{\eta \in X^0} j_{\eta, *}\mathcal{F}_\eta
$$
という $\mathcal{F}$ の有理型切断の層は準連接である。
\item $\mathcal{S}_x \subset \mathcal{O}_{X, x}$ は、任意の $x \in X$ に対して
非零因子全体である。
\item $\mathcal{K}_{X, x}$ は $\mathcal{O}_{X, x}$ の全商環であり、
これは任意の $x \in X$ について成り立つ。
\item $\mathcal{K}_X(U)$ は $\mathcal{O}_X(U)$ の全商環に等しく、
これは任意のアフィン開集合 $U \subset X$ について成り立つ。
\item $X$ の有理関数の環は $X$ 上の有理型関数の環であり、式で書けば
$R(X) = \Gamma(X, \mathcal{K}_X)$ である。
\end{enumerate}
\end{lemma}

\begin{proof}
この補題は Lemma \ref{divisors-lemma-meromorphic-weakass-finite} の特別な場合である。
実際、被約スキーム上では弱随伴点は Lemma
\ref{divisors-lemma-weakass-reduced} により一般点にほかならない。
\end{proof}

\begin{lemma}
\label{divisors-lemma-reduced-normalization}
$X$ をスキームとする。$X$ は被約であり、任意の準コンパクト開集合
$U \subset X$ は有限個の既約成分をもつと仮定する。このとき正規化射
$\nu : X^\nu \to X$ は射
$$
\underline{\Spec}_X(\mathcal{O}') \longrightarrow X
$$
である。ここで $\mathcal{O}' \subset \mathcal{K}_X$ は、有理型関数の層における
$\mathcal{O}_X$ の整閉包である。
\end{lemma}

\begin{proof}
正規化射 $\nu : X^\nu \to X$ の定義（Morphisms, Definition
\ref{morphisms-definition-normalization} を参照）と、上の Lemma
\ref{divisors-lemma-reduced-finite-irreducible} における $\mathcal{K}_X$ の記述を比較すればよい。
\end{proof}

\begin{lemma}
\label{divisors-lemma-meromorphic-functions-integral-scheme}
$X$ を一般点 $\eta$ をもつ整スキームとする。次が成り立つ。
\begin{enumerate}
\item 有理型関数の層は、$X$ の関数体（Morphisms, Definition
\ref{morphisms-definition-function-field} を参照）を値とする定数層と同型である。
\item 任意の準連接層 $\mathcal{F}$ で $X$ 上のものに対して、層
$\mathcal{K}_X(\mathcal{F})$ は $\mathcal{F}_\eta$ を値とする定数層と同型である。
\end{enumerate}
\end{lemma}

\begin{proof}
省略する。
\end{proof}

\noindent
場合によっては正則有理型切断の存在を示せる。

\begin{lemma}
\label{divisors-lemma-regular-meromorphic-section-exists}
$X$ をスキームとする。$\mathcal{L}$ を可逆 $\mathcal{O}_X$-加群とする。
次の各場合に $\mathcal{L}$ は正則有理型切断をもつ。
\begin{enumerate}
\item $X$ は整である。
\item $X$ は被約であり、任意の準コンパクト開集合は有限個の既約成分をもつ。
\item $X$ は局所 Noether であり、埋め込まれた点をもたない。
\end{enumerate}
\end{lemma}

\begin{proof}
(1) の場合、$\eta \in X$ を一般点とする。Lemma
\ref{divisors-lemma-meromorphic-functions-integral-scheme} により、
$\mathcal{K}_X$ および $\mathcal{K}_X(\mathcal{L})$ はそれぞれ
$\kappa(\eta)$ および $\mathcal{L}_\eta$ を値とする定数層である。
$\dim_{\kappa(\eta)} \mathcal{L}_\eta = 1$ なので、非零元
$s \in \mathcal{L}_\eta$ を選べる。明らかに $s$ は $\mathcal{L}$ の
正則有理型切断である。(2) の場合、非零な
$s_\eta \in \mathcal{L}_\eta$ を、一般点 $\eta$、すなわち $X$ のすべての
一般点について選ぶ。これは $\mathcal{L}_\eta$ が $1$ 次元の
$\kappa(\eta)$-ベクトル空間なので可能である。
Lemma \ref{divisors-lemma-reduced-finite-irreducible} における
$\mathcal{K}_X$ と $\mathcal{K}_X(\mathcal{L})$ の記述から直ちに、
$s = \prod s_\eta$ は $\mathcal{L}$ の正則有理型切断である。
(3) は Lemma \ref{divisors-lemma-regular-meromorphic-section-exists-noetherian} である。
\end{proof}






\section{Weil 因子}
\label{divisors-section-Weil-divisors}

\noindent
局所 Noether 整スキームに対して、Weil 因子および Weil 因子の有理同値を導入する。
スキームが準コンパクトであるとは仮定しないので、Weil 因子の定義には少し注意が
必要である。たとえば有理関数は無限個の極をもちうるので、素因子の無限和を
許さなければならない。準コンパクトスキームでは、通常どおり Weil 因子は有限和に
なる。まず、閉部分スキームの族が局所有限であることを示す際に頻繁に用いる
基本的な補題を述べる。

\begin{lemma}
\label{divisors-lemma-components-locally-finite}
$X$ を局所 Noether スキームとする。$Z \subset X$ を閉部分スキームとする。
$Z$ の既約成分全体は $X$ において局所有限である。
\end{lemma}

\begin{proof}
$U \subset X$ を準コンパクトな開部分スキームとする。すると $U$ は Noether
スキームであり、したがってその基礎位相空間は Noether である
（Properties, Lemma \ref{properties-lemma-Noetherian-topology}）。
よってすべての部分空間は Noether であり、有限個の既約成分しかもたない
（Topology, Lemma \ref{topology-lemma-Noetherian} を参照）。
\end{proof}

\noindent
$Z$ がスキーム $X$ の既約閉部分集合ならば、$Z$ の $X$ における余次元は、
局所環 $\mathcal{O}_{X, \xi}$ の次元に等しい。ここで $\xi \in Z$ は一般点である。
このことを思い出そう。Properties, Lemma
\ref{properties-lemma-codimension-local-ring} を参照。

\begin{definition}
\label{divisors-definition-Weil-divisor}
$X$ を局所 Noether 整スキームとする。
\begin{enumerate}
\item {\it 素因子}とは、整閉部分スキーム $Z \subset X$ で余次元 $1$ のものである。
\item {\it Weil 因子}とは形式和 $D = \sum n_Z Z$ であって、和が $X$ の
素因子全体を走り、族 $\{Z \mid n_Z \not = 0\}$ が局所有限であるものをいう
（Topology, Definition \ref{topology-definition-locally-finite}）。
\end{enumerate}
$X$ 上の Weil 因子全体の群を $\text{Div}(X)$ と書く。
\end{definition}

\noindent
次に、有理関数に付随する Weil 因子を定義する。そのため、素因子に沿う
有理関数の零点位数を用いる。これは次のように定義される。

\begin{definition}
\label{divisors-definition-order-vanishing}
$X$ を局所 Noether 整スキームとする。$f \in R(X)^*$ とする。
各素因子 $Z \subset X$ に対して、{\it $f$ の $Z$ に沿う零点位数}を整数
$$
\text{ord}_Z(f) = \text{ord}_{\mathcal{O}_{X, \xi}}(f)
$$
として定義する。右辺は Algebra, Definition \ref{algebra-definition-ord} の概念であり、
$\xi$ は $Z$ の一般点である。
\end{definition}

\noindent
$f, g \in R(X)^*$ に対して
$$
\text{ord}_Z(fg) = \text{ord}_Z(f) + \text{ord}_Z(g).
$$
であることに注意する。もちろん $\text{ord}_Z(f) < 0$ となることもある。
この場合、$f$ は $Z$ に沿って {\it 極}をもつといい、
$-\text{ord}_Z(f) > 0$ を {\it $f$ の $Z$ に沿う極の位数}という。
条件 $\text{ord}_Z(f) \geq 0$ は、条件 $f \in \mathcal{O}_{X, \xi}$ と
{\bf 同値ではない}ことが重要である。ただし、局所環
$\mathcal{O}_{X, \xi}$ が離散付値環である場合を除く。

\begin{lemma}
\label{divisors-lemma-divisor-locally-finite}
$X$ を局所 Noether 整スキームとする。$f \in R(X)^*$ とする。族
$$
\{Z \subset X \mid Z\text{ a prime divisor with generic point }\xi
\text{ and }f\text{ not in }\mathcal{O}_{X, \xi}\}
$$
および
$$
\{Z \subset X \mid Z \text{ a prime divisor and }\text{ord}_Z(f) \not = 0\}
$$
は $X$ において局所有限である。
\end{lemma}

\begin{proof}
空でない開部分スキーム $U \subset X$ で、$f$ が
$\Gamma(U, \mathcal{O}_X^*)$ の切断に対応するものが存在する。
したがって補題の集合に現れうる素因子はすべて $X \setminus U$ の既約成分である。
よって Lemma \ref{divisors-lemma-components-locally-finite} が望む結果を与える。
\end{proof}

\noindent
この補題により、次の定義を置くことができる。

\begin{definition}
\label{divisors-definition-principal-divisor}
$X$ を局所 Noether 整スキームとする。$f \in R(X)^*$ とする。
{\it $f$ に付随する主 Weil 因子}とは Weil 因子
$$
\text{div}(f) = \text{div}_X(f) = \sum \text{ord}_Z(f) [Z]
$$
である。ここで和は素因子全体を走り、$\text{ord}_Z(f)$ は
Definition \ref{divisors-definition-order-vanishing} のとおりである。
Lemma \ref{divisors-lemma-divisor-locally-finite} によりこれは意味をもつ。
\end{definition}

\begin{lemma}
\label{divisors-lemma-div-additive}
$X$ を局所 Noether 整スキームとする。$f, g \in R(X)^*$ とする。このとき
$$
\text{div}_X(fg) = \text{div}_X(f) + \text{div}_X(g)
$$
が $X$ 上の Weil 因子として成り立つ。
\end{lemma}

\begin{proof}
$\text{ord}$ 関数の加法性から明らかである。
\end{proof}

\noindent
上の補題から、主 Weil 因子全体は Weil 因子全体の群の部分群をなす。
これにより次の定義に至る。

\begin{definition}
\label{divisors-definition-class-group}
$X$ を局所 Noether 整スキームとする。$X$ の {\it Weil 因子類群}とは、
Weil 因子の群を主 Weil 因子の部分群で割った商である。記号は
$\text{Cl}(X)$ とする。
\end{definition}

\noindent
構成により完全複体
\begin{equation}
\label{divisors-equation-Weil-divisor-class}
R(X)^* \xrightarrow{\text{div}} \text{Div}(X) \to \text{Cl}(X) \to 0
\end{equation}
を得る。これは $\text{Cl}(X)$ の表示と考えられる。次に Weil 因子類群と
Picard 群を関係づける。





\section{可逆加群に付随する Weil 因子類}
\label{divisors-section-c1}

\noindent
この節では Section \ref{divisors-section-Weil-divisors} とまったく同じ段階をたどり、
局所 Noether 整スキーム上の標準写像
$\Pic(X) \to \text{Cl}(X)$
を定義する。

\medskip\noindent
$X$ をスキームとする。$\mathcal{L}$ を可逆 $\mathcal{O}_X$-加群とする。
$\xi \in X$ を点とする。$s_\xi, s'_\xi \in \mathcal{L}_\xi$ が
$\mathcal{L}_\xi$ を $\mathcal{O}_{X, \xi}$-加群として生成するならば、
単元 $u \in \mathcal{O}_{X, \xi}^*$ で $s_\xi = u s'_\xi$ となるものが存在する。
有理型切断の層 $\mathcal{K}_X(\mathcal{L})$、すなわち $\mathcal{L}$ の
有理型切断の層の $x$ における茎は
$\mathcal{K}_{X, x} \otimes_{\mathcal{O}_{X, x}} \mathcal{L}_x$ に等しい。
したがって、任意の有理型切断 $s$ で $\mathcal{L}$ のものの $x$ における
茎での像は、$s = fs_\xi$ と書ける。ここで
$f \in \mathcal{K}_{X, x}$ である。以下ではこれを
$f = s/s_\xi$ と略記する。また、$X$ が整ならば
$\mathcal{K}_{X, x} = R(X)$ は $X$ の関数体に等しいので
$s/s_\xi \in R(X)$ である。$s$ が正則有理型切断ならば、実際には
$s/s_\xi \in R(X)^*$ である。整スキーム上では、正則有理型切断であることは
非零有理型切断であることと同じである。最後に、$s/s_\xi$ は $s_\xi$ の選択に、
局所環 $\mathcal{O}_{X, x}$ の単元による乗法を除いて依存しない。
以上を合わせると、次の定義には意味がある。

\begin{definition}
\label{divisors-definition-order-vanishing-meromorphic}
$X$ を局所 Noether 整スキームとする。
$\mathcal{L}$ を可逆 $\mathcal{O}_X$-加群とする。
$s \in \Gamma(X, \mathcal{K}_X(\mathcal{L}))$ を $\mathcal{L}$ の
正則有理型切断とする。各素因子 $Z \subset X$ に対して、
{\it $s$ の $Z$ に沿う零点位数}を整数
$$
\text{ord}_{Z, \mathcal{L}}(s)
= \text{ord}_{\mathcal{O}_{X, \xi}}(s/s_\xi)
$$
として定義する。右辺は Algebra, Definition \ref{algebra-definition-ord} の概念であり、
$\xi \in Z$ は一般点、$s_\xi \in \mathcal{L}_\xi$ は生成元である。
\end{definition}

\noindent
主因子の場合と同様に、次の補題が成り立つ。

\begin{lemma}
\label{divisors-lemma-divisor-meromorphic-locally-finite}
$X$ を局所 Noether 整スキームとする。$\mathcal{L}$ を可逆
$\mathcal{O}_X$-加群とする。$s \in \mathcal{K}_X(\mathcal{L})$ を
$\mathcal{L}$ の正則（すなわち非零）有理型切断とする。このとき集合
$$
\{Z \subset X \mid Z \text{ a prime divisor with generic point }\xi
\text{ and }s\text{ not in }\mathcal{L}_\xi\}
$$
および
$$
\{Z \subset X \mid Z \text{ is a prime divisor and }
\text{ord}_{Z, \mathcal{L}}(s) \not = 0\}
$$
は $X$ において局所有限である。
\end{lemma}

\begin{proof}
空でない開部分スキーム $U \subset X$ で、$s$ が
$\Gamma(U, \mathcal{L})$ の切断に対応し、その切断が $\mathcal{L}$ を
$U$ 上生成するものが存在する。したがって補題の集合に現れうる素因子は
すべて $X \setminus U$ の既約成分である。よって Lemma
\ref{divisors-lemma-components-locally-finite} が望む結果を与える。
\end{proof}

\begin{lemma}
\label{divisors-lemma-divisor-meromorphic-well-defined}
$X$ を局所 Noether 整スキームとする。
$\mathcal{L}$ を可逆 $\mathcal{O}_X$-加群とする。
$s, s' \in \mathcal{K}_X(\mathcal{L})$ を $\mathcal{L}$ の
非零有理型切断とする。このとき $f = s/s'$ は $R(X)^*$ の元であり、
$$
\sum \text{ord}_{Z, \mathcal{L}}(s)[Z]
=
\sum \text{ord}_{Z, \mathcal{L}}(s')[Z]
+
\text{div}(f)
$$
が Weil 因子として成り立つ。
\end{lemma}

\begin{proof}
定義から明らかである。Lemma
\ref{divisors-lemma-divisor-meromorphic-locally-finite} により、これらの和が実際に
Weil 因子であることに注意する。
\end{proof}

\begin{definition}
\label{divisors-definition-divisor-invertible-sheaf}
$X$ を局所 Noether 整スキームとする。
$\mathcal{L}$ を可逆 $\mathcal{O}_X$-加群とする。
\begin{enumerate}
\item 任意の非零有理型切断 $s$ で $\mathcal{L}$ のものに対して、
{\it $s$ に付随する Weil 因子}を
$$
\text{div}_\mathcal{L}(s) =
\sum \text{ord}_{Z, \mathcal{L}}(s) [Z] \in \text{Div}(X)
$$
と定義する。ここで和は素因子全体を走る。
\item {\it $\mathcal{L}$ に付随する Weil 因子類}を、
$\text{div}_\mathcal{L}(s)$ の $\text{Cl}(X)$ における像として定義する。
ここで $s$ は $\mathcal{L}$ の任意の非零有理型切断で、$X$ 上のものである。
Lemma \ref{divisors-lemma-divisor-meromorphic-well-defined} により、これは良定義である。
\end{enumerate}
\end{definition}

\noindent
予想どおり、この構成は可逆加群に関して加法的である。

\begin{lemma}
\label{divisors-lemma-c1-additive}
$X$ を局所 Noether 整スキームとする。
$\mathcal{L}$、$\mathcal{N}$ を可逆 $\mathcal{O}_X$-加群とする。
$s$ および $t$ を、それぞれ $\mathcal{L}$ および $\mathcal{N}$ の
非零有理型切断とする。このとき $st$ は
$\mathcal{L} \otimes \mathcal{N}$ の非零有理型切断であり、
$$
\text{div}_{\mathcal{L} \otimes \mathcal{N}}(st)
=
\text{div}_\mathcal{L}(s) + \text{div}_\mathcal{N}(t)
$$
が $\text{Div}(X)$ において成り立つ。とくに
$\mathcal{L} \otimes_{\mathcal{O}_X} \mathcal{N}$ の Weil 因子類は、
$\mathcal{L}$ と $\mathcal{N}$ の Weil 因子類の和である。
\end{lemma}

\begin{proof}
$s$ および $t$ を、それぞれ $\mathcal{L}$ および $\mathcal{N}$ の
非零有理型切断とする。このとき $st$ は
$\mathcal{L} \otimes \mathcal{N}$ の非零有理型切断である。
$Z \subset X$ を素因子とする。$\xi \in Z$ をその一般点とする。
生成元 $s_\xi \in \mathcal{L}_\xi$ および
$t_\xi \in \mathcal{N}_\xi$ を選ぶ。すると $s_\xi t_\xi$ は
$(\mathcal{L} \otimes \mathcal{N})_\xi$ の生成元である。
したがって $st/(s_\xi t_\xi) = (s/s_\xi)(t/t_\xi)$ である。よって
$$
\text{ord}_{Z, \mathcal{L} \otimes \mathcal{N}}(st)
=
\text{ord}_{Z, \mathcal{L}}(s) + \text{ord}_{Z, \mathcal{N}}(t)
$$
であることが $\text{ord}_Z$ 関数の加法性から分かる。
\end{proof}

\noindent
このようにして、可逆加群にその Weil 因子類を対応させる Abel 群の準同型
\begin{equation}
\label{divisors-equation-c1}
\Pic(X) \longrightarrow \text{Cl}(X)
\end{equation}
を得る。

\begin{lemma}
\label{divisors-lemma-normal-c1-injective}
$X$ を局所 Noether 整スキームとする。$X$ が正規ならば、写像
(\ref{divisors-equation-c1}) $\Pic(X) \to \text{Cl}(X)$ は単射である。
\end{lemma}

\begin{proof}
$\mathcal{L}$ を可逆 $\mathcal{O}_X$-加群で、付随する Weil 因子類が自明なものとする。
$s$ を $\mathcal{L}$ の正則有理型切断とする。仮定は
$\text{div}_\mathcal{L}(s) = \text{div}(f)$ が、ある
$f \in R(X)^*$ に対して成り立つことを意味する。このとき
$t = f^{-1}s$ は $\mathcal{L}$ の正則有理型切断であり、
$\text{div}_\mathcal{L}(t) = 0$ である。Lemma
\ref{divisors-lemma-divisor-meromorphic-well-defined} を参照。
$t$ が $\mathcal{L}$ の自明化を定めることを示す。これで補題の証明が終わる。
これを証明するため $X$ 上局所的に議論してよい。したがって
$X = \Spec(A)$ がアフィンであり、$\mathcal{L}$ が自明であると仮定してよい。
このとき $A$ は Noether 正規整域であり、$t$ はその分数体の元で、
$\text{ord}_{A_\mathfrak p}(t) = 0$ が高さ $1$ のすべての素イデアル
$\mathfrak p$ で $A$ のものに対して成り立つ。目標は、$t$ が $A$ の単元であることを
示すことである。高さ一の素イデアルに対して $A_\mathfrak p$ は
離散付値環である。ここでその素イデアルは $A$ のものである
（Algebra, Lemma \ref{algebra-lemma-criterion-normal}）。したがってこの条件は、
$t \in A_\mathfrak p^*$ がすべての素イデアル $\mathfrak p$ で高さ $1$
のものに対して成り立つことを意味する。
これより Algebra, Lemma
\ref{algebra-lemma-normal-domain-intersection-localizations-height-1} によって
$t \in A$ かつ $t^{-1} \in A$ となり、証明が完了する。
\end{proof}

\begin{lemma}
\label{divisors-lemma-local-rings-UFD-c1-bijective}
$X$ を局所 Noether 整スキームとする。写像
(\ref{divisors-equation-c1}) $\Pic(X) \to \text{Cl}(X)$ を考える。
次は同値である。
\begin{enumerate}
\item $X$ の局所環は UFD である。
\item $X$ は正規であり、$\Pic(X) \to \text{Cl}(X)$ は全射である。
\end{enumerate}
この場合 $\Pic(X) \to \text{Cl}(X)$ は同型である。
\end{lemma}

\begin{proof}
(1) が成り立つならば、Algebra, Lemma
\ref{algebra-lemma-UFD-normal} により $X$ は正規である。
したがって写像 (\ref{divisors-equation-c1}) は Lemma
\ref{divisors-lemma-normal-c1-injective} により単射である。さらに、各素因子
$D \subset X$ は Lemma \ref{divisors-lemma-weil-divisor-is-cartier-UFD} により
有効 Cartier 因子である。この場合、標準切断 $1_D$（$\mathcal{O}_X(D)$ のもの）
（Definition \ref{divisors-definition-invertible-sheaf-effective-Cartier-divisor}）は
$D$ に沿ってちょうど消え、$D$ の類は写像 (\ref{divisors-equation-c1}) による
$\mathcal{O}_X(D)$ の像である。したがって写像は全射でもある。

\medskip\noindent
(2) が成り立つと仮定する。素因子 $D \subset X$ を選ぶ。
(\ref{divisors-equation-c1}) は全射なので、可逆層 $\mathcal{L}$、正則有理型切断 $s$、
および $f \in R(X)^*$ で
$\text{div}_\mathcal{L}(s) + \text{div}(f) = [D]$ となるものが存在する。
言い換えると $\text{div}_\mathcal{L}(fs) = [D]$ である。
$x \in X$ とし、$A = \mathcal{O}_{X, x}$ とする。すると $A$ は分数体
$K = R(X)$ をもつ Noether 局所正規整域である。高さ $1$ の各素イデアルで
$A$ のものは
$X$ 上の素因子に対応し、すべての可逆 $\mathcal{O}_X$-加群は
$\Spec(A)$ 上で自明な可逆加群に制限される。したがって、すべての高さ $1$ の
素イデアル $\mathfrak p \subset A$ に対して、$f \in K$ で
$\text{ord}_{A_\mathfrak p}(f) = 1$ かつ、他のすべての高さ一の素イデアル
$\text{ord}_{A_{\mathfrak p'}}(f) = 0$ が成り立つ他のすべての高さ一の素イデアル
$\mathfrak p'$ があるようなものが存在する。
Algebra, Lemma
\ref{algebra-lemma-normal-domain-intersection-localizations-height-1} により
$f \in A$ である。同じように論じると、すべての元 $g \in \mathfrak p$ は、
$g = af$ と書ける。ここである $a \in A$ を取れる。したがって $A$ の高さ一の
各素イデアルは単項であり、Algebra, Lemma
\ref{algebra-lemma-characterize-UFD-height-1} により $A$ は UFD である。
\end{proof}







\section{可逆加群についてさらに}
\label{divisors-section-invertible-modules}

\noindent
この節では可逆加群のいくつかの性質を論じる。

\begin{lemma}
\label{divisors-lemma-in-image-pullback}
$\varphi : X \to Y$ をスキームの射とする。$\mathcal{L}$ を可逆
$\mathcal{O}_X$-加群とする。次を仮定する。
\begin{enumerate}
\item $X$ は局所 Noether である。
\item $Y$ は局所 Noether、整かつ正規である。
\item $\varphi$ は整（したがって空でない）ファイバーをもつ平坦射である。
\item $\varphi$ は準コンパクトまたは局所有限型である。
\item $\mathcal{L}$ の $\varphi$ の一般ファイバーへの制限は自明である。
\end{enumerate}
このとき、$\mathcal{L} \cong \varphi^*\mathcal{N}$ を満たす可逆
$\mathcal{O}_Y$-加群 $\mathcal{N}$ が存在する。
\end{lemma}

\begin{proof}
$\xi \in Y$ を一般点とする。$X_\xi$ を $\varphi$ の $\xi$ 上の
スキーム論的ファイバーとする。$\mathcal{L}_\xi$ により
$\mathcal{L}$ の $X_\xi$ への引き戻しを表す。仮定 (5) は
$\mathcal{L}_\xi$ が自明であることを意味する。自明化切断
$s \in \Gamma(X_\xi, \mathcal{L}_\xi)$ を選ぶ。Lemma
\ref{divisors-lemma-flat-over-integral-integral-fibre} により $X$ は整であることに注意する。
したがって $s$ を $\mathcal{L}$ の正則有理型切断とみなせる。
Lemma \ref{divisors-lemma-pullback-meromorphic-sections-defined} により、
$\varphi$ に対して有理型関数の引き戻しが定義されている。
$\mathcal{N} \subset \mathcal{K}_Y$ を次の $\mathcal{O}_Y$-加群とする。
開集合 $V \subset Y$ 上の切断は、有理型関数
$g \in \mathcal{K}_Y(V)$ のうち
$\varphi^*(g)s \in \mathcal{L}(\varphi^{-1}V)$ となるものである。先験的には
$\varphi^*(g)s$ は $\mathcal{K}_X(\mathcal{L})$ の
$\varphi^{-1}V$ 上の切断である。$\mathcal{N}$ は可逆
$\mathcal{O}_Y$-加群であり、写像
$$
\varphi^*\mathcal{N} \longrightarrow \mathcal{L},\quad
g \longmapsto gs
$$
は同型であると主張する。

\medskip\noindent
まず次の状況で主張を証明する。$X$ と $Y$ はアフィンであり、
$\mathcal{L}$ は自明である。$Y = \Spec(R)$、$X = \Spec(A)$ とし、
$s$ は元 $s \in A \otimes_R K$ により与えられるとする。ここで $K$ は
$R$ の分数体である。$s = a/r$ と書ける。ここで非零元
$r \in R$ と元 $a \in A$ を取れる。$s$ は一般ファイバー上で $\mathcal{L}$ を生成するので、
$s' \in A \otimes_R K$ で $ss' = 1$ となるものが存在する。したがって
$s = r'/a'$ と書ける。ここで $r' \in R$ は非零で、$a' \in A$ である。
$\mathfrak p_1, \ldots, \mathfrak p_n \subset R$ を $rr'$ 上の極小素イデアルとする。
各 $R_{\mathfrak p_i}$ は離散付値環である（Algebra, Lemmas
\ref{algebra-lemma-minimal-over-1} および
\ref{algebra-lemma-criterion-normal}）。仮定により
$\mathfrak q_i = \mathfrak p_i A$ は素イデアルである。したがって
$\mathfrak q_i A_{\mathfrak q_i}$ は一元で生成され、$A_{\mathfrak q_i}$ も
離散付値環であることが分かる（Algebra, Lemma
\ref{algebra-lemma-characterize-dvr}）。もちろん
$R_{\mathfrak p_i} \to A_{\mathfrak q_i}$ の分岐指数は $1$ である。
$e_i, e'_i \geq 0$ を $a, a'$ の $A_{\mathfrak q_i}$ における付値とする。
すると $e_i + e'_i$ は $rr'$ の $R_{\mathfrak p_i}$ における付値である。
Algebra, Lemma
\ref{algebra-lemma-normal-domain-intersection-localizations-height-1} により、$R$ において
$$
\mathfrak p_1^{(e_1 + e'_1)} \cap \ldots \cap \mathfrak p_i^{(e_n + e'_n)} =
(rr')
$$
であることに注意する。
$$
I = \mathfrak p_1^{(e_1)} \cap \ldots \cap \mathfrak p_i^{(e_n)}
\quad\text{and}\quad
I' = \mathfrak p_1^{(e'_1)} \cap \ldots \cap \mathfrak p_i^{(e'_n)}
$$
と置くと、$II' \subset (rr')$ である。Algebra, Lemmas
\ref{algebra-lemma-symbolic-power-flat-extension} および
\ref{algebra-lemma-flat-intersect-ideals} により
$$
IA =
(\mathfrak p_1^{(e_1)} \cap \ldots \cap \mathfrak p_i^{(e_n)})A =
(\mathfrak p_1A)^{(e_1)} \cap \ldots \cap (\mathfrak p_i A)^{(e_n)}
$$
であることに注意する。$I'A$ についても同様である。したがって
$a \in IA$ かつ $a' \in I'A$ である。これより
$IA \otimes_A I'A \to rr'A$ は全射である。$R \to A$ の忠実平坦性から
$I \otimes_R I' \to (rr')$ も全射であることが分かる。
したがって $II' = (rr')$ であり、$I$ と $I'$ は有限局所自由階数 $1$ である。
Algebra, Lemma \ref{algebra-lemma-product-ideals-principal} を参照。
よって $R$ 上 Zariski 局所的に $I = (g)$ および $I' = (g')$ と書け、
$gg' = rr'$ である。このとき $a = ug$ かつ $a' = u'g'$ であり、
ここで $u, u' \in A$ である。$u, u'$ は単元であると結論する。
したがって $R$ 上 Zariski 局所的に $s = ug/r$ であり、この場合の主張が従う。

\medskip\noindent
$y \in Y$ を点とする。$x \in X$ を $y$ に写るように選ぶ。前段落の結果を
$\Spec(\mathcal{O}_{X, x}) \to \Spec(\mathcal{O}_{Y, y})$ に適用できる。
これより、$g \in R(Y)^*$ で、$\mathcal{O}_{Y, y}^*$ の元による乗法を除いて
定まり、$\varphi^*(g)s$ が $\mathcal{L}_x$ を生成するものが存在する。
したがって $\varphi^*(g)s$ は $\mathcal{L}$ をある近傍 $U$ 上で生成し、この近傍は $x$ を含む。
$x'$ を $y$ 上にある別の点とし、$g' \in R(Y)^*$ を
$\varphi^*(g')s$ が $\mathcal{L}$ をある開近傍 $U'$ で生成し、その近傍が $x'$ を含むものとする。
ファイバーは既約なので、
$x''$ を $U \cap U' \cap \varphi^{-1}(\{y\})$ の点として選べる。
環準同型 $\mathcal{O}_{Y, y} \to \mathcal{O}_{X, x''}$ に対する一意性から、
$g$ と $g'$ は $\mathcal{O}_{Y, y}^*$ の元を乗法的に掛けた分だけ異なる。
したがって $\varphi^*(g)s$ は $\mathcal{L}$ を
$\varphi^{-1}(y)$ のある開近傍上で生成する。$Z \subset X$ を、点
$z \in X$ のうち $\varphi^*(g)s$ が $\mathcal{L}_z$ を生成しないもの全体とする。
上の議論から $Z$ は閉であり、$Z = \varphi^{-1}(T)$ であるような部分集合
$T \subset Y$ が存在して $y \not \in T$ である。$T$ が閉であることを示せれば、
$g$ は $\mathcal{N}$ を $\mathcal{O}_Y$-加群として、
$Y \setminus T$ という $y$ の開近傍上で生成し、証明が終わる
（詳細の一部は省略する）。

\medskip\noindent
$\varphi$ が準コンパクトならば、Morphisms, Lemma
\ref{morphisms-lemma-fpqc-quotient-topology} により $T$ は閉である。
$\varphi$ が局所有限型ならば、Morphisms, Lemma
\ref{morphisms-lemma-fppf-open} により $\varphi$ は開写像である。
このとき $Y \setminus T$ は開集合 $X \setminus Z$ の像として開である。
\end{proof}

\begin{lemma}
\label{divisors-lemma-closure-effective-cartier-divisor}
$X$ を局所 Noether スキームとする。$U \subset X$ を開集合とし、
$D \subset U$ を有効 Cartier 因子とする。すべての点で
$\mathcal{O}_{X, x}$ が UFD であり、その点が $x \in X \setminus U$ を満たすならば、
有効 Cartier 因子 $D' \subset X$ で $D = U \cap D'$ となるものが存在する。
\end{lemma}

\begin{proof}
$D' \subset X$ を射 $D \to X$ のスキーム論的像とする。
$X$ は局所 Noether なので、射 $D \to X$ は準コンパクトである。
Properties, Lemma \ref{properties-lemma-immersion-into-noetherian} を参照。
したがって、$D'$ の形成は Morphisms, Lemma
\ref{morphisms-lemma-quasi-compact-scheme-theoretic-image} により
$X$ の開集合への移行と可換である。よって $X = \Spec(A)$ がアフィンであると
仮定してよい。$I \subset A$ を $D'$ に対応するイデアルとする。
$\mathfrak p \subset A$ を $X \setminus U$ の点に対応する素イデアルとする。
証明を終えるには、$I_\mathfrak p$ が一元で生成されることを示せば十分である。
Lemma \ref{divisors-lemma-effective-Cartier-in-points} を参照。したがって $X$ を
$\Spec(A_\mathfrak p)$ で置き換えてよい。Morphisms, Lemma
\ref{morphisms-lemma-flat-base-change-scheme-theoretic-image} を参照。
言い換えると、$X$ は局所 UFD $A$ のスペクトルであると仮定してよい。
すると $A$ のすべての局所環は UFD である。Lemma
\ref{divisors-lemma-effective-Cartier-divisor-is-a-sum} により
$D = \sum a_i D_i$ であり、$D_i \subset U$ は整な有効 Cartier 因子である。
$\xi_i$ を $D_i$ の一般点とすると、それは素イデアル
$\mathfrak p_i \subset A$ で高さ $1$ のものに対応する。Lemma
\ref{divisors-lemma-effective-Cartier-codimension-1} を参照。したがって、ある素元
$\mathfrak p_i = (f_i)$ となる素元 $f_i \in A$ があり、望んだとおり
$D'$ は $\prod f_i^{a_i}$ により切り出されると結論する。
\end{proof}

\begin{lemma}
\label{divisors-lemma-extend-invertible-module}
$X$ を局所 Noether スキームとする。$U \subset X$ を開集合とし、
$\mathcal{L}$ を可逆 $\mathcal{O}_U$-加群とする。すべての局所環
$\mathcal{O}_{X, x}$ が UFD であり、その点が $x \in X \setminus U$ を満たすならば、
可逆 $\mathcal{O}_X$-加群 $\mathcal{L}'$ で
$\mathcal{L} \cong \mathcal{L}'|_U$ となるものが存在する。
\end{lemma}

\begin{proof}
$x \in X$、$x \not \in U$ を選ぶ。アフィン開近傍 $W \subset X$ で、
$\mathcal{L}|_{W \cap U}$ が $W$ 上の可逆層へ延長されるものが存在することを示す。
これより層の貼り合わせ（Sheaves, Section
\ref{sheaves-section-glueing-sheaves}）によって、$\mathcal{L}$ を真に大きい開集合
$U \cup W$ へ延長できる。$W = \Spec(A)$ をアフィン開近傍とする。
$U \cap W$ は準アフィンなので、Lemma
\ref{divisors-lemma-quasi-projective-Noetherian-pic-effective-Cartier} により
$\mathcal{L}|_{W \cap U}$ を
$\mathcal{O}(D_1) \otimes \mathcal{O}(D_2)^{\otimes -1}$ と書ける。ここで
$D_1, D_2 \subset W \cap U$ は有効 Cartier 因子である。
すると $D_1$ と $D_2$ は Lemma
\ref{divisors-lemma-closure-effective-cartier-divisor} により $W$ の有効 Cartier 因子へ
延長され、これが可逆層の延長を与える。

\medskip\noindent
$X$ が Noether ならば（実際に最もよく用いる場合である）、上の議論と
Noether 帰納法を合わせて証明が終わる。一般の場合は次のように論じる。
まず、$U$ の外の点の各局所環が整域であり、$X$ が局所 Noether なので、
$U$ の $X$ における閉包は開である。したがって $U \subset X$ は稠密かつ
スキーム論的に稠密であると仮定してよい。
$T$ を三つ組 $(U', \mathcal{L}', \alpha)$ で、$U \subset U' \subset X$ が開部分スキーム、
$\mathcal{L}'$ が可逆 $\mathcal{O}_{U'}$-加群、
$\alpha : \mathcal{L}'|_U \to \mathcal{L}$ が同型であるもの全体とする。
$T$ に半順序 $\leq$ を入れ、規則
$(U', \mathcal{L}', \alpha) \leq (U'', \mathcal{L}'', \alpha')$
が成り立つのは、$U' \subset U''$ であり、同型
$\beta : \mathcal{L}''|_{U'} \to \mathcal{L}'$ が存在して
$\alpha$ および $\alpha'$ と両立する場合、かつその場合に限るとする。
$\beta$ は存在すれば一意であることに注意する。これは $U \subset X$ が稠密だからである。
証明の前半から、任意の元 $t = (U', \mathcal{L}', \alpha)$ で $T$ に属し、
$U' \not = X$ であるものに対して、$t' \in T$ で $t' > t$ となるものが存在する。
したがって証明を終えるには Zorn の補題が適用できることを示せば十分である。
そこで全順序部分集合 $I \subset T$ を考える。$i \in I$ が三つ組
$(U_i, \mathcal{L}_i, \alpha_i)$ に対応するならば、次のようにして可逆加群
$\mathcal{L}'$ を $U' = \bigcup U_i$ 上に構成できる。
$W \subset U'$ が開かつ準コンパクトならば、$W \subset U_i$ となる
$i$ が存在し、
$$
\mathcal{L}'(W) = \mathcal{L}_i(W)
$$
と置く。遷移写像には $\beta$ たちを用いる（それらは一意なので正しく合成する）。
これにより準コンパクト開集合の基底上の可逆 $\mathcal{O}$-加群
$\mathcal{L}'$ が $U'$ 上に定まり、可逆加群を定義するにはこれで十分である
（Sheaves, Section \ref{sheaves-section-bases}）。詳細は省略する。
\end{proof}

\begin{lemma}
\label{divisors-lemma-open-subscheme-UFD}
$R$ を UFD とする。次の Picard 群は自明である。
\begin{enumerate}
\item $\Spec(R)$ およびその任意の開部分スキーム。
\item $\mathbf{A}^n_R = \Spec(R[x_1, \ldots, x_n])$ およびその任意の開部分スキーム。
\end{enumerate}
とくに、$n$ 次元アフィン空間 $\mathbf{A}^n_k$ の任意の開部分スキームの
Picard 群は自明である。ここで底は体 $k$ である。
\end{lemma}

\begin{proof}
$R$ 上の任意の多項式環は Algebra, Lemma
\ref{algebra-lemma-polynomial-ring-UFD} により UFD である。
したがって、$\Pic(U)$ が自明であることを示すのが問題である。ここで
$U \subset \Spec(R)$ は（空でない）開集合である。

\medskip\noindent
Noether の場合、Lemma \ref{divisors-lemma-extend-invertible-module} を用いて
$\Pic(\Spec(R))$ が自明であることに帰着できる。これはさらに More on Algebra,
Lemma \ref{more-algebra-lemma-UFD-Pic-trivial} で証明されている。
以下の証明は読み飛ばすことを勧める。

\medskip\noindent
$U \subset \Spec(R)$ を空でない開集合とする。非零元 $f \in R$ で
$D(f) \subset U$ となるものを選ぶ。$f$ は $R$ の素元の積なので、有限個の素元
$a_1, \ldots, a_n \in R$ で $(a_i) \not \in U$（または同値に
$U \cap V(a_i) = \emptyset$）となるものが存在する。$R$ を
$R_{a_1 \ldots a_n}$ で置き換えると、$U$ がすべての単項素イデアルを
含むと仮定してよい（UFD の局所化は UFD である。たとえば Algebra, Lemma
\ref{algebra-lemma-invert-prime-elements}）。

\medskip\noindent
$U$ がすべての単項素イデアルを含むと仮定する。
$U = \Spec(R) \setminus V(I)$ と書く。ここで $I \subset R$ はあるイデアルである。
非零な $f \in I$ を選び、素元の積として $f = a_1 \ldots a_n$ と書く。
$(a_i) \not \in V(I)$ なので、$g \in I$ で
$g \not \in (a_i)$ がすべての $i$ に対して成り立つものを選べる（素回避）。図式
$$
U' = D(f) \cup D(g) \subset U \subset \Spec(R)
$$
を考える。More on Algebra, Lemma
\ref{more-algebra-lemma-UFD-Pic-trivial} により $D(f)$ と $D(g)$ の Picard 群は
自明であり、UFD に関する初等的な議論によって写像
$R_f^* \times R_g^* \to R_{fg}^*$、$(u, v) \mapsto uv^{-1}$ は全射である。
したがって $\Pic(U')$ は自明である。

\medskip\noindent
よって、可逆加群 $\mathcal{L}$ で $U$ 上にあり、$j : U' \to U'$ による引き戻しが
自明であるものが自明であることを示せば十分である。そのためには、写像
$\mathcal{L} \to j_*j^*\mathcal{L}$ が同型であることを示せば十分である。
さらにこのため、$u \in U \setminus U'$ とする。このとき
$f, g \in \mathfrak m_u \subset \mathcal{O}_{U, u} = A$ は、共通の素因子を
もたない局所 UFD の元である。素因子を調べれば
$$
0 \to A \to A_f \times A_g \to A_{fg}
$$
が完全であることを容易に示せる。$A = \mathcal{L}_u$ である。これは
$\mathcal{L}$ が可逆だからであり、$(j_*j^*\mathcal{L})_u$ は
$A_f \times A_g \to A_{fg}$ の核であるから結論を得る。
\end{proof}

\begin{lemma}
\label{divisors-lemma-Pic-projective-space-UFD}
$R$ を UFD とする。$\mathbf{P}^n_R$ の Picard 群は $\mathbf{Z}$ である。
より正確には、同型
$$
\mathbf{Z} \longrightarrow \Pic(\mathbf{P}^n_R),\quad
m \longmapsto \mathcal{O}_{\mathbf{P}^n_R}(m)
$$
がある。とくに、射影空間 $\mathbf{P}^n_k$ の Picard 群は、底が体
$k$ ならば $\mathbf{Z}$ である。
\end{lemma}

\begin{proof}
Lemma \ref{divisors-lemma-open-subscheme-UFD} により、開集合
$D_+(T_i) \cong \mathbf{A}^n_R$ の Picard 群は自明である。
したがって $\mathcal{L}$ が $\mathbf{P}^n_R$ 上の可逆加群ならば、
開被覆
$$
\mathbf{P}^n_R = \bigcup\nolimits_{i = 0, \ldots, n} D_+(T_i)
$$
に関する $1$-コサイクルで、可逆函子の層 $\mathcal{O}^*$ に値をとるものによって与えられる。
$i \not = j$ に対して
$$
\mathcal{O}^*(D_+(T_i) \cap D_+(T_j)) =
\mathcal{O}^*(D_+(T_iT_j)) = \left(R[T_0, \ldots, T_n]_{(T_iT_j)}\right)^* =
R^* \times (T_i/T_j)^\mathbf{Z}
$$
であることに注意する。したがって、このようなコサイクル $(g_{ij})$ は単元
$$
g_{ij} = u_{ij} (T_i/T_j)^{e_{ij}}
$$
で与えられる。ここで $u_{ij} \in R^*$ および $e_{ij} \in \mathbf{Z}$ は
コサイクル条件を満たす。$D_+(T_iT_jT_k)$ 上で
$\#\{i, j, k\} = 3$ となるコサイクル条件から
$$
u_{ik} = u_{ij} u_{jk} \quad\text{and}\quad
e_{ik} = e_{ij} = e_{jk}
$$
が従う。したがって $u_{ij} = u_{i1} u_{j1}^{-1}$ は境界である。
よって可逆加群のすべての同型類は、ある $e \in \mathbf{Z}$ に対するコサイクル
$$
g_{ij} = (T_i/T_j)^e
$$
を取ることにより与えられる。$\mathcal{O}(n)$ は自明化切断
$T_i^n$ を $D_+(T_i)$ 上でもつので、$\mathcal{O}(n)$ に対応するコサイクルは
$(T_i/T_j)^n$ であり、証明が完了する。
\end{proof}






\section{正規スキーム上の Weil 因子}
\label{divisors-section-weil-divisors-normal}

\noindent
まず反射的加群の性質を論じる。

\begin{lemma}
\label{divisors-lemma-reflexive-normal}
$X$ を整な局所 Noether 正規スキームとする。$\mathcal{F}$ および $\mathcal{G}$ を連接反射的
$\mathcal{O}_X$-加群とする。このとき写像
$$
(\SheafHom_{\mathcal{O}_X}(\mathcal{F}, \mathcal{O}_X)
\otimes_{\mathcal{O}_X} \mathcal{G})^{**} \to
\SheafHom_{\mathcal{O}_X}(\mathcal{F}, \mathcal{G})
$$
は同型である。規則 $\mathcal{F}, \mathcal{G} \mapsto
(\mathcal{F} \otimes_{\mathcal{O}_X} \mathcal{G})^{**}$ は、階数 $1$ の
連接反射的 $\mathcal{O}_X$-加群の同型類の集合上に可換群則を定める。
\end{lemma}

\begin{proof}
厳密には必要でないが、証明を進める前に Remark
\ref{divisors-remark-tensor} を読むことを勧める。開部分スキーム
$j : U \to X$ を次のように選ぶ。$X \setminus U$ の各既約成分は
余次元 $\geq 2$ を $X$ においてもち、$j^*\mathcal{F}$ および
$j^*\mathcal{G}$ は有限局所自由である。Lemma
\ref{divisors-lemma-reflexive-over-normal} を参照。このとき写像
$$
\SheafHom_{\mathcal{O}_U}(j^*\mathcal{F}, \mathcal{O}_U)
\otimes_{\mathcal{O}_U} j^*\mathcal{G} \to
\SheafHom_{\mathcal{O}_U}(j^*\mathcal{F}, j^*\mathcal{G})
$$
は同型である。実際、これは局所的に調べられ、加群が有限自由なら明らかである。
補題の表示された矢印に $j^*$ を適用すると、いま同型であると示した矢印を
得ることに注意する（細部を一つ省略した）。$j^*$ は $U$ 上の連接反射的加群と
$X$ 上の連接反射的加群との間の圏同値を定めるので（Lemma
\ref{divisors-lemma-reflexive-S2-extend} および Serre の判定法 Properties, Lemma
\ref{properties-lemma-criterion-normal} による）、補題の矢印も同型であると
結論する。$\mathcal{F}$ が階数 $1$ ならば、$j^*\mathcal{F}$ は可逆
$\mathcal{O}_U$-加群であり、反射的加群
$\mathcal{F}^\vee = \SheafHom(\mathcal{F}, \mathcal{O}_X)$ の制限はその逆元である。
前と同じ仕方で
$(\mathcal{F} \otimes_{\mathcal{O}_X} \mathcal{F}^\vee)^{**} = \mathcal{O}_X$
が従う。このようにして、補題の主張で与えた群則に逆元が存在することが分かる。
\end{proof}

\begin{lemma}
\label{divisors-lemma-normal-class-group}
$X$ を整な局所 Noether 正規スキームとする。階数 $1$ の連接反射的
$\mathcal{O}_X$-加群の群は $\text{Cl}(X)$ と同型である。これは $X$ の Weil 因子類群である。
\end{lemma}

\begin{proof}
$\mathcal{F}$ を階数 $1$ の連接反射的 $\mathcal{O}_X$-加群とする。開集合
$U \subset X$ を次のように選ぶ。$X \setminus U$ の各既約成分は
余次元 $\geq 2$ を $X$ においてもち、$\mathcal{F}|_U$ は可逆である。
Lemma \ref{divisors-lemma-reflexive-over-normal} を参照。$\text{Cl}(U) = \text{Cl}(X)$
であることに注意する。これは $X$ の Weil 因子類群が、その有理関数体と、
余次元 $1$ の点およびそれらの局所環だけに依存するからである。
そこで $\mathcal{F}$ の Weil 因子類を、$\mathcal{F}|_U$ の
$\text{Cl}(U)$ における Weil 因子類として定義できる。これが $U$ の選択に
依存しないことの確認は省略する。

\medskip\noindent
$\text{Cl}'(X)$ により、階数 $1$ の連接反射的
$\mathcal{O}_X$-加群の同型類の集合を表す。上の構成は群準同型
$$
\text{Cl}'(X) \longrightarrow \text{Cl}(X)
$$
を与える。実際、$\mathcal{F}, \mathcal{G}$ を
$\text{Cl}'(X)$ の任意の二元とすると、両方に使える $U$ を選べ、可逆加群を
その Weil 因子類へ送る対応 (\ref{divisors-equation-c1}) は準同型だからである。
$\mathcal{F}$ がこの写像の核に属するなら、$\mathcal{F}|_U$ は自明である
（Lemma \ref{divisors-lemma-normal-c1-injective}）。したがって Lemma
\ref{divisors-lemma-reflexive-S2-extend} および Serre の判定法 Properties, Lemma
\ref{properties-lemma-criterion-normal} により $\mathcal{F}$ も自明である。
証明を終えるには、この写像が全射であることを調べれば十分である。

\medskip\noindent
$D = \sum n_Z Z$ を $X$ 上の Weil 因子とする。次を主張する。開集合
$U \subset X$ が存在して、$X \setminus U$ の各既約成分は
余次元 $\geq 2$ を $X$ においてもち、$Z|_U$ は
$n_Z \not = 0$ に対して有効 Cartier 因子である。この主張を証明するには $X$ がアフィンであると
仮定してよい。すると $D = n_1 Z_1 + \ldots + n_r Z_r$ は有限和で、
$Z_1, \ldots, Z_r$ は相異なると仮定してよい。$Z_i \cap Z_j$ を
$i \not = j$ に対して取り除いた後、$Z_1, \ldots, Z_r$ は互いに素であると
仮定してよい。これにより、単一の素因子 $Z$ で $X$ 上にある場合に帰着する。
$X$ は Properties, Lemma \ref{properties-lemma-criterion-normal} により
$(R_1)$ なので、局所環 $\mathcal{O}_{X, \xi}$ は離散付値環である。ここで
$\xi$ は $Z$ の一般点である。一様化元
$f \in \mathcal{O}_{X, \xi}$ を取る。開近傍 $V \subset X$ で $\xi$ を含み、
$f$ がある元 $f \in \mathcal{O}_X(V)$ の像となるものを取る。$V$ を縮小すれば、
$Z \cap V = V(f)$ がスキーム論的に成り立つと仮定してよい。これは $\xi$ における
局所環で成り立つからである。この場合
$$
U = X \setminus (Z \setminus V) = (X \setminus Z) \cup V
$$
と取れば望む開集合を得て、主張が証明される。

\medskip\noindent
$D$ の因子類が像に属することを示すため、
$D = \sum_{n_Z < 0} n_Z Z - \sum_{n_Z > 0} (-n_Z) Z$ と書ける。
上で構成した写像の加法性により、$n_Z \leq 0$ がすべての素因子 $Z$ に対して
成り立つと仮定してよいし、そうする
（望むならこの段階を避けてもよい）。$U \subset X$ を上の主張のものとする。
$U$ が準コンパクトなら、$D|_U = -n_1 Z_1 - \ldots - n_r Z_r$ と書く。ここで
$Z_i$ は相異なる素因子で、$n_i > 0$ である。そして可逆
$\mathcal{O}_U$-加群
$$
\mathcal{L} =
\mathcal{I}_1^{n_1} \ldots \mathcal{I}_r^{n_r} \subset \mathcal{O}_U
$$
を考える。ここで $\mathcal{I}_i$ は $Z_i$ のイデアル層である。
これは $U$ の選び方と Lemma \ref{divisors-lemma-sum-effective-Cartier-divisors} により
可逆である。また $\text{div}_\mathcal{L}(1) = D|_U$ である。Lemma
\ref{divisors-lemma-reflexive-S2-extend} により $\mathcal{L} = \mathcal{F}|_U$ である。ここで
階数 $1$ の連接反射的 $\mathcal{O}_X$-加群 $\mathcal{F}$ を取れるので、
$D$ は写像の像に属することが分かる。

\medskip\noindent
$U$ が準コンパクトでないなら、上の表示式によって局所的に
$\mathcal{L} \subset \mathcal{O}_U$ を定義する。この構成が貼り合わさることを
読者が示せば、前とまったく同様に証明が終わる。細部は省略する。
\end{proof}

\begin{lemma}
\label{divisors-lemma-structure-sheaf-Xs}
$X$ を整な局所 Noether 正規スキームとする。$\mathcal{F}$ を階数 1 の連接反射的
$\mathcal{O}_X$-加群とする。$s \in \Gamma(X, \mathcal{F})$ とする。さらに
$$
U = \{x \in X \mid s : \mathcal{O}_{X, x} \to \mathcal{F}_x
\text{ is an isomorphism}\}
$$
とする。このとき $j : U \to X$ は $X$ の開部分スキームであり、
$$
j_*\mathcal{O}_U =
\colim (\mathcal{O}_X \xrightarrow{s} \mathcal{F}
\xrightarrow{s} \mathcal{F}^{[2]}
\xrightarrow{s} \mathcal{F}^{[3]}
\xrightarrow{s} \ldots)
$$
である。ここで $\mathcal{F}^{[1]} = \mathcal{F}$ とし、帰納的に
$\mathcal{F}^{[n + 1]} =
(\mathcal{F} \otimes_{\mathcal{O}_X} \mathcal{F}^{[n]})^{**}$ とする。
\end{lemma}

\begin{proof}
集合 $U$ は Modules, Lemmas
\ref{modules-lemma-finite-type-surjective-on-stalk} および
\ref{modules-lemma-finite-type-to-coherent-injective-on-stalk} により開である。
$j$ は Properties, Lemma \ref{properties-lemma-immersion-into-noetherian} により
準コンパクトであることに注意する。最後の主張を証明するには、任意の
準コンパクト開集合 $W \subset X$ に対して同型
$$
\colim \Gamma(W, \mathcal{F}^{[n]})
\longrightarrow
\Gamma(U \cap W, \mathcal{O}_U)
$$
が存在し、これが制限写像と両立する $\mathcal{O}_X(W)$-加群の同型であることを
示せば十分である。両立性の確認は省略する。$X$ を $W$ で置き換え、上の式を
Hom によって書き直すと、同型
$$
\colim \Hom_{\mathcal{O}_X}(\mathcal{O}_X, \mathcal{F}^{[n]})
\longrightarrow
\Hom_{\mathcal{O}_U}(\mathcal{O}_U, \mathcal{O}_U)
$$
を構成すれば十分であることが分かる。開集合 $V \subset X$ を次のように選ぶ。
$X \setminus V$ の各既約成分は余次元 $\geq 2$ を $X$ においてもち、
$\mathcal{F}|_V$ は可逆である。Lemma \ref{divisors-lemma-reflexive-over-normal} を参照。
このとき制限は、階数 $1$ の連接反射的加群について $X$ 上と $V$ 上との間、
また階数 $1$ の連接反射的加群について $U$ 上と $V \cap U$ 上との間の
圏同値を定める。Lemma \ref{divisors-lemma-reflexive-S2-extend} および Serre の判定法
Properties, Lemma \ref{properties-lemma-criterion-normal} を参照。
したがって同型
$$
\colim \Gamma(V, (\mathcal{F}|_V)^{\otimes n}) \longrightarrow
\Gamma(V \cap U, \mathcal{O}_U)
$$
を構成すれば十分である。$\mathcal{F}|_V$ は可逆であり、$U \cap V$ は
$s|_V$ がこの可逆加群を生成する点全体に等しいので、これは Properties, Lemma
\ref{properties-lemma-invert-s-sections} の特別な場合である
（写像には明示的な公式もある）。
\end{proof}

\begin{lemma}
\label{divisors-lemma-Xs-codim-complement}
仮定と記法は Lemma \ref{divisors-lemma-structure-sheaf-Xs} のとおりとする。
$s$ が非零なら、$X \setminus U$ の各既約成分は余次元 $1$ を $X$ においてもつ。
\end{lemma}

\begin{proof}
$\xi \in X$ を、既約成分 $Z$ で $X \setminus U$ に含まれるものの一般点とする。
$X$ を $\xi$ の開近傍で置き換えると、$Z = X \setminus U$ が既約であると
仮定してよい。$s : \mathcal{O}_U \to \mathcal{F}|_U$ は同型なので、
$Z$ の $X$ における余次元が $\geq 2$ ならば、
$s : \mathcal{O}_X \to \mathcal{F}$ は Lemma
\ref{divisors-lemma-reflexive-S2-extend} および Serre の判定法 Properties, Lemma
\ref{properties-lemma-criterion-normal} により同型である。これは
$Z = \emptyset$ を意味し、矛盾である。
\end{proof}

\begin{remark}
\label{divisors-remark-structure-sheaf-Xs}
$A$ を Noether 正規整域とする。$M$ を階数 $1$ の有限反射的
$A$-加群とする。$s \in M$ を非零とする。$\mathfrak p_1, \ldots, \mathfrak p_r$
を高さ $1$ の素イデアルとする。それらは $A$ に属し、$M/As$ の台に含まれる。
このとき Lemma \ref{divisors-lemma-structure-sheaf-Xs} の開集合 $U$ は
$$
U = \Spec(A) \setminus
\left(V(\mathfrak p_1) \cup \ldots \cup V(\mathfrak p_r)\right)
$$
である。これは Lemma \ref{divisors-lemma-Xs-codim-complement} による。さらに
$M^{[n]}$ が $M \otimes_A \ldots \otimes_A M$ の反射包
（$n$ 個の因子）を表すなら、
$$
\Gamma(U, \mathcal{O}_U) = \colim M^{[n]}
$$
である。これは Lemma \ref{divisors-lemma-structure-sheaf-Xs} による。
\end{remark}

\begin{lemma}
\label{divisors-lemma-affine-Xs}
仮定と記法は Lemma \ref{divisors-lemma-structure-sheaf-Xs} のとおりとする。
次は同値である。
\begin{enumerate}
\item 包含射 $j : U \to X$ はアフィンである。
\item すべての $x \in X \setminus U$ に対して、ある $n > 0$ が存在し、
$s^n \in \mathfrak m_x \mathcal{F}^{[n]}_x$ である。
\end{enumerate}
\end{lemma}

\begin{proof}
(1) を仮定する。$x \in X \setminus U$ に対し、逆像 $U_x$ を考える。これは $U$ の
標準射 $f_x : \Spec(\mathcal{O}_{X, x}) \to X$ による逆像で、アフィンであり、
$x$ を含まない。したがって $\mathfrak m_x \Gamma(U_x, \mathcal{O}_{U_x})$ は
単位イデアルである。とくに
$$
1 = \sum f_i g_i
$$
と書ける。ここで $f_i \in \mathfrak m_x$ かつ
$g_i \in \Gamma(U_x, \mathcal{O}_{U_x})$ である。Lemma
\ref{divisors-lemma-structure-sheaf-Xs} により
$\Gamma(U_x, \mathcal{O}_{U_x}) = \colim \mathcal{F}^{[n]}_x$ であり、
遷移写像は $s$ による乗法で与えられる。したがって、ある $n > 0$ に対して
$$
s^n = \sum f_i t_i
$$
となる。ここである $t_i = s^ng_i \in \mathcal{F}^{[n]}_x$ を取れる。
ゆえに (2) が成り立つ。

\medskip\noindent
逆に (2) を仮定する。$j$ がアフィンであることの証明は $X$ 上局所的である。
Morphisms, Lemma \ref{morphisms-lemma-characterize-affine} を参照。したがって
$X$ はアフィンであると仮定してよいし、そうする。目標は $U$ が
アフィンであることを示すことである。Cohomology of Schemes, Lemma
\ref{coherent-lemma-affine-if-quasi-affine} により、
$H^p(U, \mathcal{O}_U) = 0$ が $p > 0$ に対して成り立つことを示せば十分である。
$H^p(U, \mathcal{O}_U) = H^0(X, R^pj_*\mathcal{O}_U)$ であり
（Cohomology of Schemes, Lemma
\ref{coherent-lemma-quasi-coherence-higher-direct-images-application}）、
$R^pj_*\mathcal{O}_U$ は準連接なので（Cohomology of Schemes, Lemma
\ref{coherent-lemma-quasi-coherence-higher-direct-images}）、茎
$(R^pj_*\mathcal{O}_U)_x$ が点 $x \in X$ において零であることを示せば十分である。
基底変換図式
$$
\xymatrix{
U_x \ar[d]_{j_x} \ar[r] & U \ar[d]^j \\
\Spec(\mathcal{O}_{X, x}) \ar[r] & X
}
$$
を考える。Cohomology of Schemes, Lemma
\ref{coherent-lemma-flat-base-change-cohomology} により
$(R^pj_*\mathcal{O}_U)_x = R^pj_{x, *}\mathcal{O}_{U_x}$ である。
したがって $X$ は閉点 $x$ をもつ局所スキームであると仮定してよく、
$U$ がアフィンであることを示さなければならない
（これは上記の参照により望む消滅と同値である）。とくに
$d = \dim(X)$ は有限である（Algebra, Proposition
\ref{algebra-proposition-dimension}）。$x \in U$ ならば $U = X$ であり、
結果は明らかである。$d = 0$ かつ $x \not \in U$ ならば $U = \emptyset$ であり、
結果は明らかである。そこで $d > 0$ かつ $x \not \in U$ と仮定する。
$j_*\mathcal{O}_U = \colim \mathcal{F}^{[n]}$ なので、仮定は
$$
1 = \sum f_i g_i
$$
と書けることを意味する。ここで、ある $n > 0$ に対して
$f_i \in \mathfrak m_x$、$g_i \in \mathcal{O}(U)$ である。
$d$ に関する帰納法により、$D(f_i) \cap U$ はすべての $i$ に対して
アフィンである。実際、いま述べた議論全体で
$X$ を $D(f_i)$ に置き換えると、次元が $d$ より真に小さい
Noether 局所環に行き着く。したがって望むとおり、Properties, Lemma
\ref{properties-lemma-characterize-affine} により $U$ はアフィンである。
\end{proof}

\section{相対 Proj}
\label{divisors-section-relative-proj}

\noindent
相対 Proj に関するいくつかの結果を述べる。まず非常に基本的な結果から始める。
相対 Proj は常に基底上分離的であることを思い出そう。Constructions, Lemma
\ref{constructions-lemma-relative-proj-separated} を参照。

\begin{lemma}
\label{divisors-lemma-relative-proj-quasi-compact}
$S$ をスキームとする。$\mathcal{A}$ を準連接次数付き
$\mathcal{O}_S$-代数とする。
$p : X = \underline{\text{Proj}}_S(\mathcal{A}) \to S$ を
$\mathcal{A}$ の相対 Proj とする。次のいずれかが成り立つとする。
\begin{enumerate}
\item $\mathcal{A}$ は $\mathcal{A}_0$-代数の層として有限型である。
\item $\mathcal{A}$ は $\mathcal{A}_1$ により $\mathcal{A}_0$-代数として生成され、
$\mathcal{A}_1$ は有限型 $\mathcal{A}_0$-加群である。
\item 有限型準連接 $\mathcal{A}_0$-部分加群
$\mathcal{F} \subset \mathcal{A}_{+}$ が存在し、
$\mathcal{A}_{+}/\mathcal{F}\mathcal{A}$ は
$\mathcal{A}/\mathcal{F}\mathcal{A}$ の局所冪零イデアル層である。
\end{enumerate}
このとき $p$ は準コンパクトである。
\end{lemma}

\begin{proof}
問題は基底上局所的である。Schemes, Lemma
\ref{schemes-lemma-quasi-compact-affine} を参照。したがって $S$ は
アフィンであると仮定してよい。$S = \Spec(R)$ とし、$\mathcal{A}$ は
次数付き $R$-代数 $A$ に対応するとする。このとき
$X = \text{Proj}(A)$ である。Constructions, Section
\ref{constructions-section-relative-proj-via-glueing} を参照。
(1) の場合、必要ならさらに局所化した後、$A$ は斉次元
$f_1, \ldots, f_n \in A_{+}$ により $A_0$ 上生成されると仮定してよい。このとき
Algebra, Lemma \ref{algebra-lemma-S-plus-generated} により
$A_{+} = (f_1, \ldots, f_n)$ である。(3) の場合、
$\mathcal{F} = \widetilde{M}$ であることが分かる。ここで
$A_0$ 上有限型の加群 $M \subset A_{+}$ を取れる。
$M = \sum A_0f_i$ とする。$f_i = \sum f_{i, j}$ を斉次成分への分解とする。
(3) の条件は $A_{+} \subset \sqrt{(f_{i, j})}$ を意味する。
したがって、いずれの場合も Constructions, Lemma
\ref{constructions-lemma-proj-quasi-compact} により
$\text{Proj}(A)$ は準コンパクトであると結論する。最後に (2) は (1) から従う。
\end{proof}

\begin{lemma}
\label{divisors-lemma-relative-proj-finite-type}
$S$ をスキームとする。$\mathcal{A}$ を準連接次数付き
$\mathcal{O}_S$-代数とする。
$p : X = \underline{\text{Proj}}_S(\mathcal{A}) \to S$ を
$\mathcal{A}$ の相対 Proj とする。$\mathcal{A}$ が
$\mathcal{O}_S$-代数の層として有限型ならば、$p$ は有限型であり、
$\mathcal{O}_X(d)$ は有限型 $\mathcal{O}_X$-加群である。
\end{lemma}

\begin{proof}
仮定から $p$ は準コンパクトである。Lemma
\ref{divisors-lemma-relative-proj-quasi-compact} を参照。したがって
$p$ が局所有限型であることを示せば十分である。ゆえに問題は基底と対象上で
局所的である。Morphisms, Lemma
\ref{morphisms-lemma-locally-finite-type-characterize} を参照。
$S = \Spec(R)$ とし、$\mathcal{A}$ は次数付き $R$-代数 $A$ に対応するとする。
$S$ 上でさらに局所化すれば、$A$ は有限型 $R$-代数であると仮定してよい。
スキーム $X$ は、環 $A_{(f)}$ のスペクトルを貼り合わせて構成される。ここで
$f \in A_{+}$ は斉次元である。それぞれは Algebra, Lemma
\ref{algebra-lemma-dehomogenize-finite-type} の (1) により $R$ 上有限型である。
したがって $\text{Proj}(A)$ は $R$ 上有限型である。
$\mathcal{O}_X(d)$ に関する主張には Algebra, Lemma
\ref{algebra-lemma-dehomogenize-finite-type} の (2) を用いる。
\end{proof}

\begin{lemma}
\label{divisors-lemma-relative-proj-universally-closed}
$S$ をスキームとする。$\mathcal{A}$ を準連接次数付き
$\mathcal{O}_S$-代数とする。
$p : X = \underline{\text{Proj}}_S(\mathcal{A}) \to S$ を
$\mathcal{A}$ の相対 Proj とする。$\mathcal{O}_S \to \mathcal{A}_0$
が整な代数写像\footnote{言い換えると、$\mathcal{O}_S$ の
$\mathcal{A}_0$ における整閉包（Morphisms, Definition
\ref{morphisms-definition-integral-closure} を参照）が
$\mathcal{A}_0$ に等しい。}であり、$\mathcal{A}$ が
$\mathcal{A}_0$-代数として有限型ならば、$p$ は普遍閉である。
\end{lemma}

\begin{proof}
問題は基底上局所的である。したがって $X = \Spec(R)$ はアフィンであると
仮定してよい。$\mathcal{A}$ を準連接 $\mathcal{O}_X$-代数で、次数付き
$R$-代数 $A$ に付随するものとする。仮定は $R \to A_0$ が整であり、$A$ が
$A_0$ 上有限型であることである。$X \to \Spec(R)$ を合成
$X \to \Spec(A_0) \to \Spec(R)$ として書く。$R \to A_0$ は整な環写像なので、
$\Spec(A_0) \to \Spec(R)$ は普遍閉である。Morphisms, Lemma
\ref{morphisms-lemma-integral-universally-closed} を参照。
準コンパクトな（Constructions, Lemma
\ref{constructions-lemma-proj-quasi-compact} を参照）射
$$
X = \text{Proj}(A) \to \Spec(A_0)
$$
は Constructions, Lemma
\ref{constructions-lemma-proj-valuative-criterion} により付値判定法の存在部分を
満たすので、Schemes, Proposition
\ref{schemes-proposition-characterize-universally-closed} により普遍閉である。
したがって $X \to \Spec(R)$ は普遍閉射の合成として普遍閉である。
\end{proof}

\begin{lemma}
\label{divisors-lemma-relative-proj-proper}
$S$ をスキームとする。$\mathcal{A}$ を準連接次数付き
$\mathcal{O}_S$-代数とする。
$p : X = \underline{\text{Proj}}_S(\mathcal{A}) \to S$ を
$\mathcal{A}$ の相対 Proj とする。次の条件は同値である。
\begin{enumerate}
\item $\mathcal{A}_0$ は有限型 $\mathcal{O}_S$-加群であり、
$\mathcal{A}$ は $\mathcal{A}_0$-代数として有限型である。
\item $\mathcal{A}_0$ は有限型 $\mathcal{O}_S$-加群であり、
$\mathcal{A}$ は $\mathcal{O}_S$-代数として有限型である。
\end{enumerate}
これらの条件が成り立つなら、$p$ は局所射影的であり、とくに固有である。
\end{lemma}

\begin{proof}
$\mathcal{A}_0$ が有限型 $\mathcal{O}_S$-加群であると仮定する。
アフィン開集合 $U = \Spec(R) \subset X$ を、$\mathcal{A}$ が次数付き
$R$-代数 $A$ に対応し、$A_0$ が有限 $R$-加群となるように選ぶ。
条件 (1) は（必要なら $S$ 上でさらに局所化した後）$A$ が有限型
$A_0$-代数であることを意味し、条件 (2) は（必要なら $S$ 上でさらに
局所化した後）$A$ が有限型 $R$-代数であることを意味する。
したがって、これらの条件は Algebra, Lemma
\ref{algebra-lemma-compose-finite-type} により互いを含意する。

\medskip\noindent
局所射影的な射は固有である。Morphisms, Lemma
\ref{morphisms-lemma-locally-projective-proper} を参照。そこで
$S = \Spec(R)$ および $X = \text{Proj}(A)$ であり、$A_0$ は
$R$ 上有限、$A$ は $R$ 上有限型であると仮定してよい。
$X = \text{Proj}(A) \to \Spec(R)$ が射影的であることを示す。
以下の議論を読む代わりに、$X$ から $R$ 上の射影空間への閉埋入を直接構成して、
読者自身でこれを証明することを勧める。

\medskip\noindent
Lemma \ref{divisors-lemma-relative-proj-finite-type} により、$X$ は
$\Spec(R)$ 上有限型である。Constructions, Lemma
\ref{constructions-lemma-ample-on-proj} により、$\mathcal{O}_X(d)$ は
$X$ 上豊富であり、これはある $d \geq 1$ に対して成り立つ
（Properties, Section \ref{properties-section-ample} を参照）。
したがって $X \to \Spec(R)$ は準射影的である
（Morphisms, Definition \ref{morphisms-definition-quasi-projective} による）。
Morphisms, Lemma \ref{morphisms-lemma-quasi-projective-open-projective} により、
$X$ は $\Spec(R)$ 上射影的なスキームの開部分スキームと同型であると結論する。
ゆえに証明を終えるには、$X \to \Spec(R)$ が普遍閉であることを示せば十分である
（Morphisms, Lemma \ref{morphisms-lemma-image-proper-scheme-closed} を用いる）。
これは Lemma \ref{divisors-lemma-relative-proj-universally-closed} から従う。
\end{proof}

\begin{lemma}
\label{divisors-lemma-relative-proj-projective}
$S$ をスキームとする。$\mathcal{A}$ を準連接次数付き
$\mathcal{O}_S$-代数とする。
$p : X = \underline{\text{Proj}}_S(\mathcal{A}) \to S$ を
$\mathcal{A}$ の相対 Proj とする。$\mathcal{A}$ が
$\mathcal{A}_1$ により $\mathcal{A}_0$ 上生成され、$\mathcal{A}_1$ が有限型
$\mathcal{O}_S$-加群ならば、$p$ は射影的である。
\end{lemma}

\begin{proof}
実際、次数付き $\mathcal{O}_S$-代数写像
$$
\text{Sym}_{\mathcal{O}_X}^*(\mathcal{A}_1)
\longrightarrow
\mathcal{A}
$$
に付随する射は閉埋入 $X \to \mathbf{P}(\mathcal{A}_1)$ である。
Constructions, Lemma
\ref{constructions-lemma-surjective-generated-degree-1-map-relative-proj} を参照。
\end{proof}

\begin{lemma}
\label{divisors-lemma-relative-proj-flat}
$S$ をスキームとする。$\mathcal{A}$ を準連接次数付き
$\mathcal{O}_S$-代数とする。
$p : X = \underline{\text{Proj}}_S(\mathcal{A}) \to S$ を
$\mathcal{A}$ の相対 Proj とする。$\mathcal{A}_d$ が平坦
$\mathcal{O}_S$-加群であり、これが $d \gg 0$ に対して成り立つならば、$p$ は平坦であり、
$\mathcal{O}_X(d)$ は $S$ 上平坦である。
\end{lemma}

\begin{proof}
$X$ の $S$ 上の平坦性は、アフィン局所的に次の主張に帰着する。
$R$ を環、$A$ を次数付き $R$-代数で、$A_d$ が $R$ 上平坦であり、
これが $d \gg 0$ に対して成り立つものとする。また $f \in A_d$ とし、
これはある $d > 0$ に対して取る。このとき $A_{(f)}$ は $R$ 上平坦である。
$A_{(f)} = \colim A_{nd}$ であり、遷移写像は $f$ による乗法で与えられるので、
これは Algebra, Lemma \ref{algebra-lemma-colimit-flat} から従う。
$\mathcal{O}_X(d)$ の $S$ 上の平坦性についても同様に論じる。
\end{proof}

\begin{lemma}
\label{divisors-lemma-relative-proj-finite-presentation}
$S$ をスキームとする。$\mathcal{A}$ を準連接次数付き
$\mathcal{O}_S$-代数とする。
$p : X = \underline{\text{Proj}}_S(\mathcal{A}) \to S$ を
$\mathcal{A}$ の相対 Proj とする。$\mathcal{A}$ が有限表示
$\mathcal{O}_S$-代数ならば、$p$ は有限表示であり、
$\mathcal{O}_X(d)$ は有限表示 $\mathcal{O}_X$-加群である。
\end{lemma}

\begin{proof}
アフィン局所的に、これは次の主張に帰着する。$R$ を環、$A$ を有限表示次数付き
$R$-代数とする。このとき $\text{Proj}(A) \to \Spec(R)$ は有限表示であり、
$\mathcal{O}_{\text{Proj}(A)}(d)$ は有限表示
$\mathcal{O}_{\text{Proj}(A)}$-加群である。有限表示という条件から、表示
$$
A = R[X_1, \ldots, X_n]/(F_1, \ldots, F_m)
$$
を選べる。ここで $R[X_1, \ldots, X_n]$ は重み
$d_i$ を $X_i$ に与えて次数付けした多項式環であり、$F_1, \ldots, F_m$ は次数
$e_j$ の斉次多項式である。$R_0 \subset R$ を、多項式
$F_1, \ldots, F_m$ の係数により生成される部分環とする。
$A_0 = R_0[X_1, \ldots, X_n]/(F_1, \ldots, F_m)$ と置く。
構成により $A = A_0 \otimes_{R_0} R$ である。したがって Constructions, Lemma
\ref{constructions-lemma-base-change-map-proj} により、
$X_0 = \text{Proj}(A_0)$ の $R_0$ 上で結果を証明すれば十分である。
Lemma \ref{divisors-lemma-relative-proj-finite-type} により、$X_0$ は $R_0$ 上有限型であり、
$\mathcal{O}_{X_0}(d)$ は有限型準連接 $\mathcal{O}_{X_0}$-加群である。
$R_0$ は Noether なので（有限生成 $\mathbf{Z}$-代数であるから）、
$X_0$ は $R_0$ 上有限表示であり（Morphisms, Lemma
\ref{morphisms-lemma-noetherian-finite-type-finite-presentation}）、
$\mathcal{O}_{X_0}(d)$ は Cohomology of Schemes, Lemma
\ref{coherent-lemma-coherent-Noetherian} により有限表示である。
これで証明が完了する。
\end{proof}

\section{相対 Proj の閉部分スキーム}
\label{divisors-section-closed-in-relative-proj}

\noindent
相対 Proj の閉部分スキームに関するいくつかの補助補題を述べる。

\begin{lemma}
\label{divisors-lemma-closed-subscheme-proj}
$S$ をスキームとする。$\mathcal{A}$ を準連接次数付き
$\mathcal{O}_S$-代数とする。
$p : X = \underline{\text{Proj}}_S(\mathcal{A}) \to S$ を
$\mathcal{A}$ の相対 Proj とする。$i : Z \to X$ を閉部分スキームとする。
$\mathcal{I} \subset \mathcal{A}$ により、標準写像
$$
\mathcal{A}
\longrightarrow
\bigoplus\nolimits_{d \geq 0} p_*\left((i_*\mathcal{O}_Z)(d)\right).
$$
の核を表す。$p$ が準コンパクトならば、同型
$Z = \underline{\text{Proj}}_S(\mathcal{A}/\mathcal{I})$ が存在する。
\end{lemma}

\begin{proof}
射 $p$ は Constructions, Lemma
\ref{constructions-lemma-relative-proj-separated} により分離的である。
$p$ は準コンパクトなので、$p_*$ は準連接加群を準連接加群へ移す。
Schemes, Lemma \ref{schemes-lemma-push-forward-quasi-coherent} を参照。
したがって $\mathcal{I}$ は準連接 $\mathcal{O}_S$-加群である。
とくに $\mathcal{B} = \mathcal{A}/\mathcal{I}$ は準連接次数付き
$\mathcal{O}_S$-代数である。関手性による射
$Z' = \underline{\text{Proj}}_S(\mathcal{B}) \to
\underline{\text{Proj}}_S(\mathcal{A})$ は至る所で定義され、閉埋入である。
Constructions, Lemma
\ref{constructions-lemma-surjective-graded-rings-map-relative-proj} を参照。
したがって、$Z = Z'$ であることを $X$ の閉部分スキームとして証明すれば十分である。

\medskip\noindent
以上を踏まえると、問題は基底上局所的であるから、
$S = \Spec(R)$ であり、$X = \text{Proj}(A)$ であると仮定してよい。ここで
ある次数付き $R$-代数 $A$ を取る。
$\mathcal{I} = \widetilde{I}$ であると仮定する。ここで $I \subset A$ は次数付きイデアルである。
Constructions, Lemma \ref{constructions-lemma-proj-quasi-compact} により、
$f_0, \ldots, f_n \in A_{+}$ が存在して
$A_{+} \subset \sqrt{(f_0, \ldots, f_n)}$ となる。言い換えると
$X = \bigcup D_{+}(f_i)$ である。したがって
$Z \cap D_{+}(f_i) = Z' \cap D_{+}(f_i)$ を各 $i$ に対し調べれば十分である。
番号を付け直して $i = 0$ と仮定してよい。
$Z \cap D_{+}(f_0)$ および $Z' \cap D_{+}(f_0)$ はそれぞれ
イデアル $J$ および $J'$ により切り出されるとする。これらは $A_{(f_0)}$ のイデアルである。

\medskip\noindent
包含 $J' \subset J$ を示す。
$d$ を $\deg(f_0), \ldots, \deg(f_n)$ の最小公倍数とする。
各ねじり $\mathcal{O}_X(nd)$ は可逆であり、
$f_i^{nd/\deg(f_i)}$ により $D_{+}(f_i)$ 上で自明化されること、また任意の
準連接加群 $\mathcal{F}$ で $X$ 上にあるものに対し乗法写像
$\mathcal{O}_X(nd) \otimes_{\mathcal{O}_X} \mathcal{F}(m)
\to \mathcal{F}(nd + m)$ は同型であることに注意する。
Constructions, Lemma \ref{constructions-lemma-when-invertible} を参照。
$J'$ は元 $g/f_0^e$ により生成されるイデアルであることに注意する。ここで
$g \in I$ は次数 $e\deg(f_0)$ の斉次元である
（Constructions, Lemma
\ref{constructions-lemma-surjective-graded-rings-map-proj} の証明を参照）。
もちろん、$g$ を $f_0^lg$ で置き換えれば、適当な $l$ に対して
常に $d | e$ と仮定できる。このとき $g$ は
$\mathcal{O}_X(e\deg(f_0))$ の $Z$ への制限の切断として消えるので、
$g/f_0^d$ は $J$ の元である。したがって $J' \subset J$ である。

\medskip\noindent
逆に $g/f_0^e \in J$ と仮定する。再び $d | e$ と仮定してよい。
$i \in \{1, \ldots, n\}$ を取る。このとき $Z \cap D_{+}(f_i)$ は
あるイデアル $J_i \subset A_{(f_i)}$ により切り出される。さらに
$$
J \cdot A_{(f_0f_i)} = J_i \cdot A_{(f_0f_i)}.
$$
右辺は $J_i$ の
$f_0^{\deg(f_i)}/f_i^{\deg(f_0)}$ に関する局所化である。したがって、
十分可除なある $e_i \gg 0$ に対して
$$
f_0^{e_i}g/f_i^{(e_i + e)\deg(f_0)/\deg(f_i)} \in J_i
$$
となる。これにより $f_0^{\max(e_i)}g$ は $I$ の元であることが分かる。
実際、その各アフィン開集合 $D_{+}(f_i)$ への制限は閉部分スキーム
$Z \cap D_{+}(f_i)$ 上で消える。よって $g/f_0^e \in J'$ であり、
望むとおり $J \subset J'$ と結論する。
\end{proof}

\begin{example}
\label{divisors-example-closed-subscheme-of-proj}
$A$ を次数付き環とする。$X = \text{Proj}(A)$ および
$S = \Spec(A_0)$ とする。次数付きイデアル $I \subset A$ が与えられると、
Constructions, Lemma
\ref{constructions-lemma-surjective-graded-rings-map-proj} により閉部分スキーム
$V_+(I) = \text{Proj}(A/I) \to X$ を得る。Lemma
\ref{divisors-lemma-closed-subscheme-proj} の結果を翻訳すると、$X$ が準コンパクトならば、
任意の閉部分スキーム $Z$ は $V_+(I(Z))$ の形であることが分かる。ここで
次数付きイデアル $I(Z) \subset A$ は規則
$$
I(Z) = \Ker(A \longrightarrow
\bigoplus\nolimits_{n \geq 0} \Gamma(Z, \mathcal{O}_Z(n)))
$$
で与えられる。そこで次の二つの自然な問題を問える。
\begin{enumerate}
\item どのイデアル $I$ が $I(Z)$ の形であるか。
\item 操作 $I \mapsto I(V_+(I))$ を記述できるか。
\end{enumerate}
$A$ が Noether であるときにこれらへ答える。

\medskip\noindent
まず $A$ は $A_1$ により $A_0$ 上生成されると仮定する。この場合、
任意のイデアル $I \subset A$ に対し、写像
$A/I \to \bigoplus \Gamma(\text{Proj}(A/I), \mathcal{O})$ の核は
$A/I$ のねじれ元全体である。Cohomology of Schemes, Proposition
\ref{coherent-proposition-coherent-modules-on-proj} を参照。
したがって
$$
I(V_+(I)) = \{x \in A \mid A_n x \subset I\text{ for some }n \geq 0\}
$$
と結論する。右辺のイデアルは $I$ の飽和と呼ばれることがある。これは (2) への答えであり、
(1) への答えは、イデアルが $I(Z)$ の形であるための必要十分条件は、それが
飽和していること、すなわち自身の飽和に等しいことである。

\medskip\noindent
$A$ が一般の Noether 次数付き環ならば、Cohomology of Schemes, Proposition
\ref{coherent-proposition-coherent-modules-on-proj-general} を用いる。すると
$d$ を $A$ の $A_0$ 上の生成元の次数の最小公倍数としたとき
$$
I(V_+(I)) = \{x \in A \mid (Ax)_{nd} \subset I\text{ for all }n \gg 0\}
$$
を得る。これは $I$ の飽和と異なり得る。そのようになるのは $d \not = 1$ の場合である。たとえば
$A = \mathbf{Q}[x, y]$ で $\deg(x) = 2$、$\deg(y) = 3$ とする。このとき
$d = 6$ である。$I = (y^2)$ とする。このとき $y \in I(V_+(I))$ である。
実際、任意の斉次元 $f \in A$ で $6 | \deg(fy)$ を満たすものに対して
$y | f$ であり、したがって $fy \in I$ である。これより
$I(V_+(I)) = (y)$ であるが、$x^n y \not \in I$ がすべての $n$ に対して
成り立つ。よって $I(V_+(I))$ は飽和に等しくない。
\end{example}

\begin{lemma}
\label{divisors-lemma-equation-codim-1-in-projective-space}
$R$ を Noether UFD とする。$Z \subset \mathbf{P}^n_R$ を閉部分スキームで、
埋め込まれた点をもたず、$Z$ の各既約成分が余次元 $1$ を
$\mathbf{P}^n_R$ においてもつものとする。このとき
イデアル $I(Z) \subset R[T_0, \ldots, T_n]$ で $Z$ に対応するものは単項である。
\end{lemma}

\begin{proof}
$X = \mathbf{P}^n_R$ の局所環は UFD であることに注意する。実際、$X$ は
$\mathbf{A}^n_R$ と同型なアフィン部分により被覆され、
$R[x_1, \ldots, x_n]$ は UFD だからである
（Algebra, Lemma \ref{algebra-lemma-polynomial-ring-UFD}）。
したがって $Z$ は Lemma \ref{divisors-lemma-codimension-1-is-effective-Cartier} により
有効 Cartier 因子である。$\mathcal{I} \subset \mathcal{O}_X$ を $Z$ に対応する
準連接イデアル層とする。同型 $\mathcal{O}(m) \to \mathcal{I}$ を選ぶ。ここで
ある $m \in \mathbf{Z}$ を取る。Lemma
\ref{divisors-lemma-Pic-projective-space-UFD} を参照。このとき合成
$$
\mathcal{O}_X(m) \to \mathcal{I} \to \mathcal{O}_X
$$
は非零である。したがって $m \leq 0$ であり、対応する
$\mathcal{O}_X(m)^{\otimes -1} = \mathcal{O}_X(-m)$ の切断は、
ある $F \in R[T_0, \ldots, T_n]$ により与えられ、その次数は $-m$ である。
Cohomology of Schemes, Lemma
\ref{coherent-lemma-cohomology-projective-space-over-ring} を参照。
したがって $i$ 番目の標準開集合 $U_i = D_+(T_i)$ 上で、閉部分スキーム
$Z \cap U_i$ はイデアル
$$
(F(T_0/T_i, \ldots, T_n/T_i)) \subset R[T_0/T_i, \ldots, T_n/T_i]
$$
により切り出される。ゆえに次数付きイデアルの斉次元
$I(Z) = \Ker(R[T_0, \ldots, T_n] \to \bigoplus \Gamma(\mathcal{O}_Z(m)))$
は、次を満たす斉次多項式 $G$ 全体である。
$$
G(T_0/T_i, \ldots, T_n/T_i) \in (F(T_0/T_i, \ldots, T_n/T_i))
$$
これは $i = 0, \ldots, n$ に対して成り立つ。分母を払うと、
ある $e_i \geq 0$ が存在して
$$
T_i^{e_i}G \in (F)
$$
となることが分かる。これは $i = 0, \ldots, n$ に対して成り立つ。
$R$ が UFD なので $R[T_0, \ldots, T_n]$ も UFD である。
$F | T_0^{e_0}G$ および $F | T_1^{e_1}G$ から $F | G$ が従う。
実際、$T_0^{e_0}$ と $T_1^{e_1}$ は共通因子をもたない。
したがって $I(Z) = (F)$ である。
\end{proof}

\noindent
閉部分スキームが局所有限個の方程式により切り出される場合、
$\mathcal{A}$ の有限型イデアル層によってそれを定義できる。

\begin{lemma}
\label{divisors-lemma-closed-subscheme-proj-finite}
$S$ を準コンパクトかつ準分離的なスキームとする。$\mathcal{A}$ を準連接次数付き
$\mathcal{O}_S$-代数とする。
$p : X = \underline{\text{Proj}}_S(\mathcal{A}) \to S$ を
$\mathcal{A}$ の相対 Proj とする。$i : Z \to X$ を閉部分スキームとする。
$p$ が準コンパクトで $i$ が有限表示ならば、ある $d > 0$ と有限型準連接
$\mathcal{O}_S$-部分加群 $\mathcal{F} \subset \mathcal{A}_d$ が存在して、
$Z = \underline{\text{Proj}}_S(\mathcal{A}/\mathcal{F}\mathcal{A})$ となる。
\end{lemma}

\begin{proof}
Lemma \ref{divisors-lemma-closed-subscheme-proj} により、準連接次数付きイデアル層
$\mathcal{I} \subset \mathcal{A}$ が存在して
$Z = \underline{\text{Proj}}(\mathcal{A}/\mathcal{I})$ となる。
$S$ は準コンパクトなので、有限アフィン開被覆
$S = U_1 \cup \ldots \cup U_n$ を選べる。$U_i = \Spec(R_i)$ とする。
$\mathcal{A}|_{U_i}$ は次数付き $R_i$-代数 $A_i$ に対応し、
$\mathcal{I}|_{U_i}$ は次数付きイデアル $I_i \subset A_i$ に対応するとする。
$p^{-1}(U_i) = \text{Proj}(A_i)$ が $R_i$ 上のスキームとして成り立つことに注意する。
$p$ は準コンパクトなので、有限個の斉次元 $f_{i, j} \in A_{i, +}$ を
$p^{-1}(U_i) = \bigcup_j D_{+}(f_{i, j})$ となるように選べる。
$Z \to X$ に関する条件は、$Z$ の $\mathcal{O}_X$ におけるイデアル層が
有限型であることを意味する。Morphisms, Lemma
\ref{morphisms-lemma-closed-immersion-finite-presentation} を参照。
したがって有限個の斉次元 $h_{i, j, k} \in I_i \cap A_{i, +}$ を、
$Z \cap D_{+}(f_{i, j})$ のイデアルが元
$h_{i, j, k}/f_{i, j}^{e_{i, j, k}}$ により生成されるように取れる。
$d > 0$ を、すべての整数 $\deg(f_{i, j})$ および
$\deg(h_{i, j, k})$ の公倍数として選ぶ。Properties, Lemma
\ref{properties-lemma-quasi-coherent-colimit-finite-type} により、有限型準連接
$\mathcal{F} \subset \mathcal{I}_d$ が存在して、すべての局所切断
$$
h_{i, j, k}f_{i, j}^{(d - \deg(h_{i, j, k}))/\deg(f_{i, j})}
$$
は $\mathcal{F}$ の切断である。構成により $\mathcal{F}$ が解となる。
\end{proof}

\noindent
Lemma \ref{divisors-lemma-closed-subscheme-proj-finite} の次の版を、Lemma
\ref{divisors-lemma-composition-admissible-blowups} の証明で用いる。

\begin{lemma}
\label{divisors-lemma-closed-subscheme-proj-finite-type}
$S$ を準コンパクトかつ準分離的なスキームとする。$\mathcal{A}$ を準連接次数付き
$\mathcal{O}_S$-代数とする。
$p : X = \underline{\text{Proj}}_S(\mathcal{A}) \to S$ を
$\mathcal{A}$ の相対 Proj とする。$i : Z \to X$ を閉部分スキームとする。
$U \subset S$ を開集合とする。次を仮定する。
\begin{enumerate}
\item $p$ は準コンパクトである。
\item $i$ は有限表示である。
\item $U \cap p(i(Z)) = \emptyset$ である。
\item $U$ は準コンパクトである。
\item $\mathcal{A}_n$ は有限型 $\mathcal{O}_S$-加群であり、これはすべての
$n$ に対して成り立つ。
\end{enumerate}
このとき、ある $d > 0$ と有限型準連接
$\mathcal{O}_S$-部分加群 $\mathcal{F} \subset \mathcal{A}_d$ が存在して、
(a) $Z = \underline{\text{Proj}}_S(\mathcal{A}/\mathcal{F}\mathcal{A})$ であり、
(b) $\mathcal{A}_d/\mathcal{F}$ の台は $U$ と交わらない。
\end{lemma}

\begin{proof}
$\mathcal{I} \subset \mathcal{A}$ を、Lemma
\ref{divisors-lemma-closed-subscheme-proj} で構成した準連接次数付きイデアル層とする。
$U_i$、$R_i$、$A_i$、$I_i$、$f_{i, j}$、$h_{i, j, k}$ および $d$ を、
Lemma \ref{divisors-lemma-closed-subscheme-proj-finite} の証明で構成したものとする。
$U \cap p(i(Z)) = \emptyset$ なので、
$\mathcal{I}_d|_U = \mathcal{A}_d|_U$ である
（$\mathcal{I}$ を核として構成したことによる）。$U$ は準コンパクトなので、
有限アフィン開被覆 $U = W_1 \cup \ldots \cup W_m$ を選べる。
$\mathcal{A}_d$ は有限型なので、有限個の切断
$g_{t, s} \in \mathcal{A}_d(W_t)$ を
$\mathcal{A}_d|_{W_t} = \mathcal{I}_d|_{W_t}$ の
$\mathcal{O}_{W_t}$-加群としての生成元に取れる。証明を終えるため、
Properties, Lemma \ref{properties-lemma-quasi-coherent-colimit-finite-type} により
有限型 $\mathcal{F} \subset \mathcal{I}_d$ が存在して、すべての局所切断
$$
h_{i, j, k}f_{i, j}^{(d - \deg(h_{i, j, k}))/\deg(f_{i, j})}
\quad\text{and}\quad
g_{t, s}
$$
が $\mathcal{F}$ の切断となることに注意する。構成により
$\mathcal{F}$ が解となる。
\end{proof}

\begin{lemma}
\label{divisors-lemma-conormal-sheaf-section-projective-bundle}
$X$ をスキームとする。$\mathcal{E}$ を準連接
$\mathcal{O}_X$-加群とする。次の間に全単射がある。
$$
\left\{
\begin{matrix}
\text{sections }\sigma\text{ of the } \\
\text{morphism } \mathbf{P}(\mathcal{E}) \to X
\end{matrix}
\right\}
\leftrightarrow
\left\{
\begin{matrix}
\text{surjections }\mathcal{E} \to \mathcal{L}\text{ where} \\
\mathcal{L}\text{ is an invertible }\mathcal{O}_X\text{-module}
\end{matrix}
\right\}
$$
この場合 $\sigma$ は閉埋入であり、標準同型
$$
\Ker(\mathcal{E} \to \mathcal{L})
\otimes_{\mathcal{O}_X} \mathcal{L}^{\otimes -1}
\longrightarrow
\mathcal{C}_{\sigma(X)/\mathbf{P}(\mathcal{E})}
$$
がある。全単射と同型はいずれも基底変換と両立する。
\end{lemma}

\begin{proof}
$\pi : \mathbf{P}(\mathcal{E}) \to X$ は $\mathcal{E}$ 上の対称代数の
相対 Proj であることを思い出そう。Constructions, Definition
\ref{constructions-definition-projective-bundle} を参照。
したがって切断 $\sigma$ の記述は、Constructions, Lemma
\ref{constructions-lemma-apply-relative} における
$\mathbf{P}(\mathcal{E})$ の点関手の記述から直ちに従う。$\pi$ は分離的なので、
任意の切断は閉埋入である（Constructions, Lemma
\ref{constructions-lemma-relative-proj-separated} および Schemes, Lemma
\ref{schemes-lemma-section-immersion}）。
$U \subset X$ をアフィン開集合とし、$k \in \mathcal{E}(U)$ および
$s \in \mathcal{E}(U)$ を局所切断として、$k$ は $\mathcal{L}$ において
零へ写り、$s$ は生成元 $\overline{s}$ へ写るものとする。この生成元は
$\mathcal{L}$ のものである。
このとき $f = k/s$ は $\mathcal{O}_{\mathbf{P}(\mathcal{E})}$ の切断であり、
開近傍 $D_+(s)$ 上で定義される。これは $s(U)$ の
$\pi^{-1}(U)$ における開近傍である。
さらに $k$ は $\mathcal{L}$ において零へ写るので、$f$ は
$s(U)$ のイデアル層の切断であることが分かる。このイデアル層は
$\pi^{-1}(U)$ におけるものである。そこで像 $\overline{f}$ を $f$ から取り、
$\mathcal{C}_{\sigma(X)/\mathbf{P}(\mathcal{E})}(U)$ において取れる。
次を主張する。(1) 像 $\overline{f}$ は切断 $k$ と $\overline{s}$ のみに依存し、
$s$ の選択に依存しない。(2) この方法で $U$ 上の同型を得る（下記参照）。
しかし (1) と (2) が分かれば、構成は基底変換 $U' \to U$ と両立する。
ここで $U'$ はアフィンである。これにより、これらの局所写像は貼り合わさり、
任意の基底変換と両立することが分かる。

\medskip\noindent
(1) と (2) を証明するため、何が起きているかを明示する。すなわち
$U = \Spec(A)$ とし、$\mathcal{E} \to \mathcal{L}$ は
$A$-加群の写像 $M \to N$ に対応するとする。このとき
$k \in K = \Ker(M \to N)$ であり、$s \in M$ は生成元
$\overline{s}$ へ写る。この生成元は $N$ のものである。したがって
$M = K \oplus A s$ であり、
$$
\text{Sym}(M) = \text{Sym}(K)[s]
$$
となる。同一視 $\text{Sym}(K) \to \text{Sym}(M)_{(s)}$ を、
規則 $g \mapsto g/s^n$ によって考える。ここで
$g \in \text{Sym}^n(K)$ である。これは同型
$D_+(s) = \Spec(\text{Sym}(K))$ を与え、そのもとで $\sigma$ は環写像
$\text{Sym}(K) \to A$ に対応し、この写像は $K$ を零へ写す。
この同型を通じて、$K$ に対応する準連接加群は
$\mathcal{C}_{\sigma(U)/D_+(s)}$ と同一視され、(2) が証明される。
最後に、$s' = k' + s$ であると仮定する。ここである $k' \in K$ を取る。このとき
$$
k/s' = (k/s) (s/s') = (k/s) (s'/s)^{-1} = (k/s) (1 + k'/s)^{-1}
$$
が $\sigma(U)$ の $D_+(s)$ における開近傍上で成り立つ。
したがって $s'/s$ は $1$ に制限され、この制限は $\sigma(U)$ 上であり、
$k/s'$ は $k/s$ と同じ余法層の元へ写ることが分かる。これで (1) が証明される。
\end{proof}

\section{ブローアップ}
\label{divisors-section-blowing-up}

\noindent
ブローアップは代数幾何学における重要な道具である。

\begin{definition}
\label{divisors-definition-blow-up}
$X$ をスキームとする。$\mathcal{I} \subset \mathcal{O}_X$ を
準連接イデアル層とし、$Z \subset X$ を $\mathcal{I}$ に対応する閉部分スキームとする。
Schemes, Definition \ref{schemes-definition-immersion} を参照。
{\it $X$ の $Z$ に沿うブローアップ}、または
{\it $X$ のイデアル層 $\mathcal{I}$ におけるブローアップ}とは、射
$$
b :
\underline{\text{Proj}}_X
\left(\bigoplus\nolimits_{n \geq 0} \mathcal{I}^n\right)
\longrightarrow
X
$$
のことである。ブローアップの{\it 例外因子}とは逆像
$b^{-1}(Z)$ である。$Z$ をブローアップの{\it 中心}と呼ぶこともある。
\end{definition}

\noindent
後に、例外因子が有効 Cartier 因子であることを見る。さらに、ブローアップは、
$X$ 上で $Z$ の逆像が有効 Cartier 因子となるような「最小の」スキームとして
特徴づけられる。

\medskip\noindent
$b : X' \to X$ が $X$ の中心 $Z$ におけるブローアップであるとき、
構造層の捩りをしばしば $\mathcal{O}_{X'}(n)$ と書く。これらは可逆
$\mathcal{O}_{X'}$-加群であり、また
$\mathcal{O}_{X'}(n) = \mathcal{O}_{X'}(1)^{\otimes n}$
であることに注意する。実際、$X'$ は準連接次数付き
$\mathcal{O}_X$-代数の相対 Proj であり、この代数は次数 $1$ で生成される。
Constructions, Lemma
\ref{constructions-lemma-apply-relative} を参照。
$\mathcal{O}_{X'}(1)$ は $b$-相対的に非常に豊富である。
$b$ が有限型であるとは限らず、準コンパクトでさえないことがあるにもかかわらず、
$X'$ には $\mathbf{P}(\mathcal{I})$ への閉埋入が備わっているからである。
Morphisms, Example \ref{morphisms-example-very-ample} を参照。

\begin{lemma}
\label{divisors-lemma-blowing-up-affine}
$X$ をスキームとする。$\mathcal{I} \subset \mathcal{O}_X$ を
準連接イデアル層とする。$U = \Spec(A)$ を $X$ のアフィン開部分スキームとし、
$I \subset A$ を $\mathcal{I}|_U$ に対応するイデアルとする。
$b : X' \to X$ が $X$ の $\mathcal{I}$ におけるブローアップならば、標準同型
$$
b^{-1}(U) = \text{Proj}(\bigoplus\nolimits_{d \geq 0} I^d)
$$
があり、これは $b^{-1}(U)$ を $I$ の $A$ における Rees 代数の
斉次スペクトルと同一視する。さらに、$b^{-1}(U)$ は
アフィン・ブローアップ代数 $A[\frac{I}{a}]$ のスペクトルによる
アフィン開被覆をもつ。
\end{lemma}

\begin{proof}
第一の主張は、貼り合わせによる相対 Proj の構成から明らかである。
Constructions, Section
\ref{constructions-section-relative-proj-via-glueing} を参照。
$a \in I$ に対し、$a^{(1)}$ を、$a$ を次数 $1$ の元とみなしたものとして表す。
ここで次数付き代数は Rees 代数 $\bigoplus_{n \geq 0} I^n$ である。
これらの元は $A$ 上で Rees 代数を生成するので、
$\text{Proj}(\bigoplus_{d \geq 0} I^d)$ はアフィン開集合
$D_{+}(a^{(1)})$ で被覆される。アフィンスキーム $D_{+}(a^{(1)})$ は
アフィン・ブローアップ代数 $A' = A[\frac{I}{a}]$ のスペクトルである。
Algebra, Definition \ref{algebra-definition-blow-up} を参照。これで証明が終わる。
\end{proof}

\begin{lemma}
\label{divisors-lemma-flat-base-change-blowing-up}
\begin{slogan}
ブローアップは平坦基底変換と可換である。
\end{slogan}
$X_1 \to X_2$ をスキームの平坦射とする。$Z_2 \subset X_2$ を
閉部分スキームとし、$Z_1$ を $Z_2$ の $X_1$ における逆像とする。
$X'_i$ を $Z_i$ の $X_i$ におけるブローアップとする。このときスキームの
Cartesian 図式
$$
\xymatrix{
X_1' \ar[r] \ar[d] & X_2' \ar[d] \\
X_1 \ar[r] & X_2
}
$$
が存在する。
\end{lemma}

\begin{proof}
$\mathcal{I}_2$ を $Z_2$ の $X_2$ におけるイデアル層とする。
与えられた射を $g : X_1 \to X_2$ と表す。このときイデアル層
$\mathcal{I}_1$ は $Z_1$ のものであり、これは
$g^*\mathcal{I}_2 \to \mathcal{O}_{X_1}$ の像である
（逆像の定義による。Schemes, Definition
\ref{schemes-definition-inverse-image-closed-subscheme} を参照）。
Constructions, Lemma \ref{constructions-lemma-relative-proj-base-change} により、
$X_1 \times_{X_2} X_2'$ は $\bigoplus_{n \geq 0} g^*\mathcal{I}_2^n$ の
相対 Proj である。$g$ は平坦なので、写像
$g^*\mathcal{I}_2^n \to \mathcal{O}_{X_1}$ は単射で、その像は
$\mathcal{I}_1^n$ である。したがって
$X_1 \times_{X_2} X_2' = X_1'$ である。
\end{proof}

\begin{lemma}
\label{divisors-lemma-blowing-up-gives-effective-Cartier-divisor}
$X$ をスキームとし、$Z \subset X$ を閉部分スキームとする。
ブローアップ $b : X' \to X$ が $Z$ に沿う $X$ のものであるとき、
これは次の性質をもつ。
\begin{enumerate}
\item $b|_{b^{-1}(X \setminus Z)} : b^{-1}(X \setminus Z) \to X \setminus Z$
は同型である。
\item 例外因子 $E = b^{-1}(Z)$ は $X'$ 上の有効 Cartier 因子である。
\item 標準同型
$\mathcal{O}_{X'}(-1) = \mathcal{O}_{X'}(E)$
がある。
\end{enumerate}
\end{lemma}

\begin{proof}
ブローアップは開部分スキームへの制限と可換なので
（Lemma \ref{divisors-lemma-flat-base-change-blowing-up}）、第一の主張は
$X' = X$ であることを意味するにすぎない。ただし $Z = \emptyset$ とする。この場合、
イデアル層 $\mathcal{I} = \mathcal{O}_X$ におけるブローアップを取っており、
主張は Constructions, Example
\ref{constructions-example-trivial-proj} から従う。

\medskip\noindent
第二の主張は $X$ 上局所的なので、$X$ はアフィンであると仮定してよい。
$X = \Spec(A)$ および $Z = \Spec(A/I)$ と書く。
Lemma \ref{divisors-lemma-blowing-up-affine} により、$X'$ はアフィン・ブローアップ代数
$A' = A[\frac{I}{a}]$ のスペクトルで被覆される。
このとき $IA' = aA'$ であり、Algebra, Lemma
\ref{algebra-lemma-affine-blowup} によれば、$a$ は $A'$ の非零因子へ写る。
$Z$ の $\Spec(A')$ における逆像は
$\Spec(A'/IA') \subset \Spec(A')$ に対応するので、これで主張が証明される。

\medskip\noindent
標準写像
$\psi_{univ, 1} : b^*\mathcal{I} \to \mathcal{O}_{X'}(1)$ を考える。
Constructions, Definition
\ref{constructions-definition-relative-proj} に続く議論を参照。
これは同型 $\mathcal{I}_E \to \mathcal{O}_{X'}(1)$ を経由して分解すると主張する
（これにより最後の主張が証明される）。実際、上で述べたブローアップ代数
$A' = A[\frac{I}{a}]$ に対応するアフィン開集合上で、
$\psi_{univ, 1}$ は $A'$-加群写像
$$
I \otimes_A A'
\longrightarrow
\left(\Big(\bigoplus\nolimits_{d \geq 0} I^d\Big)_{a^{(1)}}\right)_1
$$
に対応する。ここで $a^{(1)}$ は Algebra, Definition
\ref{algebra-definition-blow-up} におけるものと同じである。
これが写像 $I \otimes_A A' \to IA' = aA'$ であることの検証は省略する。
\end{proof}

\begin{lemma}[ブローアップの普遍性]
\label{divisors-lemma-universal-property-blowing-up}
$X$ をスキームとし、$Z \subset X$ を閉部分スキームとする。
$\mathcal{C}$ を $(\Sch/X)$ の充満部分圏であって、$Y \to X$ のうち
$Z$ の逆像が $Y$ 上の有効 Cartier 因子となるものからなる部分圏とする。このとき
ブローアップ $b : X' \to X$ が $Z$ に沿う $X$ のものであれば、それは
$\mathcal{C}$ の終対象である。
\end{lemma}

\begin{proof}
Lemma \ref{divisors-lemma-blowing-up-gives-effective-Cartier-divisor} により、
$b : X' \to X$ は $\mathcal{C}$ の対象である。
$f : Y \to X$ を $\mathcal{C}$ の対象とする。一意な射
$Y \to X'$ が $X$ 上で存在することを示さなければならない。
$D = f^{-1}(Z)$ とおく。
$\mathcal{I} \subset \mathcal{O}_X$ を $Z$ のイデアル層とし、
$\mathcal{I}_D$ を $D$ のイデアル層とする。このとき
$f^*\mathcal{I} \to \mathcal{I}_D$ は、可逆 $\mathcal{O}_Y$-加群への全射である。
これは写像
$\psi : \bigoplus f^*\mathcal{I}^d \to \bigoplus \mathcal{I}_D^d$
へ拡張され、この写像は次数付き $\mathcal{O}_Y$-代数の間のものである。
（$\mathcal{I}_D^d = \mathcal{I}_D^{\otimes d}$ であることに
注意する。実際、$D$ は有効 Cartier 因子である。）
Constructions, Section \ref{constructions-section-relative-proj} の内容により、
三つ組 $(1, f : Y \to X, \psi)$ は射 $Y \to X'$ を $X$ 上で定める。制限
$$
Y \setminus D \longrightarrow X' \setminus b^{-1}(Z) = X \setminus Z
$$
は一意である。Lemma \ref{divisors-lemma-complement-effective-Cartier-divisor} により、
開部分 $Y \setminus D$ は $Y$ においてスキーム論的に稠密である。
したがって Morphisms, Lemma
\ref{morphisms-lemma-equality-of-morphisms} により射 $Y \to X'$ は一意である
（また $b$ は Constructions, Lemma
\ref{constructions-lemma-relative-proj-separated} により分離的である）。
\end{proof}

\begin{lemma}
\label{divisors-lemma-characterize-affine-blowup}
$b : X' \to X$ を、スキーム $X$ の閉部分スキーム $Z$ に沿う
ブローアップとする。$U = \Spec(A)$ を $X$ のアフィン開集合とし、
$I \subset A$ を $Z \cap U$ に対応するイデアルとする。
$a \in I$ とし、$x' \in X'$ を $U$ の点へ写る点とする。このとき、
この点 $x'$ がアフィン開集合
$U' = \Spec(A[\frac{I}{a}])$ に属するための必要十分条件は、
$a$ の $\mathcal{O}_{X', x'}$ における像が例外因子を切り出すことである。
\end{lemma}

\begin{proof}
$U'$ 上の例外因子は $a$ の像で切り出される。この像は
$A' = A[\frac{I}{a}]$ におけるものであるから、
一方の向きは明らかである。逆に、$a$ の
$\mathcal{O}_{X', x'}$ における像が $E$ を切り出すと仮定する。
$I$ の各元は $E$ を定義するイデアルの元へ写る。このイデアルは
$b^{-1}(U)$ 上のものである。したがって $I$ の元は $a$ で
$\mathcal{O}_{X', x'}$ において割り切れる。
$f \in I^n$ に対し、$f = \psi(f) a^n$ と書ける。ここで
ある $\psi(f) \in \mathcal{O}_{X', x'}$ を取る。
$a$ は $\mathcal{O}_{X', x'}$ の非零因子へ写るので、元
$\psi(f)$ はこの性質によって一意に特徴づけられることに注意する。
そこで
$$
A' \longrightarrow \mathcal{O}_{X', x'},\quad
f/a^n \longmapsto \psi(f)
$$
と定める。ここでは Algebra, Definition \ref{divisors-definition-blow-up} に続いて
与えられたブローアップ代数の記述を用いる。上の一意性により、
これは $A$-代数準同型である。この準同型は射
$\Spec(\mathcal{O}_{X', x"}) \to \Spec(A') = U'$ を与える。
ブローアップの普遍性
（Lemma \ref{divisors-lemma-universal-property-blowing-up}）により、これは
$X'$ 上の射であり、したがってもちろん $x' \in U'$ である。
\end{proof}

\begin{lemma}
\label{divisors-lemma-blow-up-effective-Cartier-divisor}
$X$ をスキームとし、$Z \subset X$ を有効 Cartier 因子とする。
$X$ の $Z$ におけるブローアップは $X$ の恒等射である。
\end{lemma}

\begin{proof}
ブローアップの普遍性
（Lemma \ref{divisors-lemma-universal-property-blowing-up}）から直ちに従う。
\end{proof}

\begin{lemma}
\label{divisors-lemma-blow-up-reduced-scheme}
$X$ をスキームとし、$\mathcal{I} \subset \mathcal{O}_X$ を
準連接イデアル層とする。$X$ が被約ならば、ブローアップ
$X'$、すなわち $X$ の $\mathcal{I}$ におけるブローアップは被約である。
\end{lemma}

\begin{proof}
Lemma \ref{divisors-lemma-blowing-up-affine} と Algebra, Lemma
\ref{algebra-lemma-blowup-reduced} を組み合わせればよい。
\end{proof}

\begin{lemma}
\label{divisors-lemma-blow-up-integral-scheme}
$X$ をスキームとし、$\mathcal{I} \subset \mathcal{O}_X$ を
零でない準連接イデアル層とする。$X$ が整ならば、ブローアップ
$X'$、すなわち $X$ の $\mathcal{I}$ におけるブローアップは整である。
\end{lemma}

\begin{proof}
Lemma \ref{divisors-lemma-blowing-up-affine} と Algebra, Lemma
\ref{algebra-lemma-blowup-domain} を組み合わせればよい。
\end{proof}

\begin{lemma}
\label{divisors-lemma-blow-up-and-irreducible-components}
$X$ をスキームとし、$Z \subset X$ を閉部分スキームとする。
$b : X' \to X$ を $X$ の $Z$ に沿うブローアップとする。このとき
$b$ は、$X'$ の既約成分の一般点の集合から、
$X$ の既約成分の一般点のうち $Z$ に属さないものの集合への全単射を誘導する。
\end{lemma}

\begin{proof}
例外因子 $E \subset X'$ は有効 Cartier 因子であり、
$X' \setminus E \to X \setminus Z$ は同型である。
Lemma \ref{divisors-lemma-blowing-up-gives-effective-Cartier-divisor} を参照。
したがって、次を示せば十分である。有効 Cartier 因子
$D \subset S$ がスキーム $S$ 上に与えられたとき、$S$ の既約成分の一般点はいずれも
$D$ に含まれない。これを見るため、$S$ をアフィン開被覆の各要素で
置き換えてよい。したがって Lemma
\ref{divisors-lemma-characterize-effective-Cartier-divisor} により、
$S = \Spec(A)$ かつ $D = V(f)$ と仮定してよい。ここで $f \in A$ は
非零因子である。このとき、$f$ がどの極小素イデアル
$\mathfrak p \subset A$ にも含まれないことを示さなければならない。
もし含まれるならば、$f$ は $R_\mathfrak p$ の極大イデアルに含まれる
非零因子へ写ることになり、Algebra, Lemma
\ref{algebra-lemma-minimal-prime-reduced-ring} に矛盾する。
\end{proof}

\begin{lemma}
\label{divisors-lemma-blow-up-pullback-effective-Cartier}
$X$ をスキームとし、$b : X' \to X$ を、ある閉部分スキームにおける
$X$ のブローアップとする。引き戻し $b^{-1}D$ は、すべての有効 Cartier 因子
$D \subset X$ に対して定義され、$b$ による有理型関数の引き戻しも定義される
（Definitions
\ref{divisors-definition-pullback-effective-Cartier-divisor} および
\ref{divisors-definition-pullback-meromorphic-sections}）。
\end{lemma}

\begin{proof}
Lemmas \ref{divisors-lemma-blowing-up-affine} および
\ref{divisors-lemma-characterize-effective-Cartier-divisor} により、これは次の代数的事実に
帰着する。$A$ を環、$I \subset A$ をイデアル、$a \in I$ とし、
$x \in A$ を非零因子とする。このとき $x$ の
$A[\frac{I}{a}]$ における像は非零因子である。実際、
$x (y/a^n) = 0$ が $A[\frac{I}{a}]$ において成り立つと仮定する。
このとき $a^mxy = 0$ が $A$ において、ある $m$ に対して成り立つ。
ゆえに $a^my = 0$ である。実際 $x$ は非零因子である。したがって望んだとおり、
$y/a^n$ は $A[\frac{I}{a}]$ において零である。
\end{proof}

\begin{lemma}
\label{divisors-lemma-blowing-up-two-ideals}
$X$ をスキームとし、
$\mathcal{I}, \mathcal{J} \subset \mathcal{O}_X$ を準連接イデアル層とする。
$b : X' \to X$ を $X$ の $\mathcal{I}$ におけるブローアップとする。
$b' : X'' \to X'$ を
$X'$ の $b^{-1}\mathcal{J} \mathcal{O}_{X'}$ におけるブローアップとする。
このとき $X'' \to X$ は、
$X$ の $\mathcal{I}\mathcal{J}$ におけるブローアップと標準的に同型である。
\end{lemma}

\begin{proof}
$E \subset X'$ を $b$ の例外因子とする。これは Lemma
\ref{divisors-lemma-blowing-up-gives-effective-Cartier-divisor} により有効 Cartier
因子である。このとき $(b')^{-1}E$ は Lemma
\ref{divisors-lemma-blow-up-pullback-effective-Cartier} により
$X''$ 上の有効 Cartier 因子である。
$E' \subset X''$ を $b'$ の例外因子とする（これも有効 Cartier 因子である）。
有効 Cartier 因子 $E'' = E' + (b')^{-1}E$ を考える。構成により
$E''$ のイデアルは
$(b \circ b')^{-1}\mathcal{I} (b \circ b')^{-1}\mathcal{J} \mathcal{O}_{X''}$
である。したがって Lemma \ref{divisors-lemma-universal-property-blowing-up} により、
$X''$ からブローアップ $c : Y \to X$ への標準射がある。このブローアップは
$X$ の $\mathcal{I}\mathcal{J}$ におけるものである。逆に、
$\mathcal{I}\mathcal{J}$ は可逆イデアルへ引き戻されるので、
$c^{-1}\mathcal{I}\mathcal{O}_Y$ は有効 Cartier 因子を定義する。
Lemma \ref{divisors-lemma-sum-closed-subschemes-effective-Cartier} を参照。
したがって Lemma \ref{divisors-lemma-universal-property-blowing-up} により、射
$c' : Y \to X'$ が $X$ 上で得られる。
すると
$(c')^{-1}b^{-1}\mathcal{J}\mathcal{O}_Y = c^{-1}\mathcal{J}\mathcal{O}_Y$
も有効 Cartier 因子を定義する。したがって射
$c'' : Y \to X''$ が $X'$ 上で得られる。この射が先に構成した射
$X'' \to Y$ の逆射であることの検証は省略する。
\end{proof}

\begin{lemma}
\label{divisors-lemma-blowing-up-projective}
$X$ をスキームとし、$\mathcal{I} \subset \mathcal{O}_X$ を
準連接イデアル層とする。$b : X' \to X$ を $X$ のイデアル層
$\mathcal{I}$ におけるブローアップとする。
$\mathcal{I}$ が有限型ならば、
\begin{enumerate}
\item $b : X' \to X$ は射影的射であり、
\item $\mathcal{O}_{X'}(1)$ は $b$-相対的に豊富な可逆層である。
\end{enumerate}
\end{lemma}

\begin{proof}
次数付き $\mathcal{O}_X$-代数の全射
$$
\text{Sym}_{\mathcal{O}_X}^*(\mathcal{I})
\longrightarrow
\bigoplus\nolimits_{d \geq 0} \mathcal{I}^d
$$
は、Constructions, Lemma
\ref{constructions-lemma-surjective-generated-degree-1-map-relative-proj} により
閉埋入
$$
X' = \underline{\text{Proj}}_X (\bigoplus\nolimits_{d \geq 0} \mathcal{I}^d)
\longrightarrow
\mathbf{P}(\mathcal{I}).
$$
を定める。したがって $b$ は射影的である。
Morphisms, Definition \ref{morphisms-definition-projective} を参照。
第二の主張は、例えば Morphisms, Lemma
\ref{morphisms-lemma-characterize-relatively-ample} における
相対的に豊富な可逆層の特徴づけから従う。詳細の一部は省略する。
\end{proof}

\begin{lemma}
\label{divisors-lemma-composition-finite-type-blowups}
\begin{slogan}
ブローアップの合成はブローアップである。
\end{slogan}
$X$ を準コンパクトかつ準分離的なスキームとする。
$Z \subset X$ を有限表示の閉部分スキームとする。
$b : X' \to X$ を中心 $Z$ のブローアップとする。
$Z' \subset X'$ を有限表示の閉部分スキームとする。
$X'' \to X'$ を中心 $Z'$ のブローアップとする。
このとき、有限表示の閉部分スキーム $Y \subset X$ であって、次を満たすものが存在する。
\begin{enumerate}
\item 集合論的に $Y = Z \cup b(Z')$ である。
\item 合成 $X'' \to X$ は、$X$ の $Y$ におけるブローアップと同型である。
\end{enumerate}
\end{lemma}

\begin{proof}
$Z \to X$ が有限表示であるという条件は、$Z$ が有限型準連接イデアル層
$\mathcal{I} \subset \mathcal{O}_X$ によって切り出されることを意味する。
Morphisms, Lemma
\ref{morphisms-lemma-closed-immersion-finite-presentation} を参照。
$\mathcal{A} = \bigoplus_{n \geq 0} \mathcal{I}^n$ と書くと、
$X' = \underline{\text{Proj}}(\mathcal{A})$ である。
Properties, Lemma
\ref{properties-lemma-quasi-coherent-finite-type-ideals} により、
$X \setminus Z$ は $X$ の準コンパクト開集合であることに注意する。
$b^{-1}(X \setminus Z) \to X \setminus Z$ は同型なので
（Lemma \ref{divisors-lemma-blowing-up-gives-effective-Cartier-divisor}）、
同じ結果から $b^{-1}(X \setminus Z) \setminus Z'$ は
$X'$ の準コンパクト開集合である。したがって
$U = X \setminus (Z \cup b(Z'))$ は $X$ の準コンパクト開集合である。
Lemma \ref{divisors-lemma-closed-subscheme-proj-finite-type} により、ある $d > 0$ と
有限型 $\mathcal{O}_X$-部分加群
$\mathcal{F} \subset \mathcal{I}^d$ が存在して、
$Z' = \underline{\text{Proj}}(\mathcal{A}/\mathcal{F}\mathcal{A})$ であり、
かつ $\mathcal{I}^d/\mathcal{F}$ の台は $X \setminus U$ に含まれる。

\medskip\noindent
$\mathcal{F} \subset \mathcal{I}^d$ は $\mathcal{O}_X$-部分加群なので、
$\mathcal{F} \subset \mathcal{I}^d \subset \mathcal{O}_X$ を
$X$ 上の有限型準連接イデアル層とみなせる。混同を避けるため、これを
$\mathcal{J} \subset \mathcal{O}_X$ と表す。
$\mathcal{I}^d / \mathcal{J}$ と $\mathcal{O}/\mathcal{I}^d$ の台は
$X \setminus U$ に含まれるので、$V(\mathcal{J})$ は
$X \setminus U$ に含まれる。逆に
$\mathcal{J} \subset \mathcal{I}^d$ なので、
$Z \subset V(\mathcal{J})$ である。
$X \setminus Z \cong X' \setminus b^{-1}(Z)$ 上で、イデアル層
$\mathcal{J}$ は $Z'$ を切り出す（下の表示式を参照）。
したがって $V(\mathcal{J})$ は $Z \cup b(Z')$ に等しい。
ゆえに集合論的にも $V(\mathcal{I}\mathcal{J}) = Z \cup b(Z')$ である。
さらに、$\mathcal{I}\mathcal{J}$ は有限型イデアル二つの積なので有限型である。
$X'' \to X$ は $X$ の $\mathcal{I}\mathcal{J}$ におけるブローアップと同型であると
主張する。これが示されれば $Y = V(\mathcal{I}\mathcal{J})$ とおくことで
補題の証明が終わる。

\medskip\noindent
まず、$X$ の $\mathcal{I}\mathcal{J}$ におけるブローアップは、
$X'$ の $b^{-1}\mathcal{J} \mathcal{O}_{X'}$ におけるブローアップと同じであることを
思い出そう。Lemma \ref{divisors-lemma-blowing-up-two-ideals} を参照。
したがって、$X'$ の $b^{-1}\mathcal{J} \mathcal{O}_{X'}$ における
ブローアップが、$X'$ の $Z'$ におけるブローアップと一致することを
示せば十分である。次を示す。
$$
b^{-1}\mathcal{J} \mathcal{O}_{X'} = \mathcal{I}_E^d \mathcal{I}_{Z'}
$$
これは $X''$ 上のイデアル層としての等式である。
$\mathcal{I}_E^d$ は有効 Cartier 因子 $dE$ を切り出すので、この等式は
望む主張を証明する。Lemmas
\ref{divisors-lemma-blow-up-effective-Cartier-divisor} および
\ref{divisors-lemma-blowing-up-two-ideals} を用いればよい。

\medskip\noindent
表示したイデアルの等式を見るため、局所的に議論してよい。
Lemma \ref{divisors-lemma-blowing-up-affine} と同じ記号 $A$, $I$, $a \in I$ を用いると、
$\mathcal{F}$ は $R$-部分加群 $M \subset I^d$ に対応し、これは
イデアル $J \subset R$ へ同型に写る。
$Z' = \underline{\text{Proj}}(\mathcal{A}/\mathcal{F}\mathcal{A})$
という条件は、$Z' \cap \Spec(A[\frac{I}{a}])$ が元
$m/a^d$, $m \in M$ で生成されるイデアルによって切り出されることを意味する。
元 $m \in M$ が関数 $f \in J$ に対応するとする。このとき、
アフィン・ブローアップ代数 $A' = A[\frac{I}{a}]$ において
$f = (a^dm)/a^d = a^d (m/a^d)$ である。したがって等式が成り立つ。
\end{proof}

\section{狭義変換}
\label{divisors-section-strict-transform}

\noindent
この節では、ブローアップによる狭義変換を簡単に論じる。
$S$ をスキームとし、$Z \subset S$ を閉部分スキームとする。
$b : S' \to S$ を $S$ の $Z$ におけるブローアップとし、
$E \subset S'$ を例外因子 $E = b^{-1}Z$ と表す。以下ではしばしば
$X$ という $S$ 上のスキームを考え、Cartesian 図式
$$
\xymatrix{
\text{pr}_{S'}^{-1}E \ar[r] \ar[d] &
X \times_S S' \ar[r]_-{\text{pr}_X} \ar[d]_{\text{pr}_{S'}} &
X \ar[d]^f \\
E \ar[r] & S' \ar[r] & S
}
$$
を作る。$E$ は有効 Cartier 因子なので
（Lemma \ref{divisors-lemma-blowing-up-gives-effective-Cartier-divisor}）、
$\text{pr}_{S'}^{-1}E \subset X \times_S S'$ は局所単項である
（Lemma \ref{divisors-lemma-pullback-locally-principal}）。
したがって、$\text{pr}_{S'}^{-1}E$ の
$X \times_S S'$ における補集合は逆コンパクトである
（Lemma \ref{divisors-lemma-complement-locally-principal-closed-subscheme}）。
従って、準連接 $\mathcal{O}_{X \times_S S'}$-加群 $\mathcal{G}$ に対し、
$\text{pr}_{S'}^{-1}E$ 上に台をもつ切断の部分層は準連接部分加群である。
Properties, Lemma
\ref{properties-lemma-sections-supported-on-closed-subset} を参照。
$\mathcal{G}$ が準連接代数層、例えば
$\mathcal{G} = \mathcal{O}_{X \times_S S'}$ であれば、この部分層は
$\mathcal{G}$ のイデアルである。

\begin{definition}
\label{divisors-definition-strict-transform}
上と同じ $Z \subset S$ および $f : X \to S$ のもとで、次のように定める。
\begin{enumerate}
\item 準連接 $\mathcal{O}_X$-加群 $\mathcal{F}$ が与えられたとき、
{\it 狭義変換}、すなわち $\mathcal{F}$ の $S$ の $Z$ における
ブローアップに関する狭義変換とは、商 $\mathcal{F}'$ のことである。
これは $\text{pr}_X^*\mathcal{F}$ を
$\text{pr}_{S'}^{-1}E$ 上に台をもつ切断の部分加群で割ったものである。
\item $X$ の{\it 狭義変換}とは、閉部分スキーム
$X' \subset X \times_S S'$ のことである。これは
$\mathcal{O}_{X \times_S S'}$ の切断からなる準連接イデアルであって、
$\text{pr}_{S'}^{-1}E$ 上に台をもつものによって切り出される。
\end{enumerate}
\end{definition}

\noindent
ブローアップに沿って狭義変換を取る操作は、ブローアップに用いた
閉部分スキームに依存し、射 $S' \to S$ だけには依存しないことに注意する。
この概念はしばしば $S$ の閉部分スキームに用いられる。
$X$ の狭義変換は実際には $X$ のブローアップである。

\begin{lemma}
\label{divisors-lemma-strict-transform}
Definition \ref{divisors-definition-strict-transform} の状況において、次が成り立つ。
\begin{enumerate}
\item 狭義変換 $X'$、すなわち $X$ の狭義変換は、$X$ の閉部分スキーム
$f^{-1}Z$ における $X$ のブローアップである。
\item 準連接 $\mathcal{O}_X$-加群 $\mathcal{F}$ に対し、狭義変換
$\mathcal{F}'$ は、$X' \to X \times_S S'$ に沿う順像と標準的に同型である。
ここで順像するのは $\mathcal{F}$ のブローアップ $X' \to X$ に関する
狭義変換である。
\end{enumerate}
\end{lemma}

\begin{proof}
$X'' \to X$ を $X$ の $f^{-1}Z$ におけるブローアップとする。
ブローアップの普遍性
（Lemma \ref{divisors-lemma-universal-property-blowing-up}）により、可換図式
$$
\xymatrix{
X'' \ar[r] \ar[d] & X \ar[d] \\
S' \ar[r] & S
}
$$
が存在し、したがって射 $X'' \to X \times_S S'$ が得られる。
第一の主張は、この射が像 $X'$ をもつ閉埋入であるということである。
問題は $X$ 上局所的である。そこで $X$ と $S$ はアフィンであると仮定してよい。
$S = \Spec(A)$、$X = \Spec(B)$ と書き、$Z$ はイデアル
$I \subset A$ によって切り出されるとする。$J = IB$ とおく。写像
$B \otimes_A \bigoplus_{n \geq 0} I^n \to \bigoplus_{n \geq 0} J^n$
は閉埋入 $X'' \to X \times_S S'$ を定める。Constructions, Lemmas
\ref{constructions-lemma-base-change-map-proj} および
\ref{constructions-lemma-surjective-graded-rings-generated-degree-1-map-proj}
を参照。この射が、上で普遍性から構成したものと同じであることの検証は省略する。
$a \in I$ を取る。これはアフィン開集合
$\Spec(A[\frac{I}{a}]) \subset S'$ に対応する。
Lemma \ref{divisors-lemma-blowing-up-affine} を参照。
$\Spec(A[\frac{I}{a}])$ の $X'$ における逆像は、$X$ の狭義変換において
$$
B' = (B \otimes_A A[\textstyle{\frac{I}{a}}])/a\text{-power-torsion}
$$
のスペクトルである。Properties, Lemma
\ref{properties-lemma-sections-supported-on-closed-subset} を参照。
他方、$b \in J$ を $a$ の像とすると、
$\Spec(B[\frac{J}{b}])$ は $\Spec(A[\frac{I}{a}])$ の
$X''$ における逆像である。
Algebra, Lemma \ref{algebra-lemma-blowup-base-change} により、開集合
$\Spec(B[\frac{J}{b}])$ は開部分スキーム
$\text{pr}_{S'}^{-1}(\Spec(A[\frac{I}{a}]))$ へ同型に写る。
これは $X'$ の開部分スキームである。
したがって $X'' \to X'$ は同型である。

\medskip\noindent
上の記号のもとで、$\mathcal{F}$ が $B$-加群 $N$ に対応するとする。
$\mathcal{F}$ の狭義変換は
$B \otimes_A A[\frac{I}{a}]$-加群
$$
N' = (N \otimes_A A[\textstyle{\frac{I}{a}}])/a\text{-power-torsion}
$$
に対応する。Properties, Lemma
\ref{properties-lemma-sections-supported-on-closed-subset} を参照。
$\mathcal{F}$ の $X$ の $f^{-1}Z$ におけるブローアップに関する狭義変換は、
$B[\frac{J}{b}]$-加群
$N \otimes_B B[\frac{J}{b}]/b\text{-power-torsion}$ に対応する。
上とまったく同じ方法で、これら二つの加群が同型であることが証明される。
詳細は省略する。
\end{proof}

\begin{lemma}
\label{divisors-lemma-strict-transform-flat}
Definition \ref{divisors-definition-strict-transform} の状況において、次が成り立つ。
\begin{enumerate}
\item $X$ が $S$ 上、$Z$ 上にあるすべての点で平坦ならば、
$X$ の狭義変換は基底変換 $X \times_S S'$ に等しい。
\item $\mathcal{F}$ を準連接 $\mathcal{O}_X$-加群とする。
$\mathcal{F}$ が $S$ 上、$Z$ 上にあるすべての点で平坦ならば、
狭義変換 $\mathcal{F}'$、すなわち $\mathcal{F}$ の狭義変換は引き戻し
$\text{pr}_X^*\mathcal{F}$ に等しい。
\end{enumerate}
\end{lemma}

\begin{proof}
(2) を証明すればよい。これは $S$ 上のスキームの狭義変換の定義により
(1) を含意するからである。問題は $X$ 上局所的である。
したがって $S = \Spec(A)$、$X = \Spec(B)$ と仮定し、
$\mathcal{F}$ は $B$-加群 $N$ に対応するとしてよい。このとき
$\mathcal{F}'$ は開集合 $\Spec(B \otimes_A A[\frac{I}{a}])$ 上で
加群に対応する。この開集合は $X \times_S S'$ のものであり、その加群は
$$
N' = (N \otimes_A A[\textstyle{\frac{I}{a}}])/a\text{-power-torsion}
$$
に対応する。Properties, Lemma
\ref{properties-lemma-sections-supported-on-closed-subset} を参照。
そこで、$a$-冪捩れが零であることを示さなければならない。これは
$N \otimes_A A[\frac{I}{a}]$ の捩れである。
$y \in N \otimes_A A[\frac{I}{a}]$ が
$a^n y = 0$ を満たすとする。
$\mathfrak q \subset B$ が素イデアルで $a \not \in \mathfrak q$ ならば、
$y$ は $(N \otimes_A A[\frac{I}{a}])_\mathfrak q$ において零へ写る。
他方、$a \in \mathfrak q$ ならば $N_\mathfrak q$ は平坦
$A$-加群であり、
$N_\mathfrak q \otimes_A A[\frac{I}{a}]
=(N \otimes_A A[\frac{I}{a}])_\mathfrak q$
には $a$-冪捩れがない
（$A[\frac{I}{a}]$ 自体にもないからである）。
従ってこの局所化においても $y$ は零へ写る。Algebra, Lemma
\ref{algebra-lemma-characterize-zero-local} により $y$ は零である。
\end{proof}

\begin{lemma}
\label{divisors-lemma-strict-transform-affine}
$S$ をスキームとし、$Z \subset S$ を閉部分スキームとする。
$b : S' \to S$ を $Z$ における $S$ のブローアップとする。
$g : X \to Y$ を $S$ 上のスキームのアフィン射とする。
$\mathcal{F}$ を $X$ 上の準連接層とする。
$g' : X \times_S S' \to Y \times_S S'$ を $g$ の基底変換とする。
$\mathcal{F}'$ を $\mathcal{F}$ の $b$ に関する狭義変換とする。
このとき $g'_*\mathcal{F}'$ は $g_*\mathcal{F}$ の狭義変換である。
\end{lemma}

\begin{proof}
Cohomology of Schemes, Lemma
\ref{coherent-lemma-affine-base-change} により
$g'_*\text{pr}_X^*\mathcal{F} = \text{pr}_Y^*g_*\mathcal{F}$ である。
$\mathcal{K} \subset \text{pr}_X^*\mathcal{F}$ を、
$Z$ の $X \times_S S'$ における逆像に台をもつ切断の部分層とする。
Properties, Lemma
\ref{properties-lemma-push-sections-supported-on-closed-subset} により、
順像 $g'_*\mathcal{K}$ は、$\text{pr}_Y^*g_*\mathcal{F}$ の切断の
部分層である。その台は $Z$ の $Y \times_S S'$ における逆像に含まれる。
$g'$ はアフィンなので
（Morphisms, Lemma \ref{morphisms-lemma-base-change-affine}）、
$g'_*$ は完全であり、これで結論が従う。
\end{proof}

\begin{lemma}
\label{divisors-lemma-strict-transform-different-centers}
$S$ をスキームとし、$Z \subset S$ を閉部分スキームとする。
$D \subset S$ を有効 Cartier 因子とする。
$Z' \subset S$ を、$Z$ と $D$ のイデアル層の積によって切り出される
閉部分スキームとする。$S' \to S$ を $S$ の $Z$ におけるブローアップとする。
\begin{enumerate}
\item $S$ の $Z'$ におけるブローアップは $S' \to S$ と同型である。
\item $f : X \to S$ をスキームの射とし、$\mathcal{F}$ を準連接
$\mathcal{O}_X$-加群とする。$\mathcal{F}$ が $f^{-1}D$ 上に台をもつ
零でない局所切断をもたないならば、$\mathcal{F}$ の
$Z$ におけるブローアップに関する狭義変換は、$\mathcal{F}$ の
$S$ の $Z'$ におけるブローアップに関する狭義変換と一致する。
\end{enumerate}
\end{lemma}

\begin{proof}
第一の主張は Lemmas \ref{divisors-lemma-blowing-up-two-ideals} および
\ref{divisors-lemma-blow-up-effective-Cartier-divisor} を組み合わせれば従う。
Lemma \ref{divisors-lemma-blowing-up-affine} を用いると、第二の主張は次の代数的問題に
翻訳される。$A$ を環、$I \subset A$ をイデアル、
$x \in A$ を非零因子、$a \in I$ とする。$M$ を $A$-加群であって、
$x$-捩れが零であるものとする。示すべきことは、$a$-冪捩れが
$M \otimes_A A[\frac{I}{a}]$ において
$xa$-冪捩れに等しいことである。
その理由は、写像 $A \to A[\frac{I}{a}]$ の核と余核が
$a$-冪捩れであり、従って $a$ を可逆化するとこの写像が同型になるからである。
ゆえに $M \to M \otimes_A A[\frac{I}{a}]$ の核と余核も
$a$-冪捩れである。これから結果が従う。
\end{proof}

\begin{lemma}
\label{divisors-lemma-strict-transform-composition-blowups}
$S$ をスキームとし、$Z \subset S$ を閉部分スキームとする。
$b : S' \to S$ を中心 $Z$ のブローアップとする。$Z' \subset S'$ を
閉部分スキームとし、$S'' \to S'$ を中心 $Z'$ のブローアップとする。
$Y \subset S$ を閉部分スキームであって、集合論的に
$Y = Z \cup b(Z')$ であり、合成 $S'' \to S$ が
$S$ の $Y$ におけるブローアップと同型になるものとする。
この状況で、任意のスキーム $X$ とその基底 $S$、および
$\mathcal{F} \in \QCoh(\mathcal{O}_X)$ が与えられたとき、次が成り立つ。
\begin{enumerate}
\item $\mathcal{F}$ の $S$ の $Y$ におけるブローアップに関する狭義変換は、
ブローアップ $S'' \to S'$ の $Z'$ における狭義変換に等しい。
後者で狭義変換されるのは、$\mathcal{F}$ のブローアップ
$S' \to S$、すなわち $S$ の $Z$ におけるブローアップに関する狭義変換である。
\item $X$ の $S$ の $Y$ におけるブローアップに関する狭義変換は、
ブローアップ $S'' \to S'$ の $Z'$ における狭義変換に等しい。
後者で狭義変換されるのは、$X$ のブローアップ
$S' \to S$、すなわち $S$ の $Z$ におけるブローアップに関する狭義変換である。
\end{enumerate}
\end{lemma}

\begin{proof}
$\mathcal{F}'$ を、$\mathcal{F}$ のブローアップ $S' \to S$、
すなわち $S$ の $Z$ におけるブローアップに関する狭義変換とする。
$\mathcal{F}''$ を、$\mathcal{F}'$ のブローアップ $S'' \to S'$、
すなわち $S'$ の $Z'$ におけるブローアップに関する狭義変換とする。
$\mathcal{G}$ を、$\mathcal{F}$ のブローアップ $S'' \to S$、
すなわち $S$ の $Y$ におけるブローアップに関する狭義変換とする。
射にも次の記号を付ける。
$$
\xymatrix{
X \times_S S'' \ar[r]_q \ar[d]^{f''} &
X \times_S S' \ar[r]_p \ar[d]^{f'} &
X \ar[d]^f \\
S'' \ar[r] & S' \ar[r] & S
}
$$
定義により全射 $p^*\mathcal{F} \to \mathcal{F}'$ と全射
$q^*\mathcal{F}' \to \mathcal{F}''$ があり、$q^*$ の右完全性により
これらは全射 $(p \circ q)^*\mathcal{F} \to \mathcal{F}''$ を与える。
また全射 $(p \circ q)^*\mathcal{F} \to \mathcal{G}$ もある。
従って、これら二つの全射が同じ核をもつことを示せば十分である。

\medskip\noindent
全射 $p^*\mathcal{F} \to \mathcal{F}'$ の核は
$(f \circ p)^{-1}Z$ 上に台をもつので、この写像は補集合の点で同型である。
従って $q^*p^*\mathcal{F} \to q^*\mathcal{F}'$ の核は
$(f \circ p \circ q)^{-1}Z$ 上に台をもつ。
$q^*\mathcal{F}' \to \mathcal{F}''$ の核は
$(f' \circ q)^{-1}Z'$ 上に台をもつ。
合わせると、$(p \circ q)^*\mathcal{F} \to \mathcal{F}''$ の核は
$(f \circ p \circ q)^{-1}Z \cup (f' \circ q)^{-1}Z' =
(f \circ p \circ q)^{-1}Y$
上に台をもつ。$\mathcal{G}$ の構成により、分解
$(p \circ q)^*\mathcal{F} \to \mathcal{F}'' \to \mathcal{G}$
が得られる。証明を終えるには、$\mathcal{F}''$ が
$(f \circ p \circ q)^{-1}(Y) =
(f \circ p \circ q)^{-1}Z \cup (f' \circ q)^{-1}Z'$
上に台をもつ零でない局所切断をもたないことを示せば十分である。
これは、$\mathcal{F}'$ に、$X \times_S S'$ 上で、かつ $S'$ 上で、
閉部分スキーム $Z'$ および有効 Cartier 因子 $b^{-1}Z$ とともに
Lemma \ref{divisors-lemma-strict-transform-different-centers} を適用すれば従う。
\end{proof}

\begin{lemma}
\label{divisors-lemma-strict-transform-universally-injective}
Definition \ref{divisors-definition-strict-transform} の状況で、完全列
$$
0 \to \mathcal{F}_1 \to \mathcal{F}_2 \to \mathcal{F}_3 \to 0
$$
が $X$ 上の準連接層からなり、任意の基底変換 $T \to S$ の後も
完全であると仮定する。このとき、$\mathcal{F}_i'$ の、任意のブローアップ
$S' \to S$ に関する狭義変換も短完全列
$0 \to \mathcal{F}'_1 \to \mathcal{F}'_2 \to \mathcal{F}'_3 \to 0$
をなす。
\end{lemma}

\begin{proof}
$S$ と $X$ 上で局所化し、両方がアフィンであると仮定してよい。
このとき $\mathcal{F}_i$ を $S$ へ順像してよい。
Lemma \ref{divisors-lemma-strict-transform-affine} を参照。
ブローアップは射 $1 : S \to S$ であると仮定してよい。この射は
有効 Cartier 因子 $D \subset S$ に付随する。このとき代数への翻訳は次である。
$A$ を環とし、
$0 \to M_1 \to M_2 \to M_3 \to 0$ を $A$-加群の普遍完全列とする。
$a\in A$ とする。このとき列
$$
0 \to
M_1/a\text{-power torsion} \to
M_2/a\text{-power torsion} \to
M_3/a\text{-power torsion} \to 0
$$
も完全である。実際、最後の写像の全射性と最初の写像の単射性は直ちに分かる。
問題は中央における完全性である。
$x \in M_2$ が $M_3/a\text{-power torsion}$ において零へ写ると仮定する。
このとき $y = a^n x \in M_1$ であり、これはある $n$ に対して成り立つ。
すると $y$ は $M_2/a^nM_2$ において零へ写る。
$M_1 \to M_2$ は普遍単射なので、$y$ は
$M_1/a^nM_1$ において零へ写る。従って
$y = a^n z$ であり、これはある $z \in M_1$ に対して成り立つ。
ゆえに $a^n(x - y) = 0$ である。
従って $y$ は $x$ の $M_2/a\text{-power torsion}$ における類へ写る。
これが望む結論である。
\end{proof}

\section{許容ブローアップ}
\label{divisors-section-admissible-blowups}

\noindent
ブローアップをもう少し制御するために、次の標準的な用語を導入する。

\begin{definition}
\label{divisors-definition-admissible-blowup}
$X$ をスキームとする。$U \subset X$ を開部分スキームとする。射
$X' \to X$ が {\it $U$-許容ブローアップ}であるとは、有限表示の
閉埋め込み $Z \to X$ であって、$Z$ が $U$ と交わらず、
$X'$ が $X$ の $Z$ に沿うブローアップと同型になるものが存在することをいう。
\end{definition}

\noindent
$Z \to X$ が有限表示であるための必要十分条件は、イデアル層
$\mathcal{I}_Z \subset \mathcal{O}_X$ が有限型であることであった。
Morphisms, Lemma
\ref{morphisms-lemma-closed-immersion-finite-presentation} を参照。
特に、$U$-許容ブローアップは射影的射である。
Lemma \ref{divisors-lemma-blowing-up-projective} を参照。
同じ射を生じる中心が複数あり得ることに注意せよ。
従って、要請しているのは、$U$ と交わらず $X'$ を生じる何らかの中心が
存在することだけである。
最後に、射 $b : X' \to X$ は $U$ 上の同型なので
(Lemma \ref{divisors-lemma-blowing-up-gives-effective-Cartier-divisor} を参照)、
しばしば記法を濫用して、$U$ を $X'$ の開部分スキームでもあるとみなす。

\begin{lemma}
\label{divisors-lemma-composition-admissible-blowups}
\begin{slogan}
許容ブローアップは合成について安定である。
\end{slogan}
$X$ を準コンパクトかつ準分離なスキームとする。
$U \subset X$ を準コンパクトな開部分スキームとする。
$b : X' \to X$ を $U$-許容ブローアップとする。
$X'' \to X'$ を $U$-許容ブローアップとする。
このとき合成 $X'' \to X$ は $U$-許容ブローアップである。
\end{lemma}

\begin{proof}
より精密な Lemma \ref{divisors-lemma-composition-finite-type-blowups} から直ちに従う。
\end{proof}

\begin{lemma}
\label{divisors-lemma-extend-admissible-blowups}
$X$ を準コンパクトかつ準分離なスキームとする。
$U, V \subset X$ を準コンパクトな開部分スキームとする。
$b : V' \to V$ を $U \cap V$-許容ブローアップとする。
このとき、$U$-許容ブローアップ $X' \to X$ であって、
その $V$ への制限が $V'$ であるものが存在する。
\end{lemma}

\begin{proof}
$\mathcal{I} \subset \mathcal{O}_V$ を、$V(\mathcal{I})$ が
$U \cap V$ と交わらず、$V'$ が $V$ の $\mathcal{I}$ に沿う
ブローアップと同型になるような有限型準連接イデアル層とする。
$\mathcal{I}' \subset \mathcal{O}_{U \cup V}$ を、その $U$ への制限が
$\mathcal{O}_U$ であり、その $V$ への制限が $\mathcal{I}$ であるような
準連接イデアル層とする
(Sheaves, Section \ref{sheaves-section-glueing-sheaves} を参照)。
Properties, Lemma \ref{properties-lemma-extend} により、
有限型準連接イデアル層 $\mathcal{J} \subset \mathcal{O}_X$ であって、
その $U \cup V$ への制限が $\mathcal{I}'$ であるものが存在する。
これで補題が従う。
\end{proof}

\begin{lemma}
\label{divisors-lemma-dominate-admissible-blowups}
$X$ を準コンパクトかつ準分離なスキームとする。
$U \subset X$ を準コンパクトな開部分スキームとする。
$b_i : X_i \to X$, $i = 1, \ldots, n$ を $U$-許容ブローアップとする。
次を満たす $U$-許容ブローアップ $b : X' \to X$ が存在する：
(a) $b$ は $X' \to X_i \to X$ と分解し、これは各
$i = 1, \ldots, n$ に対して成り立ち、
(b) 各射 $X' \to X_i$ は $U$-許容ブローアップである。
\end{lemma}

\begin{proof}
$\mathcal{I}_i \subset \mathcal{O}_X$ を、$V(\mathcal{I}_i)$ が
$U$ と交わらず、$X_i$ が $X$ の $\mathcal{I}_i$ に沿う
ブローアップと同型になるような有限型準連接イデアル層とする。
$\mathcal{I} = \mathcal{I}_1 \cdot \ldots \cdot \mathcal{I}_n$
と置き、$X'$ を $X$ の $\mathcal{I}$ に沿うブローアップとする。
すると Lemma \ref{divisors-lemma-blowing-up-two-ideals} により、
$X' \to X$ は $b_i$ を経由して分解する。
\end{proof}

\begin{lemma}
\label{divisors-lemma-separate-disjoint-opens-by-blowing-up}
\begin{slogan}
ブローアップによって既約成分を分離する。
\end{slogan}
$X$ を準コンパクトかつ準分離なスキームとする。
$U, V$ を $X$ の、互いに交わらない準コンパクトな開部分スキームとする。
このとき、$U \cup V$-許容ブローアップ $b : X' \to X$ であって、
$X'$ が開部分スキームの非交和
$X' = X'_1 \amalg X'_2$ となり、$b^{-1}(U) \subset X'_1$ および
$b^{-1}(V) \subset X'_2$ を満たすものが存在する。
\end{lemma}

\begin{proof}
有限型準連接イデアル層 $\mathcal{I}$ および、それぞれ
$\mathcal{J}$ を、$X \setminus U = V(\mathcal{I})$ および、それぞれ
$X \setminus V = V(\mathcal{J})$ となるように選ぶ。
Properties, Lemma
\ref{properties-lemma-quasi-coherent-finite-type-ideals} を参照。
このとき $V(\mathcal{I}\mathcal{J}) = X$ が集合論的に成り立つので、
$\mathcal{I}\mathcal{J}$ は局所冪零イデアル層である。
$\mathcal{I}$ と $\mathcal{J}$ は有限型であり、$X$ は準コンパクトなので、
$n > 0$ で $\mathcal{I}^n \mathcal{J}^n = 0$ を満たすものが存在する。
$\mathcal{I}$ を $\mathcal{I}^n$ で、$\mathcal{J}$ を $\mathcal{J}^n$ で
置き換えるものとする。従って $\mathcal{I} \mathcal{J} = 0$ である。
$b : X' \to X$ を $\mathcal{I} + \mathcal{J}$ に沿うブローアップとする。
これは $U \cup V$-許容である。実際、
$V(\mathcal{I} + \mathcal{J}) = X \setminus U \cup V$ である。
$X'$ が開部分スキームの非交和
$X' = X'_1 \amalg X'_2$ であって、
$b^{-1}\mathcal{I}|_{X'_2} = 0$ および
$b^{-1}\mathcal{J}|_{X'_1} = 0$ を満たすことを示す。
これにより補題が証明される。

\medskip\noindent
Lemma \ref{divisors-lemma-blowing-up-affine} におけるブローアップの記述を用いる。
$U = \Spec(A) \subset X$ がアフィン開であり、
$\mathcal{I}|_U$ および、それぞれ $\mathcal{J}|_U$ が、有限生成イデアル
$I \subset A$ および、それぞれ $J \subset A$ に対応すると仮定する。
このとき
$$
b^{-1}(U) = \text{Proj}(A \oplus (I + J) \oplus (I + J)^2 \oplus \ldots)
$$
これはアフィン開部分集合 $A[\frac{I + J}{x}]$ および
$A[\frac{I + J}{y}]$ によって被覆される。ここで $x \in I$ および
$y \in J$ である。
$x \in I$ は $A[\frac{I + J}{x}]$ における非零因子であり、
$IJ = 0$ なので、$J A[\frac{I + J}{x}] = 0$ を得る。
$y \in J$ は $A[\frac{I + J}{y}]$ における非零因子であり、
$IJ = 0$ なので、$I A[\frac{I + J}{y}] = 0$ を得る。さらに、
$$
\Spec(A[\textstyle{\frac{I + J}{x}}]) \cap
\Spec(A[\textstyle{\frac{I + J}{y}}]) =
\Spec(A[\textstyle{\frac{I + J}{xy}}]) = \emptyset
$$
である。なぜなら $xy$ は非零因子であると同時に零だからである。
従って $b^{-1}(U)$ は、開部分スキーム $U_1$
（標準開 $\Spec(A[\frac{I + J}{x}])$ を $x \in I$ にわたって
合併したもの）と、開部分スキーム $U_2$
（アフィン開 $\Spec(A[\frac{I + J}{y}])$ を $y \in J$ にわたって
合併したもの）との非交和である。
$b^{-1}\mathcal{I}\mathcal{O}_{X'}$ は $U_2$ 上で零に制限され、
$b^{-1}\mathcal{I}\mathcal{O}_{X'}$ は $U_1$ 上で零に制限されることを
既に見た。これらの開部分スキームが大域的な開部分スキーム
$X'_1$ および $X'_2$ に貼り合わさることの確認は省略する。
\end{proof}

\begin{lemma}
\label{divisors-lemma-blowing-up-denominators}
$X$ を局所 Noether スキームとする。
$\mathcal{L}$ を可逆 $\mathcal{O}_X$-加群とする。
$s$ を $\mathcal{L}$ の正則有理型切断とする。
$U \subset X$ を、$s$ が $\mathcal{L}$ の $U$ 上の切断に対応するような
最大の開部分スキームとする。
ブローアップ $b : X' \to X$ であって、$s$ の分母イデアルに沿うものは
$U$-許容である。有効 Cartier 因子 $D \subset X'$ と同型
$$
b^*\mathcal{L} = \mathcal{O}_{X'}(D - E),
$$
が存在する。ここで $E \subset X'$ は例外因子であり、この同型を介して
有理型切断 $b^*s$ は有理型切断
$1_D \otimes (1_E)^{-1}$ に対応する。
\end{lemma}

\begin{proof}
Definition
\ref{divisors-definition-regular-meromorphic-ideal-denominators}
における分母イデアルの定義から、$b$ が $U$-許容ブローアップであることが
直ちに分かる。記法 $1_D$, $1_E$, および $\mathcal{O}_{X'}(D - E)$
については Definition
\ref{divisors-definition-invertible-sheaf-effective-Cartier-divisor}
を参照せよ。引き戻し $b^*s$ は Lemmas
\ref{divisors-lemma-blow-up-pullback-effective-Cartier} および
\ref{divisors-lemma-meromorphic-sections-pullback} によって定義される。
従って補題の主張は意味をなす。
最後の主張は、$b^*s$ が $b^*\mathcal{L}(E)$ の大域正則切断であり、
その零スキームが $D$ である、と言い換えられる。
これは $D$ を一意に定めるので、補題を証明するには
$X$ と $X'$ 上でアフィン局所的に議論してよい。
$X = \Spec(A)$ がアフィンであり、$\mathcal{L} = \mathcal{O}_X$ であると
仮定する。このとき $s$ は正則有理型関数であり、さらに縮小して
$s = a'/a$ かつ $a', a \in A$ が非零因子であると仮定してよい。
このとき $s$ の分母イデアルはイデアル
$I = \{x \in A \mid xa' \in aA\}$ に対応する。
$X'$ はアフィン・ブローアップ代数 $A' = A[\frac{I}{x}]$ の
スペクトルによって被覆されることを思い出そう。ここで $x \in I$ である
(Lemma \ref{divisors-lemma-blowing-up-affine})。$x \in I$ を固定し、
$xa' = a a''$ と書く。ここで $a'' \in A$ である。
因子 $E \subset X'$ は、$x \in A'$ によって $A'$ のスペクトル上で
切り出される。従って $1/x$ は $\mathcal{O}_{X'}(E)$ の
$\Spec(A')$ 上での生成元である。
最後に、$A'$ の全商環において $a'/a = a''/x$ である。
従って $b^*s = a'/a$ は $\mathcal{O}_{X'}(E)$ の正則切断へ制限され、
$\Spec(A')$ 上では $a''/x$ によって与えられる。これで証明が終わる。
（因子 $D \cap \Spec(A')$ は $a''$ の $A'$ における像によって
切り出される。）
\end{proof}

\section{ブローアップと平坦性}
\label{divisors-section-blowup-flat}

\noindent
More on Algebra, Section \ref{more-algebra-section-blowup-flat}
で始めた議論を続ける。More on Flatness, Section
\ref{flat-section-blowup-flat} でさらに結果を証明する。


\begin{lemma}
\label{divisors-lemma-strict-transform-blowup-fitting-ideal}
$S$ をスキームとする。$\mathcal{F}$ を有限型準連接
$\mathcal{O}_S$-加群とする。$Z_k \subset S$ を
$\text{Fit}_k(\mathcal{F})$ によって切り出される閉部分スキームとする。
Section \ref{divisors-section-fitting-ideals} を参照。
$S' \to S$ を $S$ の $Z_k$ に沿うブローアップとし、
$\mathcal{F}'$ を $\mathcal{F}$ の狭義変換とする。
このとき $\mathcal{F}'$ は局所的に $\leq k$ 個の切断によって生成できる。
\end{lemma}

\begin{proof}
$\mathcal{F}'$ が局所的に $\leq k$ 個の切断によって生成できるための
必要十分条件は $\text{Fit}_k(\mathcal{F}') = \mathcal{O}_{S'}$ で
あったことを思い出そう。
Lemma \ref{divisors-lemma-fitting-ideal-generate-locally} を参照。
従って、この補題は More on Algebra, Lemma
\ref{more-algebra-lemma-blowup-fitting-ideal} の翻訳である。
\end{proof}

\begin{lemma}
\label{divisors-lemma-strict-transform-blowup-fitting-ideal-locally-free}
$S$ をスキームとする。$\mathcal{F}$ を有限型準連接
$\mathcal{O}_S$-加群とする。$Z_k \subset S$ を
$\text{Fit}_k(\mathcal{F})$ によって切り出される閉部分スキームとする。
Section \ref{divisors-section-fitting-ideals} を参照。
$\mathcal{F}$ が階数 $k$ の局所自由加群であり、それが
$S \setminus Z_k$ 上で成り立つと仮定する。
$S' \to S$ を $S$ の $Z_k$ に沿うブローアップとし、
$\mathcal{F}'$ を $\mathcal{F}$ の狭義変換とする。
このとき $\mathcal{F}'$ は階数 $k$ の局所自由加群である。
\end{lemma}

\begin{proof}
More on Algebra, Lemma
\ref{more-algebra-lemma-blowup-fitting-ideal-locally-free}
の翻訳である。
\end{proof}

\begin{lemma}
\label{divisors-lemma-blowup-fitting-ideal}
$X$ をスキームとする。$\mathcal{F}$ を有限表示
$\mathcal{O}_X$-加群とする。$U \subset X$ をスキーム論的に稠密な開で、
$\mathcal{F}|_U$ が定階数 $r$ の有限局所自由加群となるものとする。
このとき次が成り立つ。
\begin{enumerate}
\item $b : X' \to X$、すなわち $X$ の第 $r$ Fitting イデアル
（$\mathcal{F}$ のもの）に沿うブローアップは $U$-許容である。
\item 狭義変換 $\mathcal{F}'$、すなわち $\mathcal{F}$ の
$b$ に関する狭義変換は階数 $r$ の局所自由加群である。
\item 核 $\mathcal{K}$、すなわち全射
$b^*\mathcal{F} \to \mathcal{F}'$ の核は
有限表示であり、$\mathcal{K}|_U = 0$ である。
\item $b^*\mathcal{F}$ と $\mathcal{K}$ は完全
$\mathcal{O}_{X'}$-加群であり、Tor 次元は $\leq 1$ である。
\end{enumerate}
\end{lemma}

\begin{proof}
Lemma \ref{divisors-lemma-fitting-ideal-of-finitely-presented} により、イデアル
$\text{Fit}_r(\mathcal{F})$ は有限型であり、
Lemma \ref{divisors-lemma-fitting-ideal-finite-locally-free} により、その $U$
への制限は $\mathcal{O}_U$ に等しい。従って
$b : X' \to X$ は $U$-許容である。
Definition \ref{divisors-definition-admissible-blowup} を参照。

\medskip\noindent
Lemma \ref{divisors-lemma-fitting-ideal-finite-locally-free} により、
$\text{Fit}_{r - 1}(\mathcal{F})$ の $U$ への制限は零である。
$U$ はスキーム論的に稠密なので、
$\text{Fit}_{r - 1}(\mathcal{F}) = 0$ が $X$ 全体で成り立つと分かる。
従って Lemma \ref{divisors-lemma-fitting-ideal-finite-locally-free} により、
$\mathcal{F}$ は階数 $r$ の局所自由加群であり、それが成り立つのは
第 $r$ Fitting イデアル（$\mathcal{F}$ のもの）によって
切り出される部分スキームの補集合上である
（この補集合は $U$ より大きいことがあり得るので、論証にはこの段階が
必要であった）。従って
Lemma \ref{divisors-lemma-strict-transform-blowup-fitting-ideal-locally-free}
により、狭義変換
$$
b^*\mathcal{F} \longrightarrow \mathcal{F}'
$$
は階数 $r$ の局所自由加群である。この写像の核 $\mathcal{K}$ は
ブローアップ $b$ の例外因子上に台をもち、従って
$\mathcal{K}|_U = 0$ である。最後に、$\mathcal{F}'$ は有限局所自由であり、
表示された矢印は全射なので、$X'$ 上で局所的に
$b^*\mathcal{F}$ を $\mathcal{K}$ と $\mathcal{F}'$ の直和として書ける。
$b^*\mathcal{F}'$ は有限表示なので
(Modules, Lemma \ref{modules-lemma-pullback-finite-presentation})、
$\mathcal{K}$ についても同じことが成り立つ。

\medskip\noindent
Tor 次元に関する主張は More on Algebra, Lemma
\ref{more-algebra-lemma-fitting-ideals-and-pd1} から従う。
\end{proof}




\section{修正}
\label{divisors-section-modifications}

\noindent
この節では、「修正の後ではこれこれが成り立つ」という型の結果を集める。
後に、修正はブローアップによって支配されることを見る
(More on Flatness, Lemma
\ref{flat-lemma-dominate-modification-by-blowup})。

\begin{lemma}
\label{divisors-lemma-filter-after-modification}
$X$ を整スキームとする。$\mathcal{E}$ を有限局所自由
$\mathcal{O}_X$-加群とする。修正 $f : X' \to X$ であって、
$f^*\mathcal{E}$ が、その逐次商が可逆 $\mathcal{O}_{X'}$-加群である
フィルトレーションをもつものが存在する。
\end{lemma}

\begin{proof}
$r$ を $\mathcal{E}$ の階数とし、これに関する帰納法で証明する。
$r = 1$ または $r = 0$ なら補題は明らかである。$r > 1$ と仮定する。
$P = \mathbf{P}(\mathcal{E})$ とし、その構造射を
$\pi : P \to X$ とする。Constructions, Section
\ref{constructions-section-projective-bundle} を参照。
このとき $\pi$ は固有である
(Lemma \ref{divisors-lemma-relative-proj-proper})。
標準的な全射
$$
\pi^*\mathcal{E} \to \mathcal{O}_P(1)
$$
があり、その核は階数 $r - 1$ の有限局所自由加群である。
空でない開部分スキーム $U \subset X$ であって、
$\mathcal{E}|_U \cong \mathcal{O}_U^{\oplus r}$ となるものを選ぶ。
このとき $P_U = \pi^{-1}(U)$ は $\mathbf{P}^{r - 1}_U$ と同型である。
特に、切断 $s : U \to P_U$ であって $\pi$ の切断であるものが存在する。
$X' \subset P$ を射 $U \to P_U \to P$ のスキーム論的像とする。
このとき $X'$ は整であり
(Morphisms, Lemma
\ref{morphisms-lemma-scheme-theoretic-image-reduced})、
射 $f = \pi|_{X'} : X' \to X$ は固有であり
(Morphisms, Lemmas
\ref{morphisms-lemma-closed-immersion-proper} および
\ref{morphisms-lemma-composition-proper})、
$f^{-1}(U) \to U$ は同型である。従って $f$ は修正である
(Morphisms, Definition \ref{morphisms-definition-modification})。
構成により、引き戻し $f^*\mathcal{E}$ は二段階フィルトレーションをもつ。
その商は $\mathcal{O}_P(1)|_{X'}$ に等しいので可逆であり、その部分加群
$\mathcal{E}'$ は階数 $r - 1$ の局所自由加群である。
帰納法により、修正 $g : X'' \to X'$ であって、
$g^*\mathcal{E}'$ が補題の主張にあるようなフィルトレーションをもつものを
見いだせる。従って $f \circ g : X'' \to X$ が求める修正である。
\end{proof}

\begin{lemma}
\label{divisors-lemma-extend-rational-map-after-modification}
$S$ をスキームとする。$X$, $Y$ を $S$ 上のスキームとする。
$X$ は Noether であり、$Y$ は $S$ 上固有であると仮定する。
$S$-有理写像 $f : U \to Y$、すなわち $X$ から $Y$ への写像が与えられた
とする。このとき射 $p : X' \to X$ と $S$-射 $f' : X' \to Y$ であって、
次を満たすものが存在する。
\begin{enumerate}
\item $p$ は固有であり、$p^{-1}(U) \to U$ は同型である。
\item $f'|_{p^{-1}(U)}$ は $f \circ p|_{p^{-1}(U)}$ に等しい。
\end{enumerate}
\end{lemma}

\begin{proof}
$j : U \to X$ で包含射を表す。$X' \subset Y \times_S X$ を
$(f, j) : U \to Y \times_S X$ のスキーム論的像とする
(Morphisms, Definition
\ref{morphisms-definition-scheme-theoretic-image})。
射影 $g : Y \times_S X \to X$ は固有である
(Morphisms, Lemma \ref{morphisms-lemma-base-change-proper})。
合成 $p : X' \to X$、すなわち $X' \to Y \times_S X$ と $g$ の合成は
固有である
(Morphisms, Lemmas \ref{morphisms-lemma-closed-immersion-proper} および
\ref{morphisms-lemma-composition-proper})。
$g$ は分離的であり、$U \subset X$ は逆コンパクトである
（$X$ が Noether だからである）ので、
Morphisms, Lemma
\ref{morphisms-lemma-scheme-theoretic-image-of-partial-section}
により $p^{-1}(U) \to U$ は同型である。
一方、合成 $f' : X' \to Y$、すなわち $X' \to Y \times_S X$ と
射影 $Y \times_S X \to Y$ の合成は、$f$ と $p^{-1}(U)$ 上で一致する。
\end{proof}
